\input{preamble}

% OK, start here.
%
\begin{document}

\title{Descente}

\maketitle

\phantomsection
\label{section-phantom}

\tableofcontents

\section{Introduction}
\label{section-introduction}

\noindent
Dans le chapitre consacré aux topologies sur les schémas
(voir Topologies, section \ref{topologies-section-introduction}), nous avons
introduit les recouvrements de Zariski, étales, fppf, lisses, syntomiques et
fpqc de schémas. Dans le présent chapitre, nous examinons les structures sur
les schémas qui peuvent être descendues au moyen de tels recouvrements. Voir, par
exemple, \cite{Gr-I}, \cite{Gr-II}, \cite{Gr-III}, \cite{Gr-IV}, \cite{Gr-V}
et \cite{Gr-VI}. Il s'agit également d'introduire les notions de descente, de
donnée de descente et de donnée de descente effective dans le cadre moins
formel des questions de descente pour les faisceaux quasi-cohérents, les
schémas, etc. La notion formelle de champ sur un site est étudiée dans le
chapitre consacré aux champs (voir Champs, section
\ref{stacks-section-introduction}).

\section{Données de descente pour les faisceaux quasi-cohérents}
\label{section-equivalence}

\noindent
Dans ce chapitre, nous adoptons la convention suivant laquelle les
projections $\text{pr}_i : X \times \ldots \times X \to X$ sont numérotées à
partir de $i = 0$. Ainsi, nous avons
$\text{pr}_0, \text{pr}_1 : X \times X  \to X$,
$\text{pr}_0, \text{pr}_1, \text{pr}_2 : X \times X \times X  \to X$, etc.


\begin{definition}
\label{definition-descent-datum-quasi-coherent}
Soit $S$ un schéma. Soit $\{f_i : S_i \to S\}_{i \in I}$ une famille de
morphismes de but $S$.
\begin{enumerate}
\item Une {\it donnée de descente $(\mathcal{F}_i, \varphi_{ij})$ pour les
faisceaux quasi-cohérents} relativement à la famille donnée est constituée
d'un faisceau quasi-cohérent $\mathcal{F}_i$ sur $S_i$ pour chaque $i \in I$
et d'un isomorphisme de $\mathcal{O}_{S_i \times_S S_j}$-modules
quasi-cohérents
$\varphi_{ij} : \text{pr}_0^*\mathcal{F}_i \to \text{pr}_1^*\mathcal{F}_j$
pour chaque paire $(i, j) \in I^2$,
tel que, pour tout triplet d'indices $(i, j, k) \in I^3$, le diagramme
$$
\xymatrix{
\text{pr}_0^*\mathcal{F}_i \ar[rd]_{\text{pr}_{01}^*\varphi_{ij}}
\ar[rr]_{\text{pr}_{02}^*\varphi_{ik}} & &
\text{pr}_2^*\mathcal{F}_k \\
& \text{pr}_1^*\mathcal{F}_j \ar[ru]_{\text{pr}_{12}^*\varphi_{jk}} &
}
$$
de $\mathcal{O}_{S_i \times_S S_j \times_S S_k}$-modules soit commutatif.
C'est la {\it condition de cocycle}.
\item Un {\it morphisme
$\psi : (\mathcal{F}_i, \varphi_{ij}) \to
(\mathcal{F}'_i, \varphi'_{ij})$ de données de descente} est une famille
$\psi = (\psi_i)_{i\in I}$ de morphismes de $\mathcal{O}_{S_i}$-modules
$\psi_i : \mathcal{F}_i \to \mathcal{F}'_i$ telle que tous les diagrammes
$$
\xymatrix{
\text{pr}_0^*\mathcal{F}_i \ar[r]_{\varphi_{ij}} \ar[d]_{\text{pr}_0^*\psi_i}
& \text{pr}_1^*\mathcal{F}_j \ar[d]^{\text{pr}_1^*\psi_j} \\
\text{pr}_0^*\mathcal{F}'_i \ar[r]^{\varphi'_{ij}} &
\text{pr}_1^*\mathcal{F}'_j \\
}
$$
soient commutatifs.
\end{enumerate}
\end{definition}

\noindent
Un bon exemple à garder à l'esprit est le suivant. Supposons que
$S = \bigcup S_i$ soit un recouvrement ouvert. Nous avons alors rencontré
les données de descente pour les faisceaux d'ensembles dans Faisceaux,
section \ref{sheaves-section-glueing-sheaves}, où nous les avons appelées
``données de recollement pour les faisceaux d'ensembles relativement au
recouvrement donné''. Nous avons en outre démontré que la catégorie des
données de recollement est équivalente à la catégorie des faisceaux sur $S$.
Nous établirons l'analogue dans le cadre ci-dessus lorsque
$\{S_i \to S\}_{i\in I}$ est un recouvrement fpqc.

\medskip\noindent
Dans le cas extrême où le recouvrement $\{S \to S\}$ est donné par
$\text{id}_S$, une donnée de descente est nécessairement de la forme
$(\mathcal{F}, \text{id}_\mathcal{F})$. La condition de cocycle garantit que
l'identité de $\mathcal{F}$ est le seul morphisme admis dans ce cas. Le lemme
suivant montre notamment qu'à tout faisceau quasi-cohérent de
$\mathcal{O}_S$-modules est associée une unique donnée de descente relativement
à toute famille donnée.

\begin{lemma}
\label{lemma-refine-descent-datum}
Soient $\mathcal{U} = \{U_i \to U\}_{i \in I}$ et
$\mathcal{V} = \{V_j \to V\}_{j \in J}$ des familles de morphismes de
schémas à but fixé. Soit
$(g, \alpha : I \to J, (g_i)) : \mathcal{U} \to \mathcal{V}$ un morphisme
de familles de morphismes à but fixé, voir Sites, définition
\ref{sites-definition-morphism-coverings}. Soit
$(\mathcal{F}_j, \varphi_{jj'})$ une donnée de descente pour les faisceaux
quasi-cohérents relativement à la famille $\{V_j \to V\}_{j \in J}$. Alors :
\begin{enumerate}
\item le système
$$
\left(g_i^*\mathcal{F}_{\alpha(i)},
(g_i \times g_{i'})^*\varphi_{\alpha(i)\alpha(i')}\right)
$$
est une donnée de descente relativement à la famille
$\{U_i \to U\}_{i \in I}$ ;
\item cette construction est
fonctorielle en la donnée de descente $(\mathcal{F}_j, \varphi_{jj'})$ ;
\item étant donné un second morphisme
$(g', \alpha' : I \to J, (g'_i))$ de familles de morphismes à but fixé tel
que $g = g'$, il existe un isomorphisme fonctoriel de données de descente
$$
(g_i^*\mathcal{F}_{\alpha(i)},
(g_i \times g_{i'})^*\varphi_{\alpha(i)\alpha(i')})
\cong
((g'_i)^*\mathcal{F}_{\alpha'(i)},
(g'_i \times g'_{i'})^*\varphi_{\alpha'(i)\alpha'(i')}).
$$
\end{enumerate}
\end{lemma}

\begin{proof}
Omis. Indication : les morphismes
$g_i^*\mathcal{F}_{\alpha(i)} \to (g'_i)^*\mathcal{F}_{\alpha'(i)}$ qui
donnent l'isomorphisme des données de descente en (3) sont les images
inverses des morphismes $\varphi_{\alpha(i)\alpha'(i)}$ par
$(g_i, g'_i) : U_i \to V_{\alpha(i)} \times_V V_{\alpha'(i)}$.
\end{proof}

\noindent
Toute famille $\mathcal{U} = \{S_i \to S\}_{i \in I}$ raffine d'une unique
manière le recouvrement trivial $\{S \to S\}$. Si $\mathcal{F}$ est un
faisceau quasi-cohérent sur $S$, nous notons simplement
$(\mathcal{F}|_{S_i}, can)$ la donnée de descente relativement à
$\mathcal{U}$ obtenue par le procédé ci-dessus.

\begin{definition}
\label{definition-descent-datum-effective-quasi-coherent}
Soit $S$ un schéma. Soit $\{S_i \to S\}_{i \in I}$ une famille de morphismes
de but $S$.
\begin{enumerate}
\item Soit $\mathcal{F}$ un $\mathcal{O}_S$-module quasi-cohérent. Nous
appelons l'unique donnée de descente sur $\mathcal{F}$ relativement au
recouvrement $\{S \to S\}$ la {\it donnée de descente triviale}.
\item L'image inverse de la donnée de descente triviale sur
$\{S_i \to S\}$ est appelée la {\it donnée de descente canonique}. Notation :
$(\mathcal{F}|_{S_i}, can)$.
\item Une donnée de descente $(\mathcal{F}_i, \varphi_{ij})$ pour les
faisceaux quasi-cohérents relativement au recouvrement donné est dite
{\it effective} s'il existe un faisceau quasi-cohérent $\mathcal{F}$ sur $S$
tel que $(\mathcal{F}_i, \varphi_{ij})$ soit isomorphe à
$(\mathcal{F}|_{S_i}, can)$.
\end{enumerate}
\end{definition}

\begin{lemma}
\label{lemma-zariski-descent-effective}
Soit $S$ un schéma et soit $S = \bigcup U_i$ un recouvrement ouvert. Toute
donnée de descente sur des faisceaux quasi-cohérents relativement à la famille
$\mathcal{U} = \{U_i \to S\}$ est effective. De plus, le foncteur de la
catégorie des $\mathcal{O}_S$-modules quasi-cohérents vers la catégorie des
données de descente relativement à $\mathcal{U}$ est pleinement fidèle.
\end{lemma}

\begin{proof}
Cela résulte immédiatement de Faisceaux, section
\ref{sheaves-section-glueing-sheaves}, et du fait que la quasi-cohérence est
une propriété locale, voir Modules, définition
\ref{modules-definition-quasi-coherent}.
\end{proof}

\noindent
Pour aller plus loin, nous devons d'abord étudier le cas des modules sur des
anneaux.










\section{Descente pour les modules}
\label{section-descent-modules}

\noindent
Soit $R \to A$ un morphisme d'anneaux. D'après Simplicial, exemple
\ref{simplicial-example-push-outs-simplicial-object}, il donne naissance à une
$R$-algèbre cosimpliciale
$$
\xymatrix{
A
\ar@<1ex>[r]
\ar@<-1ex>[r]
&
A \otimes_R A
\ar@<0ex>[l]
\ar@<2ex>[r]
\ar@<0ex>[r]
\ar@<-2ex>[r]
&
A \otimes_R A \otimes_R A
\ar@<1ex>[l]
\ar@<-1ex>[l]
}
$$
Notons-la $(A/R)_\bullet$, de sorte que $(A/R)_n$ est le produit tensoriel de
$(n + 1)$ copies de $A$ sur $R$. Étant donné une application
$\varphi : [n] \to [m]$, le morphisme de $R$-algèbres
$(A/R)_\bullet(\varphi)$ est l'application
$$
a_0 \otimes \ldots \otimes a_n
\longmapsto
\prod\nolimits_{\varphi(i) = 0} a_i
\otimes
\prod\nolimits_{\varphi(i) = 1} a_i
\otimes \ldots \otimes
\prod\nolimits_{\varphi(i) = m} a_i
$$
où, par convention, le produit vide vaut $1$. Les premiers morphismes, avec
les notations de Simplicial, section
\ref{simplicial-section-cosimplicial-object}, sont donc
$$
\begin{matrix}
\delta^1_0 & : & a_0 & \mapsto & 1 \otimes a_0 \\
\delta^1_1 & : & a_0 & \mapsto & a_0 \otimes 1 \\
\sigma^0_0 & : & a_0 \otimes a_1 & \mapsto & a_0a_1 \\
\delta^2_0 & : & a_0 \otimes a_1 & \mapsto & 1 \otimes a_0 \otimes a_1 \\
\delta^2_1 & : & a_0 \otimes a_1 & \mapsto & a_0 \otimes 1 \otimes a_1 \\
\delta^2_2 & : & a_0 \otimes a_1 & \mapsto & a_0 \otimes a_1 \otimes 1 \\
\sigma^1_0 & : & a_0 \otimes a_1 \otimes a_2 & \mapsto & a_0a_1 \otimes a_2 \\
\sigma^1_1 & : & a_0 \otimes a_1 \otimes a_2 & \mapsto & a_0 \otimes a_1a_2
\end{matrix}
$$
et ainsi de suite.

\medskip\noindent
Un $R$-module $M$ donne naissance au $(A/R)_\bullet$-module cosimplicial
$(A/R)_\bullet \otimes_R M$. Autrement dit,
$M_n = (A/R)_n \otimes_R M$, et les morphismes de $R$-algèbres
$(A/R)_n \to (A/R)_m$ servent à définir les morphismes correspondants sur
$M \otimes_R (A/R)_\bullet$.

\medskip\noindent
L'analogue, pour les modules, d'une donnée de descente pour les faisceaux
quasi-cohérents est le suivant.

\begin{definition}
\label{definition-descent-datum-modules}
Soit $R \to A$ un morphisme d'anneaux.
\begin{enumerate}
\item Une {\it donnée de descente $(N, \varphi)$ pour les modules
relativement à $R \to A$} est constituée d'un $A$-module $N$ et d'un
isomorphisme de $A \otimes_R A$-modules
$$
\varphi : N \otimes_R A \to A \otimes_R N
$$
tel que la {\it condition de cocycle} soit satisfaite : le diagramme de
morphismes de $A \otimes_R A \otimes_R A$-modules
$$
\xymatrix{
N \otimes_R A \otimes_R A \ar[rr]_{\varphi_{02}}
\ar[rd]_{\varphi_{01}}
& &
A \otimes_R A \otimes_R N \\
& A \otimes_R N \otimes_R A \ar[ru]_{\varphi_{12}} &
}
$$
est commutatif (voir ci-dessous pour les notations).
\item Un {\it morphisme $(N, \varphi) \to (N', \varphi')$ de données de
descente} est un morphisme de $A$-modules $\psi : N \to N'$ tel que le
diagramme
$$
\xymatrix{
N \otimes_R A \ar[r]_\varphi \ar[d]_{\psi \otimes \text{id}_A} &
A \otimes_R N \ar[d]^{\text{id}_A \otimes \psi} \\
N' \otimes_R A \ar[r]^{\varphi'} &
A \otimes_R N'
}
$$
soit commutatif.
\end{enumerate}
\end{definition}

\noindent
Dans cette définition, nous posons
$\varphi_{01} = \varphi \otimes \text{id}_A$,
$\varphi_{12} = \text{id}_A \otimes \varphi$, et
$\varphi_{02}(n \otimes 1 \otimes 1) = \sum a_i \otimes 1 \otimes n_i$
si $\varphi(n \otimes 1) = \sum a_i \otimes n_i$. Ce sont trois
homomorphismes de $A \otimes_R A \otimes_R A$-modules. De manière
équivalente, nous avons
$$
\varphi_{ij}
=
\varphi \otimes_{(A/R)_1, \ (A/R)_\bullet(\tau^2_{ij})} (A/R)_2
$$
où $\tau^2_{ij} : [1] \to [2]$ est l'application
$0 \mapsto i$, $1 \mapsto j$. En effet,
$(A/R)_{\bullet}(\tau^2_{02})(a_0 \otimes a_1) =
a_0 \otimes 1 \otimes a_1$, et de même pour les autres\footnote{Remarquons
que $\tau^2_{ij} = \delta^2_k$ si
$\{i, j, k\} = [2] = \{0, 1, 2\}$, voir Simplicial, définition
\ref{simplicial-definition-face-degeneracy}.}.

\medskip\noindent
Nous avons besoin de quelques notations supplémentaires pour énoncer le lemme
suivant. Soit $(N, \varphi)$ une donnée de descente relativement à un
morphisme d'anneaux $R \to A$. Pour $n \geq 0$ et $i \in [n]$, posons
$$
N_{n, i} =
A \otimes_R
\ldots
\otimes_R A \otimes_R N \otimes_R A \otimes_R
\ldots
\otimes_R A
$$
où le facteur $N$ occupe la $i$-ième place. C'est un $(A/R)_n$-module.
Introduisons les applications $\tau^n_i : [0] \to [n]$, $0 \mapsto i$.
Alors
$$
N_{n, i} = N \otimes_{(A/R)_0, \ (A/R)_\bullet(\tau^n_i)} (A/R)_n
$$
Pour $0 \leq i \leq j \leq n$, soit
$\tau^n_{ij} : [1] \to [n]$ l'application qui envoie $0$ sur $i$ et $1$ sur
$j$. Comme ci-dessus, l'homomorphisme $\varphi$ induit des isomorphismes
$$
\varphi^n_{ij}
=
\varphi \otimes_{(A/R)_1, \ (A/R)_\bullet(\tau^n_{ij})} (A/R)_n :
N_{n, i} \longrightarrow N_{n, j}
$$
de $(A/R)_n$-modules lorsque $i < j$. Si $i = j$, nous posons
$\varphi^n_{ij} = \text{id}$. Tous ces morphismes étant des isomorphismes, ils
permettent de déplacer le facteur $N$ vers n'importe quelle place. La
condition de cocycle signifie précisément que le résultat ne dépend pas de la
manière dont on effectue ce déplacement (par exemple en composant deux de ces
isomorphismes ou en une seule étape). Enfin, pour toute application
$\beta : [n] \to [m]$, nous définissons le morphisme
$$
N_{\beta, i} : N_{n, i} \to N_{m, \beta(i)}
$$
comme l'unique application semi-linéaire relativement à l'homomorphisme
$(A/R)_\bullet(\beta)$ telle que
$$
N_{\beta, i}(1 \otimes \ldots \otimes n \otimes \ldots \otimes 1)
=
1 \otimes \ldots \otimes n \otimes \ldots \otimes 1
$$
pour tout $n \in N$. Cela suggère le lemme suivant.

\begin{lemma}
\label{lemma-descent-datum-cosimplicial}
Soit $R \to A$ un morphisme d'anneaux. À toute donnée de descente
$(N, \varphi)$, on peut associer un $(A/R)_\bullet$-module cosimplicial
$N_\bullet$\footnote{Il faudrait en réalité écrire
$(N, \varphi)_\bullet$.} en posant $N_n = N_{n, n}$ et, pour
$\beta : [n] \to [m]$, en définissant
$$
N_\bullet(\beta) = (\varphi^m_{\beta(n)m}) \circ N_{\beta, n} :
N_{n, n} \longrightarrow N_{m, m}.
$$
Cette construction est fonctorielle en la donnée de descente.
\end{lemma}

\begin{proof}
Voici les premières applications, où
$\varphi(n \otimes 1) = \sum \alpha_i \otimes x_i$ :
$$
\begin{matrix}
\delta^1_0 & : & N & \to & A \otimes N & n & \mapsto & 1 \otimes n \\
\delta^1_1 & : & N & \to & A \otimes N & n & \mapsto &
\sum \alpha_i \otimes x_i\\
\sigma^0_0 & : & A \otimes N & \to & N & a_0 \otimes n & \mapsto & a_0n \\
\delta^2_0 & : & A \otimes N & \to & A \otimes A \otimes N &
a_0 \otimes n & \mapsto & 1 \otimes a_0 \otimes n \\
\delta^2_1 & : & A \otimes N & \to & A \otimes A \otimes N &
a_0 \otimes n & \mapsto & a_0 \otimes 1 \otimes n \\
\delta^2_2 & : & A \otimes N & \to & A \otimes A \otimes N &
a_0 \otimes n & \mapsto & \sum a_0 \otimes \alpha_i \otimes x_i \\
\sigma^1_0 & : & A \otimes A \otimes N & \to & A \otimes N &
a_0 \otimes a_1 \otimes n & \mapsto & a_0a_1 \otimes n \\
\sigma^1_1 & : & A \otimes A \otimes N & \to & A \otimes N &
a_0 \otimes a_1 \otimes n & \mapsto & a_0 \otimes a_1n
\end{matrix}
$$
avec les notations de
Simplicial, section \ref{simplicial-section-cosimplicial-object}.
Vérifions d'abord les deux identités $\sigma^0_0 \circ \delta^1_0 = \text{id}$
et $\sigma^0_0 \circ \delta^1_1 = \text{id}$.
La première, $\sigma^0_0 \circ \delta^1_0 = \text{id}$, résulte immédiatement
de la description explicite des morphismes ci-dessus.
Pour démontrer la seconde, il faut utiliser la condition de cocycle
(car elle n'est pas vérifiée pour un isomorphisme arbitraire
$\varphi : N \otimes_R A \to A \otimes_R N$). Posons
$p = \sigma^0_0 \circ \delta^1_1 : N \to N$. D'après la description
des applications ci-dessus, $p$ s'identifie aussi à
$$
p = \varphi \otimes \text{id} :
N = (N \otimes_R A) \otimes_{(A \otimes_R A)} A
\longrightarrow
(A \otimes_R N) \otimes_{(A \otimes_R A)} A = N
$$
Puisque $\varphi$ est un isomorphisme, $p$ en est un également.
Écrivons $\varphi(n \otimes 1) = \sum \alpha_i \otimes x_i$, avec
$\alpha_i \in A$ et $x_i \in N$. Alors $p(n) = \sum \alpha_ix_i$.
Écrivons ensuite
$\varphi(x_i \otimes 1) = \sum \alpha_{ij} \otimes y_j$, avec
$\alpha_{ij} \in A$ et $y_j \in N$. La condition de cocycle donne
$$
\sum \alpha_i \otimes \alpha_{ij} \otimes y_j
=
\sum \alpha_i \otimes 1 \otimes x_i.
$$
Cela signifie que $p(n) = \sum \alpha_ix_i = \sum \alpha_i\alpha_{ij}y_j =
\sum \alpha_i p(x_i) = p(p(n))$. Ainsi, $p$ est un projecteur ; comme
c'est un isomorphisme, c'est l'identité.

\medskip\noindent
Pour achever de montrer que $N_\bullet$ est un module cosimplicial,
il faut vérifier les cinq types de relations de
Simplicial, remarque \ref{simplicial-remark-relations-cosimplicial}.
Les relations qui ne font intervenir que des composées de $\sigma$ sont
évidentes. Celles qui ne font intervenir que des composées de $\delta$ se
ramènent à la condition de cocycle pour $\varphi$.
On vérifie exactement comme ci-dessus les relations
$\sigma_j \circ \delta_j = \sigma_j \circ \delta_{j + 1} = \text{id}$.
Enfin, les autres relations entre composées de $\delta$ et de $\sigma$
sont vérifiées quel que soit $\varphi$.
\end{proof}

\noindent
À tout $R$-module $M$, on peut associer une donnée de descente canonique,
à savoir $(M \otimes_R A, can)$, où
$can : (M \otimes_R A) \otimes_R A \to A \otimes_R (M \otimes_R A)$
est l'application évidente :
$(m \otimes a) \otimes a' \mapsto a \otimes (m \otimes a')$.

\begin{lemma}
\label{lemma-canonical-descent-datum-cosimplicial}
Soit $R \to A$ un homomorphisme d'anneaux.
Soit $M$ un $R$-module. Le
$(A/R)_\bullet$-module cosimplicial associé à la donnée de descente canonique
est isomorphe au module cosimplicial $(A/R)_\bullet \otimes_R M$.
\end{lemma}

\begin{proof}
Démonstration omise.
\end{proof}

\begin{definition}
\label{definition-descent-datum-effective-module}
Soit $R \to A$ un homomorphisme d'anneaux.
Une donnée de descente $(N, \varphi)$ est dite {\it effective}
s'il existe un $R$-module $M$ et un isomorphisme de données de descente
de $(M \otimes_R A, can)$ vers
$(N, \varphi)$.
\end{definition}

\noindent
Soit $R \to A$ un homomorphisme d'anneaux.
Soit $(N, \varphi)$ une donnée de descente.
Considérons le complexe de cochaînes $s(N_\bullet)$ associé
à $N_\bullet$ (voir
Simplicial, section \ref{simplicial-section-dold-kan-cosimplicial}).
Il est de la forme
$$
N \to A \otimes_R N \to A \otimes_R A \otimes_R N \to \ldots
$$
On peut en décrire les différentielles.
La première est l'application
$$
n \longmapsto 1 \otimes n - \varphi(n \otimes 1).
$$
La seconde est donnée, sur les tenseurs purs, par
$$
a \otimes n \longmapsto 1 \otimes a \otimes n
- a \otimes 1 \otimes n + a \otimes \varphi(n \otimes 1).
$$
La règle se poursuit de manière évidente.

\medskip\noindent
Dans le cas particulier où
$N = A \otimes_R M$, pour tout $m \in M$,
l'élément $1 \otimes m$ appartient au noyau de la première différentielle
du complexe de cochaînes associé au module cosimplicial
$(A/R)_\bullet \otimes_R M$. On obtient donc un complexe de cochaînes augmenté
\begin{equation}
\label{equation-extended-complex}
0 \to M \to A \otimes_R M \to A \otimes_R A \otimes_R M \to \ldots
\end{equation}
On place ici $0$ en degré $-2$,
le module $M$ en degré $-1$, le module $A \otimes_R M$
en degré $0$, et ainsi de suite. Ce complexe est de la forme
$$
0 \to R \to A \to A \otimes_R A \to A \otimes_R A \otimes_R A \to \ldots
$$
lorsque $M = R$.

\begin{lemma}
\label{lemma-with-section-exact}
Supposons que $R \to A$ admette une section.
Alors, pour tout $R$-module $M$, le complexe de cochaînes augmenté
(\ref{equation-extended-complex}) est exact.
\end{lemma}

\begin{proof}
D'après
Simplicial, lemme \ref{simplicial-lemma-push-outs-simplicial-object-w-section},
l'application $R \to (A/R)_\bullet$ est un homotopisme
de $R$-algèbres cosimpliciales
(où $R$ désigne la $R$-algèbre cosimpliciale constante).
Par suite, $M \to (A/R)_\bullet \otimes_R M$ est
un homotopisme dans la catégorie des
$R$-modules cosimpliciaux, car $\otimes_R M$ est
un foncteur de la catégorie des $R$-algèbres vers celle
des $R$-modules ; voir
Simplicial, lemme \ref{simplicial-lemma-functorial-homotopy}.
L'application induite entre les complexes associés est donc
un homotopisme ; voir
Simplicial, lemme \ref{simplicial-lemma-homotopy-s-Q}.
Puisque le complexe associé au $R$-module cosimplicial constant
$M$ est le complexe
$$
\xymatrix{
M \ar[r]^0 & M \ar[r]^1 & M \ar[r]^0 & M \ar[r]^1 & M \ldots
}
$$
le résultat s'ensuit (le complexe augmenté s'obtient simplement en ajoutant
un terme $M$ au début).
\end{proof}

\begin{lemma}
\label{lemma-ff-exact}
Supposons que $R \to A$ soit fidèlement plat ; voir
Algèbre, définition \ref{algebra-definition-flat}.
Alors, pour tout $R$-module $M$, le complexe de cochaînes augmenté
(\ref{equation-extended-complex}) est exact.
\end{lemma}

\begin{proof}
Supposons que l'on puisse trouver un homomorphisme d'anneaux fidèlement plat
$R \to R'$ tel que le résultat soit vrai pour l'homomorphisme
$R' \to A' = R' \otimes_R A$. Le résultat vaut alors aussi pour $R \to A$.
En effet, pour tout $R$-module $M$, le module cosimplicial
$(M \otimes_R R') \otimes_{R'} (A'/R')_\bullet$
s'identifie au module cosimplicial $R' \otimes_R (M \otimes_R (A/R)_\bullet)$.
L'annulation de la cohomologie du complexe associé à
$(M \otimes_R R') \otimes_{R'} (A'/R')_\bullet$ entraîne donc celle
de la cohomologie du complexe associé à
$M \otimes_R (A/R)_\bullet$, par fidèle platitude de $R \to R'$.
Il en va de même pour l'annulation des groupes de cohomologie en degrés
$-1$ et $0$ du complexe augmenté (démonstration omise).

\medskip\noindent
Or une telle extension fidèlement plate existe : il suffit de prendre $R' = A$,
car l'homomorphisme d'anneaux $R' = A \to A' = A \otimes_R A$ admet la section
$a \otimes a' \mapsto aa'$, et le
lemme \ref{lemma-with-section-exact} s'applique.
\end{proof}

\noindent
Voici le lien entre ce complexe et la question de l'effectivité.

\begin{lemma}
\label{lemma-recognize-effective}
Soit $R \to A$ un homomorphisme d'anneaux fidèlement plat.
Soit $(N, \varphi)$ une donnée de descente.
Alors $(N, \varphi)$ est effective si et seulement si l'application canonique
$$
A \otimes_R H^0(s(N_\bullet)) \longrightarrow N
$$
est un isomorphisme.
\end{lemma}

\begin{proof}
Si $(N, \varphi)$ est effective, on peut écrire $N = A \otimes_R M$,
avec $\varphi = can$. On a alors $H^0(s(N_\bullet)) = M$, d'après
les lemmes \ref{lemma-canonical-descent-datum-cosimplicial}
et \ref{lemma-ff-exact}. Réciproquement, supposons que l'application du lemme
soit un isomorphisme. Posons $M = H^0(s(N_\bullet))$.
C'est un $R$-sous-module de $N$,
à savoir $M = \{n \in N \mid 1 \otimes n = \varphi(n \otimes 1)\}$.
Il reste seulement à vérifier que, par l'isomorphisme
$A \otimes_R M \to N$,
la donnée de descente canonique correspond à $\varphi$.
Nous omettons cette vérification.
\end{proof}

\begin{lemma}
\label{lemma-descent-descends}
Soit $R \to A$ un homomorphisme d'anneaux fidèlement plat, et soit $R \to R'$
un homomorphisme fidèlement plat. Posons $A' = R' \otimes_R A$.
Si toutes les données de descente relatives à $R' \to A'$ sont effectives,
il en est de même de toutes les données de descente relatives à $R \to A$.
\end{lemma}

\begin{proof}
Soit $(N, \varphi)$ une donnée de descente relative à $R \to A$.
Posons $N' = R' \otimes_R N = A' \otimes_A N$, et notons
$\varphi' = \text{id}_{R'} \otimes \varphi$ le changement de base
de la donnée de descente $\varphi$. Alors $(N', \varphi')$ est
une donnée de descente relative à $R' \to A'$ et
$H^0(s(N'_\bullet)) = R' \otimes_R H^0(s(N_\bullet))$.
De plus, l'application
$A' \otimes_{R'} H^0(s(N'_\bullet)) \to N'$
s'identifie au changement de base de l'homomorphisme de $A$-modules
$A \otimes_R H^0(s(N)) \to N$ par l'homomorphisme fidèlement plat
$A \to A'$. On conclut donc par le lemme \ref{lemma-recognize-effective}.
\end{proof}

\noindent
Voici le résultat principal de cette section.
Sa démonstration peut paraître un peu laborieuse ; pour une approche
plus conceptuelle, voir la remarque
\ref{remark-homotopy-equivalent-cosimplicial-algebras} ci-dessous.

\begin{proposition}
\label{proposition-descent-module}
\begin{slogan}
Descente effective des modules par les homomorphismes d'anneaux fidèlement plats.
\end{slogan}
Soit $R \to A$ un homomorphisme d'anneaux fidèlement plat.
Alors
\begin{enumerate}
\item toute donnée de descente de modules relative à $R \to A$
est effective ;
\item le foncteur $M \mapsto (A \otimes_R M, can)$ de la catégorie des $R$-modules
vers celle des données de descente est une équivalence ;
\item le foncteur inverse est donné par $(N, \varphi) \mapsto H^0(s(N_\bullet))$.
\end{enumerate}
\end{proposition}

\begin{proof}
Nous démontrons seulement (1) et omettons les démonstrations de (2) et (3).
Puisque $R \to A$ est fidèlement plat, il existe un changement de base
fidèlement plat $R \to R'$ tel que $R' \to A' = R' \otimes_R A$
admette une section (il suffit de prendre $R' = A$, comme dans la preuve du
lemme \ref{lemma-ff-exact}). Ainsi, par le
lemme \ref{lemma-descent-descends}, on peut supposer que
$R \to A$ admet une section, notée $\sigma : A \to R$.
Soit $(N, \varphi)$ une donnée de descente relative à $R \to A$.
Posons
$$
M = H^0(s(N_\bullet)) = \{n \in N \mid 1 \otimes n = \varphi(n \otimes 1)\}
\subset
N
$$
D'après le lemme \ref{lemma-recognize-effective}, il suffit de montrer que
$A \otimes_R M \to N$ est un isomorphisme.

\medskip\noindent
Soit $n \in N$. Écrivons
$\varphi(n \otimes 1) = \sum a_i \otimes x_i$, avec
$a_i \in A$ et $x_i \in N$. D'après le lemme \ref{lemma-descent-datum-cosimplicial},
on a $n = \sum a_i x_i$ dans $N$ (puisque
$\sigma^0_0 \circ \delta^1_1 = \text{id}$ dans tout objet cosimplicial).
Écrivons ensuite $\varphi(x_i \otimes 1) = \sum a_{ij} \otimes y_j$, avec
$a_{ij} \in A$ et $y_j \in N$.
La condition de cocycle s'écrit
$$
\sum a_i \otimes a_{ij} \otimes y_j = \sum a_i \otimes 1 \otimes x_i
$$
dans $A \otimes_R A \otimes_R N$. On en déduit deux identités :
\begin{enumerate}
\item en appliquant $\sigma$ au premier facteur $A$, on obtient
$\sum \sigma(a_i) \varphi(x_i \otimes 1) = \sum \sigma(a_i) \otimes x_i$ ;
\item en appliquant $\sigma$ au facteur $A$ du milieu, on obtient
$\sum_i a_i \otimes \sum_j \sigma(a_{ij}) y_j = \sum a_i \otimes x_i$.
\end{enumerate}
Le point (1) montre que $\sum \sigma(a_i) x_i \in M$. En l'appliquant à
$x_i$, on voit que $\sum \sigma(a_{ij})y_i \in M$ pour tout $i$.
En effectuant les multiplications dans l'égalité (2), on obtient
$\sum_i a_i (\sum_j \sigma(a_{ij}) y_j) = \sum a_i x_i = n$.
Ainsi, $A \otimes_R M \to N$ est surjectif.
Enfin, supposons que $m_i \in M$ et $\sum a_i m_i = 0$.
En appliquant $\varphi$ à
$\sum a_im_i \otimes 1$, on obtient $\sum a_i \otimes m_i = 0$.
Autrement dit, $A \otimes_R M \to N$ est injectif, ce qui achève la preuve.
\end{proof}

\begin{remark}
\label{remark-standard-covering}
Soit $R$ un anneau. Soient $f_1, \ldots, f_n\in R$ des éléments engendrant
l'idéal unité. L'anneau $A = \prod_i R_{f_i}$ est une
$R$-algèbre fidèlement plate. L'anneau cosimplicial $(A/R)_\bullet$
a pour composante de degré $n$ l'anneau
$$
\prod\nolimits_{i_0, \ldots, i_n} R_{f_{i_0}\ldots f_{i_n}}
$$
Les résultats ci-dessus redonnent donc
Algèbre, lemmes \ref{algebra-lemma-standard-covering},
\ref{algebra-lemma-cover-module} et \ref{algebra-lemma-glue-modules}.
Ils vont toutefois plus loin, grâce à l'exactitude en degrés supérieurs.
Celle-ci entraîne en effet la trivialité de la cohomologie de {\v C}ech
des faisceaux quasi-cohérents sur les schémas affines.
On obtient ainsi une seconde démonstration de
Cohomologie des schémas, lemme
\ref{coherent-lemma-cech-cohomology-quasi-coherent-trivial}.
\end{remark}

\begin{remark}
\label{remark-homotopy-equivalent-cosimplicial-algebras}
Soit $R$ un anneau. Soit $A_\bullet$ une $R$-algèbre cosimpliciale.
Dans ce cadre, une donnée de descente correspond à un
$A_\bullet$-module cosimplicial $M_\bullet$ tel que, pour tous
$n, m \geq 0$ et tout $\varphi : [n] \to [m]$,
l'application $M(\varphi) : M_n \to M_m$ induise un isomorphisme
$$
M_n \otimes_{A_n, A(\varphi)} A_m \longrightarrow M_m.
$$
Appelons un tel module cosimplicial un {\it module cartésien}.
Dans ce cadre, la démonstration de la proposition \ref{proposition-descent-module}
se décompose selon les étapes suivantes :
\begin{enumerate}
\item Si $R \to R'$ et $R \to A$ sont fidèlement plats,
les données de descente relatives à $A/R$ sont effectives dès que
celles relatives à $(R' \otimes_R A)/R'$ le sont.
\item Soit $A$ une $R$-algèbre. Les données de descente relatives à $A/R$
correspondent aux $(A/R)_\bullet$-modules cartésiens.
\item Si $R \to A$ admet une section, $(A/R)_\bullet$ est homotopiquement
équivalent à $R$, la $R$-algèbre cosimpliciale constante de valeur $R$.
\item Si $A_\bullet \to B_\bullet$ est une équivalence d'homotopie
de $R$-algèbres cosimpliciales, le foncteur
$M_\bullet \mapsto M_\bullet \otimes_{A_\bullet} B_\bullet$
induit une équivalence entre la catégorie des
$A_\bullet$-modules cartésiens et celle des $B_\bullet$-modules cartésiens.
\end{enumerate}
Pour (1), voir le lemme \ref{lemma-descent-descends}.
Le point (2) utilise le lemme \ref{lemma-descent-datum-cosimplicial}.
Le point (3) a été établi dans la preuve du lemme \ref{lemma-with-section-exact}
(il repose sur Simplicial,
lemme \ref{simplicial-lemma-push-outs-simplicial-object-w-section}).
Enfin, le point (4) est immédiat si l'on adopte le bon point de vue !
\end{remark}








\section{Descente pour les morphismes universellement injectifs}
\label{section-descent-universally-injective}

\noindent
De nombreuses constructions en géométrie algébrique font appel aux techniques
de {\it descente} : on construit des objets sur un espace donné en travaillant
d'abord sur un espace un peu plus grand qui se projette sur celui-ci,
ou l'on vérifie une propriété d'un espace ou d'un morphisme après
changement de base par un morphisme de recouvrement.
L'utilité de ces techniques dépend bien entendu de l'identification
d'une classe suffisamment large de {\it morphismes de descente effective}.
Dès les débuts de la géométrie algébrique moderne développée par Grothendieck,
on a montré que les morphismes {\it quasi-compacts} et {\it fidèlement plats}
permettent de faire descendre effectivement les objets, les morphismes
et nombre de leurs propriétés.

\medskip\noindent
Comme d'habitude, cet énoncé se ramène à une propriété des anneaux et des modules.
Pour qu'un homomorphisme $f: R \to S$ soit un morphisme de descente effective
pour les modules, Grothendieck a montré qu'il suffit que $f$ soit fidèlement plat.
Cette condition exclut toutefois de nombreux exemples naturels :
ainsi, tout homomorphisme d'anneaux scindé est un morphisme de descente effective.
Un exemple naturel intervient même dans la démonstration de la descente
fidèlement plate : pour tout homomorphisme d'anneaux $f: R \to S$,
l'homomorphisme $1_S \otimes f: S \to S \otimes_R S$
admet la multiplication pour rétraction, qu'il soit plat ou non.

\medskip\noindent
On peut donc chercher une condition naturelle sur un homomorphisme d'anneaux
qui entraîne la descente effective des modules et englobe à la fois
les homomorphismes fidèlement plats et les homomorphismes scindés.
Le lecteur sera peut-être surpris (l'auteur l'a du moins été)
d'apprendre qu'une réponse complète est connue depuis les environs de 1970 !
Il est facile de vérifier que, pour que $f: R \to S$ soit un morphisme
de descente effective pour les modules, il faut que $f$ soit
{\it universellement injectif} dans la catégorie des $R$-modules :
pour tout $R$-module $M$, l'application
$1_M  \otimes f: M \to M \otimes_R S$ doit être injective.
Or cette condition est aussi suffisante.
Par exemple, si $f$ est scindé dans la catégorie des $R$-modules
(sans l'être nécessairement dans celle des anneaux),
alors $f$ est un morphisme de descente effective pour les modules.

\medskip\noindent
L'histoire de ce résultat est assez complexe : il a d'abord été énoncé
par Olivier \cite{olivier}, qui appelait {\it purs} les morphismes
universellement injectifs, mais sans indication claire de démonstration.
On peut le déduire des travaux de Joyal et Tierney \cite{joyal-tierney} ;
cependant, à notre connaissance, la première démonstration autonome publiée
est celle de Mesablishvili \cite{mesablishvili1}.
Le premier objectif de cette section est d'exposer cette démonstration.
Hormis la correction de coquilles, nous modifions peu la présentation originale,
si ce n'est pour expliciter le lien entre la définition usuelle
d'une donnée de descente en géométrie algébrique et celle
employée dans \cite{mesablishvili1}.
La démonstration est entièrement catégorique et peut donc être formulée
dans le langage des monades, ce qui permet de l'appliquer
dans d'autres contextes ; voir \cite{janelidze-tholen}.

\medskip\noindent
Le second objectif est de rassembler des résultats sur les propriétés
des modules, des algèbres et des morphismes que l'on peut faire descendre le long
d'homomorphismes d'anneaux universellement injectifs.
Les cas des modules de type fini et des modules plats ont été traités
par Mesablishvili \cite{mesablishvili2}.


\subsection{Préliminaires de théorie des catégories}
\label{subsection-category-prelims}

\noindent
Commençons par rappeler quelques notions élémentaires de théorie des catégories
qui simplifieront l'exposé. Dans cette sous-section, on fixe une catégorie ambiante.

\medskip\noindent
Étant donnés deux morphismes $g_1, g_2: B \to C$, rappelons qu'un
{\it égalisateur} de $g_1$ et $g_2$ est un morphisme $f: A \to B$ tel que
$g_1 \circ f = g_2 \circ f$ et universel pour cette propriété.
Cela signifie que tout diagramme commutatif
$$
\xymatrix{A' \ar[rd]^e \ar@/^1.5pc/[rrd] \ar@{-->}[d] & & \\
A \ar[r]^f & B \ar@<1ex>[r]^{g_1} \ar@<-1ex>[r]_{g_2} & 
C
}
$$
dont on a omis la flèche en pointillé peut être complété de façon unique.
On dit également que le diagramme
\begin{equation}
\label{equation-equalizer}
\xymatrix{
A \ar[r]^f  & B \ar@<1ex>[r]^{g_1} \ar@<-1ex>[r]_{g_2} & C
}
\end{equation}
est un égalisateur. En inversant les flèches, on obtient la définition
d'un {\it coégalisateur}.
Voir Catégories, sections \ref{categories-section-equalizers} et
\ref{categories-section-coequalizers}.

\medskip\noindent
La propriété d'être un égalisateur étant définie par une propriété universelle,
elle n'est généralement pas préservée par l'application d'un foncteur covariant.
Comme pour les monomorphismes et les épimorphismes, on peut, dans certains cas,
remédier à cette difficulté en exhibant des scindages.

\begin{definition}
\label{definition-split-equalizer}
Un {\it égalisateur scindé} est un diagramme (\ref{equation-equalizer})
tel que $g_1 \circ f = g_2 \circ f$, pour lequel il existe des morphismes auxiliaires
$h : B \to A$ et $i : C \to B$ vérifiant
\begin{equation}
\label{equation-split-equalizer-conditions}
h \circ f = 1_A, \quad f \circ h = i \circ g_1, \quad i \circ g_2 = 1_B.
\end{equation}
\end{definition}

\noindent
Les identités entre les flèches imposent alors que (\ref{equation-equalizer})
soit un égalisateur : l'application $e$ se factorise de façon unique par $f$
en écrivant $e = f \circ (h \circ e)$.
Ainsi, l'image d'un égalisateur scindé par un foncteur covariant
est un égalisateur scindé ; son image par un foncteur contravariant
est un {\it coégalisateur scindé}, dont la définition est évidente.

\subsection{Morphismes universellement injectifs}
\label{subsection-universally-injective}

\noindent
Rappelons que $\textit{Rings}$ désigne la catégorie des anneaux commutatifs
avec $1$. Pour un objet $R$ de $\textit{Rings}$, on note $\text{Mod}_R$
la catégorie des $R$-modules.

\begin{remark}
\label{remark-reflects}
Tout foncteur $F : \mathcal{A} \to \mathcal{B}$ entre catégories abéliennes
qui est exact et envoie les objets non nuls sur des objets non nuls
reflète les injections et les surjections. En effet, l'exactitude implique
que $F$ préserve les noyaux et les conoyaux (comparer avec
Homologie, section \ref{homology-section-functors}).
Par exemple, si $f : R \to S$ est un homomorphisme d'anneaux fidèlement plat,
alors $\bullet \otimes_R S: \text{Mod}_R \to \text{Mod}_S$ possède ces propriétés.
\end{remark}

\noindent
Soit $R$ un anneau. Rappelons qu'un morphisme $f : M \to N$ de $\text{Mod}_R$
est {\it universellement injectif} si, pour tout $P \in \text{Mod}_R$,
le morphisme $f \otimes 1_P: M \otimes_R P \to N \otimes_R P$ est injectif.
Voir Algèbre, définition \ref{algebra-definition-universally-injective}.

\begin{definition}
\label{definition-universally-injective}
Un homomorphisme d'anneaux $f: R \to S$ est {\it universellement injectif}
s'il est universellement injectif comme morphisme de $\text{Mod}_R$.
\end{definition}

\begin{example}
\label{example-split-injection-universally-injective}
Toute injection scindée dans $\text{Mod}_R$ est universellement injective.
En particulier, toute injection scindée dans $\textit{Rings}$
est universellement injective.
\end{example}

\begin{example}
\label{example-cover-universally-injective}
Pour un anneau $R$ et des éléments $f_1, \ldots, f_n \in R$ engendrant l'idéal
unité, le morphisme $R \to R_{f_1} \oplus \ldots \oplus R_{f_n}$
est universellement injectif. Bien que cela résulte immédiatement du
lemme \ref{lemma-faithfully-flat-universally-injective},
il est instructif de le vérifier directement : on se ramène aussitôt
au cas où $R$ est local. L'un des $f_i$ est alors inversible,
et l'application $R \to R_{f_i}$ est donc un isomorphisme.
\end{example}

\begin{lemma}
\label{lemma-faithfully-flat-universally-injective}
Tout homomorphisme d'anneaux fidèlement plat est universellement injectif.
\end{lemma}

\begin{proof}
C'est une reformulation de
Algèbre, lemme \ref{algebra-lemma-faithfully-flat-universally-injective}.
\end{proof}

\noindent
L'observation essentielle de \cite{mesablishvili1} est que l'injectivité
universelle se reformule utilement en termes de scindage, grâce à la
construction usuelle d'un cogénérateur injectif dans $\text{Mod}_R$.

\begin{definition}
\label{definition-C}
Soit $R$ un anneau. On définit le foncteur contravariant
{\it $C$} $ : \text{Mod}_R \to \text{Mod}_R$ par
$$
C(M) = \Hom_{\textit{Ab}}(M, \mathbf{Q}/\mathbf{Z}),
$$
l'action de $R$ sur $C(M)$ étant donnée par $rf(s) = f(rs)$.
\end{definition}

\noindent
Ce foncteur était noté $M \mapsto M^\vee$ dans
Compléments d'algèbre, section \ref{more-algebra-section-injectives-modules}.

\begin{lemma}
\label{lemma-C-is-faithful}
Pour tout anneau $R$, le foncteur $C : \text{Mod}_R \to \text{Mod}_R$
est exact et permet de reconnaître les injections et les surjections.
\end{lemma}

\begin{proof}
L'exactitude est établie dans
Compléments d'algèbre, lemme \ref{more-algebra-lemma-vee-exact}.
Les autres propriétés s'en déduisent ; voir la remarque \ref{remark-reflects}.
\end{proof}

\begin{remark}
\label{remark-adjunction}
Nous utiliserons fréquemment l'adjonction usuelle entre $\Hom$ et
le produit tensoriel, sous la forme de l'isomorphisme naturel de foncteurs
contravariants
\begin{equation}
\label{equation-adjunction}
C(\bullet_1 \otimes_R \bullet_2) \cong \Hom_R(\bullet_1, C(\bullet_2)): 
\text{Mod}_R \times \text{Mod}_R \to \text{Mod}_R
\end{equation}
qui associe à $f: M_1 \otimes_R M_2 \to \mathbf{Q}/\mathbf{Z}$
l'application $m_1 \mapsto 
(m_2 \mapsto f(m_1 \otimes m_2))$. Voir
Algèbre, lemme \ref{algebra-lemma-hom-from-tensor-product-variant}.
Il en résulte que, si
$$
\xymatrix@C=9pc{
C(M) \ar@<1ex>[r] \ar@<-1ex>[r] & C(N) \ar[r] & C(P)
}
$$
est un diagramme de coégalisateur scindé dans $\text{Mod}_R$,
il en est de même de
$$
\xymatrix@C=9pc{
C(M \otimes_R Q) \ar@<1ex>[r] \ar@<-1ex>[r] & C(N \otimes_R Q) \ar[r] & C(P 
\otimes_R Q)
}
$$
pour tout $Q \in \text{Mod}_R$.
\end{remark}

\begin{lemma}
\label{lemma-split-surjection}
Soit $R$ un anneau. Un morphisme $f: M \to N$ de $\text{Mod}_R$
est universellement injectif si et seulement si
$C(f): C(N) \to C(M)$ est une surjection scindée.
\end{lemma}

\begin{proof}
D'après (\ref{equation-adjunction}), pour tout $P \in \text{Mod}_R$,
on a un diagramme commutatif
$$
\xymatrix@C=9pc{
\Hom_R( P, C(N)) \ar[r]_{\Hom_R(P,C(f))} \ar[d]^{\cong} &
\Hom_R(P,C(M)) \ar[d]^{\cong} \\
C(P \otimes_R N ) \ar[r]^{C(1_{P} \otimes f)} & C(P \otimes_R M ).
}
$$
Si $f$ est universellement injectif, alors $1_{C(M)} \otimes f: C(M) \otimes_R M \to 
C(M) \otimes_R N$ est injectif ; les deux flèches horizontales du diagramme
ci-dessus sont donc surjectives pour $P = C(M)$.
On peut ainsi relever $1_{C(M)} \in \Hom_R(C(M), C(M))$
en un élément $g \in \Hom_R(C(N), C(M))$ qui scinde $C(f)$.
Réciproquement, si $C(f)$ est une surjection scindée,
les deux flèches horizontales du diagramme ci-dessus sont surjectives.
Le lemme \ref{lemma-C-is-faithful} entraîne alors que $1_{P} \otimes f$ est injectif.
\end{proof}

\begin{remark}
\label{remark-functorial-splitting}
Soit $f: M \to N$ un morphisme universellement injectif de $\text{Mod}_R$.
En choisissant un scindage $g$ de $C(f)$, on peut construire
un scindage fonctoriel de $C(1_P \otimes f)$ pour tout $P \in \text{Mod}_R$.
En effet, d'après (\ref{equation-adjunction}), il suffit de scinder
$\Hom_R(P, 
C(f))$ fonctoriellement en $P$,
ce que réalise l'application $g \circ \bullet$.
\end{remark}


\subsection{Descente pour les modules et leurs morphismes}
\label{subsection-descent-modules-morphisms}

\noindent
Dans toute cette sous-section, on fixe un homomorphisme d'anneaux $f: R \to S$.
Comme on l'a vu dans la section \ref{section-descent-modules},
les données de descente de modules peuvent se formuler dans le langage
des algèbres cosimpliciales. Nous préférons ici une présentation plus élémentaire.

\medskip\noindent
Pour $i = 1, 2, 3$, notons $S_i$ le produit tensoriel de $i$ copies
de $S$ sur $R$.
Définissons les homomorphismes d'anneaux $\delta_0^1, \delta_1^1: S_1 \to S_2$,
$\delta_{01}^1, \delta_{02}^1, \delta_{12}^1: S_1 \to S_3$ et
$\delta_0^2, \delta_1^2, \delta_2^2: S_2 \to S_3$ par les formules
\begin{align*}
\delta^1_0  (a_0) & =  1 \otimes a_0 \\
\delta^1_1  (a_0) & = a_0 \otimes 1 \\
\delta^2_0  (a_0 \otimes a_1) & =  1 \otimes a_0 \otimes a_1 \\
\delta^2_1  (a_0 \otimes a_1) & =  a_0 \otimes 1 \otimes a_1 \\
\delta^2_2  (a_0 \otimes a_1) & =  a_0 \otimes a_1 \otimes 1 \\
\delta_{01}^1(a_0) & = 1 \otimes 1 \otimes a_0 \\
\delta_{02}^1(a_0) & = 1 \otimes a_0 \otimes 1 \\
\delta_{12}^1(a_0) & = a_0 \otimes 1 \otimes 1.
\end{align*}
Autrement dit, l'indice supérieur indique l'anneau de départ,
et l'indice inférieur les positions où insérer des facteurs égaux à 1.
(Cette notation est compatible avec celle introduite dans la
section \ref{section-descent-modules}.)

\medskip\noindent
Rappelons\footnote{Plus précisément, notre $\theta$ est ici l'inverse de
$\varphi$ dans la définition \ref{definition-descent-datum-modules}.}
que, d'après la définition \ref{definition-descent-datum-modules}, pour
$M \in \text{Mod}_S$, une {\it donnée de descente} sur $M$ relative à $f$
est un isomorphisme
$$
\theta :
M \otimes_{S,\delta^1_0} S_2
\longrightarrow
M \otimes_{S,\delta^1_1} S_2
$$
de $S_2$-modules satisfaisant la {\it condition de cocycle}
\begin{equation}
\label{equation-cocycle-condition}
(\theta \otimes \delta_2^2) \circ (\theta \otimes \delta_2^0) = (\theta \otimes 
\delta_2^1):
M \otimes_{S, \delta^1_{01}} S_3 \to M \otimes_{S,\delta^1_{12}} S_3.
\end{equation}
Notons $DD_{S/R}$ la catégorie des $S$-modules munis de données de descente
relatives à $f$.

\medskip\noindent
Par exemple, pour $M_0 \in \text{Mod}_R$, le choix d'un isomorphisme
$M \cong M_0 \otimes_R S$ fournit une donnée de descente en identifiant
naturellement $M \otimes_{S,\delta^1_0} S_2$ et
$M \otimes_{S,\delta^1_1} S_2$ à $M_0 \otimes_R S_2$.
Cette construction définit en particulier un foncteur
$f^*: \text{Mod}_R \to DD_{S/R}$.

\begin{definition}
\label{definition-effective-descent}
Le foncteur $f^*: \text{Mod}_R \to DD_{S/R}$
est appelé {\it extension des scalaires par $f$}.
On dit que $f$ est un {\it morphisme de descente pour les modules}
si $f^*$ est pleinement fidèle.
On dit que $f$ est un {\it morphisme de descente effective pour les modules}
si $f^*$ est une équivalence de catégories.
\end{definition}

\noindent
Notre but est de montrer que, lorsque $f$ est universellement injectif,
on peut utiliser $\theta$ pour retrouver $M_0$ à l'intérieur de $M$.
Cette construction repose de manière essentielle sur certains
diagrammes égalisateurs.

\begin{lemma}
\label{lemma-equalizer-M}
Pour $(M,\theta) \in DD_{S/R}$, le diagramme
\begin{equation}
\label{equation-equalizer-M}
\xymatrix@C=8pc{
M \ar[r]^{\theta \circ (1_M \otimes \delta_0^1)} &
M \otimes_{S, \delta_1^1} S_2
\ar@<1ex>[r]^{(\theta \otimes \delta_2^2) \circ (1_M \otimes \delta^2_0)}
\ar@<-1ex>[r]_{1_{M \otimes S_2} \otimes \delta^2_1} & 
M \otimes_{S, \delta_{12}^1} S_3
}
\end{equation}
est un égalisateur scindé.
\end{lemma}

\begin{proof}
Définissons les homomorphismes d'anneaux $\sigma^0_0: S_2 \to S_1$ et
$\sigma_0^1, 
\sigma_1^1: S_3 \to S_2$ par les formules
\begin{align*}
\sigma^0_0 (a_0 \otimes a_1) & = a_0a_1 \\
\sigma^1_0 (a_0 \otimes a_1 \otimes a_2) & = a_0a_1 \otimes a_2 \\
\sigma^1_1 (a_0 \otimes a_1 \otimes a_2) & = a_0 \otimes a_1a_2.
\end{align*}
On prend pour morphismes auxiliaires
$1_M \otimes \sigma_0^0: M \otimes_{S, \delta_1^1} S_2 \to M$
et $1_M \otimes \sigma_0^1: M \otimes_{S,\delta_{12}^1} S_3 \to M \otimes_{S, 
\delta_1^1} S_2$.
Parmi les compatibilités requises dans (\ref{equation-split-equalizer-conditions}),
la première s'obtient en tensorisant la condition de cocycle
(\ref{equation-cocycle-condition}) par $\sigma_1^1$ ;
les autres sont immédiates.
\end{proof}

\begin{lemma}
\label{lemma-equalizer-CM}
Pour $(M, \theta) \in DD_{S/R}$, le diagramme
\begin{equation}
\label{equation-coequalizer-CM}
\xymatrix@C=8pc{
C(M \otimes_{S, \delta_{12}^1} S_3)
\ar@<1ex>[r]^{C((\theta \otimes \delta_2^2) \circ (1_M \otimes \delta^2_0))}
\ar@<-1ex>[r]_{C(1_{M \otimes S_2} \otimes \delta^2_1)} &
C(M \otimes_{S, \delta_1^1} S_2 )
\ar[r]^{C(\theta \circ (1_M \otimes \delta_0^1))} & C(M).
}
\end{equation}
obtenu en appliquant $C$ à (\ref{equation-equalizer-M})
est un coégalisateur scindé.
\end{lemma}

\begin{proof}
Démonstration omise.
\end{proof}

\begin{lemma}
\label{lemma-equalizer-S}
Le diagramme
\begin{equation}
\label{equation-equalizer-S}
\xymatrix@C=8pc{
S_1 \ar[r]^{\delta^1_1} &
S_2 \ar@<1ex>[r]^{\delta^2_2} \ar@<-1ex>[r]_{\delta^2_1} & 
S_3
}
\end{equation}
est un égalisateur scindé.
\end{lemma}

\begin{proof}
Dans le lemme \ref{lemma-equalizer-M}, prendre $(M, \theta) = f^*(S)$.
\end{proof}

\noindent
Cela suggère la définition d'un candidat au rôle de quasi-inverse de $f^*$.

\begin{definition}
\label{definition-pushforward}
On définit le foncteur {\it $f_*$} $: DD_{S/R} \to \text{Mod}_R$
en prenant pour $f_*(M, \theta)$ le $R$-sous-module de $M$
tel que le diagramme
\begin{equation}
\label{equation-equalizer-f}
\xymatrix@C=8pc{f_*(M,\theta) \ar[r] & M \ar@<1ex>^{\theta \circ (1_M \otimes 
\delta_0^1)}[r] \ar@<-1ex>_{1_M \otimes \delta_1^1}[r] & 
M \otimes_{S, \delta_1^1} S_2 
}
\end{equation}
soit un égalisateur.
\end{definition}

\noindent
Le lemme \ref{lemma-equalizer-M} et le fait que le foncteur de restriction
$\text{Mod}_S \to \text{Mod}_R$ est adjoint à droite au foncteur
d'extension des scalaires $\bullet \otimes_R S: \text{Mod}_R \to \text{Mod}_S$
montrent que $f_*$ est adjoint à droite à $f^*$.

\medskip\noindent
Nous pouvons maintenant démontrer le lemme clé. Dans le cas fidèlement plat,
il est immédiat (voir la remarque \ref{remark-descent-lemma}),
mais le cas général demande une démonstration.

\begin{lemma}
\label{lemma-descent-lemma}
Si $f$ est universellement injectif, le diagramme
\begin{equation}
\label{equation-equalizer-f2}
\xymatrix@C=8pc{
f_*(M, \theta) \otimes_R S
\ar[r]^{\theta \circ (1_M \otimes \delta_0^1)} &
M \otimes_{S, \delta_1^1} S_2 
\ar@<1ex>[r]^{(\theta \otimes \delta_2^2) \circ (1_M \otimes \delta^2_0)}
\ar@<-1ex>[r]_{1_{M \otimes S_2} \otimes \delta^2_1} &
M \otimes_{S, \delta_{12}^1} S_3
}
\end{equation}
obtenu en tensorisant (\ref{equation-equalizer-f}) sur $R$ par $S$
est un égalisateur.
\end{lemma}

\begin{proof}
D'après le lemme \ref{lemma-split-surjection} et la remarque
\ref{remark-functorial-splitting}, l'application
$C(1_N \otimes f): C(N \otimes_R S) \to C(N)$ admet une section
fonctorielle en $N$. On obtient ainsi les flèches verticales supérieures
du diagramme commutatif
$$
\xymatrix@C=8pc{
C(M \otimes_{S, \delta_1^1} S_2)
\ar@<1ex>^{C(\theta \circ (1_M \otimes \delta_0^1))}[r]
\ar@<-1ex>_{C(1_M \otimes \delta_1^1)}[r] \ar[d] &
C(M) \ar[r]\ar[d] & C(f_*(M,\theta)) \ar@{-->}[d] \\
C(M \otimes_{S,\delta_{12}^1} S_3)
\ar@<1ex>^{C((\theta \otimes \delta_2^2) \circ (1_M \otimes \delta^2_0))}[r]
\ar@<-1ex>_{C(1_{M \otimes S_2} \otimes \delta^2_1)}[r] \ar[d] &
C(M \otimes_{S, \delta_1^1} S_2 )
\ar[r]^{C(\theta \circ (1_M \otimes \delta_0^1))}
\ar[d]^{C(1_M \otimes \delta_1^1)} &
C(M) \ar[d] \ar@{=}[dl] \\
C(M \otimes_{S, \delta_1^1} S_2)
\ar@<1ex>[r]^{C(\theta \circ (1_M \otimes \delta_0^1))}
\ar@<-1ex>[r]_{C(1_M \otimes \delta_1^1)} &
C(M) \ar[r] &
C(f_*(M,\theta))
}
$$
dans lequel les composées le long des colonnes sont les morphismes
identités. La deuxième ligne est le diagramme coégalisateur
(\ref{equation-coequalizer-CM}), d'où la flèche en pointillés.
Le carré supérieur droit fournit des morphismes auxiliaires
$C(f_*(M,\theta)) \to 
C(M)$ et $C(M) \to C(M\otimes_{S,\delta_1^1} S_2)$ qui montrent
que la première ligne est un diagramme coégalisateur scindé.
D'après la remarque \ref{remark-adjunction}, on peut tensoriser par $S$
à l'intérieur de $C$ pour obtenir le diagramme coégalisateur scindé
$$
\xymatrix@C=8pc{
C(M \otimes_{S,\delta_2^2 \circ \delta_1^1} S_3)
\ar@<1ex>^{C((\theta \otimes \delta_2^2) \circ (1_M \otimes \delta^2_0))}[r] 
\ar@<-1ex>_{C(1_{M \otimes S_2} \otimes \delta^2_1)}[r] &
C(M \otimes_{S, \delta_1^1} S_2 )
\ar[r]^{C(\theta \circ (1_M \otimes \delta_0^1))} &
C(f_*(M,\theta) \otimes_R S).
}
$$
Le lemme \ref{lemma-C-is-faithful} montre alors que
(\ref{equation-equalizer-f2}) est lui aussi un égalisateur.
\end{proof}

\begin{remark}
\label{remark-descent-lemma}
Si $f$ est une injection scindée dans $\text{Mod}_R$, on peut simplifier
la démonstration en scindant directement $f$, sans utiliser $C$.
Le cas où $f$ est fidèlement plat est encore plus simple :
la conclusion du lemme \ref{lemma-descent-lemma} est immédiate,
car la tensorisation sur $R$ par $S$ préserve tous les égalisateurs.
\end{remark}

\begin{theorem}
\label{theorem-descent}
Les conditions suivantes sont équivalentes.
\begin{enumerate}
\item[(a)] Le morphisme $f$ est un morphisme de descente pour les modules.
\item[(b)] Le morphisme $f$ est un morphisme de descente effective pour les modules.
\item[(c)] Le morphisme $f$ est universellement injectif.
\end{enumerate}
\end{theorem}

\begin{proof}
Il est clair que (b) implique (a). Montrons que (a) implique (c).
Si $f$ n'est pas universellement injectif, il existe
$M \in \text{Mod}_R$ tel que l'application
$1_M \otimes f: M \to M \otimes_R S$ ait un noyau non nul $N$.
La projection naturelle $M \to M/N$ n'est pas un isomorphisme,
mais son image dans $DD_{S/R}$ en est un.
Ainsi, $f^*$ n'est pas pleinement fidèle.

\medskip\noindent
Montrons enfin que (c) implique (b). D'après les lemmes
\ref{lemma-equalizer-M} et \ref{lemma-descent-lemma}, pour
$(M, \theta) \in DD_{S/R}$, l'application naturelle
$f^* f_*(M,\theta) \to M$ est un isomorphisme de $S$-modules.
Pour $M_0 \in \text{Mod}_R$, montrons que $M_0 \to f_*f^*M_0$
est un isomorphisme. Par adjonction, la composée
$f^*M_0 \to f^*f_*f^*M_0 \to f^*M_0$ est l'identité, et ce qui précède
montre que la deuxième flèche est un isomorphisme. Ainsi,
$f^*M_0 \to f^*f_*f^*M_0$ est un isomorphisme. Le lemme
\ref{algebra-lemma-base-change-along-universally-injective-ring-map}
du chapitre Algèbre montre que $M_0 \to f_* f^* M_0$ est un isomorphisme.
Par conséquent, $f_*$ et $f^*$ sont des foncteurs quasi-inverses,
ce qui achève la démonstration.
\end{proof}

\subsection{Descente des propriétés des modules}
\label{subsection-descent-properties-modules}

\noindent
Dans toute cette sous-section, on fixe un morphisme d'anneaux
universellement injectif $f : R \to S$, un objet $M \in \text{Mod}_R$
et un morphisme d'anneaux $R \to A$. Nous cherchons les propriétés de
$M$ ou de $A$ qui peuvent se vérifier après extension des scalaires
par $f$. Commençons par quelques résultats de \cite{mesablishvili2}.

\begin{lemma}
\label{lemma-flat-to-injective}
Si $M \in \text{Mod}_R$ est plat, alors $C(M)$ est un $R$-module injectif.
\end{lemma}

\begin{proof}
Soit $0 \to N \to P \to Q \to 0$ une suite exacte dans $\text{Mod}_R$.
Puisque $M$ est plat, la suite
$$
0 \to N \otimes_R M \to P \otimes_R M \to Q \otimes_R M \to 0
$$
est exacte. Le lemme \ref{lemma-C-is-faithful} montre que la suite
$$
0 \to C(Q \otimes_R M) \to C(P \otimes_R M) \to C(N \otimes_R M) \to 0
$$
est exacte. En utilisant (\ref{equation-adjunction}), celle-ci s'écrit
$$
0 \to \Hom_R(Q, C(M)) \to \Hom_R(P, C(M)) \to \Hom_R(N, C(M)) \to 0.
$$
Ainsi, $C(M)$ est un objet injectif de $\text{Mod}_R$.
\end{proof}

\begin{theorem}
\label{theorem-descend-module-properties}
Si $M \otimes_R S$ possède l'une des propriétés suivantes comme $S$-module :
\begin{enumerate}
\item[(a)]
être de type fini ;
\item[(b)]
être de présentation finie ;
\item[(c)]
être plat ;
\item[(d)]
être fidèlement plat ;
\item[(e)]
être projectif de type fini ;
\end{enumerate}
alors $M$ possède la même propriété comme $R$-module, et réciproquement.
\end{theorem}

\begin{proof}
Pour démontrer (a), choisissons une famille finie $\{n_i\}$ de générateurs
de $M \otimes_R S$ dans $\text{Mod}_S$.
Écrivons chaque $n_i$ sous la forme $\sum_j m_{ij} \otimes s_{ij}$,
avec $m_{ij} \in M$ et $s_{ij} \in S$.
Soit $F$ le $R$-module libre de type fini de base $e_{ij}$, et soit
$F \to M$ le morphisme de $R$-modules qui envoie $e_{ij}$ sur $m_{ij}$.
Alors $F \otimes_R S\to M \otimes_R S$ est surjectif,
donc $\Coker(F \to M) \otimes_R S$ est nul, d'où $\Coker(F \to M)$ est nul.
Ceci démontre (a).

\medskip\noindent
Pour (b), supposons que $M \otimes_R S$ soit de présentation finie.
D'après (a), $M$ est de type fini. Choisissons une surjection
$R^{\oplus n} \to M$ de noyau $K$.
La suite $K \otimes_R S \to S^{\oplus r} \to M \otimes_R S \to 0$ est exacte.
D'après le lemme \ref{algebra-lemma-extension} du chapitre Algèbre,
le noyau de $S^{\oplus r} \to M \otimes_R S$ est un $S$-module de type fini.
Il existe donc un nombre fini d'éléments $k_1, \ldots, k_t \in K$
tels que les images de $k_i \otimes 1$ dans $S^{\oplus r}$
engendrent le noyau de $S^{\oplus r} \to M \otimes_R S$.
Soit $K' \subset K$ le sous-module engendré par $k_1, \ldots, k_t$.
Alors $M' = R^{\oplus r}/K'$ est un $R$-module de présentation finie,
muni d'un morphisme $M' \to M$ tel que
$M' \otimes_R S \to M \otimes_R S$ soit un isomorphisme.
Ainsi, $M' \cong M$, comme voulu.

\medskip\noindent
Pour démontrer (c), soit $0 \to M' \to M'' \to M \to 0$
une suite exacte courte dans $\text{Mod}_R$.
Le foncteur $\bullet \otimes_R S$ étant exact à droite,
l'application $M'' \otimes_R S \to M \otimes_R S$ est surjective.
D'après le lemme \ref{lemma-C-is-faithful}, l'application
$C(M \otimes_R S) \to C(M'' \otimes_R S)$ est donc injective.
Si $M \otimes_R S$ est plat, le lemme \ref{lemma-flat-to-injective}
montre que $C(M \otimes_R S)$ est un objet injectif de $\text{Mod}_S$.
L'injection $C(M \otimes_R S) \to C(M'' \otimes_R S)$ est donc scindée
dans $\text{Mod}_S$, et par suite dans $\text{Mod}_R$.
Comme $C(M \otimes_R S) \to C(M)$ est une surjection scindée
d'après le lemme \ref{lemma-split-surjection}, il en résulte que
$C(M) \to C(M'')$ est une injection scindée dans $\text{Mod}_R$.
Autrement dit, la suite
$$
0 \to C(M) \to C(M'') \to C(M') \to 0
$$
est exacte scindée. Pour $N \in \text{Mod}_R$,
la relation (\ref{equation-adjunction}) montre que la suite
$$
0 \to C(M \otimes_R N) \to C(M'' \otimes_R N) \to C(M' \otimes_R N) \to 0
$$
est exacte scindée. D'après le lemme \ref{lemma-C-is-faithful}, la suite
$$
0 \to M' \otimes_R N \to M'' \otimes_R N \to M \otimes_R N \to 0
$$
est exacte. Il en résulte que $M$ est plat sur $R$.
En effet, en prenant pour $M'$ un module libre se surjectant sur $M$,
on obtient $\text{Tor}_1^R(M, N) = 0$ pour tout module $N$,
et on applique le lemme \ref{algebra-lemma-characterize-flat}
du chapitre Algèbre. Ceci démontre (c).

\medskip\noindent
Pour déduire (d) de (c), remarquons que si $N \in \text{Mod}_R$
et si $M \otimes_R N$ est nul, alors
$M \otimes_R S \otimes_S (N \otimes_R S) \cong (M \otimes_R N) \otimes_R 
S$ est nul. Ainsi, $N \otimes_R S$ est nul, donc $N$ est nul.

\medskip\noindent
Pour déduire (e), il suffit maintenant de rappeler que $M$
est projectif de type fini si et seulement s'il est de présentation finie
et plat. Voir le lemme \ref{algebra-lemma-finite-projective}
du chapitre Algèbre.
\end{proof}

\noindent
Il existe une variante pour les $R$-algèbres.

\begin{theorem}
\label{theorem-descend-algebra-properties}
Si $A \otimes_R S$ possède l'une des propriétés suivantes comme $S$-algèbre :
\begin{enumerate}
\item[(a)]
être de type fini ;
\item[(b)]
être de présentation finie ;
\item[(c)]
être formellement non ramifiée ;
\item[(d)]
être non ramifiée ;
\item[(e)]
être \'etale ;
\end{enumerate}
alors $A$ possède la même propriété comme $R$-algèbre,
et la réciproque est bien entendu vraie.
\end{theorem}

\begin{proof}
Pour démontrer (a), choisissons une famille finie $\{x_i\}$ de générateurs
de $A \otimes_R S$ sur $S$.
Écrivons chaque $x_i$ sous la forme $\sum_j y_{ij} \otimes s_{ij}$,
avec $y_{ij} \in A$ et $s_{ij} \in S$.
Soit $F$ la $R$-algèbre de polynômes en les variables $e_{ij}$,
et soit $F \to M$ le morphisme de $R$-algèbres qui envoie
$e_{ij}$ sur $y_{ij}$. Alors $F \otimes_R S\to A \otimes_R S$ est surjectif,
donc $\Coker(F \to A) \otimes_R S$ est nul, d'où $\Coker(F \to A)$ est nul.
Ceci démontre (a).

\medskip\noindent
Pour (b), supposons que $A \otimes_R S$ soit une $S$-algèbre de présentation finie.
D'après (a), $A$ est de type fini sur $R$.
Choisissons une surjection $R[x_1, \ldots, x_n] \to A$ de noyau $I$.
La suite $I \otimes_R S \to S[x_1, \ldots, x_n] \to A \otimes_R S \to 0$
est exacte. D'après le lemme
\ref{algebra-lemma-finite-presentation-independent} du chapitre Algèbre,
le noyau de $S[x_1, \ldots, x_n] \to A \otimes_R S$
est un idéal de type fini.
Il existe donc un nombre fini d'éléments $y_1, \ldots, y_t \in I$
tels que les images de $y_i \otimes 1$ dans $S[x_1, \ldots, x_n]$
engendrent le noyau de $S[x_1, \ldots, x_n] \to A \otimes_R S$.
Soit $I' \subset I$ l'idéal engendré par $y_1, \ldots, y_t$.
Alors $A' = R[x_1, \ldots, x_n]/I'$ est une $R$-algèbre de présentation finie,
munie d'un morphisme $A' \to A$ tel que
$A' \otimes_R S \to A \otimes_R S$ soit un isomorphisme.
Ainsi, $A' \cong A$, comme voulu.

\medskip\noindent
Pour démontrer (c), rappelons que $A$ est formellement non ramifiée sur $R$
si et seulement si le module des différentielles relatives $\Omega_{A/R}$
est nul ; voir le lemme \ref{algebra-lemma-characterize-formally-unramified}
du chapitre Algèbre ou \cite[Proposition~17.2.1]{EGA4}.
Comme $\Omega_{(A \otimes_R S)/S} = \Omega_{A/R} \otimes_R S$,
la nullité descend en vertu du théorème \ref{theorem-descent}.

\medskip\noindent
Pour déduire (d) des cas précédents, rappelons que $A$ est non ramifiée
sur $R$ si et seulement si $A$ est formellement non ramifiée
et de type fini sur $R$ ; voir le lemme
\ref{algebra-lemma-formally-unramified-unramified} du chapitre Algèbre.

\medskip\noindent
Pour démontrer (e), rappelons que, d'après le lemme
\ref{algebra-lemma-etale-flat-unramified-finite-presentation} du chapitre Algèbre
ou \cite[Th\'eor\`eme~17.6.1]{EGA4}, l'algèbre $A$ est \'etale sur $R$
si et seulement si $A$ est plate, non ramifiée et de présentation finie sur $R$.
\end{proof}

\begin{remark}
\label{remark-when-locally-split}
La situation serait plus simple si l'on disposait d'un morphisme d'anneaux
fidèlement plat $g: R \to T$ tel que $T \to S \otimes_R T$
possède une structure supplémentaire.
Par exemple, si l'on pouvait faire en sorte que $T \to S \otimes_R T$
soit scindé dans $\textit{Rings}$, toute propriété d'un module ou d'une algèbre
qui est stable par extension des scalaires et que l'on peut faire descendre le
long de morphismes fidèlement plats pourrait aussi être descendue le long de
morphismes universellement injectifs.
On pourrait naturellement chercher $g$ tel que $T$ soit non seulement
fidèlement plat, mais encore injectif dans $\text{Mod}_R$.
Cependant, même pour $R = \mathbf{Z}$, un tel morphisme ne peut exister.
\end{remark}
















\section{Descente fpqc des faisceaux quasi-cohérents}
\label{section-fpqc-descent-quasi-coherent}

\noindent
La principale application de la descente plate pour les modules
est l'énoncé correspondant de descente des faisceaux quasi-cohérents
relativement aux recouvrements fpqc.

\begin{lemma}
\label{lemma-standard-fpqc-covering}
Soit $S$ un schéma affine.
Soit $\mathcal{U} = \{f_i : U_i \to S\}_{i = 1, \ldots, n}$
un recouvrement fpqc standard de $S$ ; voir la définition
\ref{topologies-definition-standard-fpqc} du chapitre Topologies.
Toute donnée de descente sur des faisceaux quasi-cohérents
relativement à $\mathcal{U} = \{U_i \to S\}$ est effective.
De plus, le foncteur de la catégorie des $\mathcal{O}_S$-modules
quasi-cohérents vers la catégorie des données de descente
relativement à $\mathcal{U}$ est pleinement fidèle.
\end{lemma}

\begin{proof}
Il s'agit de la proposition \ref{proposition-descent-module},
reformulée en termes de schémas.
Remarquons d'abord qu'une donnée de descente $\xi$ sur des faisceaux
quasi-cohérents relativement à $\mathcal{U}$ revient exactement
à une donnée de descente $\xi'$ sur des faisceaux quasi-cohérents
relativement au recouvrement
$\mathcal{U}' = \{\coprod_{i = 1, \ldots, n} U_i \to S\}$.
De plus, l'effectivité de $\xi$ équivaut à celle de $\xi'$.
On peut donc supposer $n = 1$, c'est-à-dire $\mathcal{U} = \{U \to S\}$,
où $U$ et $S$ sont affines.
Dans ce cas, les données de descente correspondent aux données
de descente sur les modules relativement au morphisme d'anneaux
$$
\Gamma(S, \mathcal{O})
\longrightarrow
\Gamma(U, \mathcal{O}).
$$
Comme $U \to S$ est surjectif et plat, ce morphisme d'anneaux
est fidèlement plat. La proposition \ref{proposition-descent-module}
s'applique donc et donne le résultat.
\end{proof}

\begin{proposition}
\label{proposition-fpqc-descent-quasi-coherent}
Soit $S$ un schéma.
Soit $\mathcal{U} = \{\varphi_i : U_i \to S\}$ un recouvrement fpqc ;
voir la définition \ref{topologies-definition-fpqc-covering} du chapitre Topologies.
Toute donnée de descente sur des faisceaux quasi-cohérents
relativement à $\mathcal{U} = \{U_i \to S\}$ est effective.
De plus, le foncteur de la catégorie des $\mathcal{O}_S$-modules
quasi-cohérents vers la catégorie des données de descente
relativement à $\mathcal{U}$ est pleinement fidèle.
\end{proposition}

\begin{proof}
Soit $S = \bigcup_{j \in J} V_j$ un recouvrement par des ouverts affines.
Pour $j, j' \in J$, notons $V_{jj'} = V_j \cap V_{j'}$
l'intersection, qui n'est pas nécessairement affine.
Pour un ouvert $V \subset S$, notons
$\mathcal{U}_V = \{V \times_S U_i \to V\}_{i \in I}$ ;
c'est un recouvrement fpqc d'après le lemme
\ref{topologies-lemma-fpqc} du chapitre Topologies.
Par définition d'un recouvrement fpqc, il existe, pour chaque $j \in J$,
un ensemble fini $K_j$, une application $\underline{i} : K_j \to I$
et des ouverts affines $U_{\underline{i}(k), k} \subset U_{\underline{i}(k)}$,
$k \in K_j$, tels que
$\mathcal{V}_j = \{U_{\underline{i}(k), k} \to V_j\}_{k \in K_j}$
soit un recouvrement fpqc standard de $V_j$.
De plus, $\mathcal{V}_j$ est un raffinement de $\mathcal{U}_{V_j}$.
On a le diagramme
$$
\xymatrix{
\mathcal{V}_j \ar[r] \ar@{~>}[d] &
\mathcal{U}_{V_j} \ar[r] \ar@{~>}[d] &
\mathcal{U} \ar@{~>}[d] \\
V_j \ar@{=}[r] & V_j \ar[r] & S
}
$$
où les flèches horizontales supérieures sont des morphismes de familles
de morphismes de but fixé ; voir la définition
\ref{sites-definition-morphism-coverings} du chapitre Sites.

\medskip\noindent
Pour démontrer la proposition, nous établissons successivement
la fidélité, la plénitude et la surjectivité essentielle
du foncteur des faisceaux quasi-cohérents vers les données de descente.

\medskip\noindent
Fidélité. Soient $\mathcal{F}$ et $\mathcal{G}$ des faisceaux quasi-cohérents
sur $S$, et soient $a, b : \mathcal{F} \to \mathcal{G}$
des morphismes de $\mathcal{O}_S$-modules.
Supposons que $\varphi_i^*(a) = \varphi_i^*(b)$ pour tout $i$.
Soit $s \in S$. On a $s = \varphi_i(u)$ pour certains $i \in I$
et $u \in U_i$.
Puisque $\mathcal{O}_{S, s} \to \mathcal{O}_{U_i, u}$ est plat,
donc fidèlement plat par le lemme \ref{algebra-lemma-local-flat-ff}
du chapitre Algèbre, on a
$a_s = b_s : \mathcal{F}_s \to \mathcal{G}_s$. Ainsi, $a = b$.

\medskip\noindent
Pleine fidélité. Soient $\mathcal{F}$ et $\mathcal{G}$ des faisceaux
quasi-cohérents sur $S$, et soient
$a_i : \varphi_i^*\mathcal{F} \to \varphi_i^*\mathcal{G}$
des morphismes de $\mathcal{O}_{U_i}$-modules tels que
$\text{pr}_0^*a_i = \text{pr}_1^*a_j$ sur $U_i \times_U U_j$.
Par image inverse, on obtient des morphismes
$$
a_k :
\varphi_{i(k)}^*\mathcal{F}|_{U_{\underline{i}(k), k}}
\longrightarrow
\varphi_{i(k)}^*\mathcal{G}|_{U_{\underline{i}(k), k}}
$$
pour $k \in K_j$, avec les notations précédentes.
Le lemme \ref{lemma-refine-descent-datum} montre qu'ils définissent
un morphisme entre les données de descente canoniques sur $\mathcal{V}_j$.
Le lemme \ref{lemma-standard-fpqc-covering} fournit donc
d'uniques morphismes $a_j : \mathcal{F}|_{V_j} \to \mathcal{G}|_{V_j}$
correspondants. Pour montrer que $a_j|_{V_{jj'}} = a_{j'}|_{V_{jj'}}$,
on remarque que $a_j$ et $a_{j'}$ correspondent tous deux à l'image inverse
du morphisme $(a_i)_{i \in I}$ de données de descente canoniques sur
tout recouvrement raffinant à la fois $\mathcal{V}_{j, V_{jj'}}$
et $\mathcal{V}_{j', V_{jj'}}$, puis on utilise la fidélité déjà démontrée.
On peut prendre, par exemple, le recouvrement
$\mathcal{V}_{jj'} =
\{V_k \times_S V_{k'} \to V_{jj'}\}_{k \in K_j, k' \in K_{j'}}$.

\medskip\noindent
Surjectivité essentielle. Soit $\xi = (\mathcal{F}_i, \varphi_{ii'})$
une donnée de descente sur des faisceaux quasi-cohérents relativement
au recouvrement $\mathcal{U}$.
Par image inverse, elle donne des données de descente $\xi_j$
sur des faisceaux quasi-cohérents relativement aux recouvrements
$\mathcal{V}_j$ des $V_j$.
Le lemme \ref{lemma-standard-fpqc-covering} fournit des faisceaux
quasi-cohérents $\mathcal{F}_j$ sur $V_j$ dont la donnée de descente
canonique associée est isomorphe à $\xi_j$.
La pleine fidélité démontrée ci-dessus fournit des isomorphismes
$$
\phi_{jj'} :
\mathcal{F}_j|_{V_{jj'}}
\longrightarrow
\mathcal{F}_{j'}|_{V_{jj'}}
$$
correspondant à l'isomorphisme entre les images inverses
de $\xi_j$ et de $\xi_{j'}$ sur $\mathcal{V}_{jj'}$.
Pour vérifier que les applications $\phi_{jj'}$ satisfont
la condition de cocycle, on utilise la fidélité déjà démontrée
sur les intersections triples $V_{jj'j''}$.
Le lemme \ref{lemma-zariski-descent-effective} montre alors
que les faisceaux $\mathcal{F}_j$ se recollent en un faisceau
quasi-cohérent $\mathcal{F}$, comme voulu.
Il reste à vérifier que la donnée de descente canonique relative à
$\mathcal{U}$ associée à $\mathcal{F}$ est isomorphe à la donnée de descente
initiale. Cette vérification est omise.
\end{proof}









\section{Descente galoisienne des faisceaux quasi-cohérents}
\label{section-galois-descent}

\noindent
La descente galoisienne des faisceaux quasi-cohérents est un cas particulier
de leur descente fpqc. Dans cette section, nous expliquons comment
reformuler une donnée de descente galoisienne en une donnée de descente fpqc,
puis nous appliquons les résultats précédents.

\medskip\noindent
Soit $k'/k$ une extension de corps. Alors $\{\Spec(k') \to \Spec(k)\}$
est un recouvrement fpqc. Soit $X$ un schéma sur $k$.
Pour une $k$-algèbre $A$, posons $X_A = X \times_{\Spec(k)} \Spec(A)$.
Le lemme \ref{topologies-lemma-fpqc} du chapitre Topologies montre que
$\{X_{k'} \to X\}$ est un recouvrement fpqc. On a
$$
X_{k'} \times_X X_{k'} = X_{k' \otimes_k k'}
\quad\text{et}\quad
X_{k'} \times_X X_{k'} \times_X X_{k'} = X_{k' \otimes_k k' \otimes_k k'}
$$
Une donnée de descente sur des faisceaux quasi-cohérents relativement à
$\{X_{k'} \to X\}$ consiste donc en un faisceau quasi-cohérent $\mathcal{F}$
sur $X_{k'}$ et en un isomorphisme
$\varphi : \text{pr}_0^*\mathcal{F} \to \text{pr}_1^*\mathcal{F}$
sur $X_{k' \otimes_k k'}$ satisfaisant la condition de cocycle évidente
sur $X_{k' \otimes_k k' \otimes_k k'}$.
Précisons ce que cela signifie dans le cas d'une extension galoisienne.

\medskip\noindent
Soit $k'/k$ une extension galoisienne finie de groupe de Galois
$G = \text{Gal}(k'/k)$. On a des isomorphismes de $k$-algèbres
$$
k' \otimes_k k' \longrightarrow \prod\nolimits_{\sigma \in G} k',\quad
a \otimes b \longrightarrow \prod a\sigma(b)
$$
et
$$
k' \otimes_k k' \otimes_k k' \longrightarrow
\prod\nolimits_{(\sigma, \tau) \in G \times G} k',\quad
a \otimes b \otimes c \longrightarrow \prod a\sigma(b)\sigma(\tau(c))
$$
Nous choisissons ici $a\sigma(b)\sigma(\tau(c))$, et non
$a\sigma(b)\tau(c)$, afin de simplifier les formules qui suivent,
bien que ce choix ne soit pas indispensable. Pour $\sigma \in G$, notons
$$
f_\sigma = \text{id}_X \times \Spec(\sigma) :
X_{k'} \longrightarrow X_{k'}
$$
Puisque $\Spec(-)$ est contravariant, on a
$f_{\sigma \tau} = f_\tau \circ f_\sigma$, et non la composée dans l'autre ordre.
Le premier isomorphisme ci-dessus fournit une identification
$$
X_{k' \otimes_k k'} = \coprod\nolimits_{\sigma \in G} X_{k'}
$$
par laquelle $\text{pr}_0$ correspond à l'application
$$
\coprod\nolimits_{\sigma \in G} X_{k'}
\xrightarrow{\coprod \text{id}}
X_{k'}
$$
et $\text{pr}_1$ à l'application
$$
\coprod\nolimits_{\sigma \in G} X_{k'}
\xrightarrow{\coprod f_\sigma}
X_{k'}
$$
Une donnée de descente $\varphi$ sur $\mathcal{F}$ au-dessus de $X_{k'}$
correspond donc à une famille d'isomorphismes
$\varphi_\sigma : \mathcal{F} \to f_\sigma^*\mathcal{F}$.
Pour expliciter la condition de cocycle, utilisons l'identification
fournie par le second isomorphisme d'algèbres ci-dessus :
$$
X_{k' \otimes_k k' \otimes_k k'} =
\coprod\nolimits_{(\sigma, \tau) \in G \times G} X_{k'}.
$$
Par cette identification, $\text{pr}_{01}$ correspond à l'application
$$
\coprod\nolimits_{(\sigma, \tau) \in G \times G} X_{k'}
\longrightarrow
\coprod\nolimits_{\sigma \in G} X_{k'}
$$
qui envoie la composante d'indice $(\sigma, \tau)$ sur celle d'indice
$\sigma$ par le morphisme identité.
L'application $\text{pr}_{12}$ correspond à l'application
$$
\coprod\nolimits_{(\sigma, \tau) \in G \times G} X_{k'}
\longrightarrow
\coprod\nolimits_{\sigma \in G} X_{k'}
$$
qui envoie la composante d'indice $(\sigma, \tau)$ sur celle d'indice $\tau$
par le morphisme $f_\sigma$. Enfin, $\text{pr}_{02}$ correspond à l'application
$$
\coprod\nolimits_{(\sigma, \tau) \in G \times G} X_{k'}
\longrightarrow
\coprod\nolimits_{\sigma \in G} X_{k'}
$$
qui envoie la composante d'indice $(\sigma, \tau)$ sur celle d'indice
$\sigma\tau$ par le morphisme identité. La condition de cocycle
$$
\text{pr}_{02}^*\varphi = \text{pr}_{12}^*\varphi \circ \text{pr}_{01}^*\varphi
$$
se traduit donc par une condition pour chaque couple $(\sigma, \tau)$, à savoir
$$
\varphi_{\sigma\tau} = f_\sigma^*\varphi_\tau \circ \varphi_\sigma
$$
comme applications $\mathcal{F} \to f_{\sigma\tau}^*\mathcal{F}$.
(Les sources et les buts coïncident : par exemple, le but de $\varphi_\sigma$
est $f_\sigma^*\mathcal{F}$, et la source de
$f_\sigma^*\varphi_\tau$ est également $f_\sigma^*\mathcal{F}$.)

\begin{lemma}
\label{lemma-galois-descent}
Soit $k'/k$ une extension galoisienne (finie) de groupe de Galois $G$.
Soit $X$ un schéma sur $k$. La catégorie des $\mathcal{O}_X$-modules
quasi-cohérents est équivalente à la catégorie des systèmes
$(\mathcal{F}, (\varphi_\sigma)_{\sigma \in G})$ où
\begin{enumerate}
\item $\mathcal{F}$ est un module quasi-cohérent sur $X_{k'}$ ;
\item $\varphi_\sigma : \mathcal{F} \to f_\sigma^*\mathcal{F}$
est un isomorphisme de modules ;
\item $\varphi_{\sigma\tau} = f_\sigma^*\varphi_\tau \circ \varphi_\sigma$
pour tous $\sigma, \tau \in G$.
\end{enumerate}
Ici, $f_\sigma = \text{id}_X \times \Spec(\sigma) : X_{k'} \to X_{k'}$.
\end{lemma}

\begin{proof}
Comme nous l'avons vu, une donnée
$(\mathcal{F}, (\varphi_\sigma)_{\sigma \in G})$ comme dans l'énoncé
revient à une donnée de descente relativement au recouvrement fpqc
$\{X_{k'} \to X\}$. Le lemme résulte donc de la proposition
\ref{proposition-fpqc-descent-quasi-coherent}.
\end{proof}

\noindent
Voici une légère généralisation. Supposons que l'on dispose d'un morphisme
fini \'etale surjectif $X \to Y$, d'un groupe fini $G$ et d'un morphisme de groupes
$G^{opp} \to \text{Aut}_Y(X), \sigma \mapsto f_\sigma$
tel que l'application
$$
G \times X \longrightarrow X \times_Y X,\quad
(\sigma, x) \longmapsto (x, f_\sigma(x))
$$
soit un isomorphisme. Le même résultat s'applique.

\begin{lemma}
\label{lemma-galois-descent-more-general}
Soient $X \to Y$, $G$ et $f_\sigma : X \to X$ comme ci-dessus.
La catégorie des $\mathcal{O}_Y$-modules quasi-cohérents est équivalente
à la catégorie des systèmes $(\mathcal{F}, (\varphi_\sigma)_{\sigma \in G})$ où
\begin{enumerate}
\item $\mathcal{F}$ est un $\mathcal{O}_X$-module quasi-cohérent ;
\item $\varphi_\sigma : \mathcal{F} \to f_\sigma^*\mathcal{F}$
est un isomorphisme de modules ;
\item $\varphi_{\sigma\tau} = f_\sigma^*\varphi_\tau \circ \varphi_\sigma$
pour tous $\sigma, \tau \in G$.
\end{enumerate}
\end{lemma}

\begin{proof}
Puisque $X \to Y$ est fini \'etale surjectif, $\{X \to Y\}$
est un recouvrement fpqc. Comme
$G \times X \to X \times_Y X$, $(\sigma, x) \mapsto (x, f_\sigma(x))$
est un isomorphisme, l'application
$G \times G \times X \to X \times_Y X \times_Y X$,
$(\sigma, \tau, x) \mapsto (x, f_\sigma(x), f_{\sigma\tau}(x))$
en est un également. Ces identifications montrent que la catégorie
des données de l'énoncé s'identifie à celle des données de descente
sur les faisceaux quasi-cohérents relativement au recouvrement
$\{x \to Y\}$. Le lemme résulte donc de la proposition
\ref{proposition-fpqc-descent-quasi-coherent}.
\end{proof}









\section{Descente des propriétés de finitude des modules}
\label{section-descent-finiteness}

\noindent
Dans cette section, nous montrons que certaines conditions de finitude
sur un module quasi-cohérent peuvent se vérifier sur les membres d'un
recouvrement.

\begin{lemma}
\label{lemma-finite-type-descends}
Soit $X$ un schéma.
Soit $\mathcal{F}$ un $\mathcal{O}_X$-module quasi-cohérent.
Soit $\{f_i : X_i \to X\}_{i \in I}$ un recouvrement fpqc tel que
chaque $f_i^*\mathcal{F}$ soit un $\mathcal{O}_{X_i}$-module de type fini.
Alors $\mathcal{F}$ est un $\mathcal{O}_X$-module de type fini.
\end{lemma}

\begin{proof}
Démonstration omise. Pour le cas affine, voir
Algèbre, lemme \ref{algebra-lemma-descend-properties-modules}.
\end{proof}

\begin{lemma}
\label{lemma-finite-type-descends-fppf}
Soit $f : (X, \mathcal{O}_X) \to (Y, \mathcal{O}_Y)$ un morphisme
d'espaces localement annelés. Soit $\mathcal{F}$ un faisceau de $\mathcal{O}_Y$-modules.
Si
\begin{enumerate}
\item l'application continue sous-jacente à $f$ est ouverte,
\item $f$ est surjectif et plat, et
\item $f^*\mathcal{F}$ est de type fini,
\end{enumerate}
alors $\mathcal{F}$ est de type fini.
\end{lemma}

\begin{proof}
Soit $y \in Y$ un point. Choisissons un point $x \in X$ d'image $y$.
Choisissons un ouvert $x \in U \subset X$ et des éléments $s_1, \ldots, s_n$
de $f^*\mathcal{F}(U)$ qui engendrent $f^*\mathcal{F}$ sur $U$.
Puisque $f^*\mathcal{F} =
f^{-1}\mathcal{F} \otimes_{f^{-1}\mathcal{O}_Y} \mathcal{O}_X$,
nous pouvons, quitte à restreindre $U$, supposer que $s_i = \sum t_{ij} \otimes a_{ij}$
avec $t_{ij} \in f^{-1}\mathcal{F}(U)$ et $a_{ij} \in \mathcal{O}_X(U)$.
En restreignant encore $U$, nous pouvons supposer que $t_{ij}$ provient
d'une section $s_{ij} \in \mathcal{F}(V)$ pour un ouvert $V \subset Y$
tel que $f(U) \subset V$. Soit $N$ le nombre de sections $s_{ij}$ et
considérons l'application
$$
\sigma = (s_{ij}) : \mathcal{O}_V^{\oplus N} \to \mathcal{F}|_V
$$
Le choix des sections montre que $f^*\sigma|_U$ est surjective.
Ainsi, pour tout $u \in U$, l'application
$$
\sigma_{f(u)} \otimes_{\mathcal{O}_{Y, f(u)}} \mathcal{O}_{X, u} :
\mathcal{O}_{X, u}^{\oplus N}
\longrightarrow
\mathcal{F}_{f(u)} \otimes_{\mathcal{O}_{Y, f(u)}} \mathcal{O}_{X, u}
$$
est surjective. Comme $f$ est plat, l'homomorphisme local d'anneaux
$\mathcal{O}_{Y, f(u)} \to \mathcal{O}_{X, u}$ est plat, donc
fidèlement plat (Algèbre, lemme \ref{algebra-lemma-local-flat-ff}).
Par conséquent, $\sigma_{f(u)}$ est surjective. Puisque $f$ est ouvert, $f(U)$ est
un voisinage ouvert de $y$, ce qui achève la démonstration.
\end{proof}

\begin{lemma}
\label{lemma-finite-presentation-descends}
Soit $X$ un schéma.
Soit $\mathcal{F}$ un $\mathcal{O}_X$-module quasi-cohérent.
Soit $\{f_i : X_i \to X\}_{i \in I}$ un recouvrement fpqc tel que
chaque $f_i^*\mathcal{F}$ soit un $\mathcal{O}_{X_i}$-module de présentation
finie. Alors $\mathcal{F}$ est un $\mathcal{O}_X$-module
de présentation finie.
\end{lemma}

\begin{proof}
Démonstration omise. Pour le cas affine, voir
Algèbre, lemme \ref{algebra-lemma-descend-properties-modules}.
\end{proof}

\begin{lemma}
\label{lemma-locally-generated-by-r-sections-descends}
Soit $X$ un schéma.
Soit $\mathcal{F}$ un $\mathcal{O}_X$-module quasi-cohérent.
Soit $\{f_i : X_i \to X\}_{i \in I}$ un recouvrement fpqc tel que
chaque $f_i^*\mathcal{F}$ soit localement engendré par $r$ sections comme
$\mathcal{O}_{X_i}$-module. Alors $\mathcal{F}$ est localement engendré par
$r$ sections comme $\mathcal{O}_X$-module.
\end{lemma}

\begin{proof}
D'après le lemme \ref{lemma-finite-type-descends}, le module $\mathcal{F}$
est de type fini. Le lemme de Nakayama
(Algèbre, lemme \ref{algebra-lemma-NAK}) implique donc que $\mathcal{F}$
est engendré par $r$ sections au voisinage d'un point $x \in X$
si et seulement si $\dim_{\kappa(x)} \mathcal{F}_x \otimes \kappa(x) \leq r$.
Choisissons un $i$ et un point $x_i \in X_i$ d'image $x$. Alors
$\dim_{\kappa(x)} \mathcal{F}_x \otimes \kappa(x) = 
\dim_{\kappa(x_i)} (f_i^*\mathcal{F})_{x_i} \otimes \kappa(x_i)$,
et cette dimension est $\leq r$, puisque $f_i^*\mathcal{F}$ est localement engendré par $r$
sections.
\end{proof}

\begin{lemma}
\label{lemma-flat-descends}
Soit $X$ un schéma.
Soit $\mathcal{F}$ un $\mathcal{O}_X$-module quasi-cohérent.
Soit $\{f_i : X_i \to X\}_{i \in I}$ un recouvrement fpqc tel que
chaque $f_i^*\mathcal{F}$ soit un $\mathcal{O}_{X_i}$-module plat.
Alors $\mathcal{F}$ est un $\mathcal{O}_X$-module plat.
\end{lemma}

\begin{proof}
Démonstration omise. Pour le cas affine, voir
Algèbre, lemme \ref{algebra-lemma-descend-properties-modules}.
\end{proof}

\begin{lemma}
\label{lemma-finite-locally-free-descends}
Soit $X$ un schéma.
Soit $\mathcal{F}$ un $\mathcal{O}_X$-module quasi-cohérent.
Soit $\{f_i : X_i \to X\}_{i \in I}$ un recouvrement fpqc tel que
chaque $f_i^*\mathcal{F}$ soit un $\mathcal{O}_{X_i}$-module localement libre de type fini.
Alors $\mathcal{F}$ est un $\mathcal{O}_X$-module localement libre de type fini.
\end{lemma}

\begin{proof}
Cela résulte du fait qu'un faisceau quasi-cohérent est localement libre de type fini
si et seulement s'il est de présentation finie et plat, voir
Algèbre, lemme \ref{algebra-lemma-finite-projective}.
En effet, si chaque $f_i^*\mathcal{F}$ est plat et de présentation finie,
il en est de même de $\mathcal{F}$ d'après les
lemmes \ref{lemma-flat-descends} et
\ref{lemma-finite-presentation-descends}.
\end{proof}

\noindent
La définition d'un faisceau quasi-cohérent localement projectif figure dans
Propriétés, section \ref{properties-section-locally-projective}.

\begin{lemma}
\label{lemma-locally-projective-descends}
Soit $X$ un schéma.
Soit $\mathcal{F}$ un $\mathcal{O}_X$-module quasi-cohérent.
Soit $\{f_i : X_i \to X\}_{i \in I}$ un recouvrement fpqc tel que
chaque $f_i^*\mathcal{F}$ soit un $\mathcal{O}_{X_i}$-module localement projectif.
Alors $\mathcal{F}$ est un $\mathcal{O}_X$-module localement projectif.
\end{lemma}

\begin{proof}
Démonstration omise. Pour les recouvrements de Zariski, c'est
Propriétés, lemme \ref{properties-lemma-locally-projective}.
Pour le cas affine, c'est
Algèbre, théorème \ref{algebra-theorem-ffdescent-projectivity}.
\end{proof}

\begin{remark}
\label{remark-locally-free-descends}
Être localement libre est une propriété des modules quasi-cohérents qui
ne descend pas pour la topologie fpqc. En effet, supposons que
$R$ soit un anneau et que $M$ soit un $R$-module projectif qui est
une somme directe dénombrable $M = \bigoplus L_n$ de modules localement
libres de rang 1, sans être lui-même localement libre, voir
Exemples, lemme \ref{examples-lemma-projective-not-locally-free}.
Alors $M$ devient libre après le changement de base fidèlement plat
$$
R \longrightarrow
\bigoplus\nolimits_{m \geq 1}
\bigoplus\nolimits_{(i_1, \ldots, i_m) \in \mathbf{Z}^{\oplus m}}
L_1^{\otimes i_1} \otimes_R \ldots \otimes_R L_m^{\otimes i_m}
$$
Mais nous ne savons pas ce qu'il en est pour les recouvrements fppf.
Autrement dit, nous ne connaissons pas la réponse à la question suivante :
supposons que $A \to B$ soit un homomorphisme d'anneaux fidèlement
plat et de présentation finie. Soit $M$ un $A$-module
tel que $M \otimes_A B$ soit libre. Est-ce que $M$ est un
$A$-module localement libre ? Lorsque $A$ est noethérien, la réponse
est affirmative. Cela résulte des résultats de \cite{Bass}. Mais, en général,
nous ne connaissons pas la réponse. Si vous la connaissez, ou si vous disposez
d'une référence, veuillez écrire à
\href{mailto:stacks.project@gmail.com}{stacks.project@gmail.com}.
\end{remark}

\noindent
Nous ajoutons ici deux résultats liés aux précédents, mais de nature
légèrement différente.

\begin{lemma}
\label{lemma-finite-over-finite-module}
Soit $f : X \to Y$ un morphisme de schémas.
Soit $\mathcal{F}$ un $\mathcal{O}_X$-module quasi-cohérent.
Supposons que $f$ soit un morphisme fini.
Alors $\mathcal{F}$ est un $\mathcal{O}_X$-module de type fini
si et seulement si $f_*\mathcal{F}$ est un $\mathcal{O}_Y$-module de type
fini.
\end{lemma}

\begin{proof}
Comme $f$ est fini, il est affine. On se ramène donc au cas où
$f$ est le morphisme $\Spec(B) \to \Spec(A)$ défini
par un homomorphisme fini d'anneaux $A \to B$.
Alors $\mathcal{F} = \widetilde{M}$ est le faisceau de modules
associé au $B$-module $M$.
Notons que $M$ est de type fini comme $B$-module si et seulement si
$M$ est de type fini comme $A$-module, voir
Algèbre, lemme \ref{algebra-lemma-finite-module-over-finite-extension}.
En combinant cela avec
Propriétés, lemme \ref{properties-lemma-finite-type-module},
on obtient le résultat.
\end{proof}

\begin{lemma}
\label{lemma-finite-finitely-presented-module}
Soit $f : X \to Y$ un morphisme de schémas.
Soit $\mathcal{F}$ un $\mathcal{O}_X$-module quasi-cohérent.
Supposons que $f$ soit fini et de présentation finie.
Alors $\mathcal{F}$ est un $\mathcal{O}_X$-module de présentation finie
si et seulement si $f_*\mathcal{F}$ est un $\mathcal{O}_Y$-module de présentation
finie.
\end{lemma}

\begin{proof}
Comme $f$ est fini, il est affine. On se ramène donc au cas où
$f$ est le morphisme $\Spec(B) \to \Spec(A)$ défini
par un homomorphisme d'anneaux fini et de présentation finie $A \to B$.
Alors $\mathcal{F} = \widetilde{M}$ est le faisceau de modules
associé au $B$-module $M$.
Notons que $M$ est de présentation finie comme $B$-module si et seulement si
$M$ est de présentation finie comme $A$-module, voir
Algèbre, lemme \ref{algebra-lemma-finite-finitely-presented-extension}.
En combinant cela avec
Propriétés, lemme \ref{properties-lemma-finite-presentation-module},
on obtient le résultat.
\end{proof}
















\section{Faisceaux quasi-cohérents et topologies, I}
\label{section-quasi-coherent-sheaves}

\noindent
Les résultats de cette section donnent une équivalence naturelle entre
la catégorie des modules quasi-cohérents sur un schéma $S$ et la catégorie
des modules quasi-cohérents sur de nombreux sites associés à $S$
dans le chapitre consacré aux topologies.

\medskip\noindent
Soit $S$ un schéma.
Soit $\mathcal{F}$ un $\mathcal{O}_S$-module quasi-cohérent.
Considérons le foncteur
\begin{equation}
\label{equation-quasi-coherent-presheaf}
(\Sch/S)^{opp} \longrightarrow \textit{Ab},
\quad
(f : T \to S) \longmapsto \Gamma(T, f^*\mathcal{F}).
\end{equation}

\begin{lemma}
\label{lemma-sheaf-condition-holds}
Soit $S$ un schéma.
Soit $\mathcal{F}$ un $\mathcal{O}_S$-module quasi-cohérent.
Soit $\tau \in \{Zariski, \linebreak[0] \etale, \linebreak[0] smooth,
\linebreak[0] syntomic, \linebreak[0] fppf, \linebreak[0] fpqc\}$.
Le foncteur défini dans (\ref{equation-quasi-coherent-presheaf})
satisfait la condition de faisceau relativement à tout recouvrement
pour la topologie $\tau$
$\{T_i \to T\}_{i \in I}$ d'un schéma $T$ sur $S$.
\end{lemma}

\begin{proof}
Pour $\tau \in \{Zariski, \linebreak[0] \etale, \linebreak[0] smooth,
\linebreak[0] syntomic, \linebreak[0] fppf\}$, un recouvrement pour la topologie $\tau$
est aussi un recouvrement fpqc, voir
Topologies, lemmes
\ref{topologies-lemma-zariski-etale},
\ref{topologies-lemma-zariski-etale-smooth},
\ref{topologies-lemma-zariski-etale-smooth-syntomic},
\ref{topologies-lemma-zariski-etale-smooth-syntomic-fppf} et
\ref{topologies-lemma-zariski-etale-smooth-syntomic-fppf-fpqc}.
Il suffit donc de démontrer le théorème
pour un recouvrement fpqc. Supposons que $\{f_i : T_i \to T\}_{i \in I}$
soit un recouvrement fpqc, avec $f : T \to S$ donné. Supposons
donnée une famille de sections $s_i \in \Gamma(T_i , f_i^*f^*\mathcal{F})$
telle que $s_i|_{T_i \times_T T_j} = s_j|_{T_i \times_T T_j}$.
Il faut trouver la section correspondante $s \in \Gamma(T, f^*\mathcal{F})$.
On peut interpréter les $s_i$ comme une famille de morphismes
$\varphi_i : f_i^*\mathcal{O}_T = \mathcal{O}_{T_i} \to f_i^*f^*\mathcal{F}$
compatibles avec les données de descente canoniques associées aux
faisceaux quasi-cohérents $\mathcal{O}_T$ et $f^*\mathcal{F}$ sur $T$.
D'après la proposition \ref{proposition-fpqc-descent-quasi-coherent},
on peut donc les faire descendre, de manière unique,
en une application $\mathcal{O}_T \to f^*\mathcal{F}$, qui fournit
la section $s$ cherchée.
\end{proof}

\noindent
On peut en particulier donner la définition suivante.

\begin{definition}
\label{definition-structure-sheaf}
Soit $\tau \in \{Zariski, \linebreak[0] \etale, \linebreak[0]
smooth, \linebreak[0] syntomic, \linebreak[0] fppf\}$.
Soit $S$ un schéma.
Soit $\Sch_\tau$ un grand site contenant $S$.
Soit $\mathcal{F}$ un $\mathcal{O}_S$-module quasi-cohérent.
\begin{enumerate}
\item Le {\it faisceau structural du grand site $(\Sch/S)_\tau$}
est le faisceau d'anneaux $T/S \mapsto \Gamma(T, \mathcal{O}_T)$, noté
$\mathcal{O}$ ou $\mathcal{O}_S$.
\item Si $\tau = Zariski$ ou $\tau = \etale$, le
{\it faisceau structural du petit site} $S_{Zar}$ ou $S_\etale$
est le faisceau d'anneaux $T/S \mapsto \Gamma(T, \mathcal{O}_T)$,
noté $\mathcal{O}$ ou $\mathcal{O}_S$.
\item Le {\it faisceau de $\mathcal{O}$-modules associé à
$\mathcal{F}$} sur le grand site $(\Sch/S)_\tau$
est le faisceau de $\mathcal{O}$-modules
$(f : T \to S) \mapsto \Gamma(T, f^*\mathcal{F})$,
noté $\mathcal{F}^a$ (et souvent simplement $\mathcal{F}$).
\item Si $\tau = Zariski$ ou $\tau = \etale$, le
{\it faisceau de $\mathcal{O}$-modules associé à $\mathcal{F}$}
sur le petit site $S_{Zar}$ ou $S_\etale$ est le faisceau de
$\mathcal{O}$-modules $(f : T \to S) \mapsto \Gamma(T, f^*\mathcal{F})$,
noté $\mathcal{F}^a$ (et souvent simplement $\mathcal{F}$).
\end{enumerate}
\end{definition}

\noindent
Nous employons donc la même notation $\mathcal{F}^a$ dans chaque cas.
Cela ne crée pas de confusion, puisque, par définition, la règle
qui définit le faisceau $\mathcal{F}^a$ est indépendante du site
considéré.

\begin{remark}
\label{remark-Zariski-site-space}
Dans Topologies, lemme \ref{topologies-lemma-Zariski-usual},
nous avons vu que le petit site de Zariski d'un schéma $S$ est
équivalent à $S$ vu comme espace topologique, au sens où leurs
catégories de faisceaux sont naturellement équivalentes. Maintenant que $S_{Zar}$
est aussi muni d'un faisceau structural $\mathcal{O}$, les
faisceaux de modules sur le site annelé $(S_{Zar}, \mathcal{O})$
s'identifient aux faisceaux de modules sur l'espace annelé $(S, \mathcal{O}_S)$.
\end{remark}

\begin{remark}
\label{remark-change-topologies-ringed}
Soit $f : T \to S$ un morphisme de schémas.
Chacun des morphismes de sites $f_{sites}$ énumérés dans
Topologies, section \ref{topologies-section-change-topologies},
devient un morphisme de sites annelés. En effet, chacun de ces morphismes de sites
$f_{sites} : (\Sch/T)_\tau \to (\Sch/S)_{\tau'}$, ou
$f_{sites} : (\Sch/S)_\tau \to S_{\tau'}$, est défini par le foncteur
continu $S'/S \mapsto T \times_S S'/S$. Ainsi, pour $S'/S$ donné, on prend pour
$$
f_{sites}^\sharp :
\mathcal{O}(S'/S)
\longrightarrow
f_{sites, *}\mathcal{O}(S'/S) =
\mathcal{O}(T \times_S S'/T)
$$
l'homomorphisme usuel
$\text{pr}_{S'}^\sharp : \mathcal{O}(S') \to \mathcal{O}(T \times_S S')$.
De même, le morphisme
$i_f : \Sh(T_\tau) \to \Sh((\Sch/S)_\tau)$,
pour $\tau \in \{Zar, \etale\}$, voir
Topologies, lemmes \ref{topologies-lemma-put-in-T} et
\ref{topologies-lemma-put-in-T-etale},
devient un morphisme de topos annelés, puisque $i_f^{-1}\mathcal{O} = \mathcal{O}$.
Voici quelques cas particuliers :
\begin{enumerate}
\item Le morphisme de grands sites
$f_{big} : (\Sch/X)_{fppf} \to (\Sch/Y)_{fppf}$
devient un morphisme de sites annelés
$$
(f_{big}, f_{big}^\sharp) :
((\Sch/X)_{fppf}, \mathcal{O}_X)
\longrightarrow
((\Sch/Y)_{fppf}, \mathcal{O}_Y)
$$
au sens de Modules sur les sites, définition \ref{sites-modules-definition-ringed-site}.
Il en va de même pour les grands sites syntomique, lisse, \'etale et de Zariski.
\item Le morphisme de petits sites
$f_{small} : X_\etale \to Y_\etale$
devient un morphisme de sites annelés
$$
(f_{small}, f_{small}^\sharp) :
(X_\etale, \mathcal{O}_X)
\longrightarrow
(Y_\etale, \mathcal{O}_Y)
$$
au sens de Modules sur les sites, définition \ref{sites-modules-definition-ringed-site}.
Il en va de même pour le petit site de Zariski.
\end{enumerate}
\end{remark}

\noindent
Soit $S$ un schéma. Étant donné un $\mathcal{O}$-module sur, par exemple,
$(\Sch/S)_{Zar}$, son image inverse sur, par exemple, $(\Sch/S)_{fppf}$
est simplement le faisceau associé pour la topologie fppf. Pour comparer
les grands et les petits sites, on dispose du résultat suivant.

\begin{lemma}
\label{lemma-compare-sites}
Soit $S$ un schéma. Notons
$$
\begin{matrix}
\text{id}_{\tau, Zar} & : & (\Sch/S)_\tau \to S_{Zar}, &
\tau \in \{Zar, \etale, smooth, syntomic, fppf\} \\
\text{id}_{\tau, \etale} & : &
(\Sch/S)_\tau \to S_\etale, &
\tau \in \{\etale, smooth, syntomic, fppf\} \\
\text{id}_{small, \etale, Zar} & : & S_\etale \to S_{Zar},
\end{matrix}
$$
les morphismes de sites annelés de la
remarque \ref{remark-change-topologies-ringed}.
Soit $\mathcal{F}$ un faisceau de $\mathcal{O}_S$-modules,
que l'on considère comme un faisceau de $\mathcal{O}$-modules sur $S_{Zar}$. Alors
\begin{enumerate}
\item $(\text{id}_{\tau, Zar})^*\mathcal{F}$ est le faisceau associé pour la topologie $\tau$
au faisceau de Zariski
$$
(f : T \to S) \longmapsto \Gamma(T, f^*\mathcal{F})
$$
sur $(\Sch/S)_\tau$, et
\item $(\text{id}_{small, \etale, Zar})^*\mathcal{F}$ est le
faisceau associé pour la topologie \'etale au faisceau de Zariski
$$
(f : T \to S) \longmapsto \Gamma(T, f^*\mathcal{F})
$$
sur $S_\etale$.
\end{enumerate}
Soit $\mathcal{G}$ un faisceau de $\mathcal{O}$-modules
sur $S_\etale$. Alors
\begin{enumerate}
\item[(3)] $(\text{id}_{\tau, \etale})^*\mathcal{G}$ est le
faisceau associé pour la topologie $\tau$ au faisceau pour la topologie \'etale
$$
(f : T \to S) \longmapsto \Gamma(T, f_{small}^*\mathcal{G})
$$
où $f_{small} : T_\etale \to S_\etale$
est le morphisme de petits sites annelés pour la topologie \'etale de la
remarque \ref{remark-change-topologies-ringed}.
\end{enumerate}
\end{lemma}

\begin{proof}
Démonstration de (1). Remarquons d'abord que le résultat est vrai pour $\tau = Zar$,
car on dispose alors du morphisme de topos
$i_f : \Sh(T_{Zar}) \to \Sh((\Sch/S)_{Zar})$
tel que $\text{id}_{\tau, Zar} \circ i_f = f_{small}$ comme morphismes
$T_{Zar} \to S_{Zar}$, voir
Topologies, lemmes \ref{topologies-lemma-put-in-T} et
\ref{topologies-lemma-morphism-big-small}.
Par transitivité de l'image inverse (voir
Modules sur les sites,
lemme \ref{sites-modules-lemma-push-pull-composition-modules}),
on obtient
$i_f^*(\text{id}_{\tau, Zar})^*\mathcal{F} = f_{small}^*\mathcal{F}$,
comme souhaité. D'après la remarque précédant ce lemme,
$(\text{id}_{\tau, Zar})^*\mathcal{F}$ est donc le faisceau associé pour la topologie $\tau$
au préfaisceau $T \mapsto \Gamma(T, f^*\mathcal{F})$.

\medskip\noindent
La démonstration de (3) est identique à celle de (1), à ceci près qu'elle
utilise
Topologies, lemmes \ref{topologies-lemma-put-in-T-etale} et
\ref{topologies-lemma-morphism-big-small-etale}.
Nous omettons la démonstration de (2).
\end{proof}

\begin{remark}
\label{remark-change-topologies-ringed-sites}
La remarque \ref{remark-change-topologies-ringed}
et le
lemme \ref{lemma-compare-sites}
ont les applications suivantes :
\begin{enumerate}
\item Soit $S$ un schéma.
La construction $\mathcal{F} \mapsto \mathcal{F}^a$
est l'image inverse par le morphisme de sites annelés
$\text{id}_{\tau, Zar} : ((\Sch/S)_\tau, \mathcal{O})
\to (S_{Zar}, \mathcal{O})$
ou par le morphisme
$\text{id}_{small, \etale, Zar} :
(S_\etale, \mathcal{O}) \to (S_{Zar}, \mathcal{O})$.
\item Soit $f : X \to Y$ un morphisme de schémas.
Pour chacun des morphismes $f_{sites}$ de sites annelés de la
remarque \ref{remark-change-topologies-ringed},
on a
$$
(f^*\mathcal{F})^a = f_{sites}^*\mathcal{F}^a.
$$
Cela résulte de (1) et de la compatibilité de l'image inverse avec
la composition des morphismes de sites annelés, voir
Modules sur les sites,
lemme \ref{sites-modules-lemma-push-pull-composition-modules}.
\end{enumerate}
\end{remark}

\begin{lemma}
\label{lemma-quasi-coherent-gives-quasi-coherent}
Soit $S$ un schéma.
Soit $\mathcal{F}$ un $\mathcal{O}_S$-module quasi-cohérent.
Soit $\tau \in \{Zariski, \linebreak[0] \etale, \linebreak[0]
smooth, \linebreak[0] syntomic, \linebreak[0] fppf\}$.
\begin{enumerate}
\item Le faisceau $\mathcal{F}^a$ est un $\mathcal{O}$-module
quasi-cohérent sur $(\Sch/S)_\tau$, au sens de
Modules sur les sites, définition \ref{sites-modules-definition-site-local}.
\item Si $\tau = Zariski$ ou $\tau = \etale$, alors le faisceau
$\mathcal{F}^a$ est un $\mathcal{O}$-module quasi-cohérent sur
$S_{Zar}$ ou $S_\etale$, au sens de
Modules sur les sites, définition \ref{sites-modules-definition-site-local}.
\end{enumerate}
\end{lemma}

\begin{proof}
Soit $\{S_i \to S\}$ un recouvrement de Zariski tel que l'on ait des
suites exactes
$$
\bigoplus\nolimits_{k \in K_i} \mathcal{O}_{S_i} \longrightarrow
\bigoplus\nolimits_{j \in J_i} \mathcal{O}_{S_i} \longrightarrow
\mathcal{F}|_{S_i} \longrightarrow 0
$$
pour certains ensembles d'indices $K_i$ et $J_i$. Cela est possible par
la définition d'un faisceau quasi-cohérent sur un espace annelé
(voir Modules, définition \ref{modules-definition-quasi-coherent}).

\medskip\noindent
Preuve de (1). Soit $\tau \in \{Zariski, \linebreak[0] fppf, \linebreak[0]
\etale, \linebreak[0] smooth, \linebreak[0] syntomic\}$.
Il est clair que $\mathcal{F}^a|_{(\Sch/S_i)_\tau}$ s'insère aussi dans
une suite exacte
$$
\bigoplus\nolimits_{k \in K_i} \mathcal{O}|_{(\Sch/S_i)_\tau}
\longrightarrow
\bigoplus\nolimits_{j \in J_i} \mathcal{O}|_{(\Sch/S_i)_\tau}
\longrightarrow
\mathcal{F}^a|_{(\Sch/S_i)_\tau} \longrightarrow 0
$$
Ainsi $\mathcal{F}^a$ est quasi-cohérent par Modules sur les sites,
lemme \ref{sites-modules-lemma-local-final-object}.

\medskip\noindent
Preuve de (2). Soit $\tau = \etale$. Il est clair que
$\mathcal{F}^a|_{(S_i)_\etale}$ s'insère aussi dans une suite exacte
$$
\bigoplus\nolimits_{k \in K_i} \mathcal{O}|_{(S_i)_\etale}
\longrightarrow
\bigoplus\nolimits_{j \in J_i} \mathcal{O}|_{(S_i)_\etale}
\longrightarrow
\mathcal{F}^a|_{(S_i)_\etale} \longrightarrow 0
$$
Ainsi $\mathcal{F}^a$ est quasi-cohérent par Modules sur les sites,
lemme \ref{sites-modules-lemma-local-final-object}.
Le cas $\tau = Zariski$ est analogue (il est même tautologique, puisque
les topos annelés correspondants coïncident).
\end{proof}

\begin{lemma}
\label{lemma-fully-faithful-associated}
Soit $S$ un schéma.
Soit $\tau \in \{Zariski, \linebreak[0] \etale, \linebreak[0]
smooth, \linebreak[0] syntomic, \linebreak[0] fppf\}$.
Chacun des foncteurs $\mathcal{F} \mapsto \mathcal{F}^a$ de la
définition \ref{definition-structure-sheaf}
$$
\QCoh(\mathcal{O}_S) \to \QCoh((\Sch/S)_\tau, \mathcal{O})
\quad\text{ou}\quad
\QCoh(\mathcal{O}_S) \to \QCoh(S_\tau, \mathcal{O})
$$
est pleinement fidèle.
\end{lemma}

\begin{proof}
Le lemme \ref{lemma-quasi-coherent-gives-quasi-coherent} montre que l'on
obtient bien les foncteurs indiqués. On identifie les
$\mathcal{O}_S$-modules sur $S$ aux modules sur
$(S_{Zar}, \mathcal{O}_S)$. Le foncteur
$\mathcal{F} \mapsto \mathcal{F}^a$, défini sur les modules
quasi-cohérents $\mathcal{F}$, est l'image inverse par un morphisme $f$
de sites annelés; voir la remarque
\ref{remark-change-topologies-ringed-sites}. Dans chaque cas, le foncteur
$f_*$ est la restriction suivant le foncteur d'inclusion
$S_{Zar} \to S_\tau$ ou $S_{Zar} \to (\Sch/S)_\tau$ (voir, dans
Topologies, section \ref{topologies-section-change-topologies}, la
construction de ces morphismes de sites). En combinant cela avec la
description de $f^*\mathcal{F} = \mathcal{F}^a$, on obtient
$f_*f^*\mathcal{F} = \mathcal{F}$ lorsque $\mathcal{F}$ est
quasi-cohérent. Par conséquent,
$$
\Hom_\mathcal{O}(\mathcal{F}^a, \mathcal{G}^a) =
\Hom_\mathcal{O}(f^*\mathcal{F}, f^*\mathcal{G}) =
\Hom_{\mathcal{O}_S}(\mathcal{F}, f_*f^*\mathcal{G}) =
\Hom_{\mathcal{O}_S}(\mathcal{F}, \mathcal{G})
$$
ce qu'il fallait démontrer.
\end{proof}

\begin{proposition}
\label{proposition-equivalence-quasi-coherent}
Soit $S$ un schéma.
Soit $\tau \in \{Zariski, \linebreak[0] \etale, \linebreak[0]
smooth, \linebreak[0] syntomic, \linebreak[0] fppf\}$.
\begin{enumerate}
\item Le foncteur $\mathcal{F} \mapsto \mathcal{F}^a$ définit une
équivalence de catégories
$$
\QCoh(\mathcal{O}_S)
\longrightarrow
\QCoh((\Sch/S)_\tau, \mathcal{O})
$$
entre la catégorie des faisceaux quasi-cohérents sur $S$ et celle des
$\mathcal{O}$-modules quasi-cohérents sur le grand site de $S$ pour la
topologie $\tau$.
\item Si $\tau = Zariski$ ou $\tau = \etale$, le foncteur
$\mathcal{F} \mapsto \mathcal{F}^a$ définit une équivalence de catégories
$$
\QCoh(\mathcal{O}_S)
\longrightarrow
\QCoh(S_\tau, \mathcal{O})
$$
entre la catégorie des faisceaux quasi-cohérents sur $S$ et celle des
$\mathcal{O}$-modules quasi-cohérents sur le petit site de $S$ pour la
topologie $\tau$.
\end{enumerate}
\end{proposition}

\begin{proof}
Le lemme \ref{lemma-quasi-coherent-gives-quasi-coherent} montre que le
foncteur est bien défini. Il est pleinement fidèle par le lemme
\ref{lemma-fully-faithful-associated}. Pour achever la preuve, nous allons
montrer qu'un $\mathcal{O}$-module quasi-cohérent sur $(\Sch/S)_\tau$
fournit une donnée de descente pour des faisceaux quasi-cohérents
relativement à un recouvrement de $S$ pour la topologie $\tau$. Une fois
cette donnée
construite, la proposition \ref{proposition-fpqc-descent-quasi-coherent}
donnera le faisceau quasi-cohérent correspondant sur $S$.

\medskip\noindent
Soit $\mathcal{G}$ un $\mathcal{O}$-module quasi-cohérent sur le grand
site de $S$ pour la topologie $\tau$. D'après Modules sur les sites,
définition \ref{sites-modules-definition-site-local}, il existe un
recouvrement $\{S_i \to S\}_{i \in I}$ de $S$ pour la topologie $\tau$,
tel que chacune des restrictions
$\mathcal{G}|_{(\Sch/S_i)_\tau}$ admette une présentation globale
$$
\bigoplus\nolimits_{k \in K_i} \mathcal{O}|_{(\Sch/S_i)_\tau}
\longrightarrow
\bigoplus\nolimits_{j \in J_i} \mathcal{O}|_{(\Sch/S_i)_\tau}
\longrightarrow
\mathcal{G}|_{(\Sch/S_i)_\tau} \longrightarrow 0
$$
pour certains ensembles d'indices $J_i$ et $K_i$. Nous affirmons qu'il
en résulte que $\mathcal{G}|_{(\Sch/S_i)_\tau}$ est de la forme
$\mathcal{F}_i^a$ pour un faisceau quasi-cohérent $\mathcal{F}_i$ sur
$S_i$. Cela est clair
pour les sommes directes
$\bigoplus\nolimits_{k \in K_i} \mathcal{O}|_{(\Sch/S_i)_\tau}$ et
$\bigoplus\nolimits_{j \in J_i} \mathcal{O}|_{(\Sch/S_i)_\tau}$.
Ainsi $\mathcal{G}|_{(\Sch/S_i)_\tau}$ est le conoyau d'un morphisme
$\varphi : \mathcal{K}_i^a \to \mathcal{L}_i^a$, où
$\mathcal{K}_i$ et $\mathcal{L}_i$ sont des faisceaux quasi-cohérents
sur $S_i$. Par la pleine fidélité de $(\ )^a$, on a
$\varphi = \phi^a$ pour un morphisme de faisceaux quasi-cohérents
$\phi : \mathcal{K}_i \to \mathcal{L}_i$ sur $S_i$. On a donc bien
$\mathcal{G}|_{(\Sch/S_i)_\tau} \cong \Coker(\phi)^a$.

\medskip\noindent
Comme $\mathcal{G}$ est défini sur toute la catégorie $(\Sch/S)_\tau$,
on a
$$
(\text{pr}_0^*\mathcal{F}_i)^a
\cong
\mathcal{G}|_{(\Sch/(S_i \times_S S_j))_\tau}
\cong
(\text{pr}_1^*\mathcal{F})^a
$$
comme $\mathcal{O}$-modules sur $(\Sch/(S_i \times_S S_j))_\tau$.
La pleine fidélité donne alors des isomorphismes canoniques
$$
\phi_{ij} :
\text{pr}_0^*\mathcal{F}_i
\longrightarrow
\text{pr}_1^*\mathcal{F}_j
$$
de modules quasi-cohérents sur $S_i \times_S S_j$. Nous omettons la
vérification de la condition de cocycle. Cette condition étant satisfaite,
l'effectivité de la descente des faisceaux quasi-cohérents pour le
recouvrement $\{S_i \to S\}$ (proposition
\ref{proposition-fpqc-descent-quasi-coherent}) fournit un faisceau
quasi-cohérent $\mathcal{F}$ sur $S$, muni d'isomorphismes
$\mathcal{F}|_{S_i} \cong \mathcal{F}_i$ compatibles avec la donnée de
descente. Autrement dit, on dispose d'isomorphismes de
$\mathcal{O}$-modules
$$
\phi_i :
\mathcal{F}^a|_{(\Sch/S_i)_\tau}
\longrightarrow
\mathcal{G}|_{(\Sch/S_i)_\tau}
$$
qui coïncident sur $S_i \times_S S_j$. Comme
$\SheafHom_\mathcal{O}(\mathcal{F}^a, \mathcal{G})$ est un faisceau
(Modules sur les sites, lemme \ref{sites-modules-lemma-internal-hom}),
ces isomorphismes se recollent en un morphisme de $\mathcal{O}$-modules
$\mathcal{F}^a \to \mathcal{G}$ qui redonne les isomorphismes $\phi_i$
ci-dessus. Ce morphisme est un isomorphisme, ce qui
achève la preuve.

\medskip\noindent
Le cas des sites $S_\etale$ et $S_{Zar}$ se démontre exactement de la
même manière.
\end{proof}

\begin{lemma}
\label{lemma-equivalence-quasi-coherent-properties}
Soit $S$ un schéma.
Soit $\tau \in \{Zariski, \linebreak[0] \etale, \linebreak[0]
smooth, \linebreak[0] syntomic, \linebreak[0] fppf\}$.
Soit $\mathcal{P}$ l'une des propriétés de modules\footnote{La liste est
la suivante : libre, libre de rang fini, engendré par ses sections globales,
engendré par $r$ sections globales, engendré par un nombre fini de sections
globales, admettant une présentation globale, admettant une présentation
globale finie, localement libre, localement libre de rang fini, localement
engendré par des sections, localement engendré par $r$ sections, de type
fini, de présentation finie, cohérent ou plat.} définies dans Modules sur
les sites, définitions \ref{sites-modules-definition-global},
\ref{sites-modules-definition-site-local} et
\ref{sites-modules-definition-flat}. Les équivalences de catégories
$$
\QCoh(\mathcal{O}_S)
\longrightarrow
\QCoh((\Sch/S)_\tau, \mathcal{O})
\quad\text{et}\quad
\QCoh(\mathcal{O}_S)
\longrightarrow
\QCoh(S_\tau, \mathcal{O})
$$
définies par $\mathcal{F} \mapsto \mathcal{F}^a$ dans la proposition
\ref{proposition-equivalence-quasi-coherent} vérifient
$$
\mathcal{F}\text{ vérifie }\mathcal{P}
\Leftrightarrow
\mathcal{F}^a\text{ vérifie }\mathcal{P}\text{ comme }\mathcal{O}\text{-module}
$$
sauf, peut-être, lorsque $\mathcal{P}$ est ``localement libre'' ou
``cohérent''. Si $\mathcal{P}=$``cohérent'',
l'équivalence vaut pour
$\QCoh(\mathcal{O}_S) \to \QCoh(S_\tau, \mathcal{O})$ lorsque $S$ est
localement noethérien et que $\tau$ est Zariski ou \'etale.
\end{lemma}

\begin{proof}
C'est immédiat pour les propriétés globales, c'est-à-dire celles de
Modules sur les sites, définition \ref{sites-modules-definition-global}.
Pour les propriétés locales, le lemme
\ref{sites-modules-lemma-local-final-object} de Modules sur les sites
permet de traduire ``$\mathcal{F}^a$ vérifie $\mathcal{P}$'' en une
propriété portant sur les membres d'un recouvrement de $X$. Le résultat
découle alors des lemmes \ref{lemma-finite-type-descends},
\ref{lemma-finite-presentation-descends},
\ref{lemma-locally-generated-by-r-sections-descends},
\ref{lemma-flat-descends} et
\ref{lemma-finite-locally-free-descends}. Pour un module quasi-cohérent,
être cohérent revient à être de type fini sur un schéma localement
noethérien (voir Cohomologie des schémas, lemme
\ref{coherent-lemma-coherent-Noetherian}); on se ramène donc au cas des
modules de type fini (détails omis).
\end{proof}










\section{Cohomologie des modules quasi-cohérents et topologies}
\label{section-quasi-coherent-cohomology}

\noindent
Dans cette section, nous montrons que la cohomologie des modules
quasi-cohérents est indépendante du choix de la topologie.

\begin{lemma}
\label{lemma-standard-covering-Cech}
Soit $S$ un schéma. Supposons que
\begin{enumerate}
\item[(a)] $\tau \in \{Zariski, \linebreak[0] fppf, \linebreak[0]
\etale, \linebreak[0] smooth, \linebreak[0] syntomic\}$ et
$\mathcal{C} = (\Sch/S)_\tau$, ou que
\item[(b)] $\tau = \etale$ et $\mathcal{C} = S_\etale$, ou que
\item[(c)] $\tau = Zariski$ et $\mathcal{C} = S_{Zar}$.
\end{enumerate}
Soit $\mathcal{F}$ un faisceau abélien sur $\mathcal{C}$.
Soit $U \in \Ob(\mathcal{C})$ affine.
Soit $\mathcal{U} = \{U_i \to U\}_{i = 1, \ldots, n}$ un recouvrement
affine standard pour la topologie $\tau$ dans $\mathcal{C}$. Alors
\begin{enumerate}
\item $\mathcal{V} = \{\coprod_{i = 1, \ldots, n} U_i \to U\}$ est un
recouvrement de $U$ pour la topologie $\tau$,
\item $\mathcal{U}$ est un raffinement de $\mathcal{V}$, et
\item l'application induite entre les complexes de {\v C}ech
(Cohomologie sur les sites,
équation (\ref{sites-cohomology-equation-map-cech-complexes}))
$$
\check{\mathcal{C}}^\bullet(\mathcal{V}, \mathcal{F})
\longrightarrow
\check{\mathcal{C}}^\bullet(\mathcal{U}, \mathcal{F})
$$
est un isomorphisme de complexes.
\end{enumerate}
\end{lemma}

\begin{proof}
Cela résulte de l'égalité
$$
(\coprod\nolimits_{i_0 = 1, \ldots, n} U_{i_0}) \times_U
\ldots \times_U
(\coprod\nolimits_{i_p = 1, \ldots, n} U_{i_p})
=
\coprod\nolimits_{i_0, \ldots, i_p \in \{1, \ldots, n\}}
U_{i_0} \times_U \ldots \times_U U_{i_p}
$$
et du fait que
$\mathcal{F}(\coprod_a V_a) = \prod_a \mathcal{F}(V_a)$, puisque les
réunions disjointes sont des recouvrements pour la topologie $\tau$.
\end{proof}

\begin{lemma}
\label{lemma-standard-covering-Cech-quasi-coherent}
Soit $S$ un schéma. Soit $\mathcal{F}$ un faisceau quasi-cohérent sur
$S$. Soient $\tau$, $\mathcal{C}$, $U$ et $\mathcal{U}$ comme dans le
lemme \ref{lemma-standard-covering-Cech}. Il existe alors un isomorphisme
de complexes
$$
\check{\mathcal{C}}^\bullet(\mathcal{U}, \mathcal{F}^a)
\cong
s((A/R)_\bullet \otimes_R M)
$$
(voir la section \ref{section-descent-modules}), où
$R = \Gamma(U, \mathcal{O}_U)$, $M = \Gamma(U, \mathcal{F}^a)$ et
$R \to A$ est un homomorphisme d'anneaux fidèlement plat. En particulier,
$$
\check{H}^p(\mathcal{U}, \mathcal{F}^a) = 0
$$
pour tout $p \geq 1$.
\end{lemma}

\begin{proof}
Le lemme \ref{lemma-standard-covering-Cech} montre que
$\check{\mathcal{C}}^\bullet(\mathcal{U}, \mathcal{F}^a)$ est isomorphe
à $\check{\mathcal{C}}^\bullet(\mathcal{V}, \mathcal{F}^a)$, où
$\mathcal{V} = \{V \to U\}$ et
$V = \coprod_{i = 1, \ldots n} U_i$ est encore affine. Posons
$A = \Gamma(V, \mathcal{O}_V)$. Puisque $\{V \to U\}$ est un
recouvrement pour la topologie $\tau$, l'homomorphisme $R \to A$ est fidèlement plat.
D'autre part, la définition de $\mathcal{F}^a$ donne pour terme de degré
$p$ du complexe $\check{\mathcal{C}}^p(\mathcal{V}, \mathcal{F}^a)$ :
$$
\Gamma(V \times_U \ldots \times_U V, \mathcal{F}^a)
=
\Gamma(\Spec(A \otimes_R \ldots \otimes_R A), \mathcal{F}^a)
=
A \otimes_R \ldots \otimes_R A \otimes_R M
$$
Nous omettons la vérification que les applications du complexe de
{\v C}ech coïncident avec celles du complexe
$s((A/R)_\bullet \otimes_R M)$. L'annulation de la cohomologie est le
lemme \ref{lemma-ff-exact}.
\end{proof}

\begin{proposition}
\label{proposition-same-cohomology-quasi-coherent}
\begin{slogan}
La cohomologie des faisceaux quasi-cohérents est la même quelle que soit
la topologie choisie.
\end{slogan}
Soit $S$ un schéma. Soit $\mathcal{F}$ un faisceau quasi-cohérent sur
$S$. Soit $\tau \in \{Zariski, \linebreak[0] \etale, \linebreak[0]
smooth, \linebreak[0] syntomic, \linebreak[0] fppf\}$.
\begin{enumerate}
\item Il existe un isomorphisme canonique
$$
H^q(S, \mathcal{F}) = H^q((\Sch/S)_\tau, \mathcal{F}^a).
$$
\item Il existe des isomorphismes canoniques
$$
H^q(S, \mathcal{F}) =
H^q(S_{Zar}, \mathcal{F}^a) =
H^q(S_\etale, \mathcal{F}^a).
$$
\end{enumerate}
\end{proposition}

\begin{proof}
Pour $q = 0$, le résultat est immédiat d'après la définition de
$\mathcal{F}^a$. Soit $\mathcal{C} = (\Sch/S)_\tau$, ou
$\mathcal{C} = S_\etale$, ou $\mathcal{C} = S_{Zar}$.

\medskip\noindent
Nous allons appliquer le lemme
\ref{sites-cohomology-lemma-cech-vanish-collection} de Cohomologie sur les
sites avec $\mathcal{F} = \mathcal{F}^a$, en prenant pour
$\mathcal{B} \subset \Ob(\mathcal{C})$ l'ensemble des schémas affines de
$\mathcal{C}$, et pour $\text{Cov} \subset \text{Cov}_\mathcal{C}$
l'ensemble des recouvrements affines standards pour la topologie $\tau$.
L'hypothèse (3) de ce lemme est satisfaite par le lemme
\ref{lemma-standard-covering-Cech-quasi-coherent}. On en conclut que
$H^p(U, \mathcal{F}^a) = 0$ pour tout objet affine $U$ de $\mathcal{C}$.

\medskip\noindent
Soit maintenant $U \in \Ob(\mathcal{C})$ un objet séparé. Notons
$f : U \to S$ le morphisme structural. Soit $U = \bigcup U_i$ un
recouvrement ouvert affine. On peut aussi le considérer comme le
recouvrement $\mathcal{U} = \{U_i \to U\}$ de $U$ pour la topologie $\tau$ dans
$\mathcal{C}$. Remarquons que
$U_{i_0} \times_U \ldots \times_U U_{i_p} =
U_{i_0} \cap \ldots \cap U_{i_p}$ est affine, puisque $U$ est supposé
séparé. D'après Cohomologie sur les sites, lemme
\ref{sites-cohomology-lemma-cech-spectral-sequence-application}, et le
résultat précédent, on a
$$
H^p(U, \mathcal{F}^a) = \check{H}^p(\mathcal{U}, \mathcal{F}^a)
= H^p(U, f^*\mathcal{F})
$$
la dernière égalité provenant de Cohomologie des schémas, lemme
\ref{coherent-lemma-cech-cohomology-quasi-coherent}. En particulier, si
$S$ est séparé, on peut prendre $U = S$ et $f = \text{id}_S$, ce qui
démontre la proposition. Nous suggérons au lecteur de passer le reste de
la preuve (ou de le réécrire afin d'en rendre l'exposé plus clair).

\medskip\noindent
Choisissons une résolution injective
$\mathcal{F} \to \mathcal{I}^\bullet$ sur $S$. Choisissons une
résolution injective $\mathcal{F}^a \to \mathcal{J}^\bullet$
sur $\mathcal{C}$.
Notons $\mathcal{J}^n|_S$ la restriction de $\mathcal{J}^n$ aux ouverts
de $S$; c'est un faisceau sur l'espace topologique $S$, car les
recouvrements ouverts sont des recouvrements pour la topologie $\tau$. On obtient un
complexe
$$
0 \to \mathcal{F} \to \mathcal{J}^0|_S \to \mathcal{J}^1|_S \to \ldots
$$
qui est exact, puisque ses sections sur tout ouvert affine $U \subset S$
forment une suite exacte (par l'annulation de
$H^p(U, \mathcal{F}^a)$ pour $p > 0$ obtenue ci-dessus). Par Catégories
dérivées, lemme \ref{derived-lemma-morphisms-lift}, il existe donc un
morphisme de complexes
$\mathcal{J}^\bullet|_S \to \mathcal{I}^\bullet$ qui induit notamment
un morphisme
$$
R\Gamma(\mathcal{C}, \mathcal{F}^a)
=
\Gamma(S, \mathcal{J}^\bullet)
\longrightarrow
\Gamma(S, \mathcal{I}^\bullet)
=
R\Gamma(S, \mathcal{F}).
$$
En cohomologie, on obtient le morphisme
$H^n(\mathcal{C}, \mathcal{F}^a) \to H^n(S, \mathcal{F})$, dont il faut
montrer qu'elle est un isomorphisme. Soit
$\mathcal{U} : S = \bigcup U_i$ un recouvrement ouvert affine, que l'on
peut aussi voir comme un recouvrement pour la topologie $\tau$. Ce qui précède fournit un
morphisme de complexes doubles
$$
\check{\mathcal{C}}^\bullet(\mathcal{U}, \mathcal{J})
=
\check{\mathcal{C}}^\bullet(\mathcal{U}, \mathcal{J}|_S)
\longrightarrow
\check{\mathcal{C}}^\bullet(\mathcal{U}, \mathcal{I}).
$$
Il induit un morphisme de suites spectrales
$$
{}^\tau\! E_2^{p, q} = \check{H}^p(\mathcal{U}, \underline{H}^q(\mathcal{F}^a))
\longrightarrow
E_2^{p, q} = \check{H}^p(\mathcal{U}, \underline{H}^q(\mathcal{F}))
$$
La première suite spectrale converge vers
$H^{p + q}(\mathcal{C}, \mathcal{F})$, et la seconde vers
$H^{p + q}(S, \mathcal{F})$. D'autre part, nous avons vu que les
morphismes induits ${}^\tau\! E_2^{p, q} \to E_2^{p, q}$ sont
bijectives (toutes les intersections sont séparées, car ce sont des
ouverts d'affines). Les morphismes
$H^n(\mathcal{C}, \mathcal{F}^a) \to H^n(S, \mathcal{F})$ sont donc elles
aussi des isomorphismes, ce qui achève la preuve.
\end{proof}

\begin{proposition}
\label{proposition-equivalence-quasi-coherent-functorial}
Soit $f : T \to S$ un morphisme de schémas.
\begin{enumerate}
\item Les équivalences de catégories de la
proposition \ref{proposition-equivalence-quasi-coherent}
sont compatibles avec l'image inverse.
Plus précisément, on a $f^*(\mathcal{G}^a) = (f^*\mathcal{G})^a$
pour tout faisceau quasi-cohérent $\mathcal{G}$ sur $S$.
\item Les équivalences de catégories de la
proposition \ref{proposition-equivalence-quasi-coherent}, partie (1),
ne sont en général {\bf pas} compatibles avec l'image directe.
\item Si $f$ est quasi-compact et quasi-séparé, et si
$\tau \in \{Zariski, \etale\}$, alors $f_*$ et $f_{small, *}$
préservent les faisceaux quasi-cohérents et le diagramme
$$
\xymatrix{
\QCoh(\mathcal{O}_T)
\ar[rr]_{f_*} \ar[d]_{\mathcal{F} \mapsto \mathcal{F}^a} & &
\QCoh(\mathcal{O}_S)
\ar[d]^{\mathcal{G} \mapsto \mathcal{G}^a} \\
\QCoh(T_\tau, \mathcal{O}) \ar[rr]^{f_{small, *}} & &
\QCoh(S_\tau, \mathcal{O})
}
$$
est commutatif, c'est-à-dire que
$f_{small, *}(\mathcal{F}^a) = (f_*\mathcal{F})^a$.
\end{enumerate}
\end{proposition}

\begin{proof}
La partie (1) résulte de la discussion de la
Remarque \ref{remark-change-topologies-ringed-sites}.
La partie (2) constitue simplement un avertissement, que l'on peut expliquer
comme suit. Tout d'abord, l'énoncé ne peut pas être formulé précisément,
car $f_*$ ne transforme pas en général les faisceaux quasi-cohérents en
faisceaux quasi-cohérents. Même lorsque c'est le cas pour $f$ (et pour tout
changement de base de $f$), la compatibilité sur les grands sites signifierait
que la formation de $f_*\mathcal{F}$ commute à tout changement de base, ce
qui n'est pas vrai en général. Un exemple explicite est l'immersion ouverte
quasi-compacte
$j : X = \mathbf{A}^2_k \setminus \{0\} \to \mathbf{A}^2_k = Y$,
où $k$ est un corps. On a $j_*\mathcal{O}_X = \mathcal{O}_Y$,
mais, après le changement de base vers $\Spec(k)$ par le morphisme $0$,
on constate que l'image directe est nulle.

\medskip\noindent
Démontrons (3) dans le cas $\tau = \etale$. Le morphisme $f$, ainsi que tout
changement de base de $f$, envoie les faisceaux quasi-cohérents
en faisceaux quasi-cohérents ; voir Schémas, lemme
\ref{schemes-lemma-push-forward-quasi-coherent}.
L'égalité $f_{small, *}(\mathcal{F}^a) = (f_*\mathcal{F})^a$
signifie que, pour tout morphisme \'etale $g : U \to S$, on a
$\Gamma(U, g^*f_*\mathcal{F}) = \Gamma(U \times_S T, (g')^*\mathcal{F})$,
où $g' : U \times_S T \to T$ est la projection. Cela résulte du
Lemme \ref{coherent-lemma-flat-base-change-cohomology} de Cohomologie des schémas.
\end{proof}

\begin{lemma}
\label{lemma-higher-direct-images-small-etale}
Soit $f : T \to S$ un morphisme quasi-compact et quasi-séparé de schémas.
Soit $\mathcal{F}$ un faisceau quasi-cohérent sur $T$. Pour la topologie
\'etale comme pour la topologie de Zariski, il existe des isomorphismes canoniques
$R^if_{small, *}(\mathcal{F}^a) = (R^if_*\mathcal{F})^a$.
\end{lemma}

\begin{proof}
Nous le démontrons pour la topologie \'etale ; nous omettons la démonstration
pour la topologie de Zariski. D'après Cohomologie des schémas, lemme
\ref{coherent-lemma-quasi-coherence-higher-direct-images},
les faisceaux $R^if_*\mathcal{F}$ sont quasi-cohérents, de sorte que l'énoncé
a un sens. Le faisceau $R^if_{small, *}\mathcal{F}^a$ est le faisceau associé
au préfaisceau
$$
U \longmapsto H^i(U \times_S T, \mathcal{F}^a)
$$
où $g : U \to S$ est un objet de $S_\etale$ ; voir Cohomologie sur les sites,
lemme \ref{sites-cohomology-lemma-higher-direct-images}.
Suivant nos conventions, le membre de droite est la cohomologie \'etale
de la restriction de $\mathcal{F}^a$ au site localisé
$T_\etale/U \times_S T$, lequel est égal à
$(U \times_S T)_\etale$. D'après la
proposition \ref{proposition-same-cohomology-quasi-coherent},
ce préfaisceau est identique au préfaisceau
$$
U \longmapsto
H^i(U \times_S T, (g')^*\mathcal{F}),
$$
où $g' : U \times_S T \to T$ est la projection. Si $U$ est affine,
ce groupe est égal à $H^0(U, R^if'_*(g')^*\mathcal{F})$ ; voir
Cohomologie des schémas, lemme
\ref{coherent-lemma-quasi-coherence-higher-direct-images-application}.
D'après Cohomologie des schémas, lemme
\ref{coherent-lemma-flat-base-change-cohomology},
ce dernier groupe est égal à $H^0(U, g^*R^if_*\mathcal{F})$, qui est la valeur
de $(R^if_*\mathcal{F})^a$ en $U$.
Ainsi, les valeurs des faisceaux de modules
$R^if_{small, *}(\mathcal{F}^a)$ et $(R^if_*\mathcal{F})^a$
sur tout objet affine de $S_\etale$ sont canoniquement isomorphes ;
il s'ensuit que ces faisceaux sont canoniquement isomorphes.
\end{proof}




\section{Faisceaux quasi-cohérents et topologies, II}
\label{section-quasi-coherent-sheaves-bis}

\noindent
Nous poursuivons la comparaison entre les modules quasi-cohérents sur un
schéma $S$ et les modules quasi-cohérents sur chacun des sites associés à
$S$ dans le chapitre sur les topologies.


\begin{lemma}
\label{lemma-compare-etale-zariski-flat}
Dans le lemme \ref{lemma-compare-sites}, le morphisme de sites annelés
$\text{id}_{small, \etale, Zar} : S_\etale \to S_{Zar}$ est plat.
\end{lemma}

\begin{proof}
Notons $\epsilon = \text{id}_{small, \etale, Zar}$, et notons
$\mathcal{O}_\etale$ et $\mathcal{O}_{Zar}$ les faisceaux structuraux sur
$S_\etale$ et $S_{Zar}$. Il faut montrer que $\mathcal{O}_\etale$ est un
$\epsilon^{-1}\mathcal{O}_{Zar}$-module plat. Rappelons qu'un morphisme
\'etale est ouvert ; voir Morphismes, lemme
\ref{morphisms-lemma-etale-open}. Il résulte de la construction de l'image
inverse des faisceaux que $\epsilon^{-1}\mathcal{O}_{Zar}$ est le faisceau
associé au préfaisceau $\mathcal{O}'$ sur $S_\etale$ qui associe à un
morphisme \'etale $f : V \to S$ l'anneau $\mathcal{O}_S(f(V))$. Si $V$ et
$U = f(V) \subset S$ sont tous deux affines, alors $V \to U$ est un
morphisme \'etale entre schémas affines et correspond donc à un homomorphisme
d'anneaux \'etale. Comme tout homomorphisme d'anneaux \'etale est plat,
l'homomorphisme
$\mathcal{O}_S(U) = \mathcal{O}'(V) \to
\mathcal{O}_\etale(V) = \mathcal{O}_V(V)$ est plat.
Enfin, pour tout morphisme \'etale $f : V \to S$, c'est-à-dire pour tout
objet de $S_\etale$, il existe un recouvrement ouvert affine
$V = \bigcup V_i$ tel que $f(V_i)$ soit un ouvert affine de $S$ pour tout
$i$\footnote{En effet, pour $y \in V$, choisissons un ouvert affine
$y \in V' \subset V$ tel que $f(V')$ soit contenu dans un ouvert affine
$U \subset S$. Choisissons ensuite un ouvert affine
$f(y) \in U' \subset f(V')$. Alors $V'' = f^{-1}(U') \subset V'$ est affine,
puisqu'il est égal à $U' \times_U V'$, et $f(V'') = U'$ est également affine.}.
Le résultat découle alors de Modules sur les sites, lemme
\ref{sites-modules-lemma-flatness-sheafification-refined}.
\end{proof}

\begin{lemma}
\label{lemma-equivalence-quasi-coherent-limits}
Soit $S$ un schéma.
Soit $\tau \in \{Zariski, \linebreak[0] \etale,
\linebreak[0] smooth, \linebreak[0] syntomic, \linebreak[0] fppf\}$.
Les foncteurs
$$
\QCoh(\mathcal{O}_S)
\longrightarrow
\textit{Mod}((\Sch/S)_\tau, \mathcal{O})
\quad\text{et}\quad
\QCoh(\mathcal{O}_S)
\longrightarrow
\textit{Mod}(S_\tau, \mathcal{O})
$$
définis par la règle $\mathcal{F} \mapsto \mathcal{F}^a$ de la
proposition \ref{proposition-equivalence-quasi-coherent}
ont les propriétés suivantes :
\begin{enumerate}
\item ils sont pleinement fidèles ;
\item ils commutent aux sommes directes ;
\item ils commutent aux colimites ;
\item ils sont exacts à droite ;
\item ils sont exacts comme foncteurs
$\QCoh(\mathcal{O}_S) \to \textit{Mod}(S_\etale, \mathcal{O})$ ;
\item ils ne sont en général {\bf pas} exacts comme foncteurs
$\QCoh(\mathcal{O}_S) \to
\textit{Mod}((\Sch/S)_\tau, \mathcal{O})$ ;
\item pour deux $\mathcal{O}_S$-modules quasi-cohérents
$\mathcal{F}$, $\mathcal{G}$, on a
$(\mathcal{F} \otimes_{\mathcal{O}_S} \mathcal{G})^a =
\mathcal{F}^a \otimes_\mathcal{O} \mathcal{G}^a$ ;
\item si $\tau = \etale$ ou $\tau = Zariski$, et si deux
$\mathcal{O}_S$-modules quasi-cohérents $\mathcal{F}$, $\mathcal{G}$ sont
donnés, tels que $\mathcal{F}$ soit de présentation finie, on a
$(\SheafHom_{\mathcal{O}_S}(\mathcal{F}, \mathcal{G}))^a =
\SheafHom_\mathcal{O}(\mathcal{F}^a, \mathcal{G}^a)$ dans
$\textit{Mod}(S_\tau, \mathcal{O})$ ;
\item pour deux $\mathcal{O}_S$-modules quasi-cohérents
$\mathcal{F}$, $\mathcal{G}$, on n'a en général {\bf pas}
$(\SheafHom_{\mathcal{O}_S}(\mathcal{F}, \mathcal{G}))^a =
\SheafHom_\mathcal{O}(\mathcal{F}^a, \mathcal{G}^a)$
dans $\textit{Mod}((\Sch/S)_\tau, \mathcal{O})$, même si
$\mathcal{F}$ est de présentation finie ;
\item pour toute suite exacte courte
$0 \to \mathcal{F}_1^a \to \mathcal{E} \to \mathcal{F}_2^a \to 0$
de $\mathcal{O}$-modules, le module $\mathcal{E}$ est
quasi-cohérent\footnote{Attention : cette formulation est trompeuse. Voir la
partie (6).}, c'est-à-dire que $\mathcal{E}$ appartient à l'image essentielle
du foncteur.
\end{enumerate}
\end{lemma}

\begin{proof}
La partie (1) a été établie dans la
proposition \ref{proposition-equivalence-quasi-coherent}.

\medskip\noindent
Nous avons vu dans Schémas, section
\ref{schemes-section-quasi-coherent}, qu'une colimite de faisceaux
quasi-cohérents sur un schéma est encore un faisceau quasi-cohérent.
En outre, dans la remarque
\ref{remark-change-topologies-ringed-sites}, nous avons vu que
$\mathcal{F} \mapsto \mathcal{F}^a$ est le foncteur d'image inverse d'un
morphisme de sites annelés ; il commute donc à toutes les colimites, par
Modules sur les sites, lemme
\ref{sites-modules-lemma-exactness-pushforward-pullback}.
Cela prouve (3), ainsi que son cas particulier (2).

\medskip\noindent
Cela montre aussi que le foncteur est exact à droite, c'est-à-dire qu'il
commute aux colimites finies, ce qui prouve (4).

\medskip\noindent
Le foncteur $\QCoh(\mathcal{O}_S) \to
\textit{Mod}(S_\etale, \mathcal{O})$,
$\mathcal{F} \mapsto \mathcal{F}^a$, est exact à gauche parce qu'un
morphisme \'etale est plat ; voir Morphismes, lemme
\ref{morphisms-lemma-etale-flat}. Cela prouve (5).

\medskip\noindent
Pour établir (6), supposons que $S = \Spec(\mathbf{Z})$.
Alors $2 : \mathcal{O}_S \to \mathcal{O}_S$ est injectif, mais le morphisme
associé de $\mathcal{O}$-modules sur $(\Sch/S)_\tau$ ne l'est pas, car
$2 : \mathbf{F}_2 \to \mathbf{F}_2$ n'est pas injectif et
$\Spec(\mathbf{F}_2)$ est un objet de $(\Sch/S)_\tau$.

\medskip\noindent
La partie (7) est vraie parce que, comme indiqué ci-dessus, le foncteur
$\mathcal{F} \mapsto \mathcal{F}^a$ est le foncteur d'image inverse d'un
morphisme de sites annelés, et qu'un tel foncteur commute aux produits
tensoriels par Modules sur les sites, lemme
\ref{sites-modules-lemma-tensor-product-pullback}.

\medskip\noindent
La partie (8) est évidente si $\tau = Zariski$, car la catégorie des
$\mathcal{O}$-modules sur $S_{Zar}$ est la même que celle des
$\mathcal{O}_S$-modules sur l'espace topologique $S$.
Si $\tau = \etale$, la partie (8) résulte du fait que, comme indiqué
ci-dessus, le foncteur $\mathcal{F} \mapsto \mathcal{F}^a$ est le foncteur
d'image inverse associé au morphisme plat de sites annelés
$(S_\etale, \mathcal{O}) \to (S_{Zar}, \mathcal{O}_S)$ ; voir le
lemme \ref{lemma-compare-etale-zariski-flat}. L'image inverse par un
morphisme plat de sites annelés commute à la formation du Hom interne dont
le premier argument est un module de présentation finie, par Modules sur
les sites, lemme
\ref{sites-modules-lemma-pullback-internal-hom}.

\medskip\noindent
Pour établir (9), supposons que $S = \Spec(\mathbf{Z})$. Posons
$\mathcal{F} = \Coker(2 : \mathcal{O}_S \to \mathcal{O}_S)$
et $\mathcal{G} = \mathcal{O}_S$.
Alors $\mathcal{F}^a = \Coker(2 : \mathcal{O} \to \mathcal{O})$
et $\mathcal{G}^a = \mathcal{O}$. Par conséquent,
$\SheafHom_\mathcal{O}(\mathcal{F}^a, \mathcal{G}^a) = \mathcal{O}[2]$
est le sous-faisceau de $2$-torsion de $\mathcal{O}$, qui n'est pas nul ;
voir la preuve de (6). En revanche, le module
$\SheafHom_{\mathcal{O}_S}(\mathcal{F}, \mathcal{G})$ est nul.

\medskip\noindent
Preuve de (10). Soit
$0 \to \mathcal{F}_1^a \to \mathcal{E} \to \mathcal{F}_2^a \to 0$
une suite exacte courte de $\mathcal{O}$-modules, où $\mathcal{F}_1$ et
$\mathcal{F}_2$ sont quasi-cohérents sur $S$. Considérons sa restriction
$$
0 \to \mathcal{F}_1 \to \mathcal{E}|_{S_{Zar}} \to \mathcal{F}_2
$$
à $S_{Zar}$. D'après la
proposition \ref{proposition-same-cohomology-quasi-coherent}, on a, pour
tout ouvert affine $U \subset S$,
$H^1(U, \mathcal{F}_1^a) = H^1(U, \mathcal{F}_1) = 0$.
La suite ci-dessus est donc également exacte à droite. D'après Schémas,
section \ref{schemes-section-quasi-coherent}, on en déduit que
$\mathcal{F} = \mathcal{E}|_{S_{Zar}}$ est quasi-cohérent.
On obtient ainsi un diagramme commutatif
$$
\xymatrix{
& \mathcal{F}_1^a \ar[r] \ar[d] &
\mathcal{F}^a \ar[r] \ar[d] &
\mathcal{F}_2^a \ar[r] \ar[d] & 0 \\
0 \ar[r] &
\mathcal{F}_1^a \ar[r] &
\mathcal{E} \ar[r] &
\mathcal{F}_2^a \ar[r] & 0
}
$$
Pour achever la preuve, il suffit de montrer que la ligne supérieure est
également exacte à droite. Notons de nouveau $U = \Spec(A) \subset S$ un
ouvert affine de $S$. Nous avons vu plus haut que
$0 \to \mathcal{F}_1(U) \to \mathcal{E}(U) \to \mathcal{F}_2(U) \to 0$
est exacte. Pour tout schéma affine $V/U$, avec $V = \Spec(B)$, le
morphisme $\mathcal{F}_1^a(V) \to \mathcal{E}(V)$ est injectif.
Par définition, on a
$\mathcal{F}_1^a(V) = \mathcal{F}_1(U) \otimes_A B$.
L'injection $\mathcal{F}_1^a(V) \to \mathcal{E}(V)$ se factorise en
$$
\mathcal{F}_1(U) \otimes_A B \to
\mathcal{E}(U) \otimes_A B \to \mathcal{E}(V)
$$
En considérant les $A$-algèbres $B$ de la forme $B = A \oplus M$, on voit
que $\mathcal{F}_1(U) \to \mathcal{E}(U)$ est universellement injectif
(voir Algèbre, Définition
\ref{algebra-definition-universally-injective}). Comme
$\mathcal{E}(U) = \mathcal{F}(U)$, on en déduit que
$\mathcal{F}_1 \to \mathcal{F}$ reste injectif après tout changement de
base, ou, de manière équivalente, que
$\mathcal{F}_1^a \to \mathcal{F}^a$ est injectif.
\end{proof}

\begin{lemma}
\label{lemma-properties-quasi-coherent}
Soit $S$ un schéma. La catégorie $\QCoh(S_\etale, \mathcal{O})$ des modules
quasi-cohérents sur $S_\etale$ possède les propriétés suivantes :
\begin{enumerate}
\item Toute somme directe de faisceaux quasi-cohérents est quasi-cohérente.
\item Toute colimite de faisceaux quasi-cohérents est quasi-cohérente.
\item Le noyau et le conoyau d'un morphisme de faisceaux quasi-cohérents
sont quasi-cohérents.
\item Étant donnée une suite exacte courte de $\mathcal{O}$-modules
$0 \to \mathcal{F}_1 \to \mathcal{F}_2 \to \mathcal{F}_3 \to 0$,
si deux de ces trois modules sont quasi-cohérents, le troisième l'est aussi.
\item Le produit tensoriel de deux $\mathcal{O}$-modules quasi-cohérents
est quasi-cohérent.
\item Étant donnés deux $\mathcal{O}$-modules quasi-cohérents
$\mathcal{F}$, $\mathcal{G}$ tels que $\mathcal{F}$ soit de présentation
finie, le Hom interne
$\SheafHom_\mathcal{O}(\mathcal{F}, \mathcal{G})$
est quasi-cohérent.
\end{enumerate}
\end{lemma}

\begin{proof}
Les faits correspondants sont vrais pour les modules quasi-cohérents sur le
schéma $S$ ; voir Schémas, section
\ref{schemes-section-quasi-coherent}. La preuve consiste à utiliser le
lemme \ref{lemma-equivalence-quasi-coherent-limits} pour les transporter à
$S_\etale$.

\medskip\noindent
Preuve de (1). Soient $\mathcal{F}_i$, $i \in I$, une famille d'objets de
$\QCoh(S_\etale, \mathcal{O})$. Écrivons
$\mathcal{F}_i = \mathcal{G}_i^a$, où les $\mathcal{G}_i$ sont des modules
quasi-cohérents sur $S$. Alors
$\bigoplus \mathcal{F}_i = (\bigoplus \mathcal{G}_i)^a$ par le lemme cité,
ce qui donne la conclusion.

\medskip\noindent
Preuve de (2). Soit
$\mathcal{I} \to \QCoh(S_\etale, \mathcal{O})$,
$i \mapsto \mathcal{F}_i$, un diagramme. Écrivons
$\mathcal{F}_i = \mathcal{G}_i^a$ ; on obtient ainsi un diagramme
$\mathcal{I} \to \QCoh(\mathcal{O}_S)$. Alors
$\colim \mathcal{F}_i = (\colim \mathcal{G}_i)^a$ par le lemme cité,
ce qui donne la conclusion.

\medskip\noindent
Preuve de (3). Soit $a : \mathcal{F} \to \mathcal{F}'$ une flèche de
$\QCoh(S_\etale, \mathcal{O})$. Écrivons $a = b^a$ pour un morphisme
$b : \mathcal{G} \to \mathcal{G}'$ de modules quasi-cohérents sur $S$.
Par le lemme cité, on a $\Ker(a) = \Ker(b)^a$ et
$\Coker(a) = \Coker(b)^a$, d'où le résultat.

\medskip\noindent
Preuve de (4). Cela résulte de (3), sauf dans le cas où l'on sait que
$\mathcal{F}_1$ et $\mathcal{F}_3$ sont quasi-cohérents. Dans ce cas,
écrivons $\mathcal{F}_1 = \mathcal{G}_1^a$ et
$\mathcal{F}_3 = \mathcal{G}_3^a$, avec $\mathcal{G}_i$ quasi-cohérent sur
$S$. On conclut par le lemme
\ref{lemma-equivalence-quasi-coherent-limits}, partie (10).

\medskip\noindent
Preuve de (5). Soient $\mathcal{F}$ et $\mathcal{F}'$ dans
$\QCoh(S_\etale, \mathcal{O})$. Écrivons
$\mathcal{F} = \mathcal{G}^a$ et
$\mathcal{F}' = (\mathcal{G}')^a$, avec $\mathcal{G}$ et $\mathcal{G}'$
quasi-cohérents sur $S$. Par le lemme cité, on a
$\mathcal{F} \otimes_\mathcal{O} \mathcal{F}' =
(\mathcal{G} \otimes_{\mathcal{O}_S} \mathcal{G}')^a$,
d'où la conclusion.

\medskip\noindent
Preuve de (6). Soient $\mathcal{F}$ et $\mathcal{G}$ dans
$\QCoh(S_\etale, \mathcal{O})$, avec $\mathcal{F}$ de présentation finie.
Écrivons $\mathcal{F} = \mathcal{H}^a$ et
$\mathcal{G} = (\mathcal{I})^a$, avec $\mathcal{H}$ et $\mathcal{I}$
quasi-cohérents sur $S$. Par le
lemme \ref{lemma-equivalence-quasi-coherent-properties}, le module
$\mathcal{H}$ est de présentation finie. Par le
lemme \ref{lemma-equivalence-quasi-coherent-limits}, partie (8), on a
$\SheafHom_\mathcal{O}(\mathcal{F}, \mathcal{G}) =
(\SheafHom_{\mathcal{O}_S}(\mathcal{H}, \mathcal{I}))^a$,
d'où la conclusion.
\end{proof}

\begin{lemma}
\label{lemma-properties-quasi-coherent-on-big}
Soit $S$ un schéma.
Soit $\tau \in \{Zariski, \linebreak[0] \etale, \linebreak[0]
smooth, \linebreak[0] syntomic, \linebreak[0] fppf\}$.
La catégorie $\QCoh((\Sch/S)_\tau, \mathcal{O})$ des modules
quasi-cohérents sur $(\Sch/S)_\tau$ possède les propriétés suivantes :
\begin{enumerate}
\item Toute somme directe de faisceaux quasi-cohérents est quasi-cohérente.
\item Toute colimite de faisceaux quasi-cohérents est quasi-cohérente.
\item Le conoyau d'un morphisme de faisceaux quasi-cohérents est
quasi-cohérent.
\item Étant donnée une suite exacte courte de $\mathcal{O}$-modules
$0 \to \mathcal{F}_1 \to \mathcal{F}_2 \to \mathcal{F}_3 \to 0$, si
$\mathcal{F}_1$ et $\mathcal{F}_3$ sont quasi-cohérents, alors
$\mathcal{F}_2$ l'est aussi.
\item Le produit tensoriel de deux $\mathcal{O}$-modules quasi-cohérents
est quasi-cohérent.
\item Étant donnés deux $\mathcal{O}$-modules quasi-cohérents
$\mathcal{F}$, $\mathcal{G}$ tels que $\mathcal{F}$ soit localement libre
de rang fini, le Hom interne
$\SheafHom_\mathcal{O}(\mathcal{F}, \mathcal{G})$
est quasi-cohérent.
\end{enumerate}
\end{lemma}

\begin{proof}
Les faits correspondants sont vrais pour les modules quasi-cohérents sur le
schéma $S$ ; voir Schémas, section
\ref{schemes-section-quasi-coherent}. La preuve consiste à utiliser le
lemme \ref{lemma-equivalence-quasi-coherent-limits} pour les transporter à
$(\Sch/S)_\tau$.

\medskip\noindent
Preuve de (1). Soient $\mathcal{F}_i$, $i \in I$, une famille d'objets de
$\QCoh((\Sch/S)_\tau, \mathcal{O})$. Écrivons
$\mathcal{F}_i = \mathcal{G}_i^a$, où les $\mathcal{G}_i$ sont des modules
quasi-cohérents sur $S$. Alors
$\bigoplus \mathcal{F}_i = (\bigoplus \mathcal{G}_i)^a$ par le lemme cité,
ce qui donne la conclusion.

\medskip\noindent
Preuve de (2). Soit
$\mathcal{I} \to \QCoh((\Sch/S)_\tau, \mathcal{O})$,
$i \mapsto \mathcal{F}_i$, un diagramme. Écrivons
$\mathcal{F}_i = \mathcal{G}_i^a$ ; on obtient ainsi un diagramme
$\mathcal{I} \to \QCoh(\mathcal{O}_S)$. Alors
$\colim \mathcal{F}_i = (\colim \mathcal{G}_i)^a$ par le lemme cité,
ce qui donne la conclusion.

\medskip\noindent
Preuve de (3). Soit $a : \mathcal{F} \to \mathcal{F}'$ une flèche de
$\QCoh((\Sch/S)_\tau, \mathcal{O})$. Écrivons $a = b^a$ pour un morphisme
$b : \mathcal{G} \to \mathcal{G}'$ de modules quasi-cohérents sur $S$.
Par le lemme cité, on a $\Coker(a) = \Coker(b)^a$ (car un conoyau est une
colimite), d'où le résultat.

\medskip\noindent
Preuve de (4). Écrivons $\mathcal{F}_1 = \mathcal{G}_1^a$ et
$\mathcal{F}_3 = \mathcal{G}_3^a$, avec $\mathcal{G}_i$ quasi-cohérent sur
$S$. On conclut par le Lemme
\ref{lemma-equivalence-quasi-coherent-limits}, partie (10).

\medskip\noindent
Preuve de (5). Soient $\mathcal{F}$ et $\mathcal{F}'$ dans
$\QCoh((\Sch/S)_\tau, \mathcal{O})$. Écrivons
$\mathcal{F} = \mathcal{G}^a$ et
$\mathcal{F}' = (\mathcal{G}')^a$, avec $\mathcal{G}$ et $\mathcal{G}'$
quasi-cohérents sur $S$. Par le lemme cité, on a
$\mathcal{F} \otimes_\mathcal{O} \mathcal{F}' =
(\mathcal{G} \otimes_{\mathcal{O}_S} \mathcal{G}')^a$,
d'où la conclusion.

\medskip\noindent
Preuve de (6). Écrivons $\mathcal{F} = \mathcal{H}^a$ pour un
$\mathcal{O}_S$-module quasi-cohérent. D'après le
lemme \ref{lemma-equivalence-quasi-coherent-properties}, le module
$\mathcal{H}$ est localement libre de rang fini. La question est locale sur
$S$ pour la topologie de Zariski (détails omis), de sorte que l'on peut
supposer $\mathcal{H} = \mathcal{O}_S^{\oplus n}$ libre de rang fini.
Alors $\mathcal{F} = \mathcal{O}^{\oplus n}$ et
$\SheafHom_\mathcal{O}(\mathcal{F}, \mathcal{G}) = \mathcal{G}^{\oplus n}$
est quasi-cohérent.
\end{proof}

\begin{example}
\label{example-internal-hom-not-qcoh}
Soit $S$ un schéma. Soient $\mathcal{F}$ et $\mathcal{G}$ des modules
quasi-cohérents sur $(\Sch/S)_\tau$, pour l'une des topologies $\tau$
considérées dans le lemme \ref{lemma-properties-quasi-coherent-on-big}.
En général, $\SheafHom_\mathcal{O}(\mathcal{F}, \mathcal{G})$ n'est pas
quasi-cohérent, même si $\mathcal{F}$ est de présentation finie.
Prenons en effet $S = \Spec(\mathbf{Z})$,
$\mathcal{F} = \Coker(2 : \mathcal{O} \to \mathcal{O})$ et
$\mathcal{G} = \mathcal{O}$. Alors
$\SheafHom_\mathcal{O}(\mathcal{F}, \mathcal{G}) = \mathcal{O}[2]$
est le sous-faisceau de $2$-torsion de $\mathcal{O}$, qui n'est pas
quasi-cohérent.
\end{example}

\begin{lemma}
\label{lemma-qc-colimits}
Soit $S$ un schéma.
\begin{enumerate}
\item La catégorie $\QCoh((\Sch/S)_{fppf}, \mathcal{O})$ admet des
colimites, et celles-ci coïncident avec les colimites dans les catégories
$\textit{Mod}((\Sch/S)_{Zar}, \mathcal{O})$,
$\textit{Mod}((\Sch/S)_\etale, \mathcal{O})$ et
$\textit{Mod}((\Sch/S)_{fppf}, \mathcal{O})$.
\item Étant donnés $\mathcal{F}, \mathcal{G}$ dans
$\QCoh((\Sch/S)_{fppf}, \mathcal{O})$, les produits tensoriels
$\mathcal{F} \otimes_\mathcal{O} \mathcal{G}$ calculés dans
$\textit{Mod}((\Sch/S)_{Zar}, \mathcal{O})$,
$\textit{Mod}((\Sch/S)_\etale, \mathcal{O})$ ou
$\textit{Mod}((\Sch/S)_{fppf}, \mathcal{O})$ coïncident, et leur valeur
commune est un objet de $\QCoh((\Sch/S)_{fppf}, \mathcal{O})$.
\item Étant donnés $\mathcal{F}, \mathcal{G}$ dans
$\QCoh((\Sch/S)_{fppf}, \mathcal{O})$, avec $\mathcal{F}$ localement libre
de rang fini (pour la topologie fppf, ou, de manière équivalente, pour la
topologie \'etale, ou encore pour la topologie de Zariski), les Hom internes
$\SheafHom_\mathcal{O}(\mathcal{F}, \mathcal{G})$ calculés dans
$\textit{Mod}((\Sch/S)_{Zar}, \mathcal{O})$,
$\textit{Mod}((\Sch/S)_\etale, \mathcal{O})$ ou
$\textit{Mod}((\Sch/S)_{fppf}, \mathcal{O})$ coïncident, et leur valeur
commune est un objet de $\QCoh((\Sch/S)_{fppf}, \mathcal{O})$.
\end{enumerate}
\end{lemma}

\begin{proof}
Ce lemme rassemble les résultats précédents sous une forme légèrement
différente. Tout d'abord, le
Lemme \ref{lemma-properties-quasi-coherent-on-big} montre déjà que le
résultat de chacune des constructions (1), (2) et (3) appartient à
$\QCoh((\Sch/S)_\tau, \mathcal{O})$. Il reste, dans chaque cas, à montrer
que ce résultat ne dépend pas de la topologie utilisée. Le point essentiel
est que l'équivalence
$\QCoh(\mathcal{O}_S) \to \QCoh((\Sch/S)_\tau, \mathcal{O})$,
$\mathcal{F} \mapsto \mathcal{F}^a$, de la
proposition \ref{proposition-equivalence-quasi-coherent}, est donnée par la
même formule, indépendamment du choix de la topologie
$\tau \in \{Zariski, \etale, fppf\}$.

\medskip\noindent
Preuve de (1). Soit
$\mathcal{I} \to \QCoh((\Sch/S)_{fppf}, \mathcal{O})$,
$i \mapsto \mathcal{F}_i$, un diagramme. Écrivons
$\mathcal{F}_i = \mathcal{G}_i^a$ ; on obtient ainsi un diagramme
$\mathcal{I} \to \QCoh(\mathcal{O}_S)$. Alors
$\colim \mathcal{F}_i = (\colim \mathcal{G}_i)^a$ dans
$\textit{Mod}((\Sch/S)_\tau, \mathcal{O})$ pour
$\tau \in \{Zariski, \etale, fppf\}$, par le
lemme \ref{lemma-equivalence-quasi-coherent-limits}. Cela prouve (1).

\medskip\noindent
Preuve de (2). Écrivons $\mathcal{F} = \mathcal{H}^a$ et
$\mathcal{G} = (\mathcal{I})^a$, avec $\mathcal{H}$ et $\mathcal{I}$
quasi-cohérents sur $S$. Alors
$\mathcal{F} \otimes_\mathcal{O} \mathcal{G} =
(\mathcal{H} \otimes_\mathcal{O} \mathcal{I})^a$ dans
$\textit{Mod}((\Sch/S)_\tau, \mathcal{O})$ pour
$\tau \in \{Zariski, \etale, fppf\}$, par le
lemme \ref{lemma-equivalence-quasi-coherent-limits}. Cela prouve (2).

\medskip\noindent
Preuve de (3). Soient $\mathcal{F}$ et $\mathcal{G}$ dans
$\QCoh((\Sch/S)_{fppf}, \mathcal{O})$. Écrivons
$\mathcal{F} = \mathcal{H}^a$, avec $\mathcal{H}$ quasi-cohérent sur $S$.
D'après le lemme \ref{lemma-equivalence-quasi-coherent-properties}, on a
\begin{align*}
&\mathcal{F}\text{ est localement libre de rang fini} \\
&\qquad\text{pour la topologie fppf} \\
& \Leftrightarrow
\mathcal{H}\text{ est localement libre de rang fini sur }S \\
& \Leftrightarrow
\mathcal{F}\text{ est localement libre de rang fini} \\
&\qquad\text{pour la topologie \'etale} \\
& \Leftrightarrow
\mathcal{H}\text{ est localement libre de rang fini sur }S \\
& \Leftrightarrow
\mathcal{F}\text{ est localement libre de rang fini} \\
&\qquad\text{pour la topologie de Zariski}
\end{align*}
Cela explique la parenthèse de la partie (3). Si ces conditions équivalentes
sont satisfaites, alors $\mathcal{H}$ est localement libre de rang fini.
La construction de $\SheafHom_\mathcal{O}(\mathcal{F}, \mathcal{G})$ dans
Modules sur les sites, section
\ref{sites-modules-section-internal-hom}, ne dépend que de $\mathcal{F}$ et
$\mathcal{G}$ comme préfaisceaux de modules ; seul le fait que le résultat
$\SheafHom$ soit un faisceau dépend du fait que $\mathcal{F}$ et
$\mathcal{G}$ soient eux-mêmes des faisceaux.
\end{proof}










\section{Modules quasi-cohérents et schémas affines}
\label{section-alternative-quasi-coherent}

\noindent
Soit $S$ un schéma\footnote{Dans cette section, comme dans Topologies,
section \ref{topologies-section-change-topologies}, nous choisissons les sites
$(\Sch/S)_\tau$ de façon qu'ils aient la même catégorie sous-jacente pour
$\tau \in \{Zariski, \etale, smooth, syntomic, fppf\}$. Les sites
$(\textit{Aff}/S)_\tau$ ont alors eux aussi la même catégorie sous-jacente.}.
Soit $\tau \in \{Zariski, \etale, smooth, syntomic, fppf\}$.
Rappelons que $(\textit{Aff}/S)_\tau$ est la sous-catégorie pleine de
$(\Sch/S)_\tau$ formée des objets affines, munie de la structure de site
pour laquelle les recouvrements sont les recouvrements $\tau$ standards.
D'après les lemmes
\ref{topologies-lemma-affine-big-site-Zariski},
\ref{topologies-lemma-affine-big-site-etale},
\ref{topologies-lemma-affine-big-site-smooth},
\ref{topologies-lemma-affine-big-site-syntomic} et
\ref{topologies-lemma-affine-big-site-fppf} de Topologies, il existe une équivalence de
topos
$g : \Sh((\textit{Aff}/S)_\tau) \to \Sh((\Sch/S)_\tau)$
dont le foncteur d'image inverse est donné par restriction.
Rappelons que $\mathcal{O}$ désigne le faisceau structural sur
$(\Sch/S)_\tau$ et notons, provisoirement et avec toute la précision voulue,
$\mathcal{O}_{\textit{Aff}}$ la restriction de $\mathcal{O}$ à
$(\textit{Aff}/S)_\tau$. On obtient alors une équivalence
\begin{equation}
\label{equation-alternative-ringed}
(\Sh((\textit{Aff}/S)_\tau), \mathcal{O}_{\textit{Aff}})
\longrightarrow
(\Sh((\Sch/S)_\tau), \mathcal{O})
\end{equation}
de topos annelés. Cela permet également de comparer les modules
quasi-cohérents.

\begin{lemma}
\label{lemma-quasi-coherent-alternative}
Soit $S$ un schéma. Soit $\mathcal{F}$ un préfaisceau de
$\mathcal{O}_{\textit{Aff}}$-modules sur $(\textit{Aff}/S)_{fppf}$.
Les conditions suivantes sont équivalentes :
\begin{enumerate}
\item pour tout morphisme $U \to U'$ de $(\textit{Aff}/S)_{fppf}$,
l'application
$\mathcal{F}(U') \otimes_{\mathcal{O}(U')} \mathcal{O}(U) \to \mathcal{F}(U)$
est un isomorphisme ;
\item $\mathcal{F}$ est un faisceau sur $(\textit{Aff}/S)_{Zar}$ et un
module quasi-cohérent sur le site annelé
$((\textit{Aff}/S)_{Zar}, \mathcal{O}_{\textit{Aff}})$ au sens de Modules
sur les sites, définition \ref{sites-modules-definition-site-local} ;
\item la même condition que (2) pour la topologie \'etale ;
\item la même condition que (2) pour la topologie lisse ;
\item la même condition que (2) pour la topologie syntomique ;
\item la même condition que (2) pour la topologie fppf ;
\item $\mathcal{F}$ correspond à un module quasi-cohérent sur
$(\Sch/S)_{Zar}$,
$(\Sch/S)_\etale$,
$(\Sch/S)_{smooth}$,
$(\Sch/S)_{syntomic}$ ou
$(\Sch/S)_{fppf}$ par l'équivalence
(\ref{equation-alternative-ringed}) ;
\item $\mathcal{F}$ provient d'un unique $\mathcal{O}_S$-module
quasi-cohérent $\mathcal{G}$ par le procédé décrit à la section
\ref{section-quasi-coherent-sheaves}.
\end{enumerate}
\end{lemma}

\begin{proof}
La notion de module quasi-cohérent étant intrinsèque (Modules sur les sites,
lemme \ref{sites-modules-lemma-special-locally-free}), l'équivalence
(\ref{equation-alternative-ringed}) induit une équivalence entre les
catégories de modules quasi-cohérents. La
proposition \ref{proposition-equivalence-quasi-coherent} affirme que la
catégorie des modules quasi-cohérents sur $\Sch/S$ est indépendante de la
topologie choisie, et que (8) équivaut à (7). Les conditions
(2)--(8) sont donc toutes équivalentes.

\medskip\noindent
Supposons satisfaites les conditions équivalentes (2)--(8), et soit
$\mathcal{G}$ comme en (8). Soit $h : U \to U' \to S$ un morphisme de
$\textit{Aff}/S$. Notons $f : U \to S$ et $f' : U' \to S$ les morphismes
structuraux, de sorte que $f = f' \circ h$. On a
$\mathcal{F}(U') = \Gamma(U', (f')^*\mathcal{G})$ et
$\mathcal{F}(U) = \Gamma(U, f^*\mathcal{G}) = \Gamma(U, h^*(f')^*\mathcal{G})$.
La condition (1) résulte donc du lemme
\ref{schemes-lemma-widetilde-pullback} de Schémas.

\medskip\noindent
Supposons (1). Pour achever la preuve, il suffit d'établir (2).
Soit $U$ un objet de $(\textit{Aff}/S)_{Zar}$. Écrivons
$U = \Spec(R)$. Un recouvrement ouvert standard
$U = U_1 \cup \ldots \cup U_n$ est donné par $U_i = D(f_i)$, où les
éléments $f_1, \ldots, f_n \in R$ engendrent l'idéal unité de $R$.
La propriété (1) donne
$$
\mathcal{F}(U_i) =
\mathcal{F}(U) \otimes_R R_{f_i} =
\mathcal{F}(U)_{f_i}
$$
et
$$
\mathcal{F}(U_i \cap U_j) =
\mathcal{F}(U) \otimes_R R_{f_if_j} =
\mathcal{F}(U)_{f_if_j}
$$
Le lemme \ref{algebra-lemma-cover-module} d'Algèbre montre alors que
$\mathcal{F}$ est un faisceau sur $(\textit{Aff}/S)_{Zar}$. Choisissons une
présentation
$$
\bigoplus\nolimits_{k \in K} R
\longrightarrow
\bigoplus\nolimits_{l \in L} R
\longrightarrow
\mathcal{F}(U)
\longrightarrow 0
$$
par des $R$-modules libres. La propriété (1) et l'exactitude à droite du
produit tensoriel montrent que, pour tout morphisme $U' \to U$ dans
$(\textit{Aff}/S)_{Zar}$, on obtient une présentation
$$
\bigoplus\nolimits_{k \in K} \mathcal{O}_{Aff}(U')
\longrightarrow
\bigoplus\nolimits_{l \in L} \mathcal{O}_{Aff}(U')
\longrightarrow
\mathcal{F}(U')
\longrightarrow 0
$$
Autrement dit, la restriction de $\mathcal{F}$ à la catégorie localisée
$(\textit{Aff}/S)_{Zar}/U$ admet une présentation
$$
\bigoplus\nolimits_{k \in K} \mathcal{O}_{Aff}|_{(\textit{Aff}/S)_{Zar}/U}
\longrightarrow
\bigoplus\nolimits_{l \in L} \mathcal{O}_{Aff}|_{(\textit{Aff}/S)_{Zar}/U}
\longrightarrow
\mathcal{F}|_{(\textit{Aff}/S)_{Zar}/U}
\longrightarrow 0
$$
Malgré cette notation peu élégante, cela achève la preuve.
\end{proof}

\noindent
Poursuivons la discussion commencée dans l'introduction de cette section.
Soit $\tau \in \{Zariski, \etale\}$. Rappelons que $S_{affine, \tau}$ est la
sous-catégorie pleine de $S_\tau$ formée des objets affines, munie de la
structure de site pour laquelle les recouvrements sont les recouvrements
$\tau$ standards. Voir les définitions
\ref{topologies-definition-big-small-Zariski} et
\ref{topologies-definition-big-small-etale} de Topologies. D'après les lemmes
\ref{topologies-lemma-alternative-zariski}, respectivement\  
\ref{topologies-lemma-alternative} de Topologies, il existe une équivalence de topos
$g : \Sh(S_{affine, \tau}) \to \Sh(S_\tau)$, dont le foncteur d'image inverse
est donné par restriction. Rappelons que $\mathcal{O}$ désigne le faisceau
structural sur $S_\tau$ et notons, provisoirement et avec toute la précision
voulue, $\mathcal{O}_{affine}$ la restriction de $\mathcal{O}$ à
$S_{affine, \tau}$. On obtient alors une équivalence
\begin{equation}
\label{equation-alternative-small-ringed}
(\Sh(S_{affine, \tau}), \mathcal{O}_{affine})
\longrightarrow
(\Sh(S_\tau), \mathcal{O})
\end{equation}
de topos annelés. Cela permet également de comparer les modules
quasi-cohérents.

\begin{lemma}
\label{lemma-quasi-coherent-alternative-small}
Soit $S$ un schéma. Soit $\tau \in \{Zariski, \etale\}$.
Soit $\mathcal{F}$ un préfaisceau de $\mathcal{O}_{affine}$-modules sur
$S_{affine, \tau}$. Les conditions suivantes sont équivalentes :
\begin{enumerate}
\item pour tout morphisme $U \to U'$ de $S_{affine, \tau}$, l'application
$\mathcal{F}(U') \otimes_{\mathcal{O}(U')} \mathcal{O}(U) \to \mathcal{F}(U)$
est un isomorphisme ;
\item $\mathcal{F}$ est un faisceau sur $S_{affine, \tau}$ et un module
quasi-cohérent sur le site annelé
$(S_{affine, \tau}, \mathcal{O}_{affine})$ au sens de Modules sur les sites,
définition \ref{sites-modules-definition-site-local} ;
\item $\mathcal{F}$ correspond à un module quasi-cohérent sur $S_\tau$ par
l'équivalence (\ref{equation-alternative-small-ringed}) ;
\item $\mathcal{F}$ provient d'un unique $\mathcal{O}_S$-module
quasi-cohérent $\mathcal{G}$ par le procédé décrit à la section
\ref{section-quasi-coherent-sheaves}.
\end{enumerate}
\end{lemma}

\begin{proof}
Démontrons l'énoncé dans le cas de la topologie \'etale.

\medskip\noindent
Supposons (1). Pour montrer que $\mathcal{F}$ est un faisceau, soit
$\mathcal{U} = \{U_i \to U\}_{i = 1, \ldots, n}$ un recouvrement de
$S_{affine, \etale}$. La condition de faisceau pour $\mathcal{F}$ et
$\mathcal{U}$, compte tenu de notre hypothèse sur $\mathcal{F}$, se ramène à
montrer que
$$
0 \to \mathcal{F}(U) \to
\prod \mathcal{F}(U) \otimes_{\mathcal{O}(U)} \mathcal{O}(U_i) \to
\prod \mathcal{F}(U) \otimes_{\mathcal{O}(U)} \mathcal{O}(U_i \times_U U_j)
$$
est exacte. En effet, $\mathcal{O}(U) \to \prod \mathcal{O}(U_i)$ est
fidèlement plat, par le lemme \ref{lemma-standard-covering-Cech} et parce que
les recouvrements de $S_{affine, \etale}$ sont, par définition, les
recouvrements \'etales standards ; on peut donc appliquer le lemme
\ref{lemma-ff-exact}.
Montrons ensuite que $\mathcal{F}$ est quasi-cohérent sur
$S_{affine, \etale}$. Pour $U$ dans $S_{affine, \etale}$, posons
$R = \mathcal{O}(U)$ et choisissons une présentation
$$
\bigoplus\nolimits_{k \in K} R
\longrightarrow
\bigoplus\nolimits_{l \in L} R
\longrightarrow
\mathcal{F}(U)
\longrightarrow 0
$$
par des $R$-modules libres. La propriété (1) et l'exactitude à droite du
produit tensoriel montrent que, pour tout morphisme $U' \to U$ dans
$S_{affine, \etale}$, on obtient une présentation
$$
\bigoplus\nolimits_{k \in K} \mathcal{O}(U')
\longrightarrow
\bigoplus\nolimits_{l \in L} \mathcal{O}(U')
\longrightarrow
\mathcal{F}(U')
\longrightarrow 0
$$
Autrement dit, la restriction de $\mathcal{F}$ à la catégorie localisée
$S_{affine, \etale}/U$ admet une présentation
$$
\bigoplus\nolimits_{k \in K} \mathcal{O}_{affine}|_{S_{affine, \etale}/U}
\longrightarrow
\bigoplus\nolimits_{l \in L} \mathcal{O}_{affine}|_{S_{affine, \etale}/U}
\longrightarrow
\mathcal{F}|_{S_{affine, \etale}/U}
\longrightarrow 0
$$
comme il le faut pour montrer que $\mathcal{F}$ est quasi-cohérent.
Malgré cette notation peu élégante, nous avons ainsi démontré que (1)
implique (2).

\medskip\noindent
La notion de module quasi-cohérent étant intrinsèque (Modules sur les sites,
lemme \ref{sites-modules-lemma-special-locally-free}), l'équivalence
(\ref{equation-alternative-small-ringed}) induit une équivalence entre les
catégories de modules quasi-cohérents. Les conditions (2) et (3) sont donc
équivalentes.

\medskip\noindent
L'équivalence de (3) et (4) résulte de la
proposition \ref{proposition-equivalence-quasi-coherent}.

\medskip\noindent
Supposons (4) et démontrons (1). Soit $\mathcal{G}$ comme en (4).
Soit $h : U \to U' \to S$ un morphisme de $S_{affine, \etale}$.
Notons $f : U \to S$ et $f' : U' \to S$ les morphismes structuraux, de sorte
que $f = f' \circ h$. On a
$\mathcal{F}(U') = \Gamma(U', (f')^*\mathcal{G})$ et
$\mathcal{F}(U) = \Gamma(U, f^*\mathcal{G}) = \Gamma(U, h^*(f')^*\mathcal{G})$.
La condition (1) résulte donc du lemme
\ref{schemes-lemma-widetilde-pullback} de Schémas.

\medskip\noindent
Nous omettons la démonstration dans le cas de la topologie de Zariski.
\end{proof}





\section{Modules parasites}
\label{section-parasitic}

\noindent
Les modules parasites sont ceux qui deviennent nuls lorsqu'on les restreint
aux schémas plats sur le schéma de base. Voici la définition précise.

\begin{definition}
\label{definition-parasitic}
Soit $S$ un schéma. Soit $\tau \in \{Zar, \etale,
smooth, syntomic, fppf\}$. Soit $\mathcal{F}$ un préfaisceau de
$\mathcal{O}$-modules sur $(\Sch/S)_\tau$.
\begin{enumerate}
\item On dit que $\mathcal{F}$ est
{\it parasite}\footnote{Cette terminologie n'est peut-être pas standard.}
si, pour tout morphisme plat $U \to S$, on a $\mathcal{F}(U) = 0$.
\item On dit que $\mathcal{F}$ est {\it parasite pour la topologie $\tau$}
si, pour tout recouvrement $\tau$,
$\{U_i \to S\}_{i \in I}$, on a $\mathcal{F}(U_i) = 0$ pour tout $i$.
\end{enumerate}
\end{definition}

\noindent
Si $\tau = fppf$, cela signifie que $\mathcal{F}|_{U_{Zar}} = 0$ chaque fois
que $U \to S$ est plat et localement de présentation finie ; les autres cas
sont analogues.

\begin{lemma}
\label{lemma-cohomology-parasitic}
Soit $S$ un schéma. Soit $\tau \in \{Zar, \etale, smooth,
syntomic, fppf\}$. Soit $\mathcal{G}$ un préfaisceau de
$\mathcal{O}$-modules sur $(\Sch/S)_\tau$.
\begin{enumerate}
\item Si $\mathcal{G}$ est parasite pour la topologie $\tau$, alors
$H^p_\tau(U, \mathcal{G}) = 0$ pour tout $U$ ouvert dans $S$,
resp.\ \'etale sur $S$,
resp.\ lisse sur $S$,
resp.\ syntomique sur $S$,
resp.\ plat et localement de présentation finie sur $S$.
\item Si $\mathcal{G}$ est parasite, alors
$H^p_\tau(U, \mathcal{G}) = 0$ pour tout $U$ plat sur $S$.
\end{enumerate}
\end{lemma}

\begin{proof}
Démontrons-le lorsque $\tau = fppf$ ; les autres cas se démontrent exactement
de la même manière. L'hypothèse signifie que $\mathcal{G}(U) = 0$ pour tout
$U \to S$ plat et localement de présentation finie. Appliquons le lemme
\ref{sites-cohomology-lemma-cech-vanish-collection} de Cohomologie sur les
sites au sous-ensemble
$\mathcal{B} \subset \Ob((\Sch/S)_{fppf})$ formé des $U \to S$ plats et
localement de présentation finie, et à la collection $\text{Cov}$ de tous les
recouvrements fppf des éléments de $\mathcal{B}$.
\end{proof}

\begin{lemma}
\label{lemma-direct-image-parasitic}
Soit $f : T \to S$ un morphisme de schémas. Pour tout
$\mathcal{O}$-module parasite sur $(\Sch/T)_\tau$, l'image directe
$f_*\mathcal{F}$ et les images directes supérieures $R^if_*\mathcal{F}$
sont des $\mathcal{O}$-modules parasites sur $(\Sch/S)_\tau$.
\end{lemma}

\begin{proof}
Rappelons que $R^if_*\mathcal{F}$ est le faisceau associé au préfaisceau
$$
U \mapsto H^i((\Sch/U \times_S T)_\tau, \mathcal{F})
$$
voir le lemme \ref{sites-cohomology-lemma-higher-direct-images} de
Cohomologie sur les sites.
Si $U \to S$ est plat, alors $U \times_S T \to T$ est plat par changement de
base. Le groupe affiché est donc nul par le
lemme \ref{lemma-cohomology-parasitic}. Si $\{U_i \to U\}$ est un
recouvrement $\tau$, alors $U_i \times_S T \to T$ est également plat.
Il est par conséquent clair que le faisceau associé au préfaisceau affiché
est nul sur les schémas $U$ plats sur $S$.
\end{proof}

\begin{lemma}
\label{lemma-quasi-coherent-and-flat-base-change}
Soit $S$ un schéma. Soit $\tau \in \{Zar, \etale\}$.
Soit $\mathcal{G}$ un faisceau de $\mathcal{O}$-modules sur
$(\Sch/S)_{fppf}$ tel que
\begin{enumerate}
\item $\mathcal{G}|_{S_\tau}$ est quasi-cohérent ;
\item pour tout morphisme plat et localement de présentation finie
$g : U \to S$, l'application canonique
$g_{\tau, small}^*(\mathcal{G}|_{S_\tau}) \to \mathcal{G}|_{U_\tau}$
est un isomorphisme.
\end{enumerate}
Alors $H^p(U, \mathcal{G}) = H^p(U, \mathcal{G}|_{U_\tau})$ pour tout $U$
plat et localement de présentation finie sur $S$.
\end{lemma}

\begin{proof}
Soit $\mathcal{F}$ l'image inverse de $\mathcal{G}|_{S_\tau}$ sur le grand
site fppf $(\Sch/S)_{fppf}$. Remarquons que $\mathcal{F}$ est quasi-cohérent.
Il existe une application de comparaison canonique
$\varphi : \mathcal{F} \to \mathcal{G}$ qui, en vertu des hypothèses (1) et
(2), induit un isomorphisme
$\mathcal{F}|_{U_\tau} \to \mathcal{G}|_{U_\tau}$ pour tout
$g : U \to S$ plat et localement de présentation finie.
Ainsi, dans les suites exactes courtes
$$
0 \to \Ker(\varphi) \to \mathcal{F} \to \Im(\varphi) \to 0
$$
et
$$
0 \to \Im(\varphi) \to \mathcal{G} \to \Coker(\varphi) \to 0
$$
les faisceaux $\Ker(\varphi)$ et $\Coker(\varphi)$ sont parasites pour la
topologie fppf. Le lemme \ref{lemma-cohomology-parasitic} montre donc que
$H^p(U, \mathcal{F}) \to H^p(U, \mathcal{G})$ est un isomorphisme lorsque
$g : U \to S$ est plat et localement de présentation finie.
Le résultat étant vrai pour $\mathcal{F}$ d'après la
proposition \ref{proposition-same-cohomology-quasi-coherent}, la preuve est
achevée.
\end{proof}













\section[Recouvrements fpqc et épimorphismes]{Les recouvrements fpqc sont des épimorphismes effectifs universels}
\label{section-fpqc-universal-effective-epimorphisms}

\noindent
Nous appliquons les résultats précédents pour démontrer un résultat
intéressant, à savoir le lemme
\ref{lemma-fpqc-universal-effective-epimorphisms}.
D'après Sites, section \ref{sites-section-representable-sheaves},
ce lemme implique que les préfaisceaux représentables sur chacun des sites
$(\Sch/S)_\tau$ sont des faisceaux pour
$\tau \in \{Zariski, fppf, \etale, smooth, syntomic\}$. Commençons
par démontrer un lemme auxiliaire.

\begin{lemma}
\label{lemma-equiv-fibre-product}
Pour un schéma $X$, notons $|X|$ son ensemble sous-jacent.
Soit $f : X \to S$ un morphisme de schémas.
Alors
$$
|X \times_S X| \to |X| \times_{|S|} |X|
$$
est surjectif.
\end{lemma}

\begin{proof}
Cela résulte immédiatement de la description des points du produit
fibré donnée dans Schémas, lemme
\ref{schemes-lemma-points-fibre-product}.
\end{proof}


\begin{lemma}
\label{lemma-universal-effective-epimorphism-affine}
Soit $\{f_i : X_i \to X\}_{i \in I}$ une famille de morphismes de schémas affines.
Les assertions suivantes sont équivalentes :
\begin{enumerate}
\item pour tout $\mathcal{O}_X$-module quasi-cohérent $\mathcal{F}$, on a
$$
\Gamma(X, \mathcal{F}) =
\text{Égalisateur}\left(
\xymatrix{
\prod\nolimits_{i \in I} \Gamma(X_i, f_i^*\mathcal{F})
\ar@<1ex>[r] \ar@<-1ex>[r] &
\prod\nolimits_{i, j \in I}
\Gamma(X_i \times_X X_j, (f_i \times f_j)^*\mathcal{F})
}
\right)
$$
\item $\{f_i : X_i \to X\}_{i \in I}$ est un épimorphisme effectif universel
(Sites, définition \ref{sites-definition-universal-effective-epimorphisms})
dans la catégorie des schémas affines.
\end{enumerate}
\end{lemma}

\begin{proof}
Supposons (2) et soit $\mathcal{F}$ un $\mathcal{O}_X$-module quasi-cohérent.
Considérons le schéma
(Constructions, section \ref{constructions-section-spec})
$$
X' = \underline{\Spec}_X(\mathcal{O}_X \oplus \mathcal{F})
$$
où $\mathcal{O}_X \oplus \mathcal{F}$ est la $\mathcal{O}_X$-algèbre dont
la multiplication est
$(f, s)(f', s') = (ff', fs' + f's)$.
Si $s_i \in \Gamma(X_i, f_i^*\mathcal{F})$ est une section,
alors $s_i$ détermine un unique élément de
$$
\Gamma(X' \times_X X_i, \mathcal{O}_{X' \times_X X_i}) =
\Gamma(X_i, \mathcal{O}_{X_i}) \oplus \Gamma(X_i, f_i^*\mathcal{F})
$$
Nous omettons la démonstration de cette égalité.
Si $(s_i)_{i \in I}$ appartient à l'égalisateur de (1), alors, en utilisant
l'égalité
$$
\Mor(T, \mathbf{A}^1_\mathbf{Z}) = \Gamma(T, \mathcal{O}_T)
$$
valable pour tout schéma $T$, on voit que ces sections définissent une famille
de morphismes $h_i : X' \times_X X_i \to \mathbf{A}^1_\mathbf{Z}$ telle que
$h_i \circ \text{pr}_1 = h_j \circ \text{pr}_2$ comme morphismes
$(X' \times_X X_i) \times_{X'} (X' \times_X X_j) \to \mathbf{A}^1_\mathbf{Z}$.
Puisque nous supposons (2), nous obtenons un morphisme
$h : X' \to \mathbf{A}^1_\mathbf{Z}$ compatible avec les morphismes $h_i$,
qui détermine à son tour un élément
$s \in \Gamma(X, \mathcal{F})$.
Nous omettons la vérification que $s$ s'envoie sur $s_i$ dans
$\Gamma(X_i, f_i^*\mathcal{F})$.

\medskip\noindent
Supposons (1). Soient $T$ un schéma affine et $h_i : X_i \to T$
une famille de morphismes telle que
$h_i \circ \text{pr}_1 = h_j \circ \text{pr}_2$ sur
$X_i \times_X X_j$ pour tous $i, j \in I$. Alors
$$
\prod h_i^\sharp :
\Gamma(T, \mathcal{O}_T)
\to
\prod \Gamma(X_i, \mathcal{O}_{X_i})
$$
prend ses valeurs dans l'égalisateur et l'on obtient un homomorphisme d'anneaux
$\Gamma(T, \mathcal{O}_T) \to \Gamma(X, \mathcal{O}_X)$ en appliquant
l'hypothèse du lemme à $\mathcal{F} = \mathcal{O}_X$.
Cet homomorphisme d'anneaux correspond à un morphisme $h : X \to T$
tel que $h_i = h \circ f_i$. Notre famille est donc un épimorphisme effectif.

\medskip\noindent
Soit $p : Y \to X$ un morphisme de schémas affines. Nous allons montrer que
les changements de base $g_i : Y_i \to Y$ des $f_i$ forment un épimorphisme
effectif, en appliquant le résultat du paragraphe précédent.
En effet, si $\mathcal{G}$ est un $\mathcal{O}_Y$-module quasi-cohérent, alors
$$
\Gamma(Y, \mathcal{G}) = \Gamma(X, p_*\mathcal{G}),\quad
\Gamma(Y_i, g_i^*\mathcal{G}) = \Gamma(X, f_i^*p_*\mathcal{G}),
$$
et
$$
\Gamma(Y_i \times_Y Y_j, (g_i \times g_j)^*\mathcal{G}) =
\Gamma(X, (f_i \times f_j)^*p_*\mathcal{G})
$$
par la formule triviale de changement de base
(Cohomologie des schémas, lemme \ref{coherent-lemma-affine-base-change}).
Ainsi, l'assertion (1) du lemme vaut pour la famille $g_i$.
\end{proof}

\begin{lemma}
\label{lemma-universal-effective-epimorphism-surjective}
Soit $\{f_i : X_i \to X\}_{i \in I}$ une famille de morphismes de schémas.
\begin{enumerate}
\item Si la famille est un épimorphisme effectif universel dans la catégorie
des schémas, alors $\coprod f_i$ est surjectif.
\item Si $X$ et les $X_i$ sont affines et si la famille est un épimorphisme
effectif universel dans la catégorie des schémas affines, alors
$\coprod f_i$ est surjectif.
\end{enumerate}
\end{lemma}

\begin{proof}
Omis. Indication : effectuer le changement de base par
$\Spec(\kappa(x)) \to X$ pour voir que tout $x \in X$ doit appartenir à
l'image.
\end{proof}

\begin{lemma}
\label{lemma-check-universal-effective-epimorphism-affine}
Soit $\{f_i : X_i \to X\}_{i \in I}$ une famille de morphismes de schémas.
Si, pour tout morphisme $Y \to X$ avec $Y$ affine, la famille des changements
de base $g_i : Y_i \to Y$ forme un épimorphisme effectif, alors la famille des
$f_i$ est un épimorphisme effectif universel dans la catégorie des schémas.
\end{lemma}

\begin{proof}
Soit $Y \to X$ un morphisme de schémas. Il faut montrer que les changements
de base $g_i : Y_i \to Y$ forment un épimorphisme effectif.
Pour cela, supposons donnés un schéma $T$ et des morphismes $h_i : Y_i \to T$
tels que $h_i \circ \text{pr}_1 = h_j \circ \text{pr}_2$ sur
$Y_i \times_Y Y_j$.
Choisissons un recouvrement ouvert affine $Y = \bigcup V_\alpha$.
Notons $V_{\alpha, i}$ l'image inverse de $V_\alpha$ dans $Y_i$.
Alors $V_{\alpha, i}  \to V_\alpha$ est le changement de base de $f_i$ par
$V_\alpha \to X$. Par hypothèse, la famille des restrictions
$h_i|_{V_{\alpha, i}}$ provient donc d'un morphisme de schémas
$h_\alpha : V_\alpha \to T$.
Nous laissons au lecteur le soin de montrer que ces morphismes coïncident sur
les intersections et définissent le morphisme voulu $Y \to T$.
Voir la discussion de Schémas, section
\ref{schemes-section-glueing-schemes}.
\end{proof}

\begin{lemma}
\label{lemma-universal-effective-epimorphism}
Soit $\{f_i : X_i \to X\}_{i \in I}$ une famille de morphismes de schémas
affines. Supposons que l'hypothèse équivalente du lemme
\ref{lemma-universal-effective-epimorphism-affine} soit satisfaite et que,
de plus, pour tout morphisme de schémas affines $Y \to X$, l'application
$$
\coprod X_i \times_X Y \longrightarrow Y
$$
soit une application submersive d'espaces topologiques
(Topologie, définition \ref{topology-definition-submersive}).
Alors notre famille de morphismes est un épimorphisme effectif universel dans
la catégorie des schémas.
\end{lemma}

\begin{proof}
D'après le lemme \ref{lemma-check-universal-effective-epimorphism-affine},
il suffit d'effectuer sur notre famille le changement de base par
$Y \to X$, avec $Y$ affine. Posons $Y_i = X_i \times_X Y$.
Soient $T$ un schéma et $h_i : Y_i \to T$ une famille de morphismes tels que
$h_i \circ \text{pr}_1 = h_j \circ \text{pr}_2$ sur $Y_i \times_Y Y_j$.
Remarquons que l'ensemble $Y$ est le coégalisateur des deux applications de
$\coprod Y_i \times_Y Y_j$ vers $\coprod Y_i$.
La surjectivité résulte du cas affine du lemme
\ref{lemma-universal-effective-epimorphism-surjective}, et l'injectivité du
lemme \ref{lemma-equiv-fibre-product}.
Il existe donc une application entre ensembles sous-jacents $h : Y \to T$
compatible avec les applications $h_i$. La seconde condition du lemme montre
que $h$ est continue. Ainsi, si $y \in Y$ et si $U \subset T$ est un voisinage
ouvert affine de $h(y)$, on peut trouver un ouvert affine $V \subset Y$ tel que
$h(V) \subset U$.
En posant $V_i = Y_i \times_Y V = X_i \times_X V$, on peut appliquer le
résultat démontré au lemme
\ref{lemma-universal-effective-epimorphism-affine} pour voir que
$h|_V : V \to U \subset T$ provient d'un unique morphisme de schémas affines
$h_V : V \to U$ qui coïncide avec $h_i|_{V_i}$ comme morphisme de schémas pour
tout $i$. En recollant ces $h_V$
(voir Schémas, section \ref{schemes-section-glueing-schemes}), on obtient le
morphisme voulu $Y \to T$.
\end{proof}

\begin{lemma}
\label{lemma-open-fpqc-covering}
Soit $\{f_i : T_i \to T\}_{i \in I}$ un recouvrement fpqc.
Supposons que, pour chaque $i$, on dispose d'un ouvert $W_i \subset T_i$ tel
que, pour tous $i, j \in I$,
$\text{pr}_0^{-1}(W_i) = \text{pr}_1^{-1}(W_j)$ comme ouverts de
$T_i \times_T T_j$. Alors il existe un unique ouvert $W \subset T$ tel que
$W_i = f_i^{-1}(W)$ pour tout $i$.
\end{lemma}

\begin{proof}
Appliquons le lemme \ref{lemma-equiv-fibre-product} à l'application
$\coprod_{i \in I} T_i \to T$.
Il en résulte qu'il existe une partie $W \subset T$ telle que
$W_i = f_i^{-1}(W)$ pour tout $i$, à savoir $W = \bigcup f_i(W_i)$.
Pour voir que $W$ est ouvert, on peut travailler localement pour la topologie
de Zariski sur $T$. On peut donc supposer que $T$ est affine.
D'après Topologies, définition
\ref{topologies-definition-fpqc-covering}, on peut choisir un recouvrement
fpqc standard $\{g_j : V_j \to T\}_{j \in J}$ qui raffine
$\{T_i \to T\}_{i \in I}$. Soient $\alpha : J \to I$ et
$h_j : V_j \to T_{\alpha(j)}$ comme dans Sites, définition
\ref{sites-definition-morphism-coverings}.
Alors $g_j^{-1}(W) = h_j^{-1}(W_{\alpha(j)})$.
On peut donc supposer que $\{f_i : T_i \to T\}$ est un recouvrement fpqc
standard. Dans ce cas, on applique Morphismes, lemme
\ref{morphisms-lemma-fpqc-quotient-topology}, au morphisme
$\coprod T_i \to T$ pour conclure que $W$ est ouvert.
\end{proof}

\begin{lemma}
\label{lemma-fpqc-universal-effective-epimorphisms}
Soit $\{T_i \to T\}$ un recouvrement fpqc, voir Topologies, définition
\ref{topologies-definition-fpqc-covering}.
Alors $\{T_i \to T\}$ est un épimorphisme effectif universel dans la catégorie
des schémas, voir Sites, définition
\ref{sites-definition-universal-effective-epimorphisms}.
Autrement dit, tout foncteur représentable sur la catégorie des schémas
vérifie la condition de faisceau pour la topologie fpqc, voir Topologies,
définition \ref{topologies-definition-sheaf-property-fpqc}.
\end{lemma}

\begin{proof}
Soit $S$ un schéma. Il faut montrer l'assertion suivante : étant donnés des
morphismes $\varphi_i : T_i \to S$ tels que
$\varphi_i|_{T_i \times_T T_j} = \varphi_j|_{T_i \times_T T_j}$, il existe un
unique morphisme $T \to S$ qui se restreint à $\varphi_i$ sur chaque $T_i$.
Autrement dit, il faut montrer que le foncteur
$h_S = \Mor_{\Sch}( - , S)$ vérifie la condition de faisceau pour la topologie
fpqc.

\medskip\noindent
Si $\{T_i \to T\}$ est un recouvrement de Zariski, cela résulte de Schémas,
lemme \ref{schemes-lemma-glue}. Topologies, lemme
\ref{topologies-lemma-sheaf-property-fpqc}, nous ramène donc au cas d'un
recouvrement $\{X \to Y\}$ donné par un unique morphisme plat surjectif de
schémas affines.

\medskip\noindent
Première démonstration. Le lemme \ref{lemma-sheaf-condition-holds} donne la
condition de faisceau pour les modules quasi-cohérents relativement à
$\{X \to Y\}$. D'après le lemme \ref{lemma-open-fpqc-covering}, le morphisme
$X \to Y$ est universellement submersif. On peut donc appliquer le lemme
\ref{lemma-universal-effective-epimorphism} pour voir que $\{X \to Y\}$ est un
épimorphisme effectif universel.

\medskip\noindent
Deuxième démonstration. Soit $R \to A$ l'homomorphisme d'anneaux fidèlement
plat correspondant à notre morphisme plat surjectif $\pi : X \to Y$.
Soit $f : X \to S$ un morphisme tel que
$f \circ \text{pr}_1 = f \circ \text{pr}_2$ comme morphismes
$X \times_Y X = \Spec(A \otimes_R A) \to S$.
D'après le lemme \ref{lemma-equiv-fibre-product}, on voit que, comme
application entre ensembles sous-jacents, $f$ est de la forme
$f = g \circ \pi$ pour une application ensembliste
$g : \Spec(R) \to S$. D'après Morphismes, lemme
\ref{morphisms-lemma-fpqc-quotient-topology}, et puisque $f$ est continue,
l'application $g$ est continue.

\medskip\noindent
Fixons $y \in Y = \Spec(R)$. Choisissons un ouvert affine $U \subset S$
contenant $g(y)$, disons $U = \Spec(B)$.
D'après ce qui précède, on peut choisir $r \in R$ tel que
$y \in D(r) \subset g^{-1}(U)$.
La restriction de $f$ à $\pi^{-1}(D(r))$, à valeurs dans $U$, correspond à un
homomorphisme d'anneaux $B \to A_r$. Les deux homomorphismes d'anneaux induits
$B \to A_r \otimes_{R_r} A_r = (A \otimes_R A)_r$ sont égaux par hypothèse
sur $f$. Remarquons que $R_r \to A_r$ est fidèlement plat.
D'après le lemme \ref{lemma-ff-exact}, l'égalisateur des deux flèches
$A_r \to A_r \otimes_{R_r} A_r$ est $R_r$.
On en conclut que $B \to A_r$ se factorise de manière unique par un
homomorphisme $B \to R_r$. Celui-ci donne à son tour un morphisme de schémas
$D(r) \to U \to S$, voir Schémas, lemme
\ref{schemes-lemma-morphism-into-affine}.

\medskip\noindent
Qu'avons-nous démontré jusqu'ici ? Nous avons montré que, pour tout idéal
premier $\mathfrak p \subset R$, il existe un ouvert affine standard
$D(r) \subset \Spec(R)$ tel que le morphisme
$f|_{\pi^{-1}(D(r))} : \pi^{-1}(D(r)) \to S$ se factorise de manière unique
par un certain morphisme de schémas $D(r) \to S$.
Nous omettons la vérification que ces morphismes se recollent pour donner le
morphisme voulu $\Spec(R) \to S$.
\end{proof}

\begin{lemma}
\label{lemma-coequalizer-fpqc-local}
Considérons des schémas $X, Y, Z$, des morphismes $a, b : X \to Y$ et un
morphisme $c : Y \to Z$ tel que $c \circ a = c \circ b$. Posons
$d = c \circ a = c \circ b$. S'il existe un recouvrement fpqc
$\{Z_i \to Z\}$ tel que
\begin{enumerate}
\item pour tout $i$, le morphisme $Y \times_{c, Z} Z_i \to Z_i$ soit le
coégalisateur de $(a, 1) : X \times_{d, Z} Z_i \to Y \times_{c, Z} Z_i$ et de
$(b, 1) : X \times_{d, Z} Z_i \to Y \times_{c, Z} Z_i$, et
\item pour tous $i$ et $i'$, le morphisme
$Y \times_{c, Z} (Z_i \times_Z Z_{i'}) \to (Z_i \times_Z Z_{i'})$ soit le
coégalisateur de
$(a, 1) : X \times_{d, Z} (Z_i \times_Z Z_{i'}) \to
Y \times_{c, Z} (Z_i \times_Z Z_{i'})$ et de
$(b, 1) : X \times_{d, Z} (Z_i \times_Z Z_{i'}) \to
Y \times_{c, Z} (Z_i \times_Z Z_{i'})$,
\end{enumerate}
alors $c$ est le coégalisateur de $a$ et $b$.
\end{lemma}

\begin{proof}
En effet, pour un schéma $T$, la donnée d'un morphisme $Z \to T$ équivaut à
celle d'une famille de morphismes $Z_i \to T$ qui coïncident sur les
intersections, d'après le lemme
\ref{lemma-fpqc-universal-effective-epimorphisms}.
\end{proof}















\section[Descente : finitude et lissité]{Descente des propriétés de finitude et de lissité des morphismes}
\label{section-descent-finiteness-morphisms}

\noindent
Dans cette section, nous montrons que plusieurs propriétés des morphismes
(être lisse, être localement de présentation finie, etc.) se descendent bien
par des morphismes fidèlement plats. Nous commençons par une version algébrique.
(Le lecteur ``noethérien'' devrait consulter le lemme
\ref{lemma-finite-type-local-source-fppf-algebra} plutôt que le lemme suivant.)

\begin{lemma}
\label{lemma-flat-finitely-presented-permanence-algebra}
Soient $R \to A \to B$ des homomorphismes d'anneaux.
Supposons que $R \to B$ soit de présentation finie et que $A \to B$ soit
fidèlement plat et de présentation finie. Alors $R \to A$ est de présentation
finie.
\end{lemma}

\begin{proof}
Considérons l'algèbre $C = B \otimes_A B$ munie des deux homomorphismes
$p, q : B \to C$ donnés par $p(b) = b \otimes 1$ et $q(b) = 1 \otimes b$.
Bien entendu, les deux composés $A \to B \to C$ sont égaux.
Comme $p : B \to C$ est plat et de présentation finie (c'est le changement de
base de $A \to B$), l'homomorphisme d'anneaux $R \to C$ est de présentation
finie (comme composé de $R \to B \to C$).

\medskip\noindent
Nous allons utiliser le critère d'Algèbre, lemme
\ref{algebra-lemma-characterize-finite-presentation}, pour montrer que
$R \to A$ est de présentation finie.
Soit $S$ une $R$-algèbre quelconque et supposons que
$S = \colim_{\lambda \in \Lambda} S_\lambda$ soit écrite comme une colimite
filtrante de $R$-algèbres. Soit $A \to S$ un homomorphisme de $R$-algèbres.
Il faut montrer que $A \to S$ se factorise par l'une des $S_\lambda$.
Considérons les anneaux $B' = S \otimes_A B$ et
$C' = S \otimes_A C = B' \otimes_S B'$.
Puisque $B$ est fidèlement plat et de présentation finie sur $A$, l'anneau
$B'$ est lui aussi fidèlement plat et de présentation finie sur $S$.
D'après Algèbre, lemme
\ref{algebra-lemma-flat-finite-presentation-limit-flat}, partie (2), appliqué
au couple $(S \to B', B')$ et au système $(S_\lambda)$, il existe un
$\lambda_0 \in \Lambda$ et une $S_{\lambda_0}$-algèbre plate et de présentation
finie $B_{\lambda_0}$ telle que
$B' = S \otimes_{S_{\lambda_0}} B_{\lambda_0}$.
Pour $\lambda \geq \lambda_0$, posons
$B_\lambda = S_\lambda \otimes_{S_{\lambda_0}} B_{\lambda_0}$ et
$C_\lambda = B_\lambda \otimes_{S_\lambda} B_\lambda$.

\medskip\noindent
Nous interrompons le fil de l'argument pour montrer que
$S_\lambda \to B_\lambda$ est fidèlement plat pour $\lambda$ assez grand.
(Cela devrait en réalité faire l'objet d'un lemme distinct ailleurs, peut-être
dans le chapitre sur les limites.)
Comme $\Spec(B_{\lambda_0}) \to \Spec(S_{\lambda_0})$ est plat et de
présentation finie, il est ouvert (voir Morphismes, lemme
\ref{morphisms-lemma-fppf-open}).
Soit $I \subset S_{\lambda_0}$ un idéal tel que
$V(I) \subset \Spec(S_{\lambda_0})$ soit le complémentaire de l'image.
Remarquons que la formation de l'image commute au changement de base.
Puisque $\Spec(B') \to \Spec(S)$ est surjectif et que
$B' = B_{\lambda_0} \otimes_{S_{\lambda_0}} S$, on a $IS = S$.
Il existe donc $\lambda \geq \lambda_0$ tel que
$IS_{\lambda} = S_\lambda$. Pour cet indice et tous les indices
$\lambda$ plus grands, le morphisme
$\Spec(B_\lambda) \to \Spec(S_\lambda)$ est surjectif.

\medskip\noindent
Par analogie avec la notation du premier paragraphe de la démonstration,
notons $p_\lambda, q_\lambda : B_\lambda \to C_\lambda$ les deux
homomorphismes canoniques. Alors
$B' = \colim_{\lambda \geq \lambda_0} B_\lambda$ et
$C' = \colim_{\lambda \geq \lambda_0} C_\lambda$.
Comme $B$ et $C$ sont de présentation finie sur $R$, il existe
(en appliquant plusieurs fois Algèbre, lemme
\ref{algebra-lemma-characterize-finite-presentation}) un
$\lambda \geq \lambda_0$ et des homomorphismes de $R$-algèbres
$B \to B_\lambda$, $C \to C_\lambda$ tels que le diagramme
$$
\xymatrix{
C \ar[rr] & &
C_\lambda \\
B \ar[rr]
\ar@<1ex>[u]^-p
\ar@<-1ex>[u]_-q
& &
B_\lambda
\ar@<1ex>[u]^-{p_\lambda}
\ar@<-1ex>[u]_-{q_\lambda}
}
$$
soit commutatif. Cela signifie que $A \to B \to B_\lambda$ prend ses valeurs
dans l'égalisateur de $p_\lambda$ et $q_\lambda$.
D'après le lemme \ref{lemma-ff-exact}, $S_\lambda$ est l'égalisateur de
$p_\lambda$ et $q_\lambda$. On obtient ainsi l'homomorphisme d'anneaux voulu
$A \to S_\lambda$, ce qui conclut.
\end{proof}

\noindent
Voici une version plus simple qui traite de la propriété d'être de type fini.

\begin{lemma}
\label{lemma-finite-type-local-source-fppf-algebra}
Soient $R \to A \to B$ des homomorphismes d'anneaux.
Supposons que $R \to B$ soit de type fini et que $A \to B$ soit fidèlement plat
et de présentation finie. Alors $R \to A$ est de type fini.
\end{lemma}

\begin{proof}
D'après Algèbre, lemme
\ref{algebra-lemma-descend-faithfully-flat-finite-presentation}, il existe un
diagramme commutatif
$$
\xymatrix{
R \ar[r] \ar@{=}[d] &
A_0 \ar[d] \ar[r] &
B_0 \ar[d] \\
R \ar[r] & A \ar[r] & B
}
$$
où $R \to A_0$ est de présentation finie, $A_0 \to B_0$ est fidèlement plat
et de présentation finie, et $B = A \otimes_{A_0} B_0$.
Puisque $R \to B$ est de type fini par hypothèse, on peut adjoindre quelques
éléments à $A_0$ et supposer que l'homomorphisme $B_0 \to B$ est surjectif.
Dans ce cas, comme $A_0 \to B_0$ est fidèlement plat et que
$$
(A_0 \to A) \otimes_{A_0} B_0 \cong (B_0 \to B)
$$
est surjectif, $A_0 \to A$ est lui aussi surjectif, ce qui conclut.
\end{proof}

\begin{lemma}
\label{lemma-flat-finitely-presented-permanence}
\begin{reference}
\cite[IV, 17.7.5 (i) and (ii)]{EGA}.
\end{reference}
Soit
$$
\xymatrix{
X \ar[rr]_f \ar[rd]_p & &
Y \ar[dl]^q \\
& S
}
$$
un diagramme commutatif de morphismes de schémas. Supposons que $f$ soit
surjectif, plat et localement de présentation finie, et que $p$ soit localement
de présentation finie (resp.\ localement de type fini).
Alors $q$ est localement de présentation finie (resp.\ localement de type fini).
\end{lemma}

\begin{proof}
La question est locale sur $S$ et sur $Y$. On peut donc supposer que $S$ et
$Y$ sont affines. Puisque $f$ est plat et localement de présentation finie,
$f$ est ouvert (Morphismes, lemme \ref{morphisms-lemma-fppf-open}).
Comme $Y$ est quasi-compact, il existe donc un nombre fini d'ouverts affines
$X_i \subset X$ tels que $Y = \bigcup f(X_i)$.
On peut manifestement remplacer $X$ par $\coprod X_i$ et supposer ainsi que
$X$ est lui aussi affine. Dans ce cas, le lemme équivaut au lemme
\ref{lemma-flat-finitely-presented-permanence-algebra}
(resp. au lemme \ref{lemma-finite-type-local-source-fppf-algebra}) ci-dessus.
\end{proof}

\noindent
Nous allons nous en servir pour améliorer certains des résultats antérieurs
sur les morphismes.

\begin{lemma}
\label{lemma-syntomic-smooth-etale-permanence}
Soit
$$
\xymatrix{
X \ar[rr]_f \ar[rd]_p & &
Y \ar[dl]^q \\
& S
}
$$
un diagramme commutatif de morphismes de schémas. Supposons que
\begin{enumerate}
\item $f$ soit surjectif et syntomique (resp.\ lisse, resp.\ \'etale),
\item $p$ soit syntomique (resp.\ lisse, resp.\ \'etale).
\end{enumerate}
Alors $q$ est syntomique (resp.\ lisse, resp.\ \'etale).
\end{lemma}

\begin{proof}
Combiner Morphismes, lemmes
\ref{morphisms-lemma-syntomic-permanence},
\ref{morphisms-lemma-smooth-permanence} et
\ref{morphisms-lemma-etale-permanence-two}, avec le lemme
\ref{lemma-flat-finitely-presented-permanence} ci-dessus.
\end{proof}

\noindent
En fait, on peut renforcer ce résultat comme suit.

\begin{lemma}
\label{lemma-smooth-permanence}
Soit
$$
\xymatrix{
X \ar[rr]_f \ar[rd]_p & &
Y \ar[dl]^q \\
& S
}
$$
un diagramme commutatif de morphismes de schémas. Supposons que
\begin{enumerate}
\item $f$ soit surjectif, plat et localement de présentation finie,
\item $p$ soit lisse (resp.\ \'etale).
\end{enumerate}
Alors $q$ est lisse (resp.\ \'etale).
\end{lemma}

\begin{proof}
Supposons (1) et que $p$ soit lisse. Le lemme
\ref{lemma-flat-finitely-presented-permanence} montre que $q$ est localement de
présentation finie. Morphismes, lemme
\ref{morphisms-lemma-flat-permanence}, montre que $q$ est plat.
Il suffit donc de montrer que les fibres de $q$ sont lisses, voir le lemme
\ref{morphisms-lemma-smooth-flat-smooth-fibres} de Morphismes.
Appliquons le lemme \ref{varieties-lemma-flat-under-smooth} de Variétés aux
morphismes plats surjectifs $X_s \to Y_s$, pour $s \in S$, et concluons.
Nous omettons la démonstration du cas \'etale.
\end{proof}

\begin{remark}
\label{remark-smooth-permanence}
Sous les hypothèses (1), avec $p$ lisse, du lemme
\ref{lemma-smooth-permanence}, le morphisme $X \to Y$ n'est pas nécessairement
lisse. Un contre-exemple est donné par $S = \Spec(k)$,
$X = \Spec(k[s])$, $Y = \Spec(k[t])$ et $f$ défini par
$t \mapsto s^2$. Voir cependant le lemme
\ref{lemma-syntomic-permanence} pour des renseignements sur la structure de
$f$.
\end{remark}

\begin{lemma}
\label{lemma-syntomic-permanence}
Soit
$$
\xymatrix{
X \ar[rr]_f \ar[rd]_p & &
Y \ar[dl]^q \\
& S
}
$$
un diagramme commutatif de morphismes de schémas. Supposons que
\begin{enumerate}
\item $f$ soit surjectif, plat et localement de présentation finie,
\item $p$ soit syntomique.
\end{enumerate}
Alors $q$ et $f$ sont tous deux syntomiques.
\end{lemma}

\begin{proof}
Le lemme \ref{lemma-flat-finitely-presented-permanence} montre que $q$
est de présentation finie. Le lemme
\ref{morphisms-lemma-flat-permanence} de Morphismes
montre que $q$ est plat.
D'après le lemme
\ref{morphisms-lemma-syntomic-locally-standard-syntomic} de Morphismes,
il suffit donc de montrer que les anneaux locaux des fibres de
$Y \to S$ et de $X \to Y$ sont des anneaux locaux d'intersection
complète. Pour cela, nous pouvons prendre la fibre de
$X \to Y \to S$ en un point $s \in S$ ; autrement dit, nous pouvons
supposer que $S$ est le spectre d'un corps. Choisissons un point
$x \in X$, d'image $y \in Y$, et considérons l'homomorphisme d'anneaux
$$
\mathcal{O}_{Y, y} \longrightarrow \mathcal{O}_{X, x}
$$
C'est un homomorphisme local plat d'anneaux locaux noethériens.
L'anneau local $\mathcal{O}_{X, x}$ est d'intersection complète.
Nous pouvons donc appliquer le résultat d'Avramov, voir le lemme
\ref{dpa-lemma-avramov} d'Algèbre à puissances divisées, et conclure que
$\mathcal{O}_{Y, y}$ et
$\mathcal{O}_{X, x}/\mathfrak m_y\mathcal{O}_{X, x}$ sont tous deux
des anneaux d'intersection complète.
\end{proof}

\noindent
Le type de lemme suivant est parfois utile.

\begin{lemma}
\label{lemma-curiosity}
Soient $X \to Y \to Z$ des morphismes de schémas.
Soit $P$ l'une des propriétés suivantes des morphismes de schémas :
être plat, être localement de type fini, être localement de présentation
finie. Supposons que $X \to Z$ possède la propriété $P$ et que
$\{X \to Y\}$ puisse être raffiné par un recouvrement fppf de $Y$.
Alors $Y \to Z$ possède la propriété $P$.
\end{lemma}

\begin{proof}
Soient $\Spec(C) \subset Z$ un ouvert affine et
$\Spec(B) \subset Y$ un ouvert affine dont l'image est contenue dans
$\Spec(C)$. L'hypothèse sur $X \to Y$ permet de trouver un recouvrement
fppf affine standard $\{\Spec(B_j) \to \Spec(B)\}$ et des relèvements
$x_j : \Spec(B_j) \to X$. Comme $\Spec(B_j)$ est quasi-compact, nous
pouvons trouver un nombre fini d'ouverts affines
$\Spec(A_i) \subset X$ au-dessus de $\Spec(B)$ tels que l'image de
chaque $x_j$ soit contenue dans la réunion $\bigcup \Spec(A_i)$.
Après avoir remplacé chaque $\Spec(B_j)$ par un recouvrement de Zariski
affine standard, nous pouvons donc supposer que nous disposons d'un
recouvrement fppf affine standard
$\{\Spec(B_i) \to \Spec(B)\}$ tel que chaque morphisme
$\Spec(B_i) \to Y$ se factorise par un ouvert affine
$\Spec(A_i) \subset X$ situé au-dessus de $\Spec(B)$. Autrement dit,
nous avons des homomorphismes d'anneaux $C \to B \to A_i \to B_i$
pour chaque $i$. Nous pouvons aussi considérer
$$
C \to B \to A = \prod A_i \to B' = \prod B_i
$$
et l'homomorphisme d'anneaux $B \to \prod B_i$ est fidèlement plat et
de présentation finie.

\medskip\noindent
Le cas $P = \text{plat}$. Dans ce cas, nous savons que $C \to A$ est plat et
nous devons montrer que $C \to B$ est plat. Soit
$N \to N' \to N''$ une suite exacte de $C$-modules. Nous voulons
montrer que la suite
$N \otimes_C B \to N' \otimes_C B \to N'' \otimes_C B$ est exacte.
Notons $H$ sa cohomologie et $H'$ celle de la suite
$N \otimes_C B' \to N' \otimes_C B' \to N'' \otimes_C B'$.
Comme $B \to B'$ est plat, nous avons $H' = H \otimes_B B'$.
D'autre part, la suite
$N \otimes_C A \to N' \otimes_C A \to N'' \otimes_C A$ est exacte,
donc d'homologie nulle. Le morphisme $H \to H'$ est par conséquent
nul, puisqu'il se factorise par le module nul. Ainsi $H' = 0$.
Puisque $B \to B'$ est fidèlement plat, nous concluons que $H = 0$,
comme souhaité.

\medskip\noindent
Le cas $P = \text{localement de type fini}$. Dans ce cas, nous savons que
$C \to A$ est de type fini et nous devons montrer que $C \to B$ est de
type fini. Puisque $B \to B'$ est de présentation finie (donc de type
fini), l'homomorphisme $A \to B'$ est de type fini ; voir le lemme
\ref{algebra-lemma-compose-finite-type} d'Algèbre. Par conséquent,
$C \to B'$ est de type fini, et nous concluons par le
Lemme \ref{lemma-finite-type-local-source-fppf-algebra}.

\medskip\noindent
Le cas $P = \text{localement de présentation finie}$. Dans ce cas, nous savons
que $C \to A$ est de présentation finie et nous devons montrer que
$C \to B$ est de présentation finie. Puisque $B \to B'$ est de
présentation finie et que $B \to A$ est de type fini, l'homomorphisme
$A \to B'$ est de présentation finie ; voir le lemme
\ref{algebra-lemma-compose-finite-type} d'Algèbre. Par conséquent,
$C \to B'$ est de présentation finie, et nous concluons par le
Lemme \ref{lemma-flat-finitely-presented-permanence-algebra}.
\end{proof}










\section{Propriétés locales des schémas}
\label{section-descending-properties}

\noindent
Il arrive souvent que l'on puisse établir qu'une certaine propriété est
vérifiée par les membres d'un recouvrement d'un schéma. Dans bien des cas,
le schéma possède alors lui aussi cette propriété. Par exemple, si $S$ est
un schéma et si $f : S' \to S$ est un morphisme plat surjectif tel que $S'$
soit réduit, alors $S$ est réduit. On peut le démontrer en considérant les
anneaux locaux et en appliquant le lemme
\ref{algebra-lemma-descent-reduced} d'Algèbre. On dit que la propriété d'être
réduit {\it se descend par les morphismes plats surjectifs}. Des résultats de
ce type sont réunis dans Algèbre, section
\ref{algebra-section-descending-properties}, et, pour les schémas, dans la
section \ref{section-variants}. Des résultats analogues sur la descente des
propriétés de morphismes figurent dans la
section \ref{section-descent-finiteness-morphisms}.

\medskip\noindent
En revanche, il existe des morphismes plats surjectifs $f : S' \to S$ pour
lesquels $S$ est réduit mais $S'$ ne l'est pas, par exemple
$\Spec(k[x]/(x^2)) \to \Spec(k)$. La propriété d'être réduit n'est donc pas
{\it préservée par les morphismes plats}. La propriété d'avoir des corps
résiduels infinis est préservée par les morphismes plats (mais ne se descend
évidemment pas par les morphismes plats surjectifs). Des résultats de ce
type sont réunis dans Algèbre, section
\ref{algebra-section-ascending-properties}.

\medskip\noindent
Enfin, on dit qu'une propriété est {\it locale pour la topologie plate}
si elle est préservée par les morphismes plats et se descend par les
morphismes plats surjectifs. Un exemple assez élémentaire est la propriété d'avoir des corps
résiduels d'une caractéristique donnée. Afin de préciser cette notion et de
la relier aux diverses topologies sur les schémas, posons la définition
formelle suivante.

\begin{definition}
\label{definition-property-local}
Soit $\mathcal{P}$ une propriété des schémas. Soit
$\tau \in \{fpqc, \linebreak[0] fppf, \linebreak[0] syntomic, \linebreak[0]
smooth, \linebreak[0] \etale, \linebreak[0] Zariski\}$. On dit que
$\mathcal{P}$ est {\it locale pour la topologie $\tau$} si, pour tout
recouvrement pour la topologie $\tau$, $\{S_i \to S\}_{i \in I}$ (voir
Topologies, section \ref{topologies-section-procedure}), on a
$$
S \text{ vérifie }\mathcal{P}
\Leftrightarrow
\text{chaque }S_i \text{ vérifie }\mathcal{P}.
$$
\end{definition}

\noindent
Bien entendu, puisque les isomorphismes sont toujours des recouvrements,
on constate (ou l'on exige) que la propriété $\mathcal{P}$ est vérifiée par
$S$ si et seulement si elle est vérifiée par tout schéma $S'$ isomorphe à
$S$. En fait, si
$\tau = fpqc, \linebreak[0] fppf, \linebreak[0] syntomic,
\linebreak[0] smooth, \linebreak[0] \etale$ ou $Zariski$, si $S$ vérifie
$\mathcal{P}$ et si $S' \to S$ est respectivement plat, plat et localement
de présentation finie, syntomique, lisse, \'etale, ou une immersion ouverte,
alors $S'$ vérifie $\mathcal{P}$. En effet, on peut toujours compléter
$\{S' \to S\}$ en un recouvrement pour la topologie $\tau$.

\medskip\noindent
Nous avons les implications suivantes :
$\mathcal{P}$ est locale pour la topologie fpqc
$\Rightarrow$
$\mathcal{P}$ est locale pour la topologie fppf
$\Rightarrow$
$\mathcal{P}$ est locale pour la topologie syntomique
$\Rightarrow$
$\mathcal{P}$ est locale pour la topologie lisse
$\Rightarrow$
$\mathcal{P}$ est locale pour la topologie \'etale
$\Rightarrow$
$\mathcal{P}$ est locale pour la topologie de Zariski.
Cela résulte des lemmes
\ref{topologies-lemma-zariski-etale},
\ref{topologies-lemma-zariski-etale-smooth},
\ref{topologies-lemma-zariski-etale-smooth-syntomic},
\ref{topologies-lemma-zariski-etale-smooth-syntomic-fppf} et
\ref{topologies-lemma-zariski-etale-smooth-syntomic-fppf-fpqc} de Topologies.

\begin{lemma}
\label{lemma-descending-properties}
Soit $\mathcal{P}$ une propriété des schémas.
Soit $\tau \in \{fpqc, \linebreak[0] fppf, \linebreak[0]
\etale, \linebreak[0] smooth, \linebreak[0] syntomic\}$. Supposons que
\begin{enumerate}
\item la propriété soit locale pour la topologie de Zariski,
\item pour tout morphisme de schémas affines $S' \to S$ qui est plat,
plat de présentation finie, \'etale, lisse ou syntomique selon que $\tau$
vaut fpqc, fppf, \'etale, smooth ou syntomic, la propriété $\mathcal{P}$
soit vérifiée par $S'$ si la propriété $\mathcal{P}$ est vérifiée par $S$, et
\item pour tout morphisme surjectif de schémas affines $S' \to S$ qui est
plat, plat de présentation finie, \'etale, lisse ou syntomique selon que
$\tau$ vaut fpqc, fppf, \'etale, smooth ou syntomic, la propriété
$\mathcal{P}$ soit vérifiée par $S$ si la propriété $\mathcal{P}$ est
vérifiée par $S'$.
\end{enumerate}
Alors $\mathcal{P}$ est locale sur la base pour la topologie $\tau$.
\end{lemma}

\begin{proof}
Cela résulte presque immédiatement de la définition d'un recouvrement pour
la topologie $\tau$ ; voir les définitions
\ref{topologies-definition-fpqc-covering},
\ref{topologies-definition-fppf-covering},
\ref{topologies-definition-etale-covering},
\ref{topologies-definition-smooth-covering} ou
\ref{topologies-definition-syntomic-covering} de Topologies, et les lemmes
\ref{topologies-lemma-fpqc-affine},
\ref{topologies-lemma-fppf-affine},
\ref{topologies-lemma-etale-affine},
\ref{topologies-lemma-smooth-affine} ou
\ref{topologies-lemma-syntomic-affine} de Topologies. Nous omettons les détails.
\end{proof}

\begin{remark}
\label{remark-descending-properties-standard}
Dans le lemme \ref{lemma-descending-properties}, si $\tau = smooth$, on
peut supposer, dans la condition (3), que le morphisme est un morphisme
lisse standard (surjectif). Il en va de même lorsque $\tau = syntomic$ ou
$\tau = \etale$.
\end{remark}





\section{Propriétés des schémas locales pour la topologie fppf}
\label{section-descending-properties-fppf}

\noindent
Dans cette section, nous donnons quelques propriétés de schémas qui sont
locales sur la base pour la topologie fppf.

\begin{lemma}
\label{lemma-Noetherian-local-fppf}
La propriété $\mathcal{P}(S) =$``$S$ est localement noethérien'' est locale
pour la topologie fppf.
\end{lemma}

\begin{proof}
Nous allons appliquer le lemme \ref{lemma-descending-properties}.
Tout d'abord, la propriété ``être localement noethérien'' est locale pour la
topologie de Zariski. Cela résulte immédiatement de la définition ; voir
Propriétés, définition \ref{properties-definition-noetherian}.
Ensuite, si $S' \to S$ est un morphisme plat de présentation finie entre
schémas affines et si $S$ est localement noethérien, alors $S'$ est
localement noethérien ; c'est le lemme
\ref{morphisms-lemma-finite-type-noetherian} de Morphismes. Enfin, si $S' \to S$ est un
morphisme plat, surjectif et de présentation finie entre schémas affines et
si $S'$ est localement noethérien, alors $S$ est localement noethérien ; cela
résulte du lemme \ref{algebra-lemma-descent-Noetherian} d'Algèbre.
Les conditions (1), (2) et (3) du lemme
\ref{lemma-descending-properties} sont donc satisfaites, ce qui conclut.
\end{proof}

\begin{lemma}
\label{lemma-Jacobson-local-fppf}
La propriété $\mathcal{P}(S) =$``$S$ est jacobsonien'' est locale pour la
topologie fppf.
\end{lemma}

\begin{proof}
Nous allons appliquer le Lemme \ref{lemma-descending-properties}.
Tout d'abord, la propriété ``être jacobsonien'' est locale pour la topologie
de Zariski ; c'est Propriétés, Lemme
\ref{properties-lemma-locally-jacobson}. Ensuite, si $S' \to S$ est un
morphisme plat de présentation finie entre schémas affines et si $S$ est
jacobsonien, alors $S'$ est jacobsonien ; c'est Morphismes, Lemme
\ref{morphisms-lemma-Jacobson-universally-Jacobson}. Enfin, supposons que
$f : S' \to S$ soit un morphisme plat, surjectif et de présentation finie
entre schémas affines, et que $S'$ soit jacobsonien ; nous devons montrer
que $S$ est jacobsonien. Écrivons $S = \Spec(A)$ et $S' = \Spec(B)$, le
morphisme $S' \to S$ étant donné par $A \to B$. Alors $A \to B$ est de
présentation finie et fidèlement plat. De plus, l'anneau $B$ est jacobsonien ;
voir Propriétés, Lemme \ref{properties-lemma-locally-jacobson}.

\medskip\noindent
D'après Algèbre, Lemme \ref{algebra-lemma-fppf-fpqf}, il existe un diagramme
$$
\xymatrix{
B \ar[rr] & & B' \\
& A \ar[ru] \ar[lu] &
}
$$
tel que $A \to B'$ soit de présentation finie, fidèlement plat et quasi-fini.
En particulier, $B \to B'$ est de type fini et Algèbre, Proposition
\ref{algebra-proposition-Jacobson-permanence}, montre que $B'$ est
jacobsonien. Nous pouvons donc supposer que $A \to B$ est quasi-fini, en plus
d'être fidèlement plat et de présentation finie.

\medskip\noindent
Supposons que $A$ ne soit pas jacobsonien et cherchons une contradiction.
D'après Algèbre, Lemme \ref{algebra-lemma-characterize-jacobson}, il existe
un idéal premier non maximal $\mathfrak p \subset A$ et un élément $f \in A$,
$f \not \in \mathfrak p$, tels que
$V(\mathfrak p) \cap D(f) = \{\mathfrak p\}$.

\medskip\noindent
On obtient une contradiction comme suit. Choisissons d'abord un idéal maximal
$\mathfrak m$ de $A$ contenant $\mathfrak p$. Choisissons un idéal premier
$\mathfrak m' \subset B$ au-dessus de $\mathfrak m$ (il en existe un parce
que $A \to B$ est fidèlement plat ; voir Algèbre, Lemme
\ref{algebra-lemma-ff-rings}). Comme $A \to B$ est plat, le théorème de
descente, voir Algèbre, Lemme \ref{algebra-lemma-flat-going-down}, fournit
un idéal premier $\mathfrak q \subset \mathfrak m'$ au-dessus de
$\mathfrak p$. En particulier, $\mathfrak q$ n'est pas maximal. Une nouvelle
application d'Algèbre, Lemme \ref{algebra-lemma-characterize-jacobson},
montre donc que l'ensemble $V(\mathfrak q) \cap D(f)$ est infini (c'est ici
que nous utilisons enfin le fait que $B$ est jacobsonien). Tous les points de
$V(\mathfrak q) \cap D(f)$ ont pour image
$V(\mathfrak p) \cap D(f) = \{\mathfrak p\}$. La fibre au-dessus de
$\mathfrak p$ est donc infinie. Cela contredit le fait que $A \to B$ est
quasi-fini ; voir Algèbre, Lemme \ref{algebra-lemma-quasi-finite}, ou, plus
explicitement, Morphismes, Lemme \ref{morphisms-lemma-quasi-finite}.
Le lemme est ainsi démontré.
\end{proof}

\begin{lemma}
\label{lemma-locally-finite-nr-irred-local-fppf}
La propriété $\mathcal{P}(S) =$``tout ouvert quasi-compact de $S$ possède
un nombre fini de composantes irréductibles'' est locale pour la topologie
fppf.
\end{lemma}

\begin{proof}
Nous allons appliquer le lemme \ref{lemma-descending-properties}.
Tout d'abord, $\mathcal{P}$ est locale pour la topologie de Zariski.
Ensuite, supposons que $T \to S$ soit un morphisme plat de présentation
finie entre schémas affines et que $S$ possède un nombre fini de composantes
irréductibles ; il en va alors de même de $T$. En effet, comme $T \to S$ est
plat, les points génériques de $T$ s'envoient sur les points génériques de
$S$ ; voir Morphismes, Lemme
\ref{morphisms-lemma-generalizations-lift-flat}. Il suffit donc de montrer
que, pour $s \in S$, la fibre $T_s$ possède un nombre fini de points
génériques. Or $T_s$ est un schéma affine de type fini sur $\kappa(s)$ ;
voir Morphismes, Lemme \ref{morphisms-lemma-base-change-finite-type}.
Ainsi, $T_s$ est noethérien et possède un nombre fini de composantes
irréductibles (Morphismes, Lemme
\ref{morphisms-lemma-finite-type-noetherian}, et Propriétés, Lemme
\ref{properties-lemma-Noetherian-irreducible-components}). Enfin, si
$T \to S$ est un morphisme plat, surjectif et de présentation finie entre
schémas affines et si $T$ possède un nombre fini de composantes
irréductibles, il en va de même de $S$. Cela résulte de Topologie, Lemme
\ref{topology-lemma-surjective-continuous-irreducible-components}.
Les conditions (1), (2) et (3) du lemme
\ref{lemma-descending-properties} sont donc satisfaites, ce qui conclut.
\end{proof}




\section{Propriétés des schémas locales pour la topologie syntomique}
\label{section-descending-properties-syntomic}

\noindent
Dans cette section, nous donnons quelques propriétés de schémas qui sont
locales sur la base pour la topologie syntomique.

\begin{lemma}
\label{lemma-Sk-local-syntomic}
La propriété $\mathcal{P}(S) =$``$S$ est localement noethérien et vérifie
$(S_k)$'' est locale pour la topologie syntomique.
\end{lemma}

\begin{proof}
Nous allons vérifier les conditions (1), (2) et (3) du lemme
\ref{lemma-descending-properties}. Comme un morphisme syntomique est plat de
présentation finie (Morphismes, Lemmes
\ref{morphisms-lemma-syntomic-flat} et
\ref{morphisms-lemma-syntomic-locally-finite-presentation}), la propriété
``être localement noethérien'' a déjà été traitée dans la preuve du
lemme \ref{lemma-Noetherian-local-fppf}. Nous l'utiliserons sans autre
mention. Tout d'abord, $\mathcal{P}$ est locale pour la topologie de Zariski.
Cela résulte immédiatement de la définition ; voir Cohomologie des schémas,
Définition \ref{coherent-definition-depth}. Ensuite, si $S' \to S$ est un
morphisme syntomique entre schémas affines, si $S$ vérifie $\mathcal{P}$,
alors $S'$ vérifie $\mathcal{P}$. C'est Algèbre, Lemme
\ref{algebra-lemma-Sk-goes-up} (appliquer Morphismes, Lemme
\ref{morphisms-lemma-syntomic-characterize}, ainsi qu'Algèbre, Définition
\ref{algebra-definition-lci} et Lemme \ref{algebra-lemma-lci-CM}).
Enfin, si $S' \to S$ est un morphisme syntomique surjectif entre schémas
affines et si $S'$ vérifie $\mathcal{P}$, alors $S$ vérifie $\mathcal{P}$ ;
c'est Algèbre, Lemme \ref{algebra-lemma-descent-Sk}.
Les conditions (1), (2) et (3) du lemme
\ref{lemma-descending-properties} sont donc satisfaites, ce qui conclut.
\end{proof}

\begin{lemma}
\label{lemma-CM-local-syntomic}
La propriété $\mathcal{P}(S) =$``$S$ est de Cohen-Macaulay'' est locale pour
la topologie syntomique.
\end{lemma}

\begin{proof}
Cela résulte du lemme \ref{lemma-Sk-local-syntomic}, puisqu'un schéma est de
Cohen-Macaulay si et seulement s'il est localement noethérien et vérifie
$(S_k)$ pour tout $k \geq 0$ ; voir Propriétés, Lemme
\ref{properties-lemma-scheme-CM-iff-all-Sk}.
\end{proof}







\section{Propriétés des schémas locales pour la topologie lisse}
\label{section-descending-properties-smooth}

\noindent
Dans cette section, nous donnons quelques propriétés de schémas qui sont
locales sur la base pour la topologie lisse.

\begin{lemma}
\label{lemma-reduced-local-smooth}
La propriété $\mathcal{P}(S) =$``$S$ est réduit'' est locale pour la topologie
lisse.
\end{lemma}

\begin{proof}
Nous allons appliquer le lemme \ref{lemma-descending-properties}.
Tout d'abord, la propriété ``être réduit'' est locale pour la topologie de
Zariski. Cela résulte immédiatement de la définition ; voir Schémas,
Définition \ref{schemes-definition-reduced}. Ensuite, si $S' \to S$ est un
morphisme lisse entre schémas affines et si $S$ est réduit, alors $S'$ est
réduit ; c'est Algèbre, Lemme \ref{algebra-lemma-reduced-goes-up}.
Enfin, si $S' \to S$ est un morphisme lisse surjectif entre schémas affines
et si $S'$ est réduit, alors $S$ est réduit ; c'est Algèbre, Lemme
\ref{algebra-lemma-descent-reduced}. Les conditions (1), (2) et (3) du
lemme \ref{lemma-descending-properties} sont donc satisfaites, ce qui conclut.
\end{proof}

\begin{lemma}
\label{lemma-normal-local-smooth}
\begin{slogan}
La normalité est locale pour la topologie lisse.
\end{slogan}
La propriété $\mathcal{P}(S) =$``$S$ est normal'' est locale pour la topologie
lisse.
\end{lemma}

\begin{proof}
Nous allons appliquer le lemme \ref{lemma-descending-properties}.
Tout d'abord, la propriété ``être normal'' est locale pour la topologie de
Zariski. Cela résulte immédiatement de la définition ; voir Propriétés,
Définition \ref{properties-definition-normal}. Ensuite, si $S' \to S$ est
un morphisme lisse entre schémas affines et si $S$ est normal, alors $S'$ est
normal ; c'est Algèbre, Lemme \ref{algebra-lemma-normal-goes-up}. Enfin, si
$S' \to S$ est un morphisme lisse surjectif entre schémas affines et si $S'$
est normal, alors $S$ est normal ; c'est Algèbre, Lemme
\ref{algebra-lemma-descent-normal}. Les conditions (1), (2) et (3) du
lemme \ref{lemma-descending-properties} sont donc satisfaites, ce qui conclut.
\end{proof}

\begin{lemma}
\label{lemma-Rk-local-smooth}
La propriété $\mathcal{P}(S) =$``$S$ est localement noethérien et vérifie
$(R_k)$'' est locale pour la topologie lisse.
\end{lemma}

\begin{proof}
Nous allons vérifier les conditions (1), (2) et (3) du lemme
\ref{lemma-descending-properties}. Comme un morphisme lisse est plat et de
présentation finie (voir les lemmes \ref{morphisms-lemma-smooth-flat} et
\ref{morphisms-lemma-smooth-locally-finite-presentation} de Morphismes), la propriété
``être localement noethérien'' a déjà été traitée dans la preuve du
lemme \ref{lemma-Noetherian-local-fppf}. Nous l'utiliserons sans autre
mention. Tout d'abord, $\mathcal{P}$ est locale pour la topologie de Zariski.
Cela résulte immédiatement de la définition \ref{properties-definition-Rk}
de Propriétés. Ensuite, si $S' \to S$ est un morphisme
lisse entre schémas affines et si $S$ vérifie $\mathcal{P}$, alors $S'$
vérifie $\mathcal{P}$. Cela résulte du lemme
\ref{algebra-lemma-Rk-goes-up} d'Algèbre (où l'on utilise le lemme
\ref{morphisms-lemma-smooth-characterize} de Morphismes, ainsi que les lemmes
\ref{algebra-lemma-base-change-smooth} et
\ref{algebra-lemma-characterize-smooth-over-field} d'Algèbre). Enfin, si $S' \to S$
est un morphisme lisse surjectif entre schémas affines et si $S'$ vérifie
$\mathcal{P}$, alors $S$ vérifie $\mathcal{P}$ ; c'est Algèbre, Lemme
\ref{algebra-lemma-descent-Rk}. Les conditions (1), (2) et (3) du
lemme \ref{lemma-descending-properties} sont donc satisfaites, ce qui conclut.
\end{proof}

\begin{lemma}
\label{lemma-regular-local-smooth}
La propriété $\mathcal{P}(S) =$``$S$ est régulier'' est locale pour la
topologie lisse.
\end{lemma}

\begin{proof}
Cela résulte du lemme \ref{lemma-Rk-local-smooth}, puisqu'un schéma localement
noethérien est régulier si et seulement s'il est localement noethérien et
vérifie $(R_k)$ pour tout $k \geq 0$.
\end{proof}

\begin{lemma}
\label{lemma-Nagata-local-smooth}
La propriété $\mathcal{P}(S) =$``$S$ est un schéma de Nagata'' est locale pour
la topologie lisse.
\end{lemma}

\begin{proof}
Nous allons vérifier les conditions (1), (2) et (3) du lemme
\ref{lemma-descending-properties}. Tout d'abord, la propriété d'être un schéma
de Nagata est locale pour la topologie de Zariski ; c'est le lemme
\ref{properties-lemma-locally-nagata} de Propriétés. Ensuite, si $S' \to S$ est un
morphisme lisse entre schémas affines et si $S$ est un schéma de Nagata,
alors $S'$ est un schéma de Nagata ; c'est le lemme
\ref{morphisms-lemma-finite-type-nagata} de Morphismes. Enfin, si $S' \to S$ est un
morphisme lisse surjectif entre schémas affines et si $S'$ est un schéma de
Nagata, alors $S$ est un schéma de Nagata ; c'est le lemme
\ref{algebra-lemma-descent-nagata} d'Algèbre. Les conditions (1), (2) et (3) du
lemme \ref{lemma-descending-properties} sont donc satisfaites, ce qui conclut.
\end{proof}





\section{Variantes sur la descente des propriétés}
\label{section-variants}

\noindent
Il est parfois possible de faire descendre des propriétés qui ne sont pas
locales. Nous réunissons dans cette section des résultats de ce type. Voir
aussi la section \ref{section-descent-finiteness-morphisms}, consacrée à la
descente de propriétés de morphismes, comme la lissité.

\begin{lemma}
\label{lemma-descend-reduced}
Si $f : X \to Y$ est un morphisme plat surjectif de schémas et si $X$ est
réduit, alors $Y$ est réduit.
\end{lemma}

\begin{proof}
Le résultat s'obtient en considérant les anneaux locaux (voir la définition
\ref{schemes-definition-reduced} de Schémas) et en appliquant le lemme
\ref{algebra-lemma-descent-reduced} d'Algèbre.
\end{proof}

\begin{lemma}
\label{lemma-descend-regular}
Soit $f : X \to Y$ un morphisme d'espaces algébriques. Si $f$ est localement
de présentation finie, plat et surjectif, et si $X$ est régulier, alors $Y$
est régulier.
\end{lemma}

\begin{proof}
Ce lemme se ramène à l'énoncé algébrique suivant : si $A \to B$ est un
homomorphisme d'anneaux fidèlement plat et de présentation finie, et si $B$
est noethérien et régulier, alors $A$ est noethérien et régulier. Nous voyons
que $A$ est noethérien d'après le lemme
\ref{algebra-lemma-descent-Noetherian} d'Algèbre, et régulier d'après le lemme
\ref{algebra-lemma-flat-under-regular} d'Algèbre.
\end{proof}









\section{Germes de schémas}
\label{section-germs}

\begin{definition}
\label{definition-germs}
Germes de schémas.
\begin{enumerate}
\item Un couple $(X, x)$ formé d'un schéma $X$ et d'un point $x \in X$ est
appelé le {\it germe de $X$ en $x$}.
\item Un {\it morphisme de germes} $f : (X, x) \to (S, s)$ est une classe
d'équivalence de morphismes de schémas $f : U \to S$ tels que $f(x) = s$,
où $U \subset X$ est un voisinage ouvert de $x$. Deux tels morphismes
$f$, $f'$ sont dits équivalents si et seulement si $f$ et $f'$ coïncident
sur un voisinage ouvert de $x$.
\item On définit la {\it composition des morphismes de germes} en composant
des représentants (cette opération est bien définie).
\end{enumerate}
\end{definition}

\noindent
Avant de poursuivre, il nous faut encore une définition.

\begin{definition}
\label{definition-etale-morphism-germs}
Soit $f : (X, x) \to (S, s)$ un morphisme de germes. On dit que $f$ est
{\it \'etale} (resp.\ {\it lisse}) s'il existe un représentant
$f : U \to S$ de $f$ qui soit un morphisme \'etale (resp.\ un morphisme
lisse) de schémas.
\end{definition}








\section{Propriétés locales des germes}
\label{section-properties-germs-local}

\begin{definition}
\label{definition-local-at-point}
Soit $\mathcal{P}$ une propriété des germes de schémas. On dit que
$\mathcal{P}$ est {\it locale pour la topologie \'etale}
(resp.\ {\it locale pour la topologie lisse}) si, pour tout morphisme de
germes \'etale (resp.\ lisse) $(U', u') \to (U, u)$, on a
$\mathcal{P}(U, u) \Leftrightarrow \mathcal{P}(U', u')$.
\end{definition}

\noindent
Soit $(X, x)$ un germe de schéma. La dimension de $X$ en $x$ est le minimum
des dimensions des voisinages ouverts de $x$ dans $X$, et tout voisinage
ouvert suffisamment petit possède cette dimension. C'est donc un invariant
de la classe d'isomorphisme du germe. On la note simplement $\dim_x(X)$.
Le lemme suivant affirme que la relation $\dim_x(X) = d$ est une propriété
locale des germes pour la topologie \'etale.

\begin{lemma}
\label{lemma-dimension-at-point-local}
Soit $f : U \to V$ un morphisme \'etale de schémas. Soient $u \in U$ et
$v = f(u)$. Alors $\dim_u(U) = \dim_v(V)$.
\end{lemma}

\begin{proof}
Dans l'énoncé, $\dim_u(U)$ désigne la dimension de $U$ en $u$, définie dans
la définition \ref{topology-definition-Krull} de Topologie, comme le minimum des
dimensions de Krull des voisinages ouverts de $u$ dans $U$. Il en va de même
pour $\dim_v(V)$.

\medskip\noindent
Montrons que $\dim_v(V) \geq \dim_u(U)$. Soit $V'$ un voisinage ouvert de
$v$ dans $V$. Il existe un voisinage ouvert $U'$ de $u$ dans $U$, contenu
dans $f^{-1}(V')$, tel que $\dim_u(U) = \dim(U')$. Soit
$Z_0 \subset Z_1 \subset \ldots \subset Z_n$ une chaîne de sous-schémas
fermés irréductibles de $U'$. Si $\xi_i \in Z_i$ est le point générique,
nous avons les spécialisations
$\xi_n \leadsto \xi_{n - 1} \leadsto \ldots \leadsto \xi_0$. Elles donnent
les spécialisations
$f(\xi_n) \leadsto f(\xi_{n - 1}) \leadsto \ldots \leadsto f(\xi_0)$ dans
$V'$. De plus, $f(\xi_j) \not = f(\xi_i)$ si $i \not = j$, car les fibres de
$f$ sont discrètes (voir le lemme
\ref{morphisms-lemma-etale-over-field} de Morphismes). Ainsi, $\dim(V') \geq n$.
L'inégalité $\dim_v(V) \geq \dim_u(U)$ en résulte formellement.

\medskip\noindent
Montrons que $\dim_u(U) \geq \dim_v(V)$. Soit $U'$ un voisinage ouvert de
$u$ dans $U$. L'ensemble $V' = f(U')$ est un voisinage ouvert de $v$ d'après
le lemme \ref{morphisms-lemma-fppf-open} de Morphismes. Ainsi,
$\dim(V') \geq \dim_v(V)$. Choisissons une chaîne
$Z_0 \subset Z_1 \subset \ldots \subset Z_n$ de sous-schémas fermés
irréductibles de $V'$. Si $\xi_i \in Z_i$ est le point générique, nous avons
les spécialisations
$\xi_n \leadsto \xi_{n - 1} \leadsto \ldots \leadsto \xi_0$. Comme
$\xi_0 \in f(U')$, il existe un point $\eta_0 \in U'$ tel que
$f(\eta_0) = \xi_0$. Considérons l'homomorphisme d'anneaux locaux
$$
\mathcal{O}_{V', \xi_0} \longrightarrow \mathcal{O}_{U', \eta_0}
$$
qui est plat et local d'après le lemme
\ref{morphisms-lemma-etale-flat} de Morphismes. Les points $\xi_i$ correspondent aux idéaux
premiers de l'anneau de gauche d'après le lemme
\ref{schemes-lemma-specialize-points} de Schémas. Le théorème de descente (voir la
section \ref{algebra-section-going-up} d'Algèbre) appliqué à l'homomorphisme affiché
fournit une suite de spécialisations
$\eta_n \leadsto \eta_{n - 1} \leadsto \ldots \leadsto \eta_0$ dans $U'$,
qui s'envoie sur la suite
$\xi_n \leadsto \xi_{n - 1} \leadsto \ldots \leadsto \xi_0$ par $f$.
Il s'ensuit que $\dim_u(U) \geq \dim_v(V)$.
\end{proof}

\noindent
Soit $(X, x)$ un germe de schéma. La classe d'isomorphisme de l'anneau local
$\mathcal{O}_{X, x}$ est un invariant du germe. Le lemme suivant affirme que
la propriété $\dim(\mathcal{O}_{X, x}) = d$ est locale pour la topologie
\'etale sur les germes.

\begin{lemma}
\label{lemma-dimension-local-ring-local}
Soit $f : U \to V$ un morphisme \'etale de schémas. Soient $u \in U$ et
$v = f(u)$. Alors
$\dim(\mathcal{O}_{U, u}) = \dim(\mathcal{O}_{V, v})$.
\end{lemma}

\begin{proof}
L'énoncé algébrique à démontrer est le suivant : si $A \to B$ est un
homomorphisme d'anneaux \'etale, si $\mathfrak q$ est un idéal premier de
$B$ situé au-dessus de $\mathfrak p \subset A$, alors
$\dim(A_{\mathfrak p}) = \dim(B_{\mathfrak q})$. C'est le lemme
\ref{more-algebra-lemma-dimension-etale-extension} de Compléments d'algèbre.
\end{proof}

\noindent
Soit $(X, x)$ un germe de schéma. La classe d'isomorphisme de l'anneau local
$\mathcal{O}_{X, x}$ est un invariant du germe. Le lemme suivant affirme que
la propriété ``$\mathcal{O}_{X, x}$ est régulier'' est locale pour la
topologie \'etale sur les germes.

\begin{lemma}
\label{lemma-regular-local-ring-local}
Soit $f : U \to V$ un morphisme \'etale de schémas. Soient $u \in U$ et
$v = f(u)$. Alors $\mathcal{O}_{U, u}$ est un anneau local régulier si et
seulement si $\mathcal{O}_{V, v}$ est un anneau local régulier.
\end{lemma}

\begin{proof}
L'énoncé algébrique à démontrer est le suivant : si $A \to B$ est un
homomorphisme d'anneaux \'etale, si $\mathfrak q$ est un idéal premier de
$B$ situé au-dessus de $\mathfrak p \subset A$, alors $A_{\mathfrak p}$ est
régulier si et seulement si $B_{\mathfrak q}$ est régulier. C'est le lemme
\ref{more-algebra-lemma-regular-etale-extension} de Compléments d'algèbre.
\end{proof}







\section{Propriétés des morphismes locales sur la cible}
\label{section-descending-properties-morphisms}

\noindent
Soit $f : X \to Y$ un morphisme de schémas. Soit $g : Y' \to Y$ un
morphisme de schémas. Notons $f' : X' \to Y'$ le changement de base de $f$
par $g$ :
$$
\xymatrix{
X' \ar[d]_{f'} \ar[r]_{g'} & X \ar[d]^f \\
Y' \ar[r]^g & Y
}
$$
Soit $\mathcal{P}$ une propriété des morphismes de schémas. On peut se
demander si (a) $\mathcal{P}(f) \Rightarrow \mathcal{P}(f')$ et si,
réciproquement, (b) $\mathcal{P}(f') \Rightarrow \mathcal{P}(f)$. Si (a) est
vérifiée chaque fois que $g$ est plat, on dit que $\mathcal{P}$ est stable
par changement de base plat. Si (b) est vérifiée chaque fois que $g$ est plat
et surjectif, on dit que $\mathcal{P}$ se descend par les changements de base
plats surjectifs. Si $\mathcal{P}$ est stable par changement de base plat et
se descend par les changements de base plats surjectifs, on dit que
$\mathcal{P}$ est locale sur la cible pour la topologie plate. Comparer avec
la discussion de la section \ref{section-descending-properties}. Cette notion
est très importante ; nous la formalisons dans la définition suivante.

\begin{definition}
\label{definition-property-morphisms-local}
Soit $\mathcal{P}$ une propriété des morphismes de schémas au-dessus d'une
base. Soit $\tau \in \{fpqc, fppf, syntomic, smooth, \etale, Zariski\}$.
On dit que $\mathcal{P}$ est {\it locale sur la base pour $\tau$}, ou
{\it locale sur la cible pour $\tau$}, ou encore {\it locale sur la base
pour la topologie $\tau$}, si, pour tout recouvrement pour la topologie
$\tau$, $\{Y_i \to Y\}_{i \in I}$ (voir la section
\ref{topologies-section-procedure} de Topologies), et pour tout morphisme de schémas
$f : X \to Y$ au-dessus de $S$, on a
$$
f \text{ vérifie }\mathcal{P}
\Leftrightarrow
\text{chaque }Y_i \times_Y X \to Y_i\text{ vérifie }\mathcal{P}.
$$
\end{definition}

\noindent
Bien entendu, puisque les isomorphismes sont toujours des recouvrements, on
constate (ou l'on exige) que la propriété $\mathcal{P}$ est vérifiée par
$X \to Y$ si et seulement si elle est vérifiée par toute flèche
$X' \to Y'$ isomorphe à $X \to Y$. Si une propriété est locale sur la cible
pour $\tau$, elle est stable par changement de base par les morphismes qui
interviennent dans les recouvrements pour $\tau$. Voici un énoncé formel.

\begin{lemma}
\label{lemma-pullback-property-local-target}
Soit $\tau \in \{fpqc, fppf, syntomic, smooth, \etale, Zariski\}$.
Soit $\mathcal{P}$ une propriété des morphismes locale sur la cible pour
$\tau$. Soit $f : X \to Y$ un morphisme vérifiant $\mathcal{P}$. Pour tout
morphisme $Y' \to Y$ qui est plat, resp.\ plat et localement de présentation
finie, resp.\ syntomique, resp.\ \'etale, resp.\ une immersion ouverte, le
changement de base $f' : Y' \times_Y X \to Y'$ de $f$ vérifie
$\mathcal{P}$.
\end{lemma}

\begin{proof}
Cela tient au fait que l'on peut compléter $Y' \to Y$ en une famille de
morphismes qui forme un recouvrement pour la topologie $\tau$.
\end{proof}

\noindent
Une conséquence simple et souvent utilisée est la suivante. Si
$f : X \to Y$ vérifie une propriété $\mathcal{P}$ locale sur la cible pour
$\tau$, et si $f(X) \subset V$ pour un sous-schéma ouvert $V \subset Y$,
alors le morphisme induit $X \to V$ vérifie lui aussi $\mathcal{P}$.
En effet, le changement de base de $f$ par $V \to Y$ est $X \to V$.

\begin{lemma}
\label{lemma-largest-open-of-the-base}
Soit $\tau \in \{fppf, syntomic, smooth, \etale\}$. Soit $\mathcal{P}$ une
propriété des morphismes locale sur la cible pour $\tau$. Pour tout morphisme
de schémas $f : X \to Y$, il existe un plus grand ouvert $W(f) \subset Y$
tel que la restriction $X_{W(f)} \to W(f)$ vérifie $\mathcal{P}$. En outre,
\begin{enumerate}
\item si $g : Y' \to Y$ est plat et localement de présentation finie,
syntomique, lisse ou \'etale, et si le changement de base
$f' : X_{Y'} \to Y'$ vérifie $\mathcal{P}$, alors $g(Y') \subset W(f)$,
\item si $g : Y' \to Y$ est plat et localement de présentation finie,
syntomique, lisse ou \'etale, alors $W(f') = g^{-1}(W(f))$, et
\item si $\{g_i : Y_i \to Y\}$ est un recouvrement pour la topologie $\tau$,
alors $g_i^{-1}(W(f)) = W(f_i)$, où $f_i$ est le changement de base de $f$
par $Y_i \to Y$.
\end{enumerate}
\end{lemma}

\begin{proof}
Considérons la réunion $W$ des images $g(Y') \subset Y$ des morphismes
$g : Y' \to Y$ ayant les propriétés suivantes :
\begin{enumerate}
\item $g$ est plat et localement de présentation finie, syntomique, lisse
ou \'etale, et
\item le changement de base $Y' \times_{g, Y} X \to Y'$ vérifie la propriété
$\mathcal{P}$.
\end{enumerate}
Un tel morphisme $g$ étant ouvert (voir le lemme
\ref{morphisms-lemma-fppf-open} de Morphismes), $W \subset Y$ est un ouvert de $Y$.
Puisque $\mathcal{P}$ est locale pour la topologie $\tau$, la restriction
$X_W \to W$ vérifie $\mathcal{P}$ : on dispose en effet d'un recouvrement
$\{Y' \to W\}$ de $W$ dont les changements de base vérifient
$\mathcal{P}$. Cela prouve l'existence et montre que $W(f)$ possède la
propriété (1). Pour (2), on a $W(f') \supset g^{-1}(W(f))$, car
$\mathcal{P}$ est stable par changement de base plat et localement de
présentation finie, syntomique, lisse ou \'etale ; voir le
lemme \ref{lemma-pullback-property-local-target}. D'autre part, si
$Y'' \subset Y'$ est un ouvert tel que $X_{Y''} \to Y''$ vérifie
$\mathcal{P}$, alors $Y'' \to Y$ se factorise par $W$ par construction,
c'est-à-dire $Y'' \subset g^{-1}(W(f))$. Cela prouve (2). L'assertion (3)
résulte de (2), puisque chaque morphisme $Y_i \to Y$ est plat et localement
de présentation finie, syntomique, lisse ou \'etale, par définition d'un
recouvrement pour la topologie $\tau$.
\end{proof}

\begin{lemma}
\label{lemma-descending-properties-morphisms}
Soit $\mathcal{P}$ une propriété des morphismes de schémas au-dessus d'une
base. Soit $\tau \in \{fpqc, fppf, \etale, lisse, syntomique\}$. Supposons que
\begin{enumerate}
\item la propriété soit stable par changement de base plat, plat et
localement de présentation finie, \'etale, lisse ou syntomique selon que
$\tau$ vaut fpqc, fppf, \'etale, smooth ou syntomic (comparer avec la
définition \ref{schemes-definition-preserved-by-base-change} de Schémas),
\item la propriété soit locale sur la base pour la topologie de Zariski,
\item pour tout morphisme surjectif de schémas affines $S' \to S$ qui est
plat, plat de présentation finie, \'etale, lisse ou syntomique selon que
$\tau$ vaut fpqc, fppf, \'etale, smooth ou syntomic, et pour tout morphisme de
schémas $f : X \to S$, la propriété $\mathcal{P}$ soit vérifiée par $f$ si
la propriété $\mathcal{P}$ est vérifiée par le changement de base
$f' : X' = S' \times_S X \to S'$.
\end{enumerate}
Alors $\mathcal{P}$ est locale sur la base pour la topologie $\tau$.
\end{lemma}

\begin{proof}
Cela résulte presque immédiatement de la définition d'un recouvrement pour
la topologie $\tau$ ; voir les définitions
\ref{topologies-definition-fpqc-covering},
\ref{topologies-definition-fppf-covering},
\ref{topologies-definition-etale-covering},
\ref{topologies-definition-smooth-covering} ou
\ref{topologies-definition-syntomic-covering} de Topologies, et les lemmes
\ref{topologies-lemma-fpqc-affine},
\ref{topologies-lemma-fppf-affine},
\ref{topologies-lemma-etale-affine},
\ref{topologies-lemma-smooth-affine} ou
\ref{topologies-lemma-syntomic-affine} de Topologies. Nous omettons les détails.
\end{proof}

\begin{remark}
\label{remark-descending-properties-morphisms-standard}
(Il s'agit d'une répétition de la remarque
\ref{remark-descending-properties-standard}.) Dans le
lemme \ref{lemma-descending-properties-morphisms}, si $\tau = smooth$, on
peut supposer, dans la condition (3), que le morphisme est un morphisme lisse
standard (surjectif). Il en va de même lorsque $\tau = syntomic$ ou
$\tau = \etale$.
\end{remark}




\section[Localité sur la cible (fpqc)]{Propriétés des morphismes locales sur la cible pour la topologie fpqc}
\label{section-descending-properties-morphisms-fpqc}

\noindent
Dans cette section, nous donnons un grand nombre de propriétés de morphismes
de schémas qui sont locales sur la base pour la topologie fpqc. En revanche,
nous montrerons dans la section
\ref{examples-section-non-descending-property-projective} d'Exemples que les propriétés
``projectif'' et ``quasi-projectif'' ne sont pas locales sur la base, même
pour la topologie de Zariski.

\begin{lemma}
\label{lemma-descending-property-quasi-compact}
La propriété $\mathcal{P}(f) =$``$f$ est quasi-compact'' est locale sur la
base pour la topologie fpqc.
\end{lemma}

\begin{proof}
Tout changement de base d'un morphisme quasi-compact est quasi-compact ; voir
le lemme \ref{schemes-lemma-quasi-compact-preserved-base-change} de Schémas.
La quasi-compacité est locale sur la base pour la topologie de Zariski ; voir
le lemme \ref{schemes-lemma-quasi-compact-affine} de Schémas. Enfin, soit
$S' \to S$ un morphisme plat surjectif de schémas affines, et soit
$f : X \to S$ un morphisme. Supposons que le changement de base
$f' : X' \to S'$ soit quasi-compact. Alors $X'$ est quasi-compact et
$X' \to X$ est surjectif. Par conséquent, $X$ est quasi-compact, donc $f$
est quasi-compact. Le lemme \ref{lemma-descending-properties-morphisms}
s'applique et conclut.
\end{proof}

\begin{lemma}
\label{lemma-descending-property-quasi-separated}
La propriété $\mathcal{P}(f) =$``$f$ est quasi-séparé'' est locale sur la
base pour la topologie fpqc.
\end{lemma}

\begin{proof}
Tout changement de base d'un morphisme quasi-séparé est quasi-séparé ; voir
le lemme \ref{schemes-lemma-separated-permanence} de Schémas. La quasi-séparation
est locale sur la base pour la topologie de Zariski (cela résulte de la
définition ou du lemme
\ref{schemes-lemma-characterize-quasi-separated} de Schémas). Enfin, soit $S' \to S$
un morphisme plat surjectif de schémas affines, et soit $f : X \to S$ un
morphisme. Supposons que le changement de base $f' : X' \to S'$ soit
quasi-séparé. Cela signifie que
$\Delta' : X' \to X'\times_{S'} X'$ est quasi-compact. Or $\Delta'$ est le
changement de base de $\Delta : X \to X \times_S X$ par $S' \to S$.
Le lemme \ref{lemma-descending-property-quasi-compact} montre donc que
$\Delta$ est quasi-compact et, par suite, que $f$ est quasi-séparé.
Le lemme \ref{lemma-descending-properties-morphisms} s'applique et conclut.
\end{proof}

\begin{lemma}
\label{lemma-descending-property-universally-closed}
La propriété $\mathcal{P}(f) =$``$f$ est universellement fermé'' est locale
sur la base pour la topologie fpqc.
\end{lemma}

\begin{proof}
Par définition, tout changement de base d'un morphisme universellement fermé
est universellement fermé. Cette propriété est locale sur la base pour la
topologie de Zariski (par définition ou d'après le lemme
\ref{morphisms-lemma-universally-closed-local-on-the-base} de Morphismes). Enfin, soit
$S' \to S$ un morphisme plat surjectif de schémas affines, et soit
$f : X \to S$ un morphisme. Supposons que le changement de base
$f' : X' \to S'$ soit universellement fermé. Soit $T \to S$ un morphisme
quelconque. Considérons le diagramme
$$
\xymatrix{
X' \ar[d] &
S' \times_S T \times_S X \ar[d] \ar[r] \ar[l] &
T \times_S X \ar[d] \\
S' &
S' \times_S T \ar[r] \ar[l] &
T
}
$$
dont les deux carrés sont cartésiens. L'hypothèse implique que la flèche
verticale médiane est fermée. Les flèches horizontales de droite sont plates,
quasi-compactes et surjectives, puisqu'elles s'obtiennent par changement de
base de $S' \to S$. Un sous-ensemble de $T$ est fermé si et seulement si son
image inverse dans $S' \times_S T$ est fermée ; voir le lemme
\ref{morphisms-lemma-fpqc-quotient-topology} de Morphismes. Une chasse au diagramme
immédiate montre que la flèche verticale de droite est également fermée ; on
en conclut que $X \to S$ est universellement fermé. Le lemme
\ref{lemma-descending-properties-morphisms} s'applique et conclut.
\end{proof}

\begin{lemma}
\label{lemma-descending-property-universally-open}
La propriété $\mathcal{P}(f) =$``$f$ est universellement ouvert'' est locale
sur la base pour la topologie fpqc.
\end{lemma}

\begin{proof}
La démonstration est identique à celle du
lemme \ref{lemma-descending-property-universally-closed}.
\end{proof}

\begin{lemma}
\label{lemma-descending-property-universally-submersive}
La propriété $\mathcal{P}(f) =$``$f$ est universellement submersif'' est
locale sur la base pour la topologie fpqc.
\end{lemma}

\begin{proof}
La démonstration est identique à celle du
lemme \ref{lemma-descending-property-universally-closed}, en utilisant qu'un
morphisme plat, quasi-compact et surjectif est universellement submersif
d'après le lemme \ref{morphisms-lemma-fpqc-quotient-topology} de Morphismes.
\end{proof}

\begin{lemma}
\label{lemma-descending-property-separated}
La propriété $\mathcal{P}(f) =$``$f$ est séparé'' est locale sur la base pour
la topologie fpqc.
\end{lemma}

\begin{proof}
Tout changement de base d'un morphisme séparé est séparé ; voir le lemme
\ref{schemes-lemma-separated-permanence} de Schémas. La séparation est locale sur
la base pour la topologie de Zariski (par définition ou d'après le lemme
\ref{schemes-lemma-characterize-separated} de Schémas). Enfin, soit $S' \to S$ un
morphisme plat surjectif de schémas affines, et soit $f : X \to S$ un
morphisme. Supposons que le changement de base $f' : X' \to S'$ soit séparé.
Cela signifie que $\Delta' : X' \to X'\times_{S'} X'$ est une immersion
fermée, donc universellement fermée. Or $\Delta'$ est le changement de base
de $\Delta : X \to X \times_S X$ par $S' \to S$. Le
lemme \ref{lemma-descending-property-universally-closed} montre donc que
$\Delta$ est universellement fermé. Comme c'est une immersion (Schémas,
lemme \ref{schemes-lemma-diagonal-immersion} de Schémas), $\Delta$ est une immersion
fermée. Ainsi $f$ est séparé. Le
lemme \ref{lemma-descending-properties-morphisms} s'applique et conclut.
\end{proof}

\begin{lemma}
\label{lemma-descending-property-surjective}
La propriété $\mathcal{P}(f) =$``$f$ est surjectif'' est locale sur la base
pour la topologie fpqc.
\end{lemma}

\begin{proof}
C'est immédiat.
\end{proof}

\begin{lemma}
\label{lemma-descending-property-dominant}
La propriété $\mathcal{P}(f) =$``$f$ est quasi-compact et dominant'' est
locale sur la base pour la topologie fpqc.
\end{lemma}

\begin{proof}
D'après le lemme \ref{morphisms-lemma-base-change-dominant} de Morphismes, les
morphismes quasi-compacts dominants sont stables par changement de base plat.
La réciproque est plus facile. En effet, la quasi-compacité est locale sur la
base pour la topologie fpqc d'après le
lemme \ref{lemma-descending-property-quasi-compact}. La dominance est
manifestement locale sur la base pour la topologie de Zariski. Enfin, soit
$S' \to S$ un morphisme surjectif de schémas (non nécessairement plat), et
soit $f : X \to S$ un morphisme. Supposons que le changement de base
$f' : X' \to S'$ soit dominant. Alors $X' \to S'\to S$ est le composé de
deux morphismes dominants, donc est dominant. Comme il s'agit aussi du composé
$X' \to X \to S$, il s'ensuit que $X\to S$ est dominant.
\end{proof}

\noindent
Les morphismes dominants ne sont pas stables par changement de base plat en
toute généralité, mais le cas particulier suivant mérite d'être signalé.
Voir aussi le lemme \ref{lemma-descending-fppf-property-dominant}.

\begin{lemma}
\label{lemma-descending-property-dominant-over-field}
Soit $E/k$ une extension de corps. Un morphisme $X \to Y$ au-dessus de $k$
est alors dominant si et seulement si le changement de base $X_E \to Y_E$
est dominant.
\end{lemma}

\begin{proof}
D'après le lemme
\ref{morphisms-lemma-scheme-over-field-universally-open} de Morphismes, le morphisme
$\Spec(E) \to \Spec(k)$ est universellement ouvert. Ainsi $Y_E \to Y$ est
ouvert. Par conséquent, si $X \to Y$ est dominant, le lemme
\ref{morphisms-lemma-open-base-change-dominant} de Morphismes montre que $X_E \to Y_E$
est dominant. Réciproquement, supposons que $X_E \to Y_E$ soit dominant.
Le morphisme $Y_E \to Y$ est surjectif d'après le lemme
\ref{morphisms-lemma-base-change-surjective} de Morphismes ; le composé
$X_E \to Y_E \to Y$ est donc dominant. Mais ce composé est aussi
$X_E \to X \to Y$ ; il s'ensuit que $X \to Y$ est dominant.
\end{proof}

\begin{lemma}
\label{lemma-descending-property-universally-injective}
La propriété $\mathcal{P}(f) =$``$f$ est universellement injectif'' est locale
sur la base pour la topologie fpqc.
\end{lemma}

\begin{proof}
Tout changement de base d'un morphisme universellement injectif est
universellement injectif (c'est formel). Cette propriété est locale sur la
base pour la topologie de Zariski, ce qui résulte immédiatement de la
définition. Enfin, soit $S' \to S$ un morphisme plat surjectif de schémas
affines, et soit $f : X \to S$ un morphisme. Supposons que le changement de
base $f' : X' \to S'$ soit universellement injectif. Soit $K$ un corps, et
soient $a, b : \Spec(K) \to X$ deux morphismes tels que
$f \circ a = f \circ b$. Puisque $S' \to S$ est surjectif, la discussion de
la section \ref{schemes-section-points} de Schémas fournit une extension de corps
$K'/K$ et un morphisme $\Spec(K')
\to S'$ tels que le diagramme suivant, formé de flèches pleines, soit commutatif :
$$
\xymatrix{
\Spec(K') \ar[rrd] \ar@{-->}[rd]_{a', b'} \ar[dd] \\
 &
X' \ar[r] \ar[d] &
S' \ar[d] \\
\Spec(K) \ar[r]^{a, b} &
X \ar[r] &
S
}
$$
Le carré étant cartésien, on obtient les deux flèches pointillées $a'$, $b'$
qui rendent le diagramme commutatif. Comme $X' \to S'$ est universellement
injectif, nous avons $a' = b'$ d'après le lemme
\ref{morphisms-lemma-universally-injective} de Morphismes. Cela impose évidemment $a = b$
(voir la discussion de la section \ref{schemes-section-points} de Schémas). Le
lemme \ref{lemma-descending-properties-morphisms} s'applique et conclut.

\medskip\noindent
On peut aussi utiliser la caractérisation d'un morphisme universellement
injectif par la surjectivité de sa diagonale ; voir le lemme
\ref{morphisms-lemma-universally-injective} de Morphismes. Le résultat découle alors du fait
que la surjectivité est locale sur la base pour la topologie fpqc ; voir le
lemme \ref{lemma-descending-property-surjective}. (Indication : utiliser que
le changement de base de la diagonale est la diagonale du changement de base.)
\end{proof}

\begin{lemma}
\label{lemma-descending-property-universal-homeomorphism}
La propriété $\mathcal{P}(f) =$``$f$ est un homéomorphisme universel'' est
locale sur la base pour la topologie fpqc.
\end{lemma}

\begin{proof}
On peut procéder exactement comme dans la démonstration du
lemme \ref{lemma-descending-property-universally-closed}. On peut aussi
observer qu'une application d'espaces topologiques est un homéomorphisme si
et seulement si elle est injective, surjective et ouverte. Un homéomorphisme
universel équivaut donc à un morphisme surjectif, universellement injectif et
universellement ouvert. Le résultat découle ainsi des lemmes
\ref{lemma-descending-property-surjective},
\ref{lemma-descending-property-universally-injective} et
\ref{lemma-descending-property-universally-open}.
\end{proof}

\begin{lemma}
\label{lemma-descending-property-locally-finite-type}
La propriété $\mathcal{P}(f) =$``$f$ est localement de type fini'' est locale
sur la base pour la topologie fpqc.
\end{lemma}

\begin{proof}
La propriété d'être localement de type fini est stable par changement de base ;
voir le lemme \ref{morphisms-lemma-base-change-finite-type} de Morphismes. Elle est
locale sur la base pour la topologie de Zariski ; voir le lemme
\ref{morphisms-lemma-locally-finite-type-characterize} de Morphismes. Enfin, soit
$S' \to S$ un morphisme plat surjectif de schémas affines, et soit
$f : X \to S$ un morphisme. Supposons que le changement de base
$f' : X' \to S'$ soit localement de type fini. Soit $U \subset X$ un ouvert
affine. Alors $U' = S' \times_S U$ est affine et de type fini sur $S'$.
Écrivons $S = \Spec(R)$, $S' = \Spec(R')$, $U = \Spec(A)$ et
$U' = \Spec(A')$. Nous savons que $R \to R'$ est fidèlement plat, que
$A' = R' \otimes_R A$ et que $R' \to A'$ est de type fini. Il faut montrer
que $R \to A$ est de type fini. C'est le lemme
\ref{algebra-lemma-finite-type-descends} d'Algèbre. Il s'ensuit que $f$ est localement
de type fini. Le lemme \ref{lemma-descending-properties-morphisms}
s'applique et conclut.
\end{proof}

\begin{lemma}
\label{lemma-descending-property-locally-finite-presentation}
La propriété $\mathcal{P}(f) =$``$f$ est localement de présentation finie''
est locale sur la base pour la topologie fpqc.
\end{lemma}

\begin{proof}
La propriété d'être localement de présentation finie est stable par
changement de base ; voir le lemme
\ref{morphisms-lemma-base-change-finite-presentation} de Morphismes. Elle est locale sur la
base pour la topologie de Zariski ; voir le lemme
\ref{morphisms-lemma-locally-finite-presentation-characterize} de Morphismes. Enfin, soit
$S' \to S$ un morphisme plat surjectif de schémas affines, et soit
$f : X \to S$ un morphisme. Supposons que le changement de base
$f' : X' \to S'$ soit localement de présentation finie. Soit $U \subset X$
un ouvert affine. Alors $U' = S' \times_S U$ est affine et de type fini sur
$S'$. Écrivons $S = \Spec(R)$, $S' = \Spec(R')$, $U = \Spec(A)$ et
$U' = \Spec(A')$. Nous savons que $R \to R'$ est fidèlement plat, que
$A' = R' \otimes_R A$ et que $R' \to A'$ est de présentation finie. Il faut
montrer que $R \to A$ est de présentation finie. C'est le lemme
\ref{algebra-lemma-finite-presentation-descends} d'Algèbre. Il s'ensuit que $f$ est
localement de présentation finie. Le
lemme \ref{lemma-descending-properties-morphisms} s'applique et conclut.
\end{proof}

\begin{lemma}
\label{lemma-descending-property-finite-type}
La propriété $\mathcal{P}(f) =$``$f$ est de type fini''
est locale sur la base pour la topologie fpqc.
\end{lemma}

\begin{proof}
Combiner les lemmes \ref{lemma-descending-property-quasi-compact}
et \ref{lemma-descending-property-locally-finite-type}.
\end{proof}

\begin{lemma}
\label{lemma-descending-property-finite-presentation}
La propriété $\mathcal{P}(f) =$``$f$ est de présentation finie''
est locale sur la base pour la topologie fpqc.
\end{lemma}

\begin{proof}
Combiner les lemmes \ref{lemma-descending-property-quasi-compact},
\ref{lemma-descending-property-quasi-separated} et
\ref{lemma-descending-property-locally-finite-presentation}.
\end{proof}

\begin{lemma}
\label{lemma-descending-property-proper}
La propriété $\mathcal{P}(f) =$``$f$ est propre''
est locale sur la base pour la topologie fpqc.
\end{lemma}

\begin{proof}
Le lemme résulte de la combinaison des
lemmes \ref{lemma-descending-property-universally-closed},
\ref{lemma-descending-property-separated}
et \ref{lemma-descending-property-finite-type}.
\end{proof}

\begin{lemma}
\label{lemma-descending-property-flat}
La propriété $\mathcal{P}(f) =$``$f$ est plat''
est locale sur la base pour la topologie fpqc.
\end{lemma}

\begin{proof}
La platitude est préservée par tout changement de base, voir
le lemme \ref{morphisms-lemma-base-change-flat} de Morphismes.
Elle est locale pour la topologie de Zariski sur la base par définition.
Enfin, soit
$S' \to S$ un morphisme plat surjectif de schémas affines,
et soit $f : X \to S$ un morphisme. Supposons que le changement de base
$f' : X' \to S'$ soit plat.
Soit $U \subset X$ un ouvert affine. Alors $U' = S' \times_S U$
est affine. Écrivons
$S = \Spec(R)$,
$S' = \Spec(R')$,
$U = \Spec(A)$ et
$U' = \Spec(A')$.
Nous savons que $R \to R'$ est fidèlement plat,
que $A' = R' \otimes_R A$ et que $R' \to A'$ est plat.
Il s'agit de montrer que $R \to A$ est plat.
Cela résulte immédiatement du lemme
\ref{algebra-lemma-flatness-descends} d'Algèbre.
Ainsi $f$ est plat.
Le lemme \ref{lemma-descending-properties-morphisms} s'applique donc et conclut.
\end{proof}

\begin{lemma}
\label{lemma-descending-property-open-immersion}
La propriété $\mathcal{P}(f) =$``$f$ est une immersion ouverte''
est locale sur la base pour la topologie fpqc.
\end{lemma}

\begin{proof}
La propriété d'être une immersion ouverte est stable par changement de base,
voir le lemme \ref{schemes-lemma-base-change-immersion} de Schémas.
Elle est locale pour la topologie de Zariski sur la base
(cela est évident).

\medskip\noindent
Soit $S' \to S$ un morphisme plat surjectif de schémas affines,
et soit $f : X \to S$ un morphisme. Supposons que le changement de base
$f' : X' \to S'$ soit une immersion ouverte. Nous affirmons que $f$ est une
immersion ouverte.
Alors $f'$ est universellement ouvert et universellement injectif.
Nous en déduisons que $f$ est universellement ouvert par le
lemme \ref{lemma-descending-property-universally-open}, et
universellement injectif par le
lemme \ref{lemma-descending-property-universally-injective}.
En particulier, $f(X) \subset S$ est ouvert. S'il est possible de montrer,
pour tout ouvert affine $U \subset f(X)$, que $f^{-1}(U) \to U$
est un isomorphisme, alors $f$ est une immersion ouverte et la preuve est
achevée. Si $U' \subset S'$ désigne l'image inverse de $U$,
alors $U' \to U$ est un morphisme fidèlement plat entre schémas affines et
$(f')^{-1}(U') \to U'$ est un isomorphisme (car $f'(X')$ contient $U'$
par le choix de $U$). Nous sommes ainsi ramenés au cas traité dans le
paragraphe suivant.

\medskip\noindent
Soit $S' \to S$ un morphisme plat surjectif de schémas affines,
soit $f : X \to S$ un morphisme, et supposons que le changement de base
$f' : X' \to S'$ soit un isomorphisme. Il faut montrer que $f$ est lui aussi
un isomorphisme. Il est clair que $f$ est surjectif, universellement injectif
et universellement ouvert (voir les arguments ci-dessus pour les deux
dernières propriétés). Ainsi $f$ est bijectif, c'est-à-dire que $f$ est un
homéomorphisme. Le morphisme $f$ est donc affine par le lemme
\ref{morphisms-lemma-homeomorphism-affine} de Morphismes. Comme
$$
\mathcal{O}(S') \to
\mathcal{O}(X') =
\mathcal{O}(S') \otimes_{\mathcal{O}(S)} \mathcal{O}(X)
$$
est un isomorphisme et que $\mathcal{O}(S) \to \mathcal{O}(S')$
est fidèlement plat, il en résulte que
$\mathcal{O}(S) \to \mathcal{O}(X)$ est un isomorphisme.
Ainsi $f$ est un isomorphisme. Cela démontre l'affirmation ci-dessus.
Le lemme \ref{lemma-descending-properties-morphisms} s'applique donc et conclut.
\end{proof}

\begin{lemma}
\label{lemma-descending-property-isomorphism}
La propriété $\mathcal{P}(f) =$``$f$ est un isomorphisme''
est locale sur la base pour la topologie fpqc.
\end{lemma}

\begin{proof}
Combiner les lemmes \ref{lemma-descending-property-surjective}
et \ref{lemma-descending-property-open-immersion}.
\end{proof}

\begin{lemma}
\label{lemma-descending-property-affine}
La propriété $\mathcal{P}(f) =$``$f$ est affine''
est locale sur la base pour la topologie fpqc.
\end{lemma}

\begin{proof}
Tout changement de base d'un morphisme affine est affine, voir
le lemme \ref{morphisms-lemma-base-change-affine} de Morphismes.
La propriété d'être affine est locale pour la topologie de Zariski sur la base,
voir le lemme \ref{morphisms-lemma-characterize-affine} de Morphismes.
Enfin, soit
$g : S' \to S$ un morphisme plat surjectif de schémas affines,
et soit $f : X \to S$ un morphisme. Supposons que le changement de base
$f' : X' \to S'$ soit affine. Autrement dit, $X'$ est affine, disons
$X' = \Spec(A')$. Écrivons aussi $S = \Spec(R)$
et $S' = \Spec(R')$. Il faut montrer que $X$ est affine.

\medskip\noindent
Les lemmes \ref{lemma-descending-property-quasi-compact}
et \ref{lemma-descending-property-separated} montrent que
$X \to S$ est séparé et quasi-compact. Par conséquent,
$f_*\mathcal{O}_X$ est un faisceau quasi-cohérent de $\mathcal{O}_S$-algèbres,
voir le lemme \ref{schemes-lemma-push-forward-quasi-coherent} de Schémas.
Ainsi $f_*\mathcal{O}_X = \widetilde{A}$ pour une certaine $R$-algèbre $A$.
Bien entendu, $A = \Gamma(X, \mathcal{O}_X)$.
De plus, par changement de base plat
(voir par exemple le lemme
\ref{coherent-lemma-flat-base-change-cohomology} de Cohomologie des schémas),
nous avons $g^*f_*\mathcal{O}_X = f'_*\mathcal{O}_{X'}$.
Autrement dit, $A' = R' \otimes_R A$.
Considérons le morphisme canonique
$$
X \longrightarrow \Spec(A)
$$
au-dessus de $S$ donné par le lemme
\ref{schemes-lemma-morphism-into-affine} de Schémas.
D'après ce qui précède, le changement de base de ce morphisme à $S'$ est un
isomorphisme. Il s'agit donc d'un isomorphisme par le
lemme \ref{lemma-descending-property-isomorphism}.
Le lemme \ref{lemma-descending-properties-morphisms} s'applique donc et conclut.
\end{proof}

\begin{lemma}
\label{lemma-descending-property-closed-immersion}
La propriété $\mathcal{P}(f) =$``$f$ est une immersion fermée''
est locale sur la base pour la topologie fpqc.
\end{lemma}

\begin{proof}
Soit $f : X \to Y$ un morphisme de schémas.
Soit $\{Y_i \to Y\}$ un recouvrement fpqc.
Supposons que chaque $f_i : Y_i \times_Y X \to Y_i$
soit une immersion fermée.
Cela implique que chaque $f_i$ est affine, voir
le lemme \ref{morphisms-lemma-closed-immersion-affine} de Morphismes.
Le lemme \ref{lemma-descending-property-affine}
permet d'en déduire que $f$ est affine. Il reste à montrer que
$\mathcal{O}_Y \to f_*\mathcal{O}_X$ est surjectif.
Pour tout $y \in Y$, il existe un $i$ et un point
$y_i \in Y_i$ d'image $y$.
Par le lemme \ref{coherent-lemma-flat-base-change-cohomology}
de Cohomologie des schémas,
le faisceau $f_{i, *}(\mathcal{O}_{Y_i \times_Y X})$
est l'image inverse de $f_*\mathcal{O}_X$.
Par hypothèse, c'est un quotient de $\mathcal{O}_{Y_i}$.
Par conséquent,
$$
\Big(
\mathcal{O}_{Y, y} \longrightarrow (f_*\mathcal{O}_X)_y
\Big)
\otimes_{\mathcal{O}_{Y, y}} \mathcal{O}_{Y_i, y_i}
$$
est surjectif. Puisque $\mathcal{O}_{Y_i, y_i}$ est fidèlement
plat sur $\mathcal{O}_{Y, y}$, il s'ensuit que
$\mathcal{O}_{Y, y} \longrightarrow (f_*\mathcal{O}_X)_y$ est surjectif,
comme voulu.
\end{proof}

\begin{lemma}
\label{lemma-descending-property-quasi-affine}
La propriété $\mathcal{P}(f) =$``$f$ est quasi-affine''
est locale sur la base pour la topologie fpqc.
\end{lemma}

\begin{proof}
Soit $f : X \to Y$ un morphisme de schémas.
Soit $\{g_i : Y_i \to Y\}$ un recouvrement fpqc.
Supposons que chaque $f_i : Y_i \times_Y X \to Y_i$
soit quasi-affine.
Alors chaque $f_i$ est quasi-compact et séparé.
Les lemmes \ref{lemma-descending-property-quasi-compact}
et \ref{lemma-descending-property-separated}
impliquent donc que $f$ est quasi-compact et séparé.
Considérons le faisceau de $\mathcal{O}_Y$-algèbres
$\mathcal{A} = f_*\mathcal{O}_X$.
Par le lemme \ref{schemes-lemma-push-forward-quasi-coherent} de Schémas,
c'est une $\mathcal{O}_Y$-algèbre quasi-cohérente.
Considérons le morphisme canonique
$$
j : X \longrightarrow \underline{\Spec}_Y(\mathcal{A})
$$
voir le lemme \ref{constructions-lemma-canonical-morphism} de Constructions.
Par changement de base plat
(voir par exemple le lemme
\ref{coherent-lemma-flat-base-change-cohomology} de Cohomologie des schémas),
nous avons $g_i^*f_*\mathcal{O}_X = f_{i, *}\mathcal{O}_{X'}$,
où les $g_i : Y_i \to Y$ sont les morphismes plats donnés.
Ainsi, le changement de base $j_i$ de $j$ par $g_i$ est le morphisme canonique
du lemme \ref{constructions-lemma-canonical-morphism} de Constructions,
pour le morphisme $f_i$. Par hypothèse et par
le lemme \ref{morphisms-lemma-characterize-quasi-affine} de Morphismes,
tous ces morphismes $j_i$ sont des immersions ouvertes quasi-compactes.
Les lemmes \ref{lemma-descending-property-quasi-compact} et
\ref{lemma-descending-property-open-immersion} montrent donc que
$j$ est une immersion ouverte quasi-compacte.
Une nouvelle application du
lemme \ref{morphisms-lemma-characterize-quasi-affine} de Morphismes
montre alors que $f$ est quasi-affine.
\end{proof}

\begin{lemma}
\label{lemma-descending-property-quasi-compact-immersion}
La propriété $\mathcal{P}(f) =$``$f$ est une immersion quasi-compacte''
est locale sur la base pour la topologie fpqc.
\end{lemma}

\begin{proof}
Soit $f : X \to Y$ un morphisme de schémas.
Soit $\{Y_i \to Y\}$ un recouvrement fpqc.
Notons $X_i = Y_i \times_Y X$ et $f_i : X_i \to Y_i$
le changement de base de $f$. Notons aussi
$q_i : Y_i \to Y$ les morphismes plats donnés.
Supposons que chaque $f_i$ soit une immersion quasi-compacte.
Par le lemme \ref{schemes-lemma-immersions-monomorphisms} de Schémas,
chaque $f_i$ est séparé.
Les lemmes \ref{lemma-descending-property-quasi-compact} et
\ref{lemma-descending-property-separated}
impliquent donc que $f$ est quasi-compact et séparé.
Soit $X \to Z \to Y$ la factorisation de $f$ par son
image schématique. Par le
lemme \ref{morphisms-lemma-quasi-compact-scheme-theoretic-image} de Morphismes,
le sous-schéma fermé $Z \subset Y$ est défini par le
faisceau quasi-cohérent d'idéaux
$\mathcal{I} = \Ker(\mathcal{O}_Y \to f_*\mathcal{O}_X)$,
car $f$ est quasi-compact. Par changement de base plat
(voir par exemple le lemme
\ref{coherent-lemma-flat-base-change-cohomology} de Cohomologie des schémas ;
nous utilisons ici le fait que $f$ est séparé),
nous voyons que $f_{i, *}\mathcal{O}_{X_i}$ est l'image inverse
$q_i^*f_*\mathcal{O}_X$.
Ainsi, $Y_i \times_Y Z$ est défini par le
faisceau quasi-cohérent d'idéaux $q_i^*\mathcal{I} =
\Ker(\mathcal{O}_{Y_i} \to f_{i, *}\mathcal{O}_{X_i})$.
Par le lemme \ref{morphisms-lemma-quasi-compact-immersion} de Morphismes,
les morphismes $X_i \to Y_i \times_Y Z$
sont des immersions ouvertes. Le
lemme \ref{lemma-descending-property-open-immersion}
montre alors que $X \to Z$ est une immersion ouverte, et donc que
$f$ est une immersion, comme voulu
(nous avons déjà vu qu'il était quasi-compact).
\end{proof}

\begin{lemma}
\label{lemma-descending-property-integral}
La propriété $\mathcal{P}(f) =$``$f$ est entier''
est locale sur la base pour la topologie fpqc.
\end{lemma}

\begin{proof}
Un morphisme entier est la même chose qu'un morphisme affine et
universellement fermé. Voir
le lemme \ref{morphisms-lemma-integral-universally-closed} de Morphismes.
Le lemme résulte donc de la combinaison des
lemmes \ref{lemma-descending-property-universally-closed}
et \ref{lemma-descending-property-affine}.
\end{proof}

\begin{lemma}
\label{lemma-descending-property-finite}
La propriété $\mathcal{P}(f) =$``$f$ est fini''
est locale sur la base pour la topologie fpqc.
\end{lemma}

\begin{proof}
Un morphisme fini est la même chose qu'un morphisme entier
qui est localement de type fini. Voir
le lemme \ref{morphisms-lemma-finite-integral} de Morphismes.
Le lemme résulte donc de la combinaison des
lemmes \ref{lemma-descending-property-locally-finite-type}
et \ref{lemma-descending-property-integral}.
\end{proof}

\begin{lemma}
\label{lemma-descending-property-quasi-finite}
Les propriétés
$\mathcal{P}(f) =$``$f$ est localement quasi-fini''
et
$\mathcal{P}(f) =$``$f$ est quasi-fini''
sont locales sur la base pour la topologie fpqc.
\end{lemma}

\begin{proof}
Soit $f : X \to S$ un morphisme de schémas, et soit $\{S_i \to S\}$
un recouvrement fpqc tel que chaque changement de base
$f_i : X_i \to S_i$ soit localement quasi-fini.
Nous avons déjà vu
(lemme \ref{lemma-descending-property-locally-finite-type})
que la propriété ``être localement de type fini'' est locale sur la base pour la
topologie fpqc ; ainsi, $f$ est localement de type fini.
Il résulte alors du
lemme \ref{morphisms-lemma-base-change-quasi-finite} de Morphismes
que $f$ est localement quasi-fini. Le cas quasi-fini en découle,
puisque nous avons déjà vu que la propriété ``être quasi-compact'' est locale
sur la base pour la topologie fpqc
(lemme \ref{lemma-descending-property-quasi-compact}).
\end{proof}

\begin{lemma}
\label{lemma-descending-property-relative-dimension-d}
La propriété $\mathcal{P}(f) =$``$f$ est localement de type fini
de dimension relative $d$'' est locale sur la base pour la topologie fpqc.
\end{lemma}

\begin{proof}
Cela résulte immédiatement du fait que la propriété d'être localement de type
fini est locale sur la base pour la topologie fpqc, et du
lemme \ref{morphisms-lemma-dimension-fibre-after-base-change} de Morphismes.
\end{proof}

\begin{lemma}
\label{lemma-descending-property-syntomic}
La propriété $\mathcal{P}(f) =$``$f$ est syntomique''
est locale sur la base pour la topologie fpqc.
\end{lemma}

\begin{proof}
Un morphisme est syntomique si et seulement s'il est localement de présentation
finie, plat, et si ses fibres sont localement d'intersection complète.
Nous avons déjà vu que la platitude et la propriété d'être localement de
présentation finie sont locales sur la base pour la topologie fpqc (lemmes
\ref{lemma-descending-property-flat} et
\ref{lemma-descending-property-locally-finite-presentation}).
Le résultat pour les morphismes syntomiques découle donc du
lemme \ref{morphisms-lemma-set-points-where-fibres-lci} de Morphismes.
\end{proof}

\begin{lemma}
\label{lemma-descending-property-smooth}
La propriété $\mathcal{P}(f) =$``$f$ est lisse''
est locale sur la base pour la topologie fpqc.
\end{lemma}

\begin{proof}
Un morphisme est lisse si et seulement s'il est localement de présentation
finie, plat, et si ses fibres sont lisses. Nous avons déjà vu que la platitude
et la propriété d'être localement de présentation finie sont locales sur la
base pour la topologie fpqc (lemmes
\ref{lemma-descending-property-flat} et
\ref{lemma-descending-property-locally-finite-presentation}).
Le résultat pour les morphismes lisses découle donc du
lemme \ref{morphisms-lemma-set-points-where-fibres-smooth} de Morphismes.
\end{proof}

\begin{lemma}
\label{lemma-descending-property-unramified}
La propriété $\mathcal{P}(f) =$``$f$ est non ramifié''
est locale sur la base pour la topologie fpqc.
La propriété $\mathcal{P}(f) =$``$f$ est G-non ramifié''
est locale sur la base pour la topologie fpqc.
\end{lemma}

\begin{proof}
Un morphisme est non ramifié (resp.\ G-non ramifié) si et seulement s'il est
localement de type fini (resp.\ de présentation finie)
et si son morphisme diagonal est une immersion ouverte (voir
le lemme \ref{morphisms-lemma-diagonal-unramified-morphism} de Morphismes).
Nous avons déjà vu que les propriétés d'être localement de type fini
(resp.\ localement de présentation finie) et d'être une immersion ouverte sont
locales sur la base pour la topologie fpqc (lemmes
\ref{lemma-descending-property-locally-finite-presentation},
\ref{lemma-descending-property-locally-finite-type} et
\ref{lemma-descending-property-open-immersion}).
Le résultat en découle formellement.
\end{proof}

\begin{lemma}
\label{lemma-descending-property-etale}
La propriété $\mathcal{P}(f) =$``$f$ est \'etale''
est locale sur la base pour la topologie fpqc.
\end{lemma}

\begin{proof}
Un morphisme est \'etale si et seulement s'il est plat et G-non ramifié.
Voir le lemme \ref{morphisms-lemma-flat-unramified-etale} de Morphismes.
Nous avons déjà vu que la platitude et la propriété d'être G-non ramifié
sont locales sur la base pour la topologie fpqc (lemmes
\ref{lemma-descending-property-flat} et
\ref{lemma-descending-property-unramified}).
Le résultat en découle.
\end{proof}

\begin{lemma}
\label{lemma-descending-property-finite-locally-free}
La propriété $\mathcal{P}(f) =$``$f$ est fini localement libre''
est locale sur la base pour la topologie fpqc.
Soit $d \geq 0$.
La propriété $\mathcal{P}(f) =$``$f$ est fini localement libre de degré $d$''
est locale sur la base pour la topologie fpqc.
\end{lemma}

\begin{proof}
La propriété d'être fini localement libre équivaut à être
fini, plat et localement de présentation finie
(voir le lemme \ref{morphisms-lemma-finite-flat} de Morphismes).
Cela résulte donc des lemmes
\ref{lemma-descending-property-finite},
\ref{lemma-descending-property-flat} et
\ref{lemma-descending-property-locally-finite-presentation}.
Si $f : Z \to U$ est fini localement libre, et si $\{U_i \to U\}$ est une
famille surjective de morphismes telle que chaque image inverse
$Z \times_U U_i \to U_i$ soit de degré $d$, alors $Z \to U$ est de degré $d$;
on peut par exemple lire le degré en un point $u \in U$ sur la fibre
$(f_*\mathcal{O}_Z)_u \otimes_{\mathcal{O}_{U, u}} \kappa(u)$.
\end{proof}

\begin{lemma}
\label{lemma-descending-property-monomorphism}
La propriété $\mathcal{P}(f) =$``$f$ est un monomorphisme''
est locale sur la base pour la topologie fpqc.
\end{lemma}

\begin{proof}
Soit $f : X \to S$ un morphisme de schémas.
Soit $\{S_i \to S\}$ un recouvrement fpqc, et supposons que chacun des
changements de base $f_i : X_i \to S_i$ de $f$ soit un monomorphisme.
Soient $a, b : T \to X$ deux morphismes tels que
$f \circ a = f \circ b$. Il faut montrer que $a = b$.
Puisque $f_i$ est un monomorphisme, nous avons $a_i = b_i$, où
$a_i, b_i : S_i \times_S T \to X_i$ sont les changements de base.
En particulier, les composés
$S_i \times_S T \to T \to X$ sont égaux.
Comme $\coprod S_i \times_S T \to T$
est un épimorphisme (voir
par ex.\ le lemme \ref{lemma-fpqc-universal-effective-epimorphisms}),
nous concluons que $a = b$.
\end{proof}

\begin{lemma}
\label{lemma-descending-property-regular-immersion}
Les propriétés
\begin{enumerate}
\item[] $\mathcal{P}(f) =$``$f$ est une immersion Koszul-régulière'',
\item[] $\mathcal{P}(f) =$``$f$ est une immersion $H_1$-régulière'', et
\item[] $\mathcal{P}(f) =$``$f$ est une immersion quasi-régulière''
\end{enumerate}
sont locales sur la base pour la topologie fpqc.
\end{lemma}

\begin{proof}
Nous allons utiliser le critère du
lemme \ref{lemma-descending-properties-morphisms}.
Par la définition \ref{divisors-definition-regular-immersion} de Diviseurs,
la propriété d'être une immersion Koszul-régulière
(resp.\ $H_1$-régulière, quasi-régulière) est locale pour la topologie de
Zariski sur la base. Par le
lemme \ref{divisors-lemma-flat-base-change-regular-immersion} de Diviseurs,
la propriété d'être une immersion Koszul-régulière
(resp.\ $H_1$-régulière, quasi-régulière)
est préservée par changement de base plat.
La dernière hypothèse (3) du
lemme \ref{lemma-descending-properties-morphisms}
se traduit par l'énoncé algébrique suivant :
soit $A \to B$ un homomorphisme fidèlement plat d'anneaux.
Soit $I \subset A$ un idéal.
Si $IB$ est localement sur $\Spec(B)$ engendré par une suite Koszul-régulière
(resp.\ $H_1$-régulière, quasi-régulière) de $B$, alors $I \subset A$
est localement sur $\Spec(A)$ engendré par une suite Koszul-régulière
(resp.\ $H_1$-régulière, quasi-régulière) de $A$. C'est le
lemme \ref{more-algebra-lemma-flat-descent-regular-ideal}
de Compléments d'algèbre.
\end{proof}





\section{Propriétés des morphismes locales sur la cible pour la topologie fppf}
\label{section-descending-properties-morphisms-fppf}

\noindent
Dans cette section, nous déterminons certaines propriétés des morphismes de
schémas dont nous n'avons pas (encore) pu montrer qu'elles sont locales sur la
base pour la topologie fpqc, mais qui le sont pour la topologie fppf.

\begin{lemma}
\label{lemma-descending-fppf-property-immersion}
La propriété $\mathcal{P}(f) =$``$f$ est une immersion''
est locale sur la base pour la topologie fppf.
\end{lemma}

\begin{proof}
La propriété d'être une immersion est stable par changement de base,
voir le lemme \ref{schemes-lemma-base-change-immersion} de Schémas.
Elle est locale pour la topologie de Zariski sur la base.
Enfin, soit
$\pi : S' \to S$ un morphisme surjectif de schémas affines,
plat et localement de présentation finie.
Notons que $\pi : S' \to S$ est ouvert par
le lemme \ref{morphisms-lemma-fppf-open} de Morphismes.
Soit $f : X \to S$ un morphisme.
Supposons que le changement de base $f' : X' \to S'$ soit une immersion.
En particulier, $f'(X') = \pi^{-1}(f(X))$ est localement fermé.
Par le lemme \ref{topology-lemma-open-morphism-quotient-topology} de Topologie,
il s'ensuit que $f(X) \subset S$
est localement fermé. Soit $Z \subset S$
le fermé $Z = \overline{f(X)} \setminus f(X)$.
Une nouvelle application du lemme
\ref{topology-lemma-open-morphism-quotient-topology} de Topologie
montre que $f'(X')$ est fermé dans $S' \setminus Z'$.
Nous pouvons donc appliquer le
lemme \ref{lemma-descending-property-closed-immersion}
au recouvrement fpqc $\{S' \setminus Z' \to S \setminus Z\}$
et conclure que $f : X \to S \setminus Z$ est une immersion fermée.
Autrement dit, $f$ est une immersion.
Le lemme \ref{lemma-descending-properties-morphisms} s'applique donc et conclut.
\end{proof}

\begin{lemma}
\label{lemma-descending-fppf-property-dominant}
La propriété $\mathcal{P}(f) =$``$f$ est dominant''
est locale sur la base pour la topologie fppf.
\end{lemma}

\begin{proof}
Par le lemme \ref{morphisms-lemma-open-base-change-dominant} de Morphismes,
les morphismes dominants sont préservés
par changement de base par des morphismes ouverts, donc
par des morphismes plats localement de présentation finie
(voir le lemme \ref{morphisms-lemma-fppf-open} de Morphismes).
La réciproque est plus facile.
En effet, la propriété d'être dominant est manifestement locale pour la
topologie de Zariski sur la base.
Soit ensuite $S' \to S$ un morphisme surjectif de schémas
(pas nécessairement plat ni localement de présentation finie),
et soit $f : X \to S$ un morphisme. Supposons que le changement de base
$f' : X' \to S'$ soit dominant. Alors $X' \to S'\to S$ est un composé
de deux morphismes dominants, donc est dominant. Comme il s'agit aussi du
composé $X' \to X \to S$, il s'ensuit que $X\to S$ est dominant.
\end{proof}






\section[Application de la descente fpqc]{Application de la descente fpqc des propriétés des morphismes}
\label{section-application-descending-properties-morphisms}

\noindent
Le lemme suivant peut sembler quelque peu anodin, mais il constitue un outil
utile pour étudier les propriétés, pour un morphisme, d'être \'etale ou non
ramifié.

\begin{lemma}
\label{lemma-flat-surjective-quasi-compact-monomorphism-isomorphism}
Soit $f : X \to Y$ un monomorphisme plat, quasi-compact et surjectif.
Alors ce morphisme est un isomorphisme.
\end{lemma}

\begin{proof}
Puisque $f$ est un morphisme plat, quasi-compact et surjectif,
$\{X \to Y\}$ est un recouvrement fpqc de $Y$.
La diagonale $\Delta : X \to X \times_Y X$ est un isomorphisme
(voir le lemme \ref{schemes-lemma-monomorphism} de Schémas).
Il s'ensuit que le changement de base de $f$ par $f$ est un
isomorphisme. Ainsi, $f$ est un isomorphisme par le
lemme \ref{lemma-descending-property-isomorphism}.
\end{proof}

\noindent
Ce lemme permet de démontrer le résultat important suivant ; nous donnons aussi
une démonstration qui évite la descente fpqc.
Nous discuterons plus en détail de ce résultat et de résultats connexes dans
la section \ref{etale-section-topological-etale} du chapitre Morphismes étales.

\begin{lemma}
\label{lemma-universally-injective-etale-open-immersion}
Tout morphisme universellement injectif et \'etale est une immersion ouverte.
\end{lemma}

\begin{proof}[Première démonstration]
Soit $f : X \to Y$ un morphisme \'etale universellement injectif.
Alors $f$ est ouvert
(voir le lemme \ref{morphisms-lemma-etale-open} de Morphismes) ;
nous pouvons donc remplacer $Y$ par $f(X)$ et supposer $f$ surjectif.
Le morphisme $f$ est alors bijectif et ouvert, donc est un homéomorphisme.
En particulier, $f$ est quasi-compact. Par le
lemme \ref{lemma-flat-surjective-quasi-compact-monomorphism-isomorphism},
il suffit donc de montrer que $f$ est un monomorphisme.
Comme $X \to Y$ est \'etale, le morphisme
$\Delta_{X/Y} : X \to X \times_Y X$ est une immersion ouverte par
le lemme \ref{morphisms-lemma-diagonal-unramified-morphism} de Morphismes
(et le lemme \ref{morphisms-lemma-flat-unramified-etale} de Morphismes).
Comme $f$ est universellement injectif, $\Delta_{X/Y}$ est aussi surjectif,
voir le lemme \ref{morphisms-lemma-universally-injective} de Morphismes.
Ainsi, $\Delta_{X/Y}$ est un isomorphisme, c'est-à-dire que
$X \to Y$ est un monomorphisme.
\end{proof}

\begin{proof}[Seconde démonstration]
Soit $f : X \to Y$ un morphisme \'etale universellement injectif.
Alors $f$ est ouvert (voir le lemme \ref{morphisms-lemma-etale-open} de Morphismes) ;
nous pouvons donc remplacer $Y$ par $f(X)$ et supposer $f$ surjectif.
Les hypothèses restant satisfaites après tout changement de base, nous en
concluons que $f$ est un homéomorphisme universel. Par conséquent, $f$ est
entier, voir le
lemme \ref{morphisms-lemma-universal-homeomorphism} de Morphismes.
Il s'ensuit que $f$ est fini par
le lemme \ref{morphisms-lemma-finite-integral} de Morphismes.
Il s'ensuit que $f$ est fini localement libre par
le lemme \ref{morphisms-lemma-finite-flat} de Morphismes.
Pour achever la démonstration, il suffit de voir que $f$ est fini localement
libre de degré $1$ (un morphisme fini localement libre de degré $1$
est un isomorphisme).
Il existe une décomposition de $Y$ en sous-schémas ouverts et fermés
$V_d$ telle que $f^{-1}(V_d) \to V_d$ soit fini localement libre de
degré $d$, voir le lemme \ref{morphisms-lemma-finite-locally-free} de Morphismes.
Si $V_d$ est non vide, nous pouvons choisir un morphisme
$\Spec(k) \to V_d \subset Y$, où $k$ est un corps algébriquement clos
(il suffit de prendre la clôture algébrique du corps résiduel d'un point de
$V_d$).
Alors $\Spec(k) \times_Y X \to \Spec(k)$ est une réunion disjointe de
copies de $\Spec(k)$, par le
lemme \ref{morphisms-lemma-etale-over-field} de Morphismes,
et parce que $k$ est algébriquement clos.
Mais, comme $f$ est universellement injectif, il ne peut y avoir qu'une seule
copie; ainsi $d = 1$, comme voulu.
\end{proof}

\noindent
La caractérisation suivante des morphismes plats universellement injectifs
permet de reformuler légèrement les hypothèses du lemme précédent.

\begin{lemma}
\label{lemma-flat-universally-injective}
Soit $f : X \to Y$ un morphisme de schémas. Notons $X^0$ l'ensemble
des points génériques des composantes irréductibles de $X$. Si
\begin{enumerate}
\item $f$ est plat et séparé,
\item pour $\xi \in X^0$, nous avons $\kappa(f(\xi)) = \kappa(\xi)$, et
\item si $\xi, \xi' \in X^0$, $\xi \not = \xi'$, alors
$f(\xi) \not = f(\xi')$,
\end{enumerate}
alors $f$ est universellement injectif.
\end{lemma}

\begin{proof}
Il faut montrer que $\Delta : X \to X \times_Y X$ est surjectif, voir le
lemme \ref{morphisms-lemma-universally-injective} de Morphismes.
Comme $X \to Y$ est séparé, l'image de $\Delta$ est fermée.
Ainsi, si $\Delta$ n'est pas surjectif, il existe un point générique
$\eta \in X \times_S X$ d'une composante irréductible de $X \times_S X$
qui n'appartient pas à l'image de $\Delta$. La projection
$\text{pr}_1 : X \times_Y X \to X$
est plate, car elle se déduit du morphisme plat $X \to Y$ par changement de
base, voir le lemme \ref{morphisms-lemma-base-change-flat} de Morphismes.
Les générisations se relèvent donc le long de $\text{pr}_1$, voir le
lemme \ref{morphisms-lemma-generalizations-lift-flat} de Morphismes.
Nous en concluons que $\xi = \text{pr}_1(\eta) \in X^0$.
Cependant, les hypothèses (2) et (3) garantissent que le schéma
$(X \times_Y X)_{f(\xi)}$ possède au plus un point pour tout $\xi \in X^0$.
Autrement dit, $\Delta(\xi) = \eta$, contradiction.
\end{proof}

\noindent
Nous pouvons donc reformuler le
lemme \ref{lemma-universally-injective-etale-open-immersion} comme suit.

\begin{lemma}
\label{lemma-characterize-open-immersion}
Soit $f : X \to Y$ un morphisme de schémas. Notons $X^0$ l'ensemble
des points génériques des composantes irréductibles de $X$. Si
\begin{enumerate}
\item $f$ est \'etale et séparé,
\item pour $\xi \in X^0$, nous avons $\kappa(f(\xi)) = \kappa(\xi)$, et
\item si $\xi, \xi' \in X^0$, $\xi \not = \xi'$, alors
$f(\xi) \not = f(\xi')$,
\end{enumerate}
alors $f$ est une immersion ouverte.
\end{lemma}

\begin{proof}
Cela résulte immédiatement des lemmes
\ref{lemma-flat-universally-injective} et
\ref{lemma-universally-injective-etale-open-immersion}.
\end{proof}

\begin{lemma}
\label{lemma-descending-property-proper-over-base}
Soit $f : X \to Y$ un morphisme de schémas localement de type fini.
Soit $Z$ un fermé de $X$. S'il existe un recouvrement fpqc
$\{Y_i \to Y\}$ tel que l'image inverse $Z_i \subset Y_i \times_Y X$
soit propre sur $Y_i$
(voir la définition \ref{coherent-definition-proper-over-base}
de Cohomologie des schémas),
alors $Z$ est propre sur $Y$.
\end{lemma}

\begin{proof}
Munissons $Z$ de la structure réduite induite de sous-schéma fermé, voir la
définition \ref{schemes-definition-reduced-induced-scheme} de Schémas.
Pour tout $i$, le changement de base $Y_i \times_Y Z$ est un sous-schéma
fermé de $Y_i \times_Y X$ dont le fermé sous-jacent est $Z_i$.
Par définition (au moyen du lemme
\ref{coherent-lemma-closed-proper-over-base} de Cohomologie des schémas),
nous en concluons que les projections $Y_i \times_Y Z \to Y_i$ sont des
morphismes propres. Le morphisme $Z \to Y$ est donc propre par le
lemme \ref{lemma-descending-property-proper}.
Ainsi, $Z$ est propre sur $Y$ par définition.
\end{proof}

\begin{lemma}
\label{lemma-descending-property-ample}
Soit $f : X \to S$ un morphisme de schémas.
Soit $\mathcal{L}$ un $\mathcal{O}_X$-module inversible.
Soit $\{g_i : S_i \to S\}_{i \in I}$ un recouvrement fpqc.
Soit $f_i : X_i \to S_i$ le changement de base de $f$, et soit
$\mathcal{L}_i$ l'image inverse de $\mathcal{L}$ sur $X_i$.
Les assertions suivantes sont équivalentes :
\begin{enumerate}
\item $\mathcal{L}$ est ample sur $X/S$, et
\item $\mathcal{L}_i$ est ample sur $X_i/S_i$
pour tout $i \in I$.
\end{enumerate}
\end{lemma}

\begin{proof}
L'implication (1) $\Rightarrow$ (2) résulte du
lemme \ref{morphisms-lemma-ample-base-change} de Morphismes.
Supposons que $\mathcal{L}_i$ soit ample sur $X_i/S_i$ pour tout $i \in I$.
Par la définition \ref{morphisms-definition-relatively-ample} de Morphismes,
cela implique que $X_i \to S_i$ est quasi-compact et, par
le lemme \ref{morphisms-lemma-relatively-ample-separated} de Morphismes,
que $X_i \to S$ est séparé.
Ainsi, $f$ est quasi-compact et séparé par les
lemmes \ref{lemma-descending-property-quasi-compact} et
\ref{lemma-descending-property-separated}.

\medskip\noindent
Cela signifie que
$\mathcal{A} = \bigoplus_{d \geq 0} f_*\mathcal{L}^{\otimes d}$
est une $\mathcal{O}_S$-algèbre graduée quasi-cohérente
(voir le lemme \ref{schemes-lemma-push-forward-quasi-coherent} de Schémas).
De plus, la formation de $\mathcal{A}$ commute aux changements de base plats par
le lemme \ref{coherent-lemma-flat-base-change-cohomology}
de Cohomologie des schémas.
En particulier, si nous posons
$\mathcal{A}_i = \bigoplus_{d \geq 0} f_{i, *}\mathcal{L}_i^{\otimes d}$,
alors $\mathcal{A}_i = g_i^*\mathcal{A}$.
Il s'ensuit que les morphismes naturels
$\psi_d : f^*\mathcal{A}_d \to \mathcal{L}^{\otimes d}$
de $\mathcal{O}_X$-modules donnent par image inverse les morphismes naturels
$\psi_{i, d} : f_i^*(\mathcal{A}_i)_d \to \mathcal{L}_i^{\otimes d}$
de $\mathcal{O}_{X_i}$-modules. Comme $\mathcal{L}_i$ est ample sur $X_i/S_i$,
pour tout point $x_i \in X_i$, il existe un $d \geq 1$ tel que
$f_i^*(\mathcal{A}_i)_d \to \mathcal{L}_i^{\otimes d}$
soit surjectif sur les germes en $x_i$. Cela résulte soit directement
de la définition d'un module relativement ample, soit de
le lemme \ref{morphisms-lemma-characterize-relatively-ample} de Morphismes.
Si $x \in X$, nous pouvons choisir un $i$ et un $x_i \in X_i$
d'image $x$. Comme $\mathcal{O}_{X, x} \to \mathcal{O}_{X_i, x_i}$
est plat, donc fidèlement plat, nous en concluons que pour tout $x \in X$,
il existe un $d \geq 1$ tel que
$f^*\mathcal{A}_d \to \mathcal{L}^{\otimes d}$
soit surjectif sur les germes en $x$.
Cela implique que l'ouvert $U(\psi) \subset X$ du
lemme \ref{constructions-lemma-invertible-map-into-relative-proj} de Constructions,
associé au morphisme
$\psi : f^*\mathcal{A} \to \bigoplus_{d \geq 0} \mathcal{L}^{\otimes d}$
d'algèbres graduées sur $\mathcal{O}_X$,
est égal à $X$. Considérons le morphisme correspondant
$$
r_{\mathcal{L}, \psi} : X \longrightarrow \underline{\text{Proj}}_S(\mathcal{A})
$$
Il résulte clairement de ce qui précède que le changement de base de
$r_{\mathcal{L}, \psi}$ à $S_i$ est le morphisme
$r_{\mathcal{L}_i, \psi_i}$, qui est une immersion ouverte par le
lemme \ref{morphisms-lemma-characterize-relatively-ample} de Morphismes.
Ainsi, $r_{\mathcal{L}, \psi}$ est une immersion ouverte par le
lemme \ref{lemma-descending-property-open-immersion},
et nous concluons que $\mathcal{L}$ est ample sur $X/S$ par
le lemme \ref{morphisms-lemma-characterize-relatively-ample} de Morphismes.
\end{proof}















\section{Propriétés des morphismes locales sur la source}
\label{section-properties-morphisms-local-source}

\noindent
Il arrive souvent que l'on puisse établir qu'un morphisme possède une certaine
propriété après l'avoir précomposé par un autre morphisme. Dans de nombreux cas,
cela implique que le morphisme initial possède lui aussi cette propriété.
Nous formalisons cette situation dans la définition suivante.

\begin{definition}
\label{definition-property-morphisms-local-source}
Soit $\mathcal{P}$ une propriété de morphismes de schémas.
Soit $\tau \in \{Zariski, \linebreak[0] fpqc, \linebreak[0] fppf, \linebreak[0]
\etale, \linebreak[0] smooth, \linebreak[0] syntomic\}$.
On dit que $\mathcal{P}$ est
{\it locale sur la source pour $\tau$}, ou
{\it locale sur la source pour la topologie $\tau$}, si pour tout
morphisme de schémas $f : X \to Y$ au-dessus de $S$, et pour tout
recouvrement pour la topologie $\tau$, $\{X_i \to X\}_{i \in I}$, on a
$$
f \text{ vérifie }\mathcal{P}
\Leftrightarrow
\text{chaque }X_i \to Y\text{ vérifie }\mathcal{P}.
$$
\end{definition}

\noindent
Bien entendu, puisque les isomorphismes sont toujours des recouvrements,
nous constatons (ou exigeons) que la propriété $\mathcal{P}$ est vérifiée par
$X \to Y$ si et seulement si elle l'est par toute flèche $X' \to Y'$
isomorphe à $X \to Y$. Si une propriété est locale sur la source pour $\tau$,
elle est préservée par précomposition avec les morphismes qui interviennent
dans les recouvrements pour la topologie $\tau$. Voici un énoncé précis.

\begin{lemma}
\label{lemma-precompose-property-local-source}
Soit $\tau \in \{fpqc, fppf, syntomic, smooth, \etale, Zariski\}$.
Soit $\mathcal{P}$ une propriété de morphismes locale sur la source pour
$\tau$. Soit $f : X \to Y$ vérifiant la propriété $\mathcal{P}$.
Pour tout morphisme $a : X' \to X$ qui est
plat, resp.\ plat et localement de présentation finie, resp.\ syntomique,
resp.\ \'etale, resp.\ une immersion ouverte, le composé
$f \circ a : X' \to Y$ vérifie la propriété $\mathcal{P}$.
\end{lemma}

\begin{proof}
Cela résulte du fait que l'on peut compléter $X' \to X$ en une famille de
morphismes qui forme un $\tau$-recouvrement.
\end{proof}

\begin{lemma}
\label{lemma-largest-open-of-the-source}
Soit $\tau \in \{fppf, syntomic, smooth, \etale\}$.
Soit $\mathcal{P}$ une propriété de morphismes locale sur la source pour
$\tau$. Pour tout morphisme de schémas $f : X \to Y$, il existe
un plus grand ouvert $W(f) \subset X$ tel que la restriction
$f|_{W(f)} : W(f) \to Y$ vérifie $\mathcal{P}$. De plus,
si $g : X' \to X$ est plat et localement de présentation finie,
syntomique, lisse ou \'etale, et si $f' = f \circ g : X' \to Y$, alors
$g^{-1}(W(f)) = W(f')$.
\end{lemma}

\begin{proof}
Considérons la réunion $W$ des images $g(X') \subset X$ des
morphismes $g : X' \to X$ qui vérifient les propriétés suivantes :
\begin{enumerate}
\item $g$ est plat et localement de présentation finie, syntomique,
lisse ou \'etale, et
\item le composé $X' \to X \to Y$ vérifie la propriété $\mathcal{P}$.
\end{enumerate}
Comme tout tel morphisme $g$ est ouvert (voir le
lemme \ref{morphisms-lemma-fppf-open} de Morphismes),
$W \subset X$ est un ouvert de $X$. Puisque $\mathcal{P}$ est locale pour la
topologie $\tau$, la restriction $f|_W : W \to Y$ vérifie la propriété
$\mathcal{P}$ : nous disposons en effet d'un recouvrement pour la topologie $\tau$,
$\{X' \to W\}$ de $W$ dont les images inverses vérifient $\mathcal{P}$.
Cela démontre l'existence de $W(f)$.
La compatibilité énoncée dans la dernière phrase résulte immédiatement de la
construction de $W(f)$.
\end{proof}

\begin{lemma}
\label{lemma-properties-morphisms-local-source}
Soit $\mathcal{P}$ une propriété de morphismes de schémas.
Soit $\tau \in \{fpqc, \linebreak[0] fppf, \linebreak[0]
\etale, \linebreak[0] smooth, \linebreak[0] syntomic\}$.
Supposons que
\begin{enumerate}
\item la propriété soit préservée par précomposition avec un morphisme
plat, plat localement de présentation finie, \'etale, lisse ou syntomique
selon que $\tau$ est fpqc, fppf, \'etale, lisse ou syntomique,
\item la propriété soit locale sur la source pour la topologie de Zariski,
\item la propriété soit locale sur la cible pour la topologie de Zariski,
\item pour tout morphisme de schémas affines $f : X \to Y$, et pour tout
morphisme surjectif de schémas affines $X' \to X$ qui est plat, plat de
présentation finie, \'etale, lisse ou syntomique selon que $\tau$ est
fpqc, fppf, \'etale, lisse ou syntomique, la propriété $\mathcal{P}$ soit
vérifiée par $f$ si la propriété $\mathcal{P}$ est vérifiée par le composé
$f' : X' \to Y$.
\end{enumerate}
Alors $\mathcal{P}$ est locale sur la source pour la topologie $\tau$.
\end{lemma}

\begin{proof}
Cela résulte presque immédiatement de la définition d'un
recouvrement pour la topologie $\tau$ ; voir les définitions
\ref{topologies-definition-fpqc-covering}
\ref{topologies-definition-fppf-covering}
\ref{topologies-definition-etale-covering}
\ref{topologies-definition-smooth-covering}, ou
\ref{topologies-definition-syntomic-covering} de Topologies,
et les lemmes
\ref{topologies-lemma-fpqc-affine},
\ref{topologies-lemma-fppf-affine},
\ref{topologies-lemma-etale-affine},
\ref{topologies-lemma-smooth-affine}, ou
\ref{topologies-lemma-syntomic-affine} de Topologies.
Détails omis. (Indication : utiliser la localité sur la source et sur le but
pour ramener la vérification de la propriété $\mathcal{P}$ au cas d'un
morphisme entre schémas affines. Appliquer ensuite (1) et (4).)
\end{proof}

\begin{remark}
\label{remark-properties-morphisms-local-source-standard}
(Cette remarque reprend les remarques
\ref{remark-descending-properties-standard}
et \ref{remark-descending-properties-morphisms-standard} ci-dessus.)
Dans le lemme \ref{lemma-properties-morphisms-local-source}, si
$\tau = smooth$, on peut supposer, dans la condition (4), que le morphisme est
un morphisme lisse standard (surjectif).
De même lorsque $\tau = syntomic$ ou $\tau = \etale$.
\end{remark}



\section[Localité sur la source (fpqc)]{Propriétés des morphismes locales sur la source pour la topologie fpqc}
\label{section-fpqc-local-source}

\noindent
Voici quelques propriétés de morphismes qui sont locales sur la source pour la
topologie fpqc.

\begin{lemma}
\label{lemma-flat-fpqc-local-source}
La propriété $\mathcal{P}(f)=$``$f$ est plat'' est locale sur la source pour la
topologie fpqc.
\end{lemma}

\begin{proof}
Puisque la platitude se définit au moyen des homomorphismes d'anneaux locaux
(voir la définition \ref{morphisms-definition-flat} de Morphismes),
il faut établir le fait algébrique suivant : supposons que
$A \to B \to C$ soient des homomorphismes locaux d'anneaux locaux, et que
$B \to C$ soit plat. Alors $A \to B$ est plat si et seulement si
$A \to C$ est plat.
Si $A \to B$ est plat, alors $A \to C$ est plat par
le lemme \ref{algebra-lemma-composition-flat} d'Algèbre.
Réciproquement, supposons $A \to C$ plat.
Notons que $B \to C$ est fidèlement plat, voir
le lemme \ref{algebra-lemma-local-flat-ff} d'Algèbre.
Ainsi, $A \to B$ est plat par
le lemme \ref{algebra-lemma-flat-permanence} d'Algèbre.
(Voir aussi le lemme \ref{morphisms-lemma-flat-permanence} de Morphismes
pour une démonstration directe.)
\end{proof}

\begin{lemma}
\label{lemma-injective-local-rings-fpqc-local-source}
La propriété
$\mathcal{P}(f : X \to Y)=$``pour tout $x \in X$, l'homomorphisme d'anneaux
locaux $\mathcal{O}_{Y, f(x)} \to \mathcal{O}_{X, x}$ est injectif''
est locale sur la source pour la topologie fpqc.
\end{lemma}

\begin{proof}
Omis. Il s'agit seulement d'une tentative (probablement malavisée) de badiner.
\end{proof}





\section{Propriétés des morphismes locales sur la source pour la topologie fppf}
\label{section-fppf-local-source}

\noindent
Voici quelques propriétés de morphismes qui sont locales sur la source pour la
topologie fppf.

\begin{lemma}
\label{lemma-locally-finite-presentation-fppf-local-source}
La propriété $\mathcal{P}(f)=$``$f$ est localement de présentation finie''
est locale sur la source pour la topologie fppf.
\end{lemma}

\begin{proof}
La propriété d'être localement de présentation finie est locale sur la source
et sur la cible pour la topologie de Zariski, voir le
lemme \ref{morphisms-lemma-locally-finite-presentation-characterize} de Morphismes.
Elle est préservée par composition, voir
le lemme \ref{morphisms-lemma-composition-finite-presentation} de Morphismes.
Cela démontre (1), (2) et (3) du
lemme \ref{lemma-properties-morphisms-local-source}.
La dernière condition (4) est le
lemme \ref{lemma-flat-finitely-presented-permanence-algebra}.
Le résultat en découle.
\end{proof}

\begin{lemma}
\label{lemma-locally-finite-type-fppf-local-source}
La propriété $\mathcal{P}(f)=$``$f$ est localement de type fini''
est locale sur la source pour la topologie fppf.
\end{lemma}

\begin{proof}
La propriété d'être localement de type fini est locale sur la source et sur la
cible pour la topologie de Zariski, voir le
lemme \ref{morphisms-lemma-locally-finite-type-characterize} de Morphismes.
Elle est préservée par composition, voir
le lemme \ref{morphisms-lemma-composition-finite-type} de Morphismes,
et tout morphisme plat localement de présentation finie est localement de type
fini, voir
le lemme \ref{morphisms-lemma-finite-presentation-finite-type} de Morphismes.
Cela démontre (1), (2) et (3) du
lemme \ref{lemma-properties-morphisms-local-source}.
La dernière condition (4) est le
lemme \ref{lemma-finite-type-local-source-fppf-algebra}.
Le résultat en découle.
\end{proof}

\begin{lemma}
\label{lemma-open-fppf-local-source}
La propriété $\mathcal{P}(f)=$``$f$ est ouvert''
est locale sur la source pour la topologie fppf.
\end{lemma}

\begin{proof}
La propriété d'être un morphisme ouvert est manifestement locale sur la source
et sur la cible pour la topologie de Zariski.
Elle est préservée par composition, voir
le lemme \ref{morphisms-lemma-composition-open} de Morphismes,
et tout morphisme plat de présentation finie est ouvert, voir
le lemme \ref{morphisms-lemma-fppf-open} de Morphismes.
Cela démontre (1), (2) et (3) du
lemme \ref{lemma-properties-morphisms-local-source}.
La dernière condition (4) résulte de
lemme \ref{morphisms-lemma-fpqc-quotient-topology} de Morphismes.
Le résultat en découle.
\end{proof}

\begin{lemma}
\label{lemma-universally-open-fppf-local-source}
La propriété $\mathcal{P}(f)=$``$f$ est universellement ouvert''
est locale sur la source pour la topologie fppf.
\end{lemma}

\begin{proof}
Soit $f : X \to Y$ un morphisme de schémas.
Soit $\{X_i \to X\}_{i \in I}$ un recouvrement fppf.
Notons $f_i : X_i \to X$ les composés. Il faut montrer que $f$ est
universellement ouvert si et seulement si chaque $f_i$ l'est.
Si $f$ est universellement ouvert, alors chaque $f_i$ l'est aussi, puisque les
morphismes $X_i \to X$ sont universellement ouverts et que les composés de
morphismes universellement ouverts sont universellement ouverts
(voir les lemmes \ref{morphisms-lemma-fppf-open}
et \ref{morphisms-lemma-composition-open} de Morphismes).
Réciproquement, supposons chaque $f_i$ universellement ouvert.
Soit $Y' \to Y$ un morphisme de schémas.
Notons $X' = Y' \times_Y X$ et $X'_i = Y' \times_Y X_i$.
Le système $\{X_i' \to X'\}_{i \in I}$ est encore un recouvrement fppf.
Les morphismes $f'_i : X_i' \to Y'$ sont ouverts par hypothèse.
Le lemme \ref{lemma-open-fppf-local-source} montre donc que
$f' : X' \to Y'$ est ouvert, comme voulu.
\end{proof}



\section{Propriétés des morphismes locales sur la source pour la topologie syntomique}
\label{section-syntomic-local-source}

\noindent
Voici quelques propriétés de morphismes qui sont locales sur la source pour la
topologie syntomique.

\begin{lemma}
\label{lemma-syntomic-syntomic-local-source}
La propriété $\mathcal{P}(f)=$``$f$ est syntomique''
est locale sur la source pour la topologie syntomique.
\end{lemma}

\begin{proof}
Combiner le lemme \ref{lemma-properties-morphisms-local-source} avec le
lemme \ref{morphisms-lemma-syntomic-characterize} de Morphismes
(localité sur la source et sur la cible pour la topologie de Zariski), le
lemme \ref{morphisms-lemma-composition-syntomic} de Morphismes
(précomposition), et le
lemme \ref{lemma-syntomic-smooth-etale-permanence} (partie (4)).
\end{proof}




\section{Propriétés des morphismes locales sur la source pour la topologie lisse}
\label{section-smooth-local-source}

\noindent
Voici quelques propriétés de morphismes qui sont locales sur la source pour la
topologie lisse. Notons aussi le résultat (à certains égards plus fort)
sur la descente de la lissité le long de morphismes plats,
lemme \ref{lemma-smooth-permanence}.

\begin{lemma}
\label{lemma-smooth-smooth-local-source}
La propriété $\mathcal{P}(f)=$``$f$ est lisse''
est locale sur la source pour la topologie lisse.
\end{lemma}

\begin{proof}
Combiner le lemme \ref{lemma-properties-morphisms-local-source} avec le
lemme \ref{morphisms-lemma-smooth-characterize} de Morphismes
(localité sur la source et sur la cible pour la topologie de Zariski), le
lemme \ref{morphisms-lemma-composition-smooth} de Morphismes
(précomposition), et le
lemme \ref{lemma-syntomic-smooth-etale-permanence} (partie (4)).
\end{proof}



\section[Localité sur la source (\'etale)]{Propriétés des morphismes locales sur la source pour la topologie \'etale}
\label{section-etale-local-source}

\noindent
Voici quelques propriétés de morphismes qui sont locales sur la source pour la
topologie \'etale.

\begin{lemma}
\label{lemma-etale-etale-local-source}
La propriété $\mathcal{P}(f)=$``$f$ est \'etale''
est locale sur la source pour la topologie \'etale.
\end{lemma}

\begin{proof}
Combiner le lemme \ref{lemma-properties-morphisms-local-source} avec le
lemme \ref{morphisms-lemma-etale-characterize} de Morphismes
(localité sur la source et sur la cible pour la topologie de Zariski), le
lemme \ref{morphisms-lemma-composition-etale} de Morphismes
(précomposition), et le
lemme \ref{lemma-syntomic-smooth-etale-permanence} (partie (4)).
\end{proof}

\begin{lemma}
\label{lemma-locally-quasi-finite-etale-local-source}
La propriété $\mathcal{P}(f)=$``$f$ est localement quasi-fini''
est locale sur la source pour la topologie \'etale.
\end{lemma}

\begin{proof}
Nous allons utiliser le
lemme \ref{lemma-properties-morphisms-local-source}.
Par le lemme \ref{morphisms-lemma-locally-quasi-finite-characterize} de
Morphismes, la propriété d'être localement quasi-fini est locale sur la source
et sur la cible pour la topologie de Zariski. Par les lemmes
\ref{morphisms-lemma-composition-quasi-finite} et
\ref{morphisms-lemma-etale-locally-quasi-finite} de Morphismes,
la précomposition d'un morphisme localement quasi-fini par un morphisme
\'etale est localement quasi-finie.
Enfin, supposons que $X \to Y$ soit un morphisme de schémas affines et que
$X' \to X$ soit un morphisme surjectif \'etale de schémas affines tel que
$X' \to Y$ soit localement quasi-fini. Alors $X' \to Y$ est de type fini, et
le lemme \ref{lemma-finite-type-local-source-fppf-algebra}
montre que $X \to Y$ est lui aussi de type fini.
De plus, par hypothèse, $X' \to Y$ a des fibres finies ; il en va donc de même
de $X \to Y$. Nous concluons que $X \to Y$ est quasi-fini par
le lemme \ref{morphisms-lemma-quasi-finite} de Morphismes.
Cela démontre la dernière hypothèse du
lemme \ref{lemma-properties-morphisms-local-source}
et achève la démonstration.
\end{proof}

\begin{lemma}
\label{lemma-unramified-etale-local-source}
La propriété $\mathcal{P}(f)=$``$f$ est non ramifié''
est locale sur la source pour la topologie \'etale.
La propriété $\mathcal{P}(f)=$``$f$ est G-non ramifié''
est locale sur la source pour la topologie \'etale.
\end{lemma}

\begin{proof}
Nous allons utiliser le
lemme \ref{lemma-properties-morphisms-local-source}.
Par le lemme \ref{morphisms-lemma-unramified-characterize} de Morphismes,
la propriété d'être non ramifié (resp.\ G-non ramifié)
est locale sur la source et sur la cible pour la topologie de Zariski. Par les
lemmes \ref{morphisms-lemma-composition-unramified} et
\ref{morphisms-lemma-etale-smooth-unramified} de Morphismes,
la précomposition d'un morphisme non ramifié (resp.\ G-non ramifié)
par un morphisme \'etale donne un morphisme non ramifié
(resp.\ G-non ramifié).
Enfin, supposons que $X \to Y$ soit un morphisme de schémas affines et que
$f : X' \to X$ soit un morphisme surjectif \'etale de schémas affines tel que
$X' \to Y$ soit non ramifié (resp.\ G-non ramifié).
Alors $X' \to Y$ est de type fini (resp.\ de présentation finie), et le
lemme \ref{lemma-finite-type-local-source-fppf-algebra}
(resp.\ le lemme \ref{lemma-flat-finitely-presented-permanence-algebra})
montre que $X \to Y$ est lui aussi de type fini
(resp.\ de présentation finie). Par
le lemme \ref{morphisms-lemma-triangle-differentials-smooth} de Morphismes,
nous avons une suite exacte courte
$$
0 \to f^*\Omega_{X/Y} \to \Omega_{X'/Y} \to \Omega_{X'/X} \to 0.
$$
Comme $X' \to Y$ est non ramifié, le terme du milieu est nul.
Puisque $f$ est fidèlement plat, il s'ensuit que $\Omega_{X/Y} = 0$.
Ainsi, $X \to Y$ est non ramifié (resp.\ G-non ramifié), voir le
lemme \ref{morphisms-lemma-unramified-omega-zero} de Morphismes.
Cela démontre la dernière hypothèse du
lemme \ref{lemma-properties-morphisms-local-source}
et achève la démonstration.
\end{proof}




\section[Localité source-cible (\'etale)]{Propriétés des morphismes locales sur la source et la cible pour la topologie \'etale}
\label{section-properties-etale-local-source-target}

\noindent
Soit $\mathcal{P}$ une propriété de morphismes de schémas. La phrase
``$\mathcal{P}$ est locale sur la source et la cible pour la topologie \'etale''
a un sens intuitif. Il s'avère toutefois que cette notion ne revient pas à
demander que $\mathcal{P}$ soit à la fois locale sur la source et locale sur
la cible pour la topologie \'etale.
Avant d'approfondir cette question, donnons deux exemples volontairement
élémentaires.

\begin{example}
\label{example-silly-one}
Considérons la propriété $\mathcal{P}$ des morphismes de schémas définie par
la règle $\mathcal{P}(X \to Y) = $``$Y$ est localement noethérien''.
Le lecteur vérifiera qu'elle est locale sur la source et locale sur la cible
pour la topologie \'etale (démonstration omise, voir le
lemme \ref{lemma-Noetherian-local-fppf}).
Mais il n'est {\bf pas} vrai que, si $f : X \to Y$ vérifie $\mathcal{P}$
et si $g : Y \to Z$ est \'etale, alors $g \circ f$ vérifie $\mathcal{P}$.
En effet, $f$ pourrait être l'identité de $Y$ et $g$ une immersion ouverte
d'un schéma localement noethérien $Y$ dans un schéma non localement
noethérien $Z$.
\end{example}

\noindent
L'exemple suivant est, en un sens, encore pire.

\begin{example}
\label{example-silly-two}
Considérons la propriété $\mathcal{P}$ des morphismes de schémas définie par
la règle $\mathcal{P}(f : X \to Y) = $``pour tout $y \in Y$ qui est une
spécialisation d'un certain $f(x)$, $x \in X$, l'anneau local
$\mathcal{O}_{Y, y}$ est noethérien''.
Vérifions qu'elle est locale sur la source et sur la cible pour la topologie
\'etale. Nous utiliserons librement le
lemme \ref{schemes-lemma-specialize-points} de Schémas.

\medskip\noindent
Localité sur la cible : soit $\{g_i : Y_i \to Y\}$ un recouvrement pour la
topologie \'etale.
Soit $f_i : X_i \to Y_i$ le changement de base de $f$, et notons
$h_i : X_i \to X$ la projection. Supposons $\mathcal{P}(f)$.
Soit $f(x_i) \leadsto y_i$ une spécialisation.
Alors $f(h_i(x_i)) \leadsto g_i(y_i)$; ainsi, $\mathcal{P}(f)$ implique que
$\mathcal{O}_{Y, g_i(y_i)}$ est noethérien.
De plus, $\mathcal{O}_{Y, g_i(y_i)} \to \mathcal{O}_{Y_i, y_i}$ est une
localisation d'un homomorphisme \'etale d'anneaux.
Par conséquent, $\mathcal{O}_{Y_i, y_i}$ est noethérien par le
lemme \ref{algebra-lemma-Noetherian-permanence} d'Algèbre.
Réciproquement, supposons $\mathcal{P}(f_i)$ pour tout $i$.
Soit $f(x) \leadsto y$ une spécialisation. Choisissons un $i$ et un
$y_i \in Y_i$ d'image $y$.
Comme $x$ peut être considéré comme un point de
$\Spec(\mathcal{O}_{Y, y}) \times_Y X$ et que
$\mathcal{O}_{Y, y} \to \mathcal{O}_{Y_i, y_i}$ est fidèlement plat,
il existe un point
$x_i \in \Spec(\mathcal{O}_{Y_i, y_i}) \times_Y X$
d'image $x$. Alors $x_i \in X_i$, et $f_i(x_i)$ se spécialise en $y_i$.
Ainsi, nous voyons que $\mathcal{O}_{Y_i, y_i}$ est noethérien par
$\mathcal{P}(f_i)$, ce qui implique que $\mathcal{O}_{Y, y}$ est noethérien par
le lemme \ref{algebra-lemma-descent-Noetherian} d'Algèbre.

\medskip\noindent
Localité sur la source : soit $\{h_i : X_i \to X\}$ un recouvrement pour la
topologie \'etale.
Soit $f_i : X_i \to Y$ le composé $f \circ h_i$.
Supposons $\mathcal{P}(f)$. Soit $f(x_i) \leadsto y$ une spécialisation.
Alors $f(h_i(x_i)) \leadsto y$; ainsi, $\mathcal{P}(f)$ implique que
$\mathcal{O}_{Y, y}$ est noethérien. Par conséquent, $\mathcal{P}(f_i)$ est
vérifiée. Réciproquement, supposons $\mathcal{P}(f_i)$ pour tout $i$.
Soit $f(x) \leadsto y$ une spécialisation. Choisissons un $i$ et un
$x_i \in X_i$ d'image $x$.
Alors $y$ est une spécialisation de $f_i(x_i) = f(x)$.
Par conséquent, $\mathcal{P}(f_i)$ implique que $\mathcal{O}_{Y, y}$ est
noethérien, comme voulu.

\medskip\noindent
Nous affirmons qu'il existe un diagramme commutatif
$$
\xymatrix{
U \ar[d]_a \ar[r]_h & V \ar[d]^b \\
X \ar[r]^f & Y
}
$$
dont les flèches verticales sont surjectives, chacune étant \'etale, tel que $h$ vérifie
$\mathcal{P}$ et que $f$ ne vérifie pas $\mathcal{P}$. En effet, posons
$$
Y =
\Spec\Big(
\mathbf{C}[x_n; n \in \mathbf{Z}]/(x_n x_m; n \not = m)
\Big)
$$
et prenons pour $X \subset Y$ le sous-schéma ouvert complémentaire du point
dont toutes les coordonnées $x_n = 0$. Posons $U = X$ et
$V = X \amalg Y$; soient $a, b$ les morphismes évidents, et soit
$h : U \to V$ l'inclusion de $U = X$ dans la première composante de $V$.
L'affirmation précédente est vraie parce que $U$ est localement noethérien,
mais que $Y$ ne l'est pas.
\end{example}

\noindent
Quelle devrait être la bonne notion pour une propriété locale sur la source et
la cible pour la topologie \'etale ? Par analogie avec la
définition \ref{morphisms-definition-property-local} de Morphismes, nous pensons
que ce devrait être la suivante.

\begin{definition}
\label{definition-local-source-target}
Soit $\mathcal{P}$ une propriété de morphismes de schémas.
Nous disons que $\mathcal{P}$ est {\it locale sur la source et la cible pour
la topologie \'etale} si
\begin{enumerate}
\item (stabilité par précomposition avec des morphismes étales)
si $f : X \to Y$ est \'etale et si $g : Y \to Z$ possède $\mathcal{P}$,
alors $g \circ f$ possède $\mathcal{P}$,
\item (stabilité par changement de base \'etale)
si $f : X \to Y$ possède $\mathcal{P}$ et si $Y' \to Y$ est \'etale, alors
le changement de base $f' : Y' \times_Y X \to Y'$ possède $\mathcal{P}$, et
\item (localité) pour un morphisme $f : X \to Y$, les assertions suivantes
sont équivalentes :
\begin{enumerate}
\item $f$ possède $\mathcal{P}$,
\item pour tout $x \in X$, il existe un diagramme commutatif
$$
\xymatrix{
U \ar[d]_a \ar[r]_h & V \ar[d]^b \\
X \ar[r]^f & Y
}
$$
dont les flèches verticales sont étales et un $u \in U$ tel que
$a(u) = x$, de sorte que $h$ possède $\mathcal{P}$.
\end{enumerate}
\end{enumerate}
\end{definition}

\noindent
Cette définition exclut en fait le comportement observé dans les
exemples \ref{example-silly-one} et \ref{example-silly-two}.
Nous la comparerons à celle de l'article \cite{DM} de Deligne et Mumford dans
la remarque \ref{remark-compare-definitions}.
En outre, une propriété locale sur la source et la cible pour la topologie
\'etale est locale sur la source et locale sur la cible pour cette topologie.
Enfin, la réciproque est presque vraie, comme nous le verrons dans le
lemme \ref{lemma-etale-local-source-target}.

\begin{lemma}
\label{lemma-local-source-target-implies}
Soit $\mathcal{P}$ une propriété de morphismes de schémas qui est locale sur
la source et la cible pour la topologie \'etale. Alors
\begin{enumerate}
\item $\mathcal{P}$ est locale sur la source pour la topologie \'etale,
\item $\mathcal{P}$ est locale sur la cible pour la topologie \'etale,
\item $\mathcal{P}$ est stable par postcomposition avec des morphismes
\'etale : si $f : X \to Y$ possède $\mathcal{P}$ et si $g : Y \to Z$ est
\'etale, alors $g \circ f$ possède $\mathcal{P}$, et
\item $\mathcal{P}$ possède la propriété de permanence suivante : étant donnés
$f : X \to Y$ et $g : Y \to Z$ de type \'etale tels que $g \circ f$ possède
$\mathcal{P}$, le morphisme $f$ possède $\mathcal{P}$.
\end{enumerate}
\end{lemma}

\begin{proof}
Donnons tous les détails.

\medskip\noindent
Preuve de (1). Soit $f : X \to Y$ un morphisme de schémas.
Soit $\{X_i \to X\}_{i \in I}$ un recouvrement de $X$ pour la topologie
\'etale. Si chaque
composé $h_i : X_i \to Y$ possède $\mathcal{P}$, alors, pour tout $x \in X$,
on peut trouver un $i \in I$ et un point $x_i \in X_i$ d'image $x$. Le
morphisme de germes $(X_i, x_i) \to (X, x)$ est alors \'etale, le morphisme
$\text{id}_Y : Y \to Y$ est \'etale, et $h_i$ est comme dans la partie (3) de
la définition \ref{definition-local-source-target}.
Nous voyons donc que $f$ possède $\mathcal{P}$.
Réciproquement, si $f$ possède $\mathcal{P}$, alors chaque $X_i \to Y$ possède
$\mathcal{P}$ d'après la partie (1) de la
définition \ref{definition-local-source-target}.

\medskip\noindent
Preuve de (2). Soit $f : X \to Y$ un morphisme de schémas.
Soit $\{Y_i \to Y\}_{i \in I}$ un recouvrement de $Y$ pour la topologie
\'etale.
Écrivons $X_i = Y_i \times_Y X$ et notons $h_i : X_i \to Y_i$ le changement
de base de $f$. Si chaque $h_i : X_i \to Y_i$ possède $\mathcal{P}$, alors,
pour tout $x \in X$, choisissons un $i \in I$ et un point $x_i \in X_i$
d'image $x$. Le morphisme de germes $(X_i, x_i) \to (X, x)$ est alors
\'etale, le morphisme $Y_i \to Y$ est \'etale, et $h_i$ est comme dans la
partie (3) de la définition \ref{definition-local-source-target}.
Nous voyons donc que $f$ possède $\mathcal{P}$.
Réciproquement, si $f$ possède $\mathcal{P}$, alors chaque $X_i \to Y_i$
possède $\mathcal{P}$ d'après la partie (2) de la
définition \ref{definition-local-source-target}.

\medskip\noindent
Preuve de (3). Supposons que $f : X \to Y$ possède $\mathcal{P}$ et que
$g : Y \to Z$ soit \'etale. Pour tout $x \in X$, on peut considérer
$(X, x) \to (X, x)$ comme un morphisme de germes \'etale, tandis que
$Y \to Z$ est un morphisme \'etale et que $h = f$ est comme dans la partie (3)
de la définition \ref{definition-local-source-target}.
Nous voyons donc que $g \circ f$ possède $\mathcal{P}$.

\medskip\noindent
Preuve de (4). Soient $f : X \to Y$ un morphisme et $g : Y \to Z$ un
morphisme \'etale tels que $g \circ f$ possède $\mathcal{P}$. La partie (2) de
la définition \ref{definition-local-source-target} montre alors que
$\text{pr}_Y : Y \times_Z X \to Y$ possède $\mathcal{P}$. Or le morphisme
$(f, 1) : X \to Y \times_Z X$ est \'etale, car c'est une section de la
projection \'etale $\text{pr}_X : Y \times_Z X \to X$ ; voir le
lemme \ref{morphisms-lemma-etale-permanence} de Morphismes.
Par conséquent, $f = \text{pr}_Y \circ (f, 1)$ possède $\mathcal{P}$ d'après
la partie (1) de la définition \ref{definition-local-source-target}.
\end{proof}

\noindent
Le lemme suivant est l'analogue du
lemme \ref{morphisms-lemma-locally-P-characterize} de Morphismes.

\begin{lemma}
\label{lemma-local-source-target-characterize}
Soit $\mathcal{P}$ une propriété de morphismes de schémas locale sur la source
et la cible pour la topologie \'etale. Soit $f : X \to Y$ un morphisme de
schémas. Les assertions suivantes sont équivalentes :
\begin{enumerate}
\item[(a)] $f$ possède la propriété $\mathcal{P}$,
\item[(b)] pour tout $x \in X$, il existe un morphisme de germes \'etale
$a : (U, u) \to (X, x)$, un morphisme \'etale $b : V \to Y$ et un morphisme
$h : U \to V$ tels que $f \circ a = b \circ h$ et que $h$ possède
$\mathcal{P}$,
\item[(c)] pour tout diagramme commutatif
$$
\xymatrix{
U \ar[d]_a \ar[r]_h & V \ar[d]^b \\
X \ar[r]^f & Y
}
$$
avec $a$, $b$ étales, le morphisme $h$ possède $\mathcal{P}$,
\item[(d)] pour un diagramme comme en (c) dans lequel $a : U \to X$ est
surjectif, le morphisme $h$ possède $\mathcal{P}$,
\item[(e)] il existe un recouvrement pour la topologie \'etale
$\{Y_i \to Y\}_{i \in I}$ tel
que chaque changement de base $Y_i \times_Y X \to Y_i$ possède
$\mathcal{P}$,
\item[(f)] il existe un recouvrement pour la topologie \'etale
$\{X_i \to X\}_{i \in I}$ tel
que chaque composé $X_i \to Y$ possède $\mathcal{P}$,
\item[(g)] il existe un recouvrement pour la topologie \'etale
$\{Y_i \to Y\}_{i \in I}$ et, pour tout $i \in I$, un recouvrement pour la
topologie \'etale
$\{X_{ij} \to Y_i \times_Y X\}_{j \in J_i}$ tels que chaque morphisme
$X_{ij} \to Y_i$ possède $\mathcal{P}$.
\end{enumerate}
\end{lemma}

\begin{proof}
L'équivalence de (a) et (b) fait partie de la
définition \ref{definition-local-source-target}.
L'équivalence de (a) et (e) est la partie (2) du
lemme \ref{lemma-local-source-target-implies}.
L'équivalence de (a) et (f) est la partie (1) du
lemme \ref{lemma-local-source-target-implies}.
Comme (a) est désormais équivalente à (e) et à (f), il s'ensuit que (a) est
équivalente à (g).

\medskip\noindent
Il est clair que (c) implique (a). Si (a) est vérifiée, alors, pour tout
diagramme comme en (c), le morphisme $f \circ a$ possède $\mathcal{P}$ d'après
la partie (1) de la définition \ref{definition-local-source-target} ; le
morphisme $h$ possède alors $\mathcal{P}$ d'après la partie (4) du
lemme \ref{lemma-local-source-target-implies}.
Ainsi, (a) et (c) sont équivalentes. Il est clair que (c) implique (d).
Pour voir que (d) implique (a), supposons donné un diagramme comme en (c),
avec $a : U \to X$ surjectif et $h$ possédant $\mathcal{P}$.
Alors $b \circ h$ possède $\mathcal{P}$ d'après la partie (3) du
lemme \ref{lemma-local-source-target-implies}.
Puisque $\{a : U \to X\}$ est un recouvrement pour la topologie \'etale,
nous concluons que
$f$ possède $\mathcal{P}$ d'après la partie (1) du
lemme \ref{lemma-local-source-target-implies}.
\end{proof}

\noindent
Il semble que le résultat du lemme suivant ne soit pas une pure formalité :
il utilise effectivement la géométrie des morphismes \'etale.

\begin{lemma}
\label{lemma-etale-local-source-target}
Soit $\mathcal{P}$ une propriété de morphismes de schémas. Supposons que
\begin{enumerate}
\item $\mathcal{P}$ soit locale sur la source pour la topologie \'etale,
\item $\mathcal{P}$ soit locale sur la cible pour la topologie \'etale, et
\item $\mathcal{P}$ soit stable par postcomposition avec les immersions
ouvertes : si $f : X \to Y$ possède $\mathcal{P}$ et si $Y \subset Z$ est un
sous-schéma ouvert, alors $X \to Z$ possède $\mathcal{P}$.
\end{enumerate}
Alors $\mathcal{P}$ est locale sur la source et la cible pour la topologie
\'etale.
\end{lemma}

\begin{proof}
Soit $\mathcal{P}$ une propriété de morphismes de schémas satisfaisant les
conditions (1), (2) et (3) du lemme. Le
lemme \ref{lemma-precompose-property-local-source} montre que $\mathcal{P}$ est
stable par précomposition avec les morphismes étales. Le
lemme \ref{lemma-pullback-property-local-target} montre que $\mathcal{P}$ est
stable par changement de base \'etale. Il suffit donc de prouver la partie (3)
de la définition \ref{definition-local-source-target}.

\medskip\noindent
Plus précisément, supposons que $f : X \to Y$ soit un morphisme de schémas
satisfaisant la partie (3)(b) de la
définition \ref{definition-local-source-target}. Autrement dit, pour tout
$x \in X$, il existe un morphisme \'etale $a_x : U_x \to X$, un point
$u_x \in U_x$ d'image $x$, un morphisme \'etale $b_x : V_x \to Y$ et un
morphisme $h_x : U_x \to V_x$ tels que $f \circ a_x = b_x \circ h_x$ et que
$h_x$ possède $\mathcal{P}$. La preuve du lemme sera achevée dès que nous
aurons montré que $f$ possède $\mathcal{P}$. Posons $U = \coprod U_x$,
$a = \coprod a_x$, $V = \coprod V_x$, $b = \coprod b_x$ et
$h = \coprod h_x$. Nous obtenons un diagramme commutatif
$$
\xymatrix{
U \ar[d]_a \ar[r]_h & V \ar[d]^b \\
X \ar[r]^f & Y
}
$$
où $a$, $b$ sont étales et où $a$ est surjectif. Remarquons que $h$
possède $\mathcal{P}$, puisque chaque $h_x$ la possède et que $\mathcal{P}$ est
locale sur la cible pour la topologie \'etale. Comme $a$ est surjectif et que
$\mathcal{P}$ est locale sur la source pour la topologie \'etale, il suffit de
prouver que $b \circ h$ possède $\mathcal{P}$. Le lemme se ramène donc à
montrer que $\mathcal{P}$ est stable par postcomposition avec un morphisme
\'etale.

\medskip\noindent
Dans la suite de la preuve, soit $f : X \to Y$ un morphisme possédant la
propriété $\mathcal{P}$, et soit $g : Y \to Z$ un morphisme \'etale.
Considérons les assertions suivantes :
\begin{enumerate}
\item[(-)] Sans hypothèse supplémentaire, $g \circ f$ possède la propriété
$\mathcal{P}$.
\item[(A)] Lorsque $Z$ est affine, $g \circ f$ possède la propriété
$\mathcal{P}$.
\item[(AA)] Lorsque $X$ et $Z$ sont affines, $g \circ f$ possède la propriété
$\mathcal{P}$.
\item[(AAA)] Lorsque $X$, $Y$ et $Z$ sont affines, $g \circ f$ possède la
propriété $\mathcal{P}$.
\end{enumerate}
Une fois (-) prouvée, la démonstration du lemme sera achevée.

\medskip\noindent
Assertion 1 : (AAA) $\Rightarrow$ (AA).
En effet, soient $f : X \to Y$, $g : Y \to Z$ comme ci-dessus, avec $X$ et
$Z$ affines. Comme $X$ est affine, donc quasi-compact, il existe un nombre fini
d'ouverts affines $Y_i \subset Y$, $i = 1, \ldots, n$, tels que
$X = \bigcup_{i = 1, \ldots, n} f^{-1}(Y_i)$. Posons
$X_i = f^{-1}(Y_i)$. D'après le
lemme \ref{lemma-pullback-property-local-target}, chacun des morphismes
$X_i \to Y_i$ possède $\mathcal{P}$. Par conséquent,
$\coprod_{i = 1, \ldots, n} X_i \to \coprod_{i = 1, \ldots, n} Y_i$ possède
$\mathcal{P}$, puisque $\mathcal{P}$ est locale sur la cible pour la topologie
\'etale. En appliquant (AAA) au morphisme
$\coprod_{i = 1, \ldots, n} X_i \to \coprod_{i = 1, \ldots, n} Y_i$ et au
morphisme \'etale $\coprod_{i = 1, \ldots, n} Y_i \to Z$, nous voyons que
$\coprod_{i = 1, \ldots, n} X_i \to Z$ possède $\mathcal{P}$. Or
$\{\coprod_{i = 1, \ldots, n} X_i \to X\}$ est un recouvrement pour la
topologie \'etale ; comme $\mathcal{P}$ est locale sur la source pour cette
topologie, nous concluons que $X \to Z$ possède $\mathcal{P}$, comme voulu.

\medskip\noindent
Assertion 2 : (AAA) $\Rightarrow$ (A).
En effet, soient $f : X \to Y$, $g : Y \to Z$ comme ci-dessus, avec $Z$
affine. Choisissons un recouvrement par des ouverts affines
$X = \bigcup X_i$. Puisque $\mathcal{P}$ est locale sur la source pour la
topologie \'etale, chaque $f|_{X_i} : X_i \to Y$ possède $\mathcal{P}$. D'après
(AA), qui découle de (AAA) par l'Assertion 1, chaque $X_i \to Z$ possède
$\mathcal{P}$, pour tout $i$. Puisque $\{X_i \to X\}$ est un recouvrement
pour la topologie \'etale et que
$\mathcal{P}$ est locale sur la source pour cette topologie, nous
concluons que $X \to Z$ possède $\mathcal{P}$.

\medskip\noindent
Assertion 3 : (AAA) $\Rightarrow$ (-).
En effet, soient $f : X \to Y$, $g : Y \to Z$ comme ci-dessus. Choisissons un
recouvrement par des ouverts affines $Z = \bigcup Z_i$. Posons
$Y_i = g^{-1}(Z_i)$ et $X_i = f^{-1}(Y_i)$. D'après le
lemme \ref{lemma-pullback-property-local-target}, chacun des morphismes
$X_i \to Y_i$ possède $\mathcal{P}$. D'après (A), qui découle de (AAA) par
l'Assertion 2, chaque $X_i \to Z_i$ possède $\mathcal{P}$, pour tout $i$.
Comme $\mathcal{P}$ est locale sur la cible et que
$X_i = (g \circ f)^{-1}(Z_i)$, nous concluons que
$X \to Z$ possède $\mathcal{P}$.

\medskip\noindent
Ainsi, pour prouver le lemme, il suffit d'établir (AAA).
Soient $f : X \to Y$ et $g : Y \to Z$ comme ci-dessus, avec $X, Y, Z$
affines. Remarquons qu'un morphisme \'etale entre schémas affines a des fibres
universellement bornées ; voir les lemmes
\ref{morphisms-lemma-etale-locally-quasi-finite} et
\ref{morphisms-lemma-locally-quasi-finite-qc-source-universally-bounded} de
Morphismes.
Nous pouvons donc raisonner par récurrence sur l'entier $n$ qui borne le degré
des fibres de $Y \to Z$. Voir le
lemme \ref{morphisms-lemma-etale-universally-bounded} de Morphismes pour une description de
cet entier dans le cas d'un morphisme \'etale.
Si $n = 1$, alors $Y \to Z$ est une immersion ouverte ; voir le
lemme \ref{lemma-universally-injective-etale-open-immersion}. Le résultat
découle alors de l'hypothèse (3) du lemme. Supposons $n > 1$.

\medskip\noindent
Considérons le diagramme commutatif suivant
$$
\xymatrix{
X \times_Z Y \ar[d] \ar[r]_{f_Y} &
Y \times_Z Y \ar[d] \ar[r]_-{\text{pr}} &
Y \ar[d] \\
X \ar[r]^f &
Y \ar[r]^g &
Z
}
$$
Remarquons que l'on a une décomposition en sous-schémas ouverts et fermés
$Y \times_Z Y = \Delta_{Y/Z}(Y) \amalg Y'$ ; voir le
lemme \ref{morphisms-lemma-diagonal-unramified-morphism} de Morphismes.
Puisqu'elle est obtenue par changement de base, la seconde projection
$\text{pr} : Y \times_Z Y \to Y$ a des fibres de degré borné par $n$ ; voir le
lemme \ref{morphisms-lemma-base-change-universally-bounded} de Morphismes.
D'autre part, $\text{pr}|_{\Delta(Y)} : \Delta(Y) \to Y$ est un isomorphisme,
et chacune de ses fibres possède exactement un point. En appliquant le
lemme \ref{morphisms-lemma-etale-universally-bounded} de Morphismes, nous
concluons que les degrés des fibres de la restriction
$\text{pr}|_{Y'} : Y' \to Y$ sont bornés par $n - 1$.
Posons $X' = f_Y^{-1}(Y')$. Considérons le diagramme
$$
\xymatrix{
X \amalg X' \ar@{=}[d] \ar[r]_-{f \amalg f'} &
\Delta(Y) \amalg Y' \ar@{=}[d] \ar[r] &
Y \ar@{=}[d] \\
X \times_Z Y \ar[r]^{f_Y} &
Y \times_Z Y \ar[r]^-{\text{pr}} &
Y
}
$$
Comme $\mathcal{P}$ est locale sur la cible pour la topologie \'etale et donc
stable par changement de base \'etale (voir le
lemme \ref{lemma-pullback-property-local-target}), nous voyons que $f_Y$
possède $\mathcal{P}$. Puisque $\mathcal{P}$ est locale sur la source pour la
topologie \'etale, $f' = f_Y|_{X'}$ possède $\mathcal{P}$. Par l'hypothèse
de récurrence, $X' \to Y$ possède $\mathcal{P}$. Comme $\mathcal{P}$ est locale
sur la source et que
$\{X \to X \times_Z Y, X' \to X \times_Y Z\}$ est un recouvrement pour la
topologie \'etale,
nous concluons que $\text{pr} \circ f_Y$ possède $\mathcal{P}$.
Remarquons que $g \circ f$ peut être considéré comme un morphisme
$g \circ f : X \to g(Y)$. Or $\text{pr} \circ f_Y$ est le changement de base
de $g \circ f : X \to g(Y)$ par le recouvrement pour la topologie \'etale
$\{Y \to g(Y)\}$ ; puisque $\mathcal{P}$ est locale sur la cible pour cette
topologie, nous concluons que $g \circ f : X \to g(Y)$ possède la
propriété $\mathcal{P}$. Enfin, en appliquant encore une fois l'hypothèse (3)
du lemme, nous concluons que $g \circ f : X \to Z$ possède la propriété
$\mathcal{P}$.
\end{proof}

\begin{remark}
\label{remark-list-local-source-target}
En utilisant le lemme \ref{lemma-etale-local-source-target} et les résultats
des sections précédentes de ce chapitre, on obtient aisément une liste de
types de morphismes dont la propriété est locale sur la source et la cible pour
la topologie \'etale. Dans chaque cas, nous indiquons le lemme qui établit la
localité sur la source, puis celui qui établit la localité sur la cible pour
cette topologie. La troisième hypothèse du
lemme \ref{lemma-etale-local-source-target} se vérifie immédiatement dans
chaque cas ; nous omettons cette vérification. Voici la liste :
\begin{enumerate}
\item plat, voir les lemmes \ref{lemma-flat-fpqc-local-source} et
\ref{lemma-descending-property-flat},
\item localement de présentation finie, voir les lemmes
\ref{lemma-locally-finite-presentation-fppf-local-source} et
\ref{lemma-descending-property-locally-finite-presentation},
\item localement de type fini, voir les lemmes
\ref{lemma-locally-finite-type-fppf-local-source} et
\ref{lemma-descending-property-locally-finite-type},
\item universellement ouvert, voir les lemmes
\ref{lemma-universally-open-fppf-local-source} et
\ref{lemma-descending-property-universally-open},
\item syntomique, voir les lemmes \ref{lemma-syntomic-syntomic-local-source} et
\ref{lemma-descending-property-syntomic},
\item lisse, voir les lemmes \ref{lemma-smooth-smooth-local-source} et
\ref{lemma-descending-property-smooth},
\item \'etale, voir les lemmes \ref{lemma-etale-etale-local-source} et
\ref{lemma-descending-property-etale},
\item localement quasi-fini, voir les lemmes
\ref{lemma-locally-quasi-finite-etale-local-source} et
\ref{lemma-descending-property-quasi-finite},
\item non ramifié, voir les lemmes \ref{lemma-unramified-etale-local-source} et
\ref{lemma-descending-property-unramified},
\item G-non ramifié, voir les lemmes
\ref{lemma-unramified-etale-local-source} et
\ref{lemma-descending-property-unramified}, et
\item tout autre type qu'il sera utile d'ajouter.
\end{enumerate}
\end{remark}

\begin{remark}
\label{remark-compare-definitions}
À ce stade, trois définitions peuvent exprimer qu'une propriété
$\mathcal{P}$ de morphismes est ``locale sur la source et la cible pour la
topologie \'etale'' :
\begin{enumerate}
\item[(ST)] $\mathcal{P}$ est locale sur la source et locale sur la cible pour
la topologie \'etale,
\item[(DM)] (définition donnée par Deligne et Mumford dans l'article
\cite[p.~100]{DM}) pour tout diagramme
$$
\xymatrix{
U \ar[d]_a \ar[r]_h & V \ar[d]^b \\
X \ar[r]^f & Y
}
$$
dont les flèches verticales sont surjectives et étales, on a
$\mathcal{P}(h) \Leftrightarrow \mathcal{P}(f)$, et
\item[(SP)] $\mathcal{P}$ est locale sur la source et la cible pour la
topologie \'etale.
\end{enumerate}
Nous avons vu dans cette section que (SP) $\Rightarrow$ (DM) $\Rightarrow$
(ST). Les exemples \ref{example-silly-one} et \ref{example-silly-two} montrent
qu'aucune des deux implications réciproques n'est vraie. Enfin, le
lemme \ref{lemma-etale-local-source-target} montre que la différence disparaît
pour les propriétés de morphismes stables par postcomposition avec les
immersions ouvertes, ce qui sera toujours le cas en pratique.
\end{remark}

\begin{lemma}
\label{lemma-etale-etale-local-source-target}
Soit $\mathcal{P}$ une propriété de morphismes de schémas locale sur la source
et la cible pour la topologie \'etale. Étant donné un diagramme commutatif
de schémas
$$
\vcenter{
\xymatrix{
X' \ar[d]_{g'} \ar[r]_{f'} & Y' \ar[d]^g \\
X \ar[r]^f & Y
}
}
\quad\text{avec les points}\quad
\vcenter{
\xymatrix{
x' \ar[d] \ar[r] & y' \ar[d] \\
x \ar[r] & y
}
}
$$
tels que $g'$ soit \'etale en $x'$ et que $g$ soit \'etale en $y'$, on a
$x \in W(f) \Leftrightarrow x' \in W(f')$, où $W(-)$ est défini dans le
lemme \ref{lemma-largest-open-of-the-source}.
\end{lemma}

\begin{proof}
Le lemme \ref{lemma-largest-open-of-the-source} s'applique, puisque
$\mathcal{P}$ est locale sur la source pour la topologie \'etale d'après le
lemme \ref{lemma-local-source-target-implies}.

\medskip\noindent
Supposons $x \in W(f)$. Soient $U' \subset X'$ et $V' \subset Y'$ des
voisinages ouverts de $x'$ et $y'$ tels que $f'(U') \subset V'$,
$g'(U') \subset W(f)$, et tels que $g'|_{U'}$ et $g|_{V'}$ soient tous deux
étales.
Alors $f \circ g'|_{U'} = g \circ f'|_{U'}$ possède $\mathcal{P}$ d'après la
propriété (1) de la définition \ref{definition-local-source-target}.
Le morphisme $f'|_{U'} : U' \to V'$ possède donc la propriété $\mathcal{P}$
d'après la partie (4) du lemme \ref{lemma-local-source-target-implies}.
Puis, par la partie (3) du lemme \ref{lemma-local-source-target-implies}, nous
concluons que $f'_{U'} : U' \to Y'$ possède $\mathcal{P}$.
Ainsi, $U' \subset W(f')$ par définition, et donc $x' \in W(f')$.

\medskip\noindent
Supposons $x' \in W(f')$. Soient $U' \subset X'$ et $V' \subset Y'$ des
voisinages ouverts de $x'$ et $y'$ tels que $f'(U') \subset V'$,
$U' \subset W(f')$, et tels que $g'|_{U'}$ et $g|_{V'}$ soient tous deux
étales.
Alors $U' \to Y'$ possède $\mathcal{P}$ par définition de $W(f')$.
Le morphisme $U' \to V'$ possède donc $\mathcal{P}$ d'après la partie (4) du
lemme \ref{lemma-local-source-target-implies}.
Ensuite, $U' \to Y$ possède $\mathcal{P}$ d'après la partie (3) du
lemme \ref{lemma-local-source-target-implies}.
Soit $U \subset X$ l'image du morphisme \'etale, donc ouvert,
$g'|_U' : U' \to X$. Alors $\{U' \to U\}$ est un recouvrement pour la
topologie \'etale, et
nous concluons que $U \to Y$ possède $\mathcal{P}$ d'après la partie (1) du
lemme \ref{lemma-local-source-target-implies}.
Ainsi, $U \subset W(f)$ par définition, et donc $x \in W(f)$.
\end{proof}

\begin{lemma}
\label{lemma-orbits}
Soit $k$ un corps. Soit $n \geq 2$. Pour $1 \leq i, j \leq n$ avec
$i \not = j$ et $d \geq 0$, notons $T_{i, j, d}$ l'automorphisme de
$\mathbf{A}^n_k$ donné en coordonnées par
$$
(x_1, \ldots, x_n) \longmapsto
(x_1, \ldots, x_{i - 1}, x_i + x_j^d, x_{i + 1}, \ldots, x_n)
$$
Soit $W \subset \mathbf{A}^n_k$ un sous-schéma ouvert non vide tel que
$T_{i, j, d}(W) = W$ pour tous $i, j, d$ comme ci-dessus. Alors soit
$W = \mathbf{A}^n_k$, soit la caractéristique de $k$ est $p > 0$ et
$\mathbf{A}^n_k \setminus W$ est un ensemble fini de points fermés dont les
coordonnées sont algébriques sur $\mathbf{F}_p$.
\end{lemma}

\begin{proof}
Pour établir le résultat, nous pouvons remplacer $k$ par une extension
quelconque. Soit $Z$ une composante irréductible de
$\mathbf{A}^n_k \setminus W$. Supposons $\dim(Z) \geq 1$, et cherchons une
contradiction. Il existe alors une extension de corps $k'/k$ et un point à
valeurs dans $k'$, $\xi = (\xi_1, \ldots, \xi_n) \in (k')^n$, de
$Z_{k'} \subset \mathbf{A}^n_{k'}$, tel qu'au moins l'un des
$x_1, \ldots, x_n$ soit transcendant sur le corps premier. Affirmation :
l'orbite de $\xi$ sous le groupe engendré par les transformations
$T_{i, j, d}$ est dense pour la topologie de Zariski dans
$\mathbf{A}^n_{k'}$. Cette affirmation donnera la contradiction cherchée.

\medskip\noindent
Si la caractéristique de $k'$ est nulle, les opérateurs $T_{i, j, 0}$
suffisent déjà, puisqu'ils transforment $\xi$ en les points
$$
(\xi_1 + a_1, \ldots, \xi_n + a_n)
$$
pour $(a_1, \ldots, a_n) \in \mathbf{Z}_{\geq 0}^n$ arbitraire. Si la
caractéristique est $p > 0$, nous pouvons supposer, après renumérotation, que
$\xi_n$ est transcendant sur $\mathbf{F}_p$. En appliquant successivement les
opérateurs $T_{i, n, d}$ pour $i < n$, nous voyons que l'orbite de $\xi$
contient les éléments
$$
(\xi_1 + P_1(\xi_n), \ldots, \xi_{n - 1} + P_{n - 1}(\xi_n), \xi_n)
$$
pour $(P_1, \ldots, P_{n - 1}) \in \mathbf{F}_p[t]$ arbitraire. La fermeture
de l'orbite pour la topologie de Zariski contient donc l'hyperplan d'équation
$x_n = \xi_n$. En répétant l'argument avec une autre coordonnée, nous
concluons que cette fermeture contient $x_i = \xi_i + P(\xi_n)$ pour tout
$P \in \mathbf{F}_p[t]$ tel que $\xi_i + P(\xi_n)$ soit transcendant sur
$\mathbf{F}_p$. Puisqu'il existe une infinité de tels $P$, l'affirmation en
résulte.

\medskip\noindent
Bien entendu, l'argument du paragraphe précédent s'applique aussi lorsque
$Z = \{z\}$ est de dimension $0$ et que les coordonnées de $z$ dans
$\kappa(z)$ ne sont pas algébriques sur $\mathbf{F}_p$. Le lemme en résulte.
\end{proof}

\begin{lemma}
\label{lemma-etale-tau-local-source-target}
Soit $\mathcal{P}$ une propriété de morphismes de schémas. Supposons que
\begin{enumerate}
\item $\mathcal{P}$ soit locale sur la source pour la topologie \'etale,
\item $\mathcal{P}$ soit locale sur la cible pour la topologie lisse,
\item $\mathcal{P}$ soit stable par postcomposition avec les immersions
ouvertes : si $f : X \to Y$ possède $\mathcal{P}$ et si $Y \subset Z$ est un
sous-schéma ouvert, alors $X \to Z$ possède $\mathcal{P}$.
\end{enumerate}
Étant donné un diagramme commutatif de schémas
$$
\vcenter{
\xymatrix{
X' \ar[d]_{g'} \ar[r]_{f'} & Y' \ar[d]^g \\
X \ar[r]^f & Y
}
}
\quad\text{avec les points}\quad
\vcenter{
\xymatrix{
x' \ar[d] \ar[r] & y' \ar[d] \\
x \ar[r] & y
}
}
$$
tels que $g$ soit lisse en $y'$ et que $X' \to X \times_Y Y'$ soit \'etale
en $x'$, on a $x \in W(f) \Leftrightarrow x' \in W(f')$, où $W(-)$ est défini
dans le lemme \ref{lemma-largest-open-of-the-source}.
\end{lemma}

\begin{proof}
Comme $\mathcal{P}$ est locale sur la source pour la topologie \'etale, nous
voyons que $x \in W(f)$ si et seulement si l'image de $x$ dans
$X \times_Y Y'$ appartient à $W(X \times_Y Y' \to Y')$. Nous pouvons donc
supposer que le diagramme du lemme est cartésien.

\medskip\noindent
Supposons $x \in W(f)$. Comme $\mathcal{P}$ est locale sur la cible pour la
topologie lisse, nous voyons que
$(g')^{-1}W(f) = W(f) \times_Y Y' \to Y'$ possède
$\mathcal{P}$. Ainsi, $(g')^{-1}W(f) \subset W(f')$. Nous concluons que
$x' \in W(f')$.

\medskip\noindent
Supposons $x' \in W(f')$. Pour tout voisinage ouvert $V' \subset Y'$ de $y'$,
on peut remplacer $Y'$ par $V'$ et $X'$ par $U' = (f')^{-1}V'$, puisque
$V' \to Y'$ est lisse et que le changement de base
$W(f') \cap U' \to V'$ de $W(f') \to Y'$ possède la propriété $\mathcal{P}$.
Nous pouvons donc supposer qu'il existe un morphisme \'etale
$Y' \to \mathbf{A}^n_Y$ au-dessus de $Y$ ; voir le
lemme \ref{morphisms-lemma-smooth-etale-over-affine-space} de Morphismes.
Considérons le diagramme
$$
\xymatrix{
X' \ar[r] \ar[d] & Y' \ar[d] \\
\mathbf{A}^n_X \ar[r]_{f_n} \ar[d] & \mathbf{A}^n_Y \ar[d] \\
X \ar[r]^f & Y
}
$$
D'après le lemme \ref{lemma-etale-local-source-target} (et puisque les
recouvrements pour la topologie \'etale sont des recouvrements pour la
topologie lisse), nous voyons que $\mathcal{P}$ est locale sur la source et la
cible pour la topologie \'etale. Le
lemme \ref{lemma-etale-etale-local-source-target} montre que
$W(f')$ est l'image inverse de l'ouvert $W(f_n) \subset \mathbf{A}^n_X$. En
particulier, $W(f_n)$ contient un point au-dessus de $x$. Après avoir remplacé
$X$ par l'image de $W(f_n)$, qui est ouverte, nous pouvons supposer que
$W(f_n) \to X$ est surjectif. Affirmation :
$W(f_n) = \mathbf{A}^n_X$. Cette affirmation implique que $f$ possède
$\mathcal{P}$, puisque $\mathcal{P}$ est locale sur la cible pour la topologie
lisse et que $\{\mathbf{A}^n_Y \to Y\}$ est un recouvrement pour cette
topologie.

\medskip\noindent
Essentiellement, l'affirmation résulte du fait que
$W(f_n) \subset \mathbf{A}^n_X$ est un ouvert ``invariant par translation''
qui rencontre toute fibre de $\mathbf{A}^n_X \to X$. Toutefois, pour produire
un argument de cette forme, il faut effectuer une localisation \'etale sur
$Y$ afin de disposer d'assez de translations, ce qui devient quelque peu
fastidieux. Nous utilisons plutôt les automorphismes du
lemme \ref{lemma-orbits} et des morphismes étales d'espaces affines.
Nous pouvons supposer $n \geq 2$. En effet, si $n = 0$, le résultat est acquis.
Si $n = 1$, considérons le diagramme
$$
\xymatrix{
\mathbf{A}^2_X \ar[r]_{f_2} \ar[d]_p & \mathbf{A}^2_Y \ar[d] \\
\mathbf{A}^1_X \ar[r]^{f_1} & \mathbf{A}^1_Y
}
$$
Nous avons $p^{-1}(W(f_1)) \subset W(f_2)$ ; voir le premier paragraphe de la
preuve. Ainsi, $W(f_2) \to X$ est encore surjectif, et nous pouvons travailler
avec $f_2$. Supposons $n \geq 2$.

\medskip\noindent
Pour tous $1 \leq i, j \leq n$ avec $i \not = j$ et $d \geq 0$, notons
$T_{i, j, d}$ l'automorphisme de $\mathbf{A}^n$ défini dans le
lemme \ref{lemma-orbits}. Nous obtenons alors un diagramme commutatif
$$
\xymatrix{
\mathbf{A}^n_X \ar[r]_{f_n} \ar[d]_{T_{i, j, d}} &
\mathbf{A}^n_Y \ar[d]^{T_{i, j, d}} \\
\mathbf{A}^n_X \ar[r]^{f_n} & \mathbf{A}^n_Y
}
$$
dont les flèches verticales sont des isomorphismes. Nous concluons que
$T_{i, j, d}(W(f_n)) = W(f_n)$. En appliquant le lemme \ref{lemma-orbits}, nous
concluons que, pour tout $x \in X$, la fibre
$W(f_n)_x \subset \mathbf{A}^n_x$ est soit $\mathbf{A}^n_x$, ce qui est le
résultat cherché, soit telle que $\kappa(x)$ est de caractéristique $p > 0$ et
que $W(f_n)_x$ est le complémentaire d'un ensemble fini
$Z_x \subset \mathbf{A}^n_x$ de points fermés. La seconde possibilité ne peut
pas se produire. Considérons en effet le morphisme
$T_p : \mathbf{A}^n \to \mathbf{A}^n$ donné par
$$
(x_1, \ldots, x_n) \mapsto (x_1 - x_1^p, \ldots, x_n - x_n^p)
$$
Comme ci-dessus, nous obtenons un diagramme commutatif
$$
\xymatrix{
\mathbf{A}^n_X \ar[r]_{f_n} \ar[d]_{T_p} &
\mathbf{A}^n_Y \ar[d]^{T_p} \\
\mathbf{A}^n_X \ar[r]^{f_n} & \mathbf{A}^n_Y
}
$$
Le morphisme $T_p : \mathbf{A}^n_X \to \mathbf{A}^n_X$ est \'etale en tout
point au-dessus de $x$, et le morphisme
$T_p : \mathbf{A}^n_Y \to \mathbf{A}^n_Y$ est \'etale en tout point au-dessus
de l'image de $x$ dans $Y$. (Détails omis ; indication : calculer les
dérivées.) Nous concluons que
$$
T_p^{-1}(W) \cap \mathbf{A}^n_x = W \cap \mathbf{A}^n_x
$$
d'après le lemme \ref{lemma-etale-etale-local-source-target} (nous avons déjà
vu que $\mathcal{P}$ est locale sur la source et la cible pour la topologie
\'etale). Puisque $T_p : \mathbf{A}^n_x \to \mathbf{A}^n_x$ est fini \'etale de
degré $p^n > 1$, si $Z_x$ n'est pas vide, il contient $T_p^{-1}(Z_x)$, qui est
strictement plus grand. Cette contradiction achève la preuve.
\end{proof}





\section[Morphismes de germes : localité étale]{Propriétés des morphismes de germes locales sur la source et la cible
pour la topologie étale}
\label{section-local-source-target-at-point}

\noindent
Dans cette section, nous examinons l'analogue des résultats de la
section \ref{section-properties-etale-local-source-target}
pour les morphismes de germes de schémas.

\begin{definition}
\label{definition-local-source-target-at-point}
Soit $\mathcal{Q}$ une propriété des morphismes de germes de schémas.
On dit que $\mathcal{Q}$ est {\it locale sur la source et la cible pour la
topologie \'etale} si, pour tout diagramme commutatif
$$
\xymatrix{
(U', u') \ar[d]_a \ar[r]_{h'} & (V', v') \ar[d]^b \\
(U, u) \ar[r]^h & (V, v)
}
$$
de germes dont les flèches verticales sont étales, on a
$\mathcal{Q}(h) \Leftrightarrow \mathcal{Q}(h')$.
\end{definition}

\begin{lemma}
\label{lemma-local-source-target-global-implies-local}
Soit $\mathcal{P}$ une propriété des morphismes de schémas
qui est locale sur la source et la cible pour la topologie \'etale.
Considérons la propriété $\mathcal{Q}$ des morphismes de germes définie par
$$
\mathcal{Q}((X, x) \to (S, s))
\Leftrightarrow
\text{il existe un représentant }U \to S
\text{ qui vérifie }\mathcal{P}
$$
Alors $\mathcal{Q}$ est locale sur la source et la cible pour la topologie
\'etale au sens de la
définition \ref{definition-local-source-target-at-point}.
\end{lemma}

\begin{proof}
Si un morphisme de germes $(X, x) \to (S, s)$ vérifie $\mathcal{Q}$,
alors il existe des voisinages ouverts arbitrairement petits
$U \subset X$ de $x$ et $V \subset S$ de $s$ tels qu'un représentant
$U \to V$ de $(X, x) \to (S, s)$ vérifie $\mathcal{P}$.
Cela résulte du lemme \ref{lemma-local-source-target-implies}. Soit
$$
\xymatrix{
(U', u') \ar[r]_{h'} \ar[d]_a & (V', v') \ar[d]^b \\
(U, u) \ar[r]^h & (V, v)
}
$$
un diagramme comme dans la
définition \ref{definition-local-source-target-at-point}.
Choisissons $U_1 \subset U$ et un représentant $h_1 : U_1 \to V$ de $h$.
Choisissons $V'_1 \subset V'$ et un représentant \'etale
$b_1 : V'_1 \to V$ de $b$
(définition \ref{definition-etale-morphism-germs}).
Choisissons $U'_1 \subset U'$ et des représentants
$a_1 : U'_1 \to U_1$ et $h'_1 : U'_1 \to V'_1$ de $a$ et $h'$,
avec $a_1$ \'etale. Après avoir rétréci $U'_1$, nous pouvons supposer que
$h_1 \circ a_1 = b_1 \circ h'_1$. D'après la remarque liminaire de la
démonstration, il s'agit de montrer que
$u' \in W(h'_1) \Leftrightarrow u \in W(h_1)$, où $W(-)$ est défini comme
au lemme \ref{lemma-largest-open-of-the-source}.
Le résultat découle donc du
lemme \ref{lemma-etale-etale-local-source-target}.
\end{proof}

\begin{lemma}
\label{lemma-local-source-target-local-implies-global}
Soit $\mathcal{P}$ une propriété des morphismes de schémas qui est
locale sur la source et la cible pour la topologie \'etale. Soit $Q$ la
propriété associée des morphismes de germes, voir le
lemme \ref{lemma-local-source-target-global-implies-local}.
Soit $f : X \to Y$ un morphisme de schémas. Les conditions suivantes
sont équivalentes :
\begin{enumerate}
\item $f$ vérifie la propriété $\mathcal{P}$, et
\item pour tout $x \in X$, le morphisme de germes
$(X, x) \to (Y, f(x))$ vérifie la propriété $\mathcal{Q}$.
\end{enumerate}
\end{lemma}

\begin{proof}
L'implication (1) $\Rightarrow$ (2) résulte directement des définitions.
L'implication (2) $\Rightarrow$ (1) résulte aussi de la partie (3) de la
définition \ref{definition-local-source-target}.
\end{proof}

\noindent
Un morphisme de germes $(X, x) \to (S, s)$ détermine un homomorphisme
bien défini d'anneaux locaux. Le lemme suivant a donc un sens.

\begin{lemma}
\label{lemma-flat-at-point}
La propriété des morphismes de germes
$$
\mathcal{P}((X, x) \to (S, s)) =
\mathcal{O}_{S, s} \to \mathcal{O}_{X, x}\text{ est plat}
$$
est locale sur la source et la cible pour la topologie \'etale.
\end{lemma}

\begin{proof}
Étant donné un diagramme comme dans la
définition \ref{definition-local-source-target-at-point}, nous obtenons le
diagramme suivant d'homomorphismes locaux d'anneaux locaux :
$$
\xymatrix{
\mathcal{O}_{U', u'} & \mathcal{O}_{V', v'} \ar[l] \\
\mathcal{O}_{U, u} \ar[u] & \mathcal{O}_{V, v} \ar[l] \ar[u]
}
$$
Les flèches verticales sont des localisations d'homomorphismes d'anneaux
étales ; en particulier, elles sont essentiellement de présentation finie,
plates et non ramifiées (voir
la section \ref{algebra-section-etale} d'Algèbre).
Ces homomorphismes verticaux sont donc fidèlement plats, voir le
lemme \ref{algebra-lemma-local-flat-ff} d'Algèbre.
Si la flèche horizontale supérieure est plate, la flèche horizontale
inférieure est plate par application du
lemme \ref{algebra-lemma-flat-permanence} d'Algèbre,
avec $R = \mathcal{O}_{V, v}$, $S = \mathcal{O}_{U, u}$ et
$M = \mathcal{O}_{U', u'}$. Si la flèche horizontale inférieure est
plate, alors l'homomorphisme d'anneaux
$$
\mathcal{O}_{V', v'} \otimes_{\mathcal{O}_{V, v}} \mathcal{O}_{U, u}
\longleftarrow
\mathcal{O}_{V', v'}
$$
est plat d'après le
lemme \ref{algebra-lemma-flat-base-change} d'Algèbre.
Enfin, l'homomorphisme d'anneaux
$$
\mathcal{O}_{U', u'}
\longleftarrow
\mathcal{O}_{V', v'} \otimes_{\mathcal{O}_{V, v}} \mathcal{O}_{U, u}
$$
est la localisation d'un homomorphisme entre des extensions de type \'etale de
$\mathcal{O}_{U, u}$ ; il est donc plat d'après le
lemme \ref{algebra-lemma-map-between-etale} d'Algèbre.
\end{proof}

\begin{lemma}
\label{lemma-etale-on-fiber}
Considérons un diagramme commutatif de morphismes de schémas
$$
\xymatrix{
U' \ar[r] \ar[d] & V' \ar[d] \\
U \ar[r] & V
}
$$
dont les flèches verticales sont étales, et un point $v' \in V'$
dont l'image est $v \in V$. Alors le morphisme entre les fibres
$U'_{v'} \to U_v$ est \'etale.
\end{lemma}

\begin{proof}
Le morphisme $U'_v \to U_v$ est \'etale, puisqu'il s'obtient par
changement de base à partir du morphisme \'etale $U' \to U$.
Le schéma $U'_v$ est un schéma sur $V'_v$. D'après le
lemme \ref{morphisms-lemma-etale-over-field} de Morphismes, le schéma
$V'_v$ est une réunion disjointe de spectres d'extensions finies séparables
de $\kappa(v)$. L'une de ces composantes est
$v' = \Spec(\kappa(v'))$. Par conséquent, $U'_{v'}$ est un sous-schéma
ouvert et fermé de $U'_v$, et le composé
$U'_{v'} \to U'_v \to U_v$ est \'etale (comme composé d'une immersion
ouverte et d'un morphisme \'etale, voir la
section \ref{morphisms-section-etale} de Morphismes).
\end{proof}

\noindent
Étant donné un morphisme de germes de schémas $(X, x) \to (S, s)$,
on peut définir sa {\it fibre} comme la classe d'isomorphisme du germe
$(U_s, x)$, où $U \to S$ est un représentant quelconque. Par abus de
notation, nous écrirons souvent simplement $(X_s, x)$.

\begin{lemma}
\label{lemma-dimension-local-ring-fibre}
Soit $d \in \{0, 1, 2, \ldots, \infty\}$.
La propriété des morphismes de germes
$$
\mathcal{P}_d((X, x) \to (S, s)) =
\text{l'anneau local }
\mathcal{O}_{X_s, x}
\text{ de la fibre est de dimension }d
$$
est locale sur la source et la cible pour la topologie \'etale.
\end{lemma}

\begin{proof}
Étant donné un diagramme comme dans la
définition \ref{definition-local-source-target-at-point}, on obtient un
morphisme \'etale entre les fibres $U'_{v'} \to U_v$ envoyant $u'$ sur $u$,
voir le lemme \ref{lemma-etale-on-fiber}. Le résultat découle donc du
lemme \ref{lemma-dimension-local-ring-local}.
\end{proof}

\begin{lemma}
\label{lemma-transcendence-degree-at-point}
Soit $r \in \{0, 1, 2, \ldots, \infty\}$.
La propriété des morphismes de germes
$$
\mathcal{P}_r((X, x) \to (S, s))
\Leftrightarrow
\text{trdeg}_{\kappa(s)} \kappa(x) = r
$$
est locale sur la source et la cible pour la topologie \'etale.
\end{lemma}

\begin{proof}
Étant donné un diagramme comme dans la
définition \ref{definition-local-source-target-at-point}, on obtient le
diagramme suivant d'homomorphismes locaux d'anneaux locaux :
$$
\xymatrix{
\mathcal{O}_{U', u'} & \mathcal{O}_{V', v'} \ar[l] \\
\mathcal{O}_{U, u} \ar[u] & \mathcal{O}_{V, v} \ar[l] \ar[u]
}
$$
Les flèches verticales sont des localisations d'homomorphismes d'anneaux
étales ; elles sont en particulier non ramifiées (voir la
section \ref{algebra-section-etale} d'Algèbre).
Par conséquent, $\kappa(u')/\kappa(u)$ et $\kappa(v')/\kappa(v)$ sont des
extensions finies séparables. On a donc
$\text{trdeg}_{\kappa(v)} \kappa(u) = \text{trdeg}_{\kappa(v')} \kappa(u)$,
ce qui démontre le lemme.
\end{proof}

\noindent
Soit $(X, x)$ un germe de schéma. La dimension de $X$ en $x$ est le minimum
des dimensions des voisinages ouverts de $x$ dans $X$, et tout voisinage
ouvert suffisamment petit a cette dimension. C'est donc un invariant de la
classe d'isomorphisme du germe. On le note simplement $\dim_x(X)$.

\begin{lemma}
\label{lemma-dimension-at-point}
Soit $d \in \{0, 1, 2, \ldots, \infty\}$.
La propriété des morphismes de germes
$$
\mathcal{P}_d((X, x) \to (S, s))
\Leftrightarrow
\dim_x (X_s) = d
$$
est locale sur la source et la cible pour la topologie \'etale.
\end{lemma}

\begin{proof}
Étant donné un diagramme comme dans la
définition \ref{definition-local-source-target-at-point}, on obtient un
morphisme \'etale entre les fibres $U'_{v'} \to U_v$ envoyant $u'$ sur $u$,
voir le lemme \ref{lemma-etale-on-fiber}. L'égalité
$\dim_u(U_v) = \dim_{u'}(U'_{v'})$ résulte alors du
lemme \ref{lemma-dimension-at-point-local}.
\end{proof}







\section{Données de descente pour les schémas au-dessus de schémas}
\label{section-descent-datum}

\noindent
La plupart des raisonnements de cette section sont formels et ne reposent que
sur la définition d'une donnée de descente. Dans
la section \ref{spaces-simplicial-section-simplicial-descent} d'Espaces
simpliciaux, nous examinerons le lien
avec les schémas simpliciaux, ce qui éclairera quelque peu la situation.

\begin{definition}
\label{definition-descent-datum}
Soit $f : X \to S$ un morphisme de schémas.
\begin{enumerate}
\item Soit $V \to X$ un schéma au-dessus de $X$.
Une {\it donnée de descente pour $V/X/S$} est un isomorphisme
$\varphi : V \times_S X \to X \times_S V$ de schémas au-dessus de
$X \times_S X$ satisfaisant la {\it condition de cocycle}, c'est-à-dire tel
que le diagramme
$$
\xymatrix{
V \times_S X \times_S X \ar[rd]^{\varphi_{01}} \ar[rr]_{\varphi_{02}} &
&
X \times_S X \times_S V\\
&
X \times_S V \times_S X \ar[ru]^{\varphi_{12}}
}
$$
soit commutatif (avec les notations évidentes).
\item On dit aussi que le couple $(V/X, \varphi)$ est
une {\it donnée de descente relativement à $X \to S$}.
\item Un {\it morphisme
$g : (V/X, \varphi) \to (V'/X, \varphi')$ de données de descente
relativement à $X \to S$} est un morphisme $g : V \to V'$ de schémas
au-dessus de $X$ tel que le diagramme
$$
\xymatrix{
V \times_S X \ar[r]_{\varphi} \ar[d]_{g \times \text{id}_X} &
X \times_S V \ar[d]^{\text{id}_X \times g} \\
V' \times_S X \ar[r]^{\varphi'} & X \times_S V'
}
$$
soit commutatif.
\end{enumerate}
\end{definition}

\noindent
La définition précédente donne lieu à toutes sortes d'identités
``miraculeuses''. Par exemple, l'image inverse de $\varphi$ par le
morphisme diagonal $\Delta : X \to X \times_S X$ peut être considéré comme
un morphisme $\Delta^*\varphi : V \to V$.
En effet, $X \times_{\Delta, X \times_S X} (V \times_S X) = V$ et
$X \times_{\Delta, X \times_S X} (X \times_S V) = V$.
En fait, $\Delta^*\varphi$ est l'identité.
Voilà un bon exercice si cette construction ne vous est pas familière.

\begin{remark}
\label{remark-easier}
Soit $X \to S$ un morphisme de schémas. Soit $(V/X, \varphi)$ une donnée
de descente relativement à $X \to S$. On peut voir l'isomorphisme
$\varphi$ comme un isomorphisme
$$
(X \times_S X) \times_{\text{pr}_0, X} V
\longrightarrow
(X \times_S X) \times_{\text{pr}_1, X} V
$$
de schémas au-dessus de $X \times_S X$. De façon informelle, on peut donc
voir $\varphi$ comme un morphisme
$\varphi : \text{pr}_0^*V \to \text{pr}_1^*V$\footnote{Malheureusement,
le sens choisi ici pour notre flèche est opposé à la convention attendue. Dans les
définitions \ref{definition-descent-datum} et
\ref{definition-descent-datum-for-family-of-morphisms}, nous aurions dû
prendre le sens opposé à celui de la
définition \ref{definition-descent-datum-quasi-coherent}, conformément au
principe général selon lequel ``fonctions'' et ``espaces'' sont
duaux.}. La condition de cocycle affirme alors que
$\text{pr}_{02}^*\varphi =
\text{pr}_{12}^*\varphi \circ \text{pr}_{01}^*\varphi$.
De ce point de vue, la situation est très analogue à celle d'une donnée de
descente sur des faisceaux quasi-cohérents.
\end{remark}

\noindent
Voici la définition lorsqu'on considère une famille de morphismes de même but.

\begin{definition}
\label{definition-descent-datum-for-family-of-morphisms}
Soit $S$ un schéma.
Soit $\{X_i \to S\}_{i \in I}$ une famille de morphismes de but $S$.
\begin{enumerate}
\item Une {\it donnée de descente $(V_i, \varphi_{ij})$ relativement à la
famille $\{X_i \to S\}$} consiste en un schéma $V_i$ au-dessus de $X_i$
pour chaque $i \in I$, et en un isomorphisme
$\varphi_{ij} : V_i \times_S X_j \to X_i \times_S V_j$
de schémas au-dessus de $X_i \times_S X_j$ pour chaque paire
$(i, j) \in I^2$, tels que, pour tout triplet d'indices
$(i, j, k) \in I^3$, le diagramme
$$
\xymatrix{
V_i \times_S X_j \times_S X_k
\ar[rd]^{\text{pr}_{01}^*\varphi_{ij}}
\ar[rr]_{\text{pr}_{02}^*\varphi_{ik}} &
&
X_i \times_S X_j \times_S V_k\\
&
X_i \times_S V_j \times_S X_k
\ar[ru]^{\text{pr}_{12}^*\varphi_{jk}}
}
$$
de schémas au-dessus de $X_i \times_S X_j \times_S X_k$ soit commutatif
(avec les notations évidentes).
\item Un {\it morphisme
$\psi : (V_i, \varphi_{ij}) \to (V'_i, \varphi'_{ij})$
de données de descente} consiste en une famille
$\psi = (\psi_i)_{i \in I}$ de morphismes de
$X_i$-schémas $\psi_i : V_i \to V'_i$ telle que tous les diagrammes
$$
\xymatrix{
V_i \times_S X_j \ar[r]_{\varphi_{ij}} \ar[d]_{\psi_i \times \text{id}} &
X_i \times_S V_j \ar[d]^{\text{id} \times \psi_j} \\
V'_i \times_S X_j \ar[r]^{\varphi'_{ij}} & X_i \times_S V'_j
}
$$
soient commutatifs.
\end{enumerate}
\end{definition}

\noindent
Cette notion apparaît naturellement, par exemple, lorsqu'on se demande si la
catégorie fibrée des courbes relatives est un champ pour la topologie fpqc
(ce n'est pas le cas -- du moins si l'on s'en tient aux schémas).

\begin{remark}
\label{remark-easier-family}
Soit $S$ un schéma.
Soit $\{X_i \to S\}_{i \in I}$ une famille de morphismes de but $S$.
Soit $(V_i, \varphi_{ij})$ une donnée de descente relativement à
$\{X_i \to S\}$. On peut voir les isomorphismes $\varphi_{ij}$ comme des
isomorphismes
$$
(X_i \times_S X_j) \times_{\text{pr}_0, X_i} V_i
\longrightarrow
(X_i \times_S X_j) \times_{\text{pr}_1, X_j} V_j
$$
de schémas au-dessus de $X_i \times_S X_j$. De façon informelle, on peut
donc voir $\varphi_{ij}$ comme un isomorphisme
$\text{pr}_0^*V_i \to \text{pr}_1^*V_j$ au-dessus de
$X_i \times_S X_j$. La condition de cocycle affirme alors que
$\text{pr}_{02}^*\varphi_{ik} =
\text{pr}_{12}^*\varphi_{jk} \circ \text{pr}_{01}^*\varphi_{ij}$.
De ce point de vue, la situation est très analogue à celle d'une donnée de
descente sur des faisceaux quasi-cohérents.
\end{remark}

\noindent
Le lemme suivant explique pourquoi nous travaillerons généralement avec une
famille réduite à un seul morphisme.

\begin{lemma}
\label{lemma-family-is-one}
Soit $S$ un schéma.
Soit $\{X_i \to S\}_{i \in I}$ une famille de morphismes de but $S$.
Posons $X = \coprod_{i \in I} X_i$ et considérons-le comme un $S$-schéma.
Il existe une équivalence canonique de catégories
$$
\begin{matrix}
\text{catégorie des données de descente } \\
\text{relativement à la famille } \{X_i \to S\}_{i \in I}
\end{matrix}
\longrightarrow
\begin{matrix}
\text{ catégorie des données de descente} \\
\text{ relativement à } X/S
\end{matrix}
$$
qui envoie $(V_i, \varphi_{ij})$ sur $(V, \varphi)$, où
$V = \coprod_{i\in I} V_i$ et $\varphi = \coprod \varphi_{ij}$.
\end{lemma}

\begin{proof}
On a $X \times_S X = \coprod_{ij} X_i \times_S X_j$, et de même pour les
produits fibrés d'ordre supérieur. Se donner un morphisme $V \to X$ équivaut
exactement à se donner une famille $V_i \to X_i$. De même, se donner une
donnée de descente $\varphi$ équivaut exactement à se donner une famille
$\varphi_{ij}$.
\end{proof}

\begin{lemma}
\label{lemma-pullback}
Changement de base des données de descente pour les schémas au-dessus de
schémas.
\begin{enumerate}
\item Soit
$$
\xymatrix{
X' \ar[r]_f \ar[d]_{a'} & X \ar[d]^a \\
S' \ar[r]^h & S
}
$$
un diagramme commutatif de morphismes de schémas.
La construction
$$
(V \to X, \varphi) \longmapsto f^*(V \to X, \varphi) = (V' \to X', \varphi')
$$
où $V' = X' \times_X V$ et où $\varphi'$ est défini comme le composé
$$
\xymatrix{
V' \times_{S'} X' \ar@{=}[r] &
(X' \times_X V) \times_{S'} X' \ar@{=}[r] &
(X' \times_{S'} X') \times_{X \times_S X} (V \times_S X)
\ar[d]^{\text{id} \times \varphi} \\
X' \times_{S'} V' \ar@{=}[r] &
X' \times_{S'} (X' \times_X V) &
(X' \times_{S'} X') \times_{X \times_S X} (X \times_S V) \ar@{=}[l]
}
$$
définit un foncteur de la catégorie des données de descente relativement à
$X \to S$ vers la catégorie des données de descente relativement à
$X' \to S'$.
\item Étant donnés deux morphismes $f_i : X' \to X$, $i = 0, 1$, qui rendent
le diagramme commutatif, les foncteurs $f_0^*$ et $f_1^*$ sont canoniquement
isomorphes.
\end{enumerate}
\end{lemma}

\begin{proof}
Nous omettons la preuve de (1), mais remarquons que le morphisme $\varphi'$
est le morphisme $(f \times f)^*\varphi$ dans les notations introduites à la
remarque \ref{remark-easier}. Pour (2), indiquons le morphisme
$f_0^*V \to f_1^*V$ qui donne l'isomorphisme de foncteurs. Comme $f_0$ et
$f_1$ s'insèrent tous deux dans le diagramme commutatif, il existe un unique
morphisme $r : X' \to X \times_S X$ tel que
$f_i = \text{pr}_i \circ r$. On prend alors
\begin{eqnarray*}
f_0^*V & = &
X' \times_{f_0, X} V \\
& = &
X' \times_{\text{pr}_0 \circ r, X} V \\
& = &
X' \times_{r, X \times_S X} (X \times_S X) \times_{\text{pr}_0, X} V \\
& \xrightarrow{\varphi} &
X' \times_{r, X \times_S X} (X \times_S X) \times_{\text{pr}_1, X} V \\
& = &
X' \times_{\text{pr}_1 \circ r, X} V \\
& = &
X' \times_{f_1, X} V \\
& = & f_1^*V
\end{eqnarray*}
Nous omettons la vérification.
\end{proof}

\begin{definition}
\label{definition-pullback-functor}
Avec $S, S', X, X', f, a, a', h$ comme dans le
lemme \ref{lemma-pullback}, le foncteur
$$
(V, \varphi) \longmapsto f^*(V, \varphi)
$$
construit dans ce lemme est appelé le {\it foncteur de changement de base}
sur les données de descente.
\end{definition}

\begin{lemma}[Changement de base des données de descente pour les schémas
au-dessus de familles]
\label{lemma-pullback-family}
Soient $\mathcal{U} = \{U_i \to S'\}_{i \in I}$ et
$\mathcal{V} = \{V_j \to S\}_{j \in J}$ des familles de morphismes de même
but. Soient $\alpha : I \to J$, $h : S' \to S$ et
$g_i : U_i \to V_{\alpha(i)}$ un morphisme de familles de morphismes de même
but, voir la définition \ref{sites-definition-morphism-coverings} de Sites.
\begin{enumerate}
\item Soit $(Y_j, \varphi_{jj'})$ une donnée de descente relativement à la
famille $\{V_j \to S'\}$. Le système
$$
\left(
g_i^*Y_{\alpha(i)},
(g_i \times g_{i'})^*\varphi_{\alpha(i)\alpha(i')}
\right)
$$
(avec les notations de la remarque \ref{remark-easier-family}) est une donnée
de descente relativement à $\mathcal{V}$.
\item Cette construction définit un foncteur entre les données de descente
relativement à $\mathcal{U}$ et les données de descente relativement à
$\mathcal{V}$.
\item Étant donnés un second $\alpha' : I \to J$, un morphisme
$h' : S' \to S$ et un morphisme de familles de morphismes de même but
$g'_i : U_i \to V_{\alpha'(i)}$, si $h = h'$, alors les deux foncteurs ainsi
obtenus entre les données de descente sont canoniquement isomorphes.
\item Ces foncteurs coïncident, via le
lemme \ref{lemma-family-is-one}, avec les foncteurs de changement de base
construits au lemme \ref{lemma-pullback}.
\end{enumerate}
\end{lemma}

\begin{proof}
Cela résulte du lemme \ref{lemma-pullback} par la correspondance du
lemme \ref{lemma-family-is-one}.
\end{proof}

\begin{definition}
\label{definition-pullback-functor-family}
Avec $\mathcal{U} = \{U_i \to S'\}_{i \in I}$,
$\mathcal{V} = \{V_j \to S\}_{j \in J}$, $\alpha : I \to J$,
$h : S' \to S$ et $g_i : U_i \to V_{\alpha(i)}$ comme dans le
lemme \ref{lemma-pullback-family}, le foncteur
$$
(Y_j, \varphi_{jj'}) \longmapsto
(g_i^*Y_{\alpha(i)}, (g_i \times g_{i'})^*\varphi_{\alpha(i)\alpha(i')})
$$
construit dans ce lemme est appelé le {\it foncteur de changement de base}
sur les données de descente.
\end{definition}

\noindent
Si $\mathcal{U}$ et $\mathcal{V}$ ont le même but $S$, et si
$\mathcal{U}$ raffine $\mathcal{V}$ (voir la
définition \ref{sites-definition-morphism-coverings} de Sites), mais qu'aucun
couple explicite $(\alpha, g_i)$ n'est donné, on peut néanmoins parler du
foncteur de changement de base : le
lemme \ref{lemma-pullback-family} montre que le choix de ce couple est
indifférent (à isomorphisme canonique près).


\begin{definition}
\label{definition-effective}
Soit $S$ un schéma.
Soit $f : X \to S$ un morphisme de schémas.
\begin{enumerate}
\item Étant donné un schéma $U$ au-dessus de $S$, on dispose de la
{\it donnée de descente triviale} de $U$ relativement à
$\text{id} : S \to S$, à savoir le morphisme identité de $U$.
\item Le lemme \ref{lemma-pullback} fournit une
{\it donnée de descente canonique} sur $X \times_S U$ relativement à
$X \to S$, obtenue par changement de base de la donnée de descente triviale
le long de $f$. Nous noterons souvent cette donnée
$(X \times_S U, can)$.
\item Une donnée de descente $(V, \varphi)$ relativement à $X/S$ est dite
{\it effective} si $(V, \varphi)$ est isomorphe à la donnée de descente
canonique $(X \times_S U, can)$ pour un certain schéma $U$ au-dessus de $S$.
\end{enumerate}
\end{definition}

\noindent
Ainsi, être effective signifie qu'il existe un schéma $U$ au-dessus de $S$
et un isomorphisme $\psi : V \to X \times_S U$ de $X$-schémas tel que
$\varphi$ soit égal au composé
$$
V \times_S X \xrightarrow{\psi \times \text{id}_X}
X \times_S U \times_S X =
X \times_S X \times_S U
\xrightarrow{\text{id}_X \times \psi^{-1}}
X \times_S V
$$

\begin{definition}
\label{definition-effective-family}
Soit $S$ un schéma.
Soit $\{X_i \to S\}$ une famille de morphismes de but $S$.
\begin{enumerate}
\item Étant donné un schéma $U$ au-dessus de $S$, on dispose d'une
{\it donnée de descente canonique} sur la famille de schémas
$X_i \times_S U$, obtenue par changement de base de la donnée de descente
triviale de $U$ relativement à $\{\text{id} : S \to S\}$.
Nous notons cette donnée de descente $(X_i \times_S U, can)$.
\item Une donnée de descente $(V_i, \varphi_{ij})$ relativement à
$\{X_i \to S\}$ est dite {\it effective} s'il existe un schéma $U$ au-dessus
de $S$ tel que $(V_i, \varphi_{ij})$ soit isomorphe à
$(X_i \times_S U, can)$.
\end{enumerate}
\end{definition}












\section{Pleine fidélité des foncteurs de changement de base}
\label{section-fully-faithful}

\noindent
Le foncteur de changement de base entre données de descente pour des
recouvrements fpqc est en fait pleinement fidèle. Autrement dit, les
morphismes de schémas satisfont à la descente fpqc. Le but de cette section
est de le démontrer. Le lecteur est plutôt invité à en chercher lui-même une
preuve. L'idée essentielle consiste à utiliser le
lemme \ref{lemma-fpqc-universal-effective-epimorphisms}.

\begin{lemma}
\label{lemma-surjective-flat-epi}
Tout morphisme surjectif et plat est un épimorphisme dans la catégorie des
schémas.
\end{lemma}

\begin{proof}
Soit $h : X' \to X$ un morphisme surjectif et plat, et soient
$a, b : X \to Y$ des morphismes tels que $a \circ h = b \circ h$.
Comme $h$ est surjectif, les morphismes $a$ et $b$ coïncident sur les
espaces topologiques sous-jacents. Choisissons $x' \in X'$ et posons
$x = h(x')$ et $y = a(x) = b(x)$. Considérons les homomorphismes d'anneaux
locaux
$$
a^\sharp_x, b^\sharp_x : \mathcal{O}_{Y, y} \to \mathcal{O}_{X, x}
$$
Ils deviennent égaux après composition avec l'homomorphisme local plat
$h^\sharp_{x'} : \mathcal{O}_{X, x} \to \mathcal{O}_{X', x'}$.
Comme tout homomorphisme local plat est fidèlement plat
(lemme \ref{algebra-lemma-local-flat-ff} d'Algèbre), on en déduit que
$h^\sharp_{x'}$ est injectif. Ainsi $a^\sharp_x = b^\sharp_x$, ce qui
entraîne $a = b$, comme voulu.
\end{proof}

\begin{lemma}
\label{lemma-ff-base-change-faithful}
Soit $h : S' \to S$ un morphisme de schémas surjectif et plat.
Le foncteur de changement de base
$$
\Sch/S \longrightarrow \Sch/S', \quad
X \longmapsto S' \times_S X
$$
est fidèle.
\end{lemma}

\begin{proof}
Soient $X_1$ et $X_2$ des schémas au-dessus de $S$.
Soient $\alpha, \beta : X_2 \to X_1$ des morphismes au-dessus de $S$.
Si les changements de base de $\alpha$ et $\beta$ sont égaux, on obtient le
diagramme commutatif
$$
\xymatrix{
X_2 \ar[d]^\alpha &
S' \times_S X_2 \ar[l] \ar[d] \ar[r] &
X_2 \ar[d]^\beta \\
X_1 &
S' \times_S X_1 \ar[l] \ar[r] &
X_1
}
$$
Il suffit donc de montrer que $S' \times_S X_2 \to X_2$ est un épimorphisme.
Puisqu'il est obtenu par changement de base à partir d'un morphisme surjectif
et plat, il est surjectif et plat (voir les lemmes
\ref{morphisms-lemma-base-change-surjective} et
\ref{morphisms-lemma-base-change-flat} de Morphismes). Le résultat découle
alors du lemme \ref{lemma-surjective-flat-epi}.
\end{proof}

\begin{lemma}
\label{lemma-faithful}
Dans la situation du lemme \ref{lemma-pullback}, supposons que
$f : X' \to X$ soit surjectif et plat. Alors le foncteur de changement de
base est fidèle.
\end{lemma}

\begin{proof}
Soient $(V_i, \varphi_i)$, $i = 1, 2$, deux données de descente pour
$X \to S$. Soient $\alpha, \beta : V_1 \to V_2$ des morphismes de données
de descente. Supposons que $f^*\alpha = f^*\beta$. Il faut montrer que
$\alpha = \beta$. Or $\alpha$ et $\beta$ sont des morphismes de schémas
au-dessus de $X$, et $f^*\alpha$ et $f^*\beta$ sont simplement les
changements de base de $\alpha$ et $\beta$, considérés comme morphismes
au-dessus de $X'$.
Le résultat découle donc du
lemme \ref{lemma-ff-base-change-faithful}.
\end{proof}

\noindent
Voici le lemme essentiel de cette section.

\begin{lemma}
\label{lemma-fully-faithful}
Dans la situation du lemme \ref{lemma-pullback}, supposons que
\begin{enumerate}
\item $\{f : X' \to X\}$ soit un recouvrement fpqc (par exemple si $f$ est
surjectif, plat et quasi-compact), et
\item $S = S'$.
\end{enumerate}
Alors le foncteur de changement de base est pleinement fidèle.
\end{lemma}

\begin{proof}
L'hypothèse (1) implique que $f$ est surjectif et plat. Le foncteur de
changement de base est donc fidèle d'après le lemme \ref{lemma-faithful}.
Soient $(V, \varphi)$ et $(W, \psi)$ deux données de descente relativement à
$X \to S$. Posons $(V', \varphi') = f^*(V, \varphi)$ et
$(W', \psi') = f^*(W, \psi)$.
Soit $\alpha' : V' \to W'$ un morphisme de données de descente pour $X'$
au-dessus de $S$. Il faut montrer qu'il existe un morphisme
$\alpha : V \to W$ de données de descente pour $X$ au-dessus de $S$ dont le
changement de base est $\alpha'$.

\medskip\noindent
Rappelons que $V'$ est le changement de base de $V$ par $f$ et que
$\varphi'$ est le changement de base de $\varphi$ par $f \times f$
(voir la remarque \ref{remark-easier}).
Par hypothèse, le diagramme
$$
\xymatrix{
V' \times_S X' \ar[r]_{\varphi'} \ar[d]_{\alpha' \times \text{id}} &
X' \times_S V' \ar[d]^{\text{id} \times \alpha'} \\
W' \times_S X' \ar[r]^{\psi'} &
X' \times_S W'
}
$$
est commutatif. Nous affirmons que les deux composés
$$
\xymatrix{
V' \times_V V' \ar[r]^-{\text{pr}_i} &
V' \ar[r]^{\alpha'} &
W' \ar[r] &
W
}
, \quad i = 0, 1
$$
sont égaux. Le lecteur gagnera à le démontrer lui-même plutôt qu'à lire la
suite de ce paragraphe. (Merci de nous écrire si vous trouvez un argument
simple et élégant.)
Soient $v_0, v_1$ des points de $V'$ ayant la même image $v \in V$.
Soit $x_i \in X'$ l'image de $v_i$, et soit $x$ le point de $X$ qui est
l'image de $v$ dans $X$. Autrement dit, $v_i = (x_i, v)$ dans
$V' = X' \times_X V$. Écrivons $\varphi(v, x) = (x, v')$ pour un certain
point $v'$ de $V$. Cela est possible parce que $\varphi$ est un morphisme
au-dessus de $X \times_S X$. Posons $v_i' = (x_i, v')$, qui est un point de
$V'$. Un calcul utilisant la définition de $\varphi'$ donne alors
$\varphi'(v_i, x_j) = (x_i, v'_j)$. Posons
$w_i = \alpha'(v_i)$ et $w'_i = \alpha'(v_i')$.
On peut écrire $w_i = (x_i, u_i)$ pour un certain point $u_i$ de $W$, et
$w_i' = (x_i, u'_i)$ pour un certain point $u_i'$ de $W$.
Notre assertion équivaut à $u_0 = u_1$. Un calcul formel à partir de la
définition de $\psi'$ (voir le lemme \ref{lemma-pullback}) montre que la
commutativité du diagramme précédent signifie que
$$
((x_i, x_j), \psi(u_i, x)) = ((x_i, x_j), (x, u'_j))
$$
comme points de
$(X' \times_S X') \times_{X \times_S X} (X \times_S W)$ pour tous
$i, j \in \{0, 1\}$. Il en résulte que
$\psi(u_0, x) = \psi(u_1, x)$, puis que $u_0 = u_1$ en appliquant
$\psi^{-1}$. Cela démontre l'assertion, car l'argument précédent est formel
et l'on peut prendre des points à valeurs dans un schéma (autrement dit,
prendre $(v_0, v_1) = \text{id}_{V' \times_V V'}$).

\medskip\noindent
Nous pouvons maintenant appliquer le
lemme \ref{lemma-fpqc-universal-effective-epimorphisms}.
En effet, $\{V' \to V\}$ est un recouvrement fpqc, puisqu'il est obtenu par
changement de base à partir du morphisme $f : X' \to X$. Par conséquent,
d'après le lemme \ref{lemma-fpqc-universal-effective-epimorphisms}, le
morphisme $\alpha' : V' \to W' \to W$ se factorise par un unique morphisme
$\alpha : V \to W$, dont le changement de base est nécessairement
$\alpha'$. Enfin, le diagramme
$$
\xymatrix{
V \times_S X \ar[r]_{\varphi} \ar[d]_{\alpha \times \text{id}} &
X \times_S V \ar[d]^{\text{id} \times \alpha} \\
W \times_S X \ar[r]^{\psi} & X \times_S W
}
$$
est commutatif, car son changement de base à $X' \times_S X'$ l'est, et le
morphisme $X' \times_S X' \to X \times_S X$ est surjectif et plat (appliquer
le lemme \ref{lemma-ff-base-change-faithful}). Ainsi, $\alpha$ est bien un
morphisme de données de descente
$(V, \varphi) \to (W, \psi)$, comme voulu.
\end{proof}

\noindent
Les deux lemmes qui suivent sont devenus superflus grâce à l'amélioration de
l'exposé précédent. Ils n'en restent pas moins vrais !

\begin{lemma}
\label{lemma-pullback-selfmap}
Soit $X \to S$ un morphisme de schémas.
Soit $f : X \to X$ un endomorphisme de $X$ au-dessus de $S$.
Dans ce cas, le changement de base par $f$ est isomorphe au foncteur identité
de la catégorie des données de descente relativement à $X \to S$.
\end{lemma}

\begin{proof}
Cela résulte immédiatement du lemme \ref{lemma-pullback}, puisque celui-ci
donne $f^* \cong \text{id}^*$.
\end{proof}

\begin{lemma}
\label{lemma-morphism-with-section-equivalence}
Soit $f : X' \to X$ un morphisme de schémas au-dessus d'un schéma de base
$S$. Supposons qu'il existe un morphisme $g : X \to X'$ au-dessus de $S$,
par exemple si $f$ admet une section. Alors le foncteur de changement de base
du lemme \ref{lemma-pullback} définit une équivalence entre les catégories de
données de descente relativement à $X/S$ et à $X'/S$.
\end{lemma}

\begin{proof}
Soit $g : X \to X'$ un morphisme au-dessus de $S$. Le
lemme \ref{lemma-pullback-selfmap} montre que les foncteurs
$f^* \circ g^* = (g \circ f)^*$ et $g^* \circ f^* = (f \circ g)^*$ sont
isomorphes aux foncteurs identité correspondants, comme voulu.
\end{proof}

\begin{lemma}
\label{lemma-morphism-source-faithfully-flat}
Soit $f : X \to X'$ un morphisme de schémas au-dessus d'un schéma de base
$S$. Supposons que $X \to S$ soit surjectif et plat. Alors le foncteur de
changement de base du lemme \ref{lemma-pullback} est un foncteur fidèle de la
catégorie des données de descente relativement à $X'/S$ vers la catégorie des
données de descente relativement à $X/S$.
\end{lemma}

\begin{proof}
On peut factoriser $X \to X'$ en
$X \to X \times_S X' \to X'$. Le premier morphisme est une section de la
première projection ; il induit donc une équivalence de catégories de données
de descente d'après le
lemme \ref{lemma-morphism-with-section-equivalence}. Le second morphisme est
surjectif et plat ; il induit donc un foncteur fidèle d'après le
lemme \ref{lemma-faithful}.
\end{proof}

\begin{lemma}
\label{lemma-morphism-source-fpqc-covering}
Soit $f : X \to X'$ un morphisme de schémas au-dessus d'un schéma de base
$S$. Supposons que $\{X \to S\}$ soit un recouvrement fpqc (par exemple si
$f$ est surjectif, plat et quasi-compact). Alors le foncteur de changement de
base du lemme \ref{lemma-pullback} est un foncteur pleinement fidèle de la
catégorie des données de descente relativement à $X'/S$ vers la catégorie des
données de descente relativement à $X/S$.
\end{lemma}

\begin{proof}
On peut factoriser $X \to X'$ en
$X \to X \times_S X' \to X'$. Le premier morphisme est une section de la
première projection ; il induit donc une équivalence de catégories de données
de descente d'après le
lemme \ref{lemma-morphism-with-section-equivalence}. Le second morphisme est un
recouvrement fpqc ; il induit donc un foncteur pleinement fidèle d'après le
lemme \ref{lemma-fully-faithful}.
\end{proof}

\begin{lemma}
\label{lemma-fpqc-refinement-coverings-fully-faithful}
Soit $S$ un schéma. Soient
$\mathcal{U} = \{U_i \to S\}_{i \in I}$ et
$\mathcal{V} = \{V_j \to S\}_{j \in J}$ des familles de morphismes de même
but $S$. Soient $\alpha : I \to J$, $\text{id} : S \to S$ et
$g_i : U_i \to V_{\alpha(i)}$ un morphisme de familles de morphismes de même
but ; voir la définition \ref{sites-definition-morphism-coverings} de Sites.
Supposons que, pour tout $j \in J$, la famille
$\{g_i : U_i \to V_j\}_{\alpha(i) = j}$ soit un recouvrement fpqc de $V_j$.
Alors le foncteur de changement de base
$$
\text{données de descente relativement à }
\mathcal{V}
\longrightarrow
\text{données de descente relativement à }
\mathcal{U}
$$
du lemme \ref{lemma-pullback-family} est pleinement fidèle.
\end{lemma}

\begin{proof}
Considérons le morphisme de schémas
$$
g :
X = \coprod\nolimits_{i \in I} U_i
\longrightarrow
Y = \coprod\nolimits_{j \in J} V_j
$$
au-dessus de $S$ qui, sur la composante d'indice $i$, prend ses valeurs dans
la composante d'indice $\alpha(i)$ au moyen du morphisme
$g_{\alpha(i)}$. Nous affirmons que $\{g : X \to Y\}$ est un recouvrement
fpqc de schémas. En effet, d'après le
lemme \ref{topologies-lemma-disjoint-union-is-fpqc-covering} de Topologies,
pour tout $j$,
le morphisme $\{\coprod_{\alpha(i) = j} U_i \to V_j\}$ est un recouvrement
fpqc. Ainsi, pour tout ouvert affine $V \subset V_j$ (que l'on peut considérer
comme un ouvert affine de $Y$), il existe un nombre fini d'ouverts affines
$W_1, \ldots, W_n \subset \coprod_{\alpha(i) = j} U_i$ (que l'on peut
considérer comme des ouverts affines de $X$) tels que
$V = \bigcup_{i = 1, \ldots, n} g(W_i)$. On dispose donc de suffisamment
d'ouverts affines de $Y$ recouverts par un nombre fini d'ouverts affines de
$X$ pour appliquer la partie (3) du
lemme \ref{topologies-lemma-recognize-fpqc-covering} de Topologies, ce qui prouve
l'assertion. Notons $DD(X/S)$, resp.\ $DD(\mathcal{U})$, la catégorie des
données de descente relativement à $X/S$, resp.\ à $\mathcal{U}$, et de même
pour $Y/S$ et $\mathcal{V}$. Considérons le diagramme
$$
\xymatrix{
DD(Y/S) \ar[r] & DD(X/S) \\
DD(\mathcal{V}) \ar[u]^{\text{lemme }\ref{lemma-family-is-one}} \ar[r] &
DD(\mathcal{U}) \ar[u]_{\text{lemme }\ref{lemma-family-is-one}}
}
$$
Ce diagramme est commutatif, comme le montre la preuve du
lemme \ref{lemma-pullback-family}. Les flèches verticales sont des
équivalences. Le résultat découle donc du lemme \ref{lemma-fully-faithful},
qui montre que la flèche horizontale supérieure du diagramme est pleinement
fidèle.
\end{proof}

\noindent
Le lemme suivant montre que, pour vérifier qu'une donnée est effective, on
peut toujours raffiner pour la topologie de Zariski la famille donnée de
morphismes de but $S$.

\begin{lemma}
\label{lemma-Zariski-refinement-coverings-equivalence}
Soit $S$ un schéma. Soient
$\mathcal{U} = \{U_i \to S\}_{i \in I}$ et
$\mathcal{V} = \{V_j \to S\}_{j \in J}$ des familles de morphismes de même
but $S$. Soient $\alpha : I \to J$, $\text{id} : S \to S$ et
$g_i : U_i \to V_{\alpha(i)}$ un morphisme de familles de morphismes de même
but ; voir la définition \ref{sites-definition-morphism-coverings} de Sites.
Supposons que, pour tout $j \in J$, la famille
$\{g_i : U_i \to V_j\}_{\alpha(i) = j}$ soit un recouvrement de Zariski (voir
la définition \ref{topologies-definition-zariski-covering} de Topologies) de
$V_j$. Alors le foncteur de
changement de base
$$
\text{données de descente relativement à }
\mathcal{V}
\longrightarrow
\text{données de descente relativement à }
\mathcal{U}
$$
du lemme \ref{lemma-pullback-family} est une équivalence de catégories.
En particulier, la catégorie des schémas au-dessus de $S$ est équivalente à
la catégorie des données de descente relativement à n'importe quel
recouvrement de Zariski de $S$.
\end{lemma}

\begin{proof}
Le foncteur est pleinement fidèle d'après le
lemme \ref{lemma-fpqc-refinement-coverings-fully-faithful}.
Indiquons comment montrer qu'il est essentiellement surjectif.
Soit $(X_i, \varphi_{ii'})$ une donnée de descente relativement à
$\mathcal{U}$. Fixons $j \in J$ et posons
$I_j = \{i \in I \mid \alpha(i) = j\}$. Pour $i, i' \in I_j$, il existe un
morphisme canonique
$$
c_{ii'} : U_i \times_{g_i, V_j, g_{i'}} U_{i'} \to U_i \times_S U_{i'}.
$$
On peut donc effectuer le changement de base de $\varphi_{ii'}$ par ce
morphisme et poser $\psi_{ii'} = c_{ii'}^*\varphi_{ii'}$ pour
$i, i' \in I_j$. On obtient ainsi une donnée de descente
$(X_i, \psi_{ii'})$ relativement au recouvrement de Zariski
$\{g_i : U_i \to V_j\}_{i \in I_j}$. Remarquons que $\psi_{ii'}$ est un
isomorphisme de l'ouvert $X_{i, U_i \times_{V_j} U_{i'}}$ de $X_i$ sur
l'ouvert correspondant de $X_{i'}$. Il résulte de la
section \ref{schemes-section-glueing-schemes} de Schémas que l'on peut recoller
$(X_i, \psi_{ii'})$ en un schéma $Y_j$ au-dessus de $V_j$. De plus, les
morphismes $\varphi_{ii'}$, pour $i \in I_j$ et $i' \in I_{j'}$, se
recollent en un morphisme
$\varphi_{jj'} : Y_j \times_S V_{j'} \to V_j \times_S Y_{j'}$ satisfaisant
la condition de cocycle (détails omis). On obtient ainsi la donnée de descente
recherchée $(Y_j, \varphi_{jj'})$ relativement à $\mathcal{V}$.
\end{proof}

\begin{lemma}
\label{lemma-refine-coverings-fully-faithful}
Soit $S$ un schéma. Soient
$\mathcal{U} = \{U_i \to S\}_{i \in I}$ et
$\mathcal{V} = \{V_j \to S\}_{j \in J}$ des recouvrements fpqc de $S$.
Si $\mathcal{U}$ est un raffinement de $\mathcal{V}$, alors le foncteur de
changement de base
$$
\text{données de descente relativement à }
\mathcal{V}
\longrightarrow
\text{données de descente relativement à }
\mathcal{U}
$$
est pleinement fidèle. En particulier, la catégorie des schémas au-dessus de
$S$ s'identifie à une sous-catégorie pleine de la catégorie des données de
descente relativement à tout recouvrement fpqc de $S$.
\end{lemma}

\begin{proof}
Considérons le recouvrement fpqc
$\mathcal{W} = \{U_i \times_S V_j \to S\}_{(i, j) \in I \times J}$ de $S$.
Il raffine à la fois $\mathcal{U}$ et $\mathcal{V}$. On a donc un diagramme
$2$-commutatif de foncteurs et de catégories
$$
\xymatrix{
DD(\mathcal{V}) \ar[rd] \ar[rr] & & DD(\mathcal{U}) \ar[ld] \\
& DD(\mathcal{W}) &
}
$$
Les notations sont celles de la preuve du
lemme \ref{lemma-fpqc-refinement-coverings-fully-faithful}, et la
commutativité résulte de la partie (3) du lemme \ref{lemma-pullback-family}.
Il suffit donc manifestement de montrer que les foncteurs
$DD(\mathcal{V}) \to DD(\mathcal{W})$ et
$DD(\mathcal{U}) \to DD(\mathcal{W})$ sont pleinement fidèles. Cela découle
du lemme \ref{lemma-fpqc-refinement-coverings-fully-faithful}, comme voulu.
\end{proof}

\begin{remark}
\label{remark-morphisms-of-schemes-satisfy-fpqc-descent}
Le lemme \ref{lemma-refine-coverings-fully-faithful} affirme que les
morphismes de schémas satisfont à la descente fpqc. Autrement dit, étant
donnés un schéma $S$ et des schémas $X$, $Y$ au-dessus de $S$, le foncteur
$$
(\Sch/S)^{opp} \longrightarrow \textit{Ensembles},
\quad
T \longmapsto \Mor_T(X_T, Y_T)
$$
satisfait à la condition de faisceau pour la topologie fpqc. Le cas le plus
simple est le suivant. Supposons que $T \to S$ soit un morphisme surjectif et
plat entre schémas affines. Soit $\psi_0 : X_T \to Y_T$ un morphisme de
schémas au-dessus de $T$, compatible avec les données de descente canoniques.
Il existe alors un unique morphisme $\psi : X \to Y$ dont le changement de
base à $T$ est $\psi_0$. Ce cas particulier découle en fait directement du
lemme \ref{lemma-fully-faithful}. Ce lemme est lui-même une conséquence
formelle des deux faits suivants : (a) le foncteur de changement de base par
un morphisme fidèlement plat est fidèle, voir le
lemme \ref{lemma-ff-base-change-faithful}, et (b) un schéma satisfait à la
condition de faisceau pour la topologie fpqc, voir le
lemme \ref{lemma-fpqc-universal-effective-epimorphisms}.
\end{remark}

\begin{lemma}
\label{lemma-effective-for-fpqc-is-local-upstairs}
Soit $X \to S$ un morphisme de schémas surjectif, quasi-compact et plat.
Soit $(V, \varphi)$ une donnée de descente relativement à $X/S$. Supposons que,
pour tout $v \in V$, il existe un sous-schéma ouvert $v \in W \subset V$ tel
que $\varphi(W \times_S X) \subset X \times_S W$ et que la donnée de descente
$(W, \varphi|_{W \times_S X})$ soit effective. Alors $(V, \varphi)$ est
effective.
\end{lemma}

\begin{proof}
Soit $V = \bigcup W_i$ un recouvrement ouvert tel que
$\varphi(W_i \times_S X) \subset X \times_S W_i$ et que la donnée de descente
$(W_i, \varphi|_{W_i \times_S X})$ soit effective. Soit $U_i \to S$ un
schéma, et soit
$\alpha_i : (X \times_S U_i, can) \to (W_i, \varphi|_{W_i \times_S X})$ un
isomorphisme de données de descente. Pour chaque paire d'indices $(i, j)$,
considérons l'ouvert
$\alpha_i^{-1}(W_i \cap W_j) \subset X \times_S U_i$. Comme toutes les
données sont compatibles avec les données de descente et que
$\{X \to S\}$ est un recouvrement fpqc, le
lemme \ref{lemma-open-fpqc-covering} fournit un ouvert
$U_{ij} \subset U_i$ tel que
$\alpha_i^{-1}(W_i \cap W_j) = X \times_S U_{ij}$. Le morphisme identité de
$W_i \cap W_j$ est compatible avec les données de descente ; il provient donc
d'un unique morphisme $\varphi_{ij} : U_{ij} \to U_{ji}$ au-dessus de $S$
(voir la remarque \ref{remark-morphisms-of-schemes-satisfy-fpqc-descent}).
Alors $(U_i, U_{ij}, \varphi_{ij})$ est une donnée de recollement au sens de
la section \ref{schemes-section-glueing-schemes} de Schémas (preuve omise). On peut
donc supposer qu'il existe un schéma $U$ au-dessus de $S$ tel que
$U_i \subset U$ soit ouvert, que $U_{ij} = U_i \cap U_j$ et que
$\varphi_{ij} = \text{id}_{U_i \cap U_j}$ ; voir le
lemme \ref{schemes-lemma-glue} de Schémas. Après changement de base à $X$, les
$\alpha_i$ fournissent l'isomorphisme recherché
$\alpha : X \times_S U \to V$.
\end{proof}







\section{Descente de types de morphismes}
\label{section-descending-types-morphisms}

\noindent
Dans la suite, nous étudions si les données de descente de schémas
relativement à un recouvrement fpqc sont effectives. La première remarque est
que ce n'est pas toujours le cas. Nous le verrons dans
l'exemple \ref{spaces-example-non-representable-descent} d'Espaces algébriques. Même les morphismes
projectifs ne satisfont pas toujours à la descente pour les recouvrements
fpqc, d'après le
lemme \ref{examples-lemma-non-effective-descent-projective} d'Exemples.

\medskip\noindent
En revanche, lorsque les schémas que l'on cherche à descendre sont
particulièrement simples, il arrive que, pour des classes entières de schémas,
les données de descente soient effectives. Nous introduisons ici une
terminologie qui décrit ce phénomène abstraitement, bien qu'elle puisse prêter
à confusion si elle est mal employée par la suite.

\begin{definition}
\label{definition-descending-types-morphisms}
Soit $\mathcal{P}$ une propriété de morphismes de schémas au-dessus d'une
base. Soit
$\tau \in \{Zariski, fpqc, fppf, \etale, lisse, syntomique\}$. On dit que
{\it les morphismes de type $\mathcal{P}$ satisfont à la descente pour les
recouvrements pour la topologie $\tau$} si, pour tout recouvrement
$\mathcal{U} : \{U_i \to S\}_{i \in I}$ pour la topologie $\tau$ (voir la
section \ref{topologies-section-procedure} de Topologies), toute donnée de descente
$(X_i, \varphi_{ij})$ relativement à $\mathcal{U}$ telle que chaque morphisme
$X_i \to U_i$ vérifie la propriété $\mathcal{P}$ est effective.
\end{definition}

\noindent
Remarquons que, dans chacun des cas considérés, nous avons déjà vu que le
foncteur des schémas au-dessus de $S$ vers les données de descente relativement
à $\mathcal{U}$ est pleinement fidèle : cela résulte du
lemme \ref{lemma-refine-coverings-fully-faithful} et du fait, établi dans
Topologies, que tout recouvrement $\tau$ est aussi un recouvrement fpqc. Nous
avons également vu que les données de descente sont toujours effectives pour
les recouvrements de Zariski
(lemme \ref{lemma-Zariski-refinement-coverings-equivalence}). Pour éviter les
raisonnements erronés, il peut être prudent de n'étudier la notion précédente
que lorsque $\mathcal{P}$ est préservée par tout changement de base, ou du
moins locale sur la base pour la topologie $\tau$ (voir la
définition \ref{definition-property-morphisms-local}). En effet, on risquerait
sinon d'invoquer $\mathcal{P}$ au cours d'une descente encore inachevée.

\medskip\noindent
Voici l'inévitable lemme qui ramène la question au cas d'un recouvrement donné
par un unique morphisme entre schémas affines.

\begin{lemma}
\label{lemma-descending-types-morphisms}
Soit $\mathcal{P}$ une propriété de morphismes de schémas au-dessus d'une
base. Soit $\tau \in \{fpqc, fppf, \etale, lisse, syntomique\}$. Supposons que
\begin{enumerate}
\item $\mathcal{P}$ soit préservée par tout changement de base (voir la
définition \ref{schemes-definition-preserved-by-base-change} de Schémas),
\item si $Y_j \to V_j$, $j = 1, \ldots, m$, vérifient $\mathcal{P}$, alors
$\coprod Y_j \to \coprod V_j$ la vérifie aussi, et
\item pour tout morphisme surjectif de schémas affines $X \to S$ qui est
plat, plat de présentation finie, \'etale, lisse ou syntomique suivant que
$\tau$ vaut fpqc, fppf, \'etale, lisse ou syntomique, toute donnée de descente
$(V, \varphi)$ relativement à $X$ au-dessus de $S$ telle que $\mathcal{P}$
soit vérifiée par $V \to X$ est effective.
\end{enumerate}
Alors les morphismes de type $\mathcal{P}$ satisfont à la descente pour les
recouvrements pour la topologie $\tau$.
\end{lemma}

\begin{proof}
Soit $S$ un schéma. Soit
$\mathcal{U} = \{\varphi_i : U_i \to S\}_{i \in I}$ un recouvrement de $S$
pour la topologie $\tau$. Soit $(X_i, \varphi_{ii'})$ une donnée de descente relativement
à $\mathcal{U}$, et supposons que chaque morphisme $X_i \to U_i$ vérifie la
propriété $\mathcal{P}$. Il faut montrer qu'il existe un schéma $X \to S$ tel
que $(X_i, \varphi_{ii'}) \cong (U_i \times_S X, can)$.

\medskip\noindent
Avant d'entamer la preuve proprement dite, remarquons que, pour toute famille
de morphismes $\mathcal{V} : \{V_j \to S\}$ et tout morphisme de familles
$\mathcal{V} \to \mathcal{U}$, si l'on effectue le changement de base de la
donnée de descente $(X_i, \varphi_{ii'})$ en une donnée de descente
$(Y_j, \varphi_{jj'})$ au-dessus de $\mathcal{V}$, chacun des morphismes
$Y_j \to V_j$ vérifie encore la propriété $\mathcal{P}$. Cela résulte de
l'hypothèse (1), selon laquelle $\mathcal{P}$ est préservée par tout changement
de base, et de la définition du changement de base (voir la
définition \ref{definition-pullback-functor-family}). Nous utiliserons ce fait
sans le rappeler davantage.

\medskip\noindent
Commençons par démontrer le lemme lorsque $S$ est affine. D'après les lemmes
\ref{topologies-lemma-fpqc-affine},
\ref{topologies-lemma-fppf-affine},
\ref{topologies-lemma-etale-affine},
\ref{topologies-lemma-smooth-affine} ou
\ref{topologies-lemma-syntomic-affine} de Topologies, il existe un recouvrement
standard pour la topologie $\tau$
$\mathcal{V} : \{V_j \to S\}_{j = 1, \ldots, m}$ qui raffine
$\mathcal{U}$. Le foncteur de changement de base
$DD(\mathcal{U}) \to DD(\mathcal{V})$ entre catégories de données de descente
est pleinement fidèle d'après le
lemme \ref{lemma-refine-coverings-fully-faithful}. Il suffit donc de montrer
que la donnée de descente au-dessus du recouvrement standard pour la topologie $\tau$
$\mathcal{V}$ est effective. L'hypothèse (2) montre que
$\coprod Y_j \to \coprod V_j$ vérifie la propriété $\mathcal{P}$. Le
lemme \ref{lemma-family-is-one} ramène alors la question au recouvrement
$\{\coprod_{j = 1, \ldots, m} V_j \to S\}$, pour lequel le résultat a été
supposé dans l'hypothèse (3) du lemme. Le lemme est donc démontré lorsque $S$
est affine.

\medskip\noindent
Supposons maintenant $S$ quelconque. Soit $V \subset S$ un ouvert affine.
Les propriétés d'un site montrent que la famille
$\mathcal{U}_V = \{V \times_S U_i \to V\}_{i \in I}$ est un recouvrement de
$V$ pour la topologie $\tau$. Notons $(X_i, \varphi_{ii'})_V$ la restriction, ou le
changement de base, de la donnée de descente donnée à $\mathcal{U}_V$. D'après
ce qui précède, on obtient un schéma $X_V$ au-dessus de $V$ dont la donnée de
descente canonique relativement à $\mathcal{U}_V$ est isomorphe à
$(X_i, \varphi_{ii'})_V$. Supposons que $V' \subset V$ soit un ouvert affine
de $V$. Alors $X_{V'}$ et $V' \times_V X_V$ ont tous deux une donnée de
descente canonique isomorphe à $(X_i, \varphi_{ii'})_{V'}$. Le
lemme \ref{lemma-refine-coverings-fully-faithful} fournit donc de nouveau un
morphisme canonique $\rho^V_{V'} : X_{V'} \to X_V$ au-dessus de $S$ qui
identifie $X_{V'}$ à l'image inverse de $V'$ dans $X_V$. Nous omettons de
vérifier que, pour des ouverts affines $V'' \subset V' \subset V$ de $S$, on a
$\rho^V_{V''} = \rho^V_{V'} \circ \rho^{V'}_{V''}$.

\medskip\noindent
D'après le lemme \ref{constructions-lemma-relative-glueing} de Constructions, les
données $(X_V, \rho^V_{V'})$ se recollent en un schéma $X \to S$. De plus, on
dispose d'isomorphismes $V \times_S X \to X_V$ qui redonnent les morphismes
$\rho^V_{V'}$. En déroulant la construction des schémas $X_V$, on obtient des
isomorphismes
$$
V \times_S U_i \times_S X
\longrightarrow
V \times_S X_i
$$
compatibles avec les morphismes $\varphi_{ii'}$ et avec la restriction aux
ouverts affines plus petits de $X$. Il en résulte que la donnée de descente
canonique sur $U_i \times_S X$ est isomorphe à la donnée de descente donnée,
ce qui conclut la preuve.
\end{proof}










\section{Descente des morphismes affines}
\label{section-affine}

\noindent
Dans cette section, nous montrons que
``les morphismes affines satisfont à la descente pour les recouvrements
pour la topologie fpqc''. En voici l'énoncé formel.

\begin{lemma}
\label{lemma-affine}
Soit $S$ un schéma. Soit $\{X_i \to S\}_{i\in I}$ un recouvrement fpqc ; voir
la définition \ref{topologies-definition-fpqc-covering} de Topologies. Soit
$(V_i/X_i, \varphi_{ij})$ une donnée de descente relativement à
$\{X_i \to S\}$. Si chaque morphisme $V_i \to X_i$ est affine, alors la
donnée de descente est effective.
\end{lemma}

\begin{proof}
La propriété d'être affine pour un morphisme de schémas est locale sur la base
et préservée par tout changement de base ; voir les lemmes
\ref{morphisms-lemma-characterize-affine} et
\ref{morphisms-lemma-base-change-affine} de Morphismes. Le
lemme \ref{lemma-descending-types-morphisms} s'applique donc, et il suffit de
démontrer l'énoncé lorsque le recouvrement fpqc est donné par un unique
morphisme plat et surjectif de schémas affines $\{X \to S\}$. Écrivons
$X = \Spec(A)$ et $S = \Spec(R)$, de sorte que $R \to A$ soit un
homomorphisme d'anneaux fidèlement plat. Soit $(V, \varphi)$ une donnée de
descente relativement à $X$ au-dessus de $S$, et supposons que $V \to X$ soit
affine. Puisque $V \to X$ est affine, on a $V = \Spec(B)$ pour une
$A$-algèbre $B$ (voir la
définition \ref{morphisms-definition-affine} de Morphismes). L'isomorphisme $\varphi$
correspond à un isomorphisme d'anneaux
$$
\varphi^\sharp :
B \otimes_R A \longleftarrow A \otimes_R B
$$
d'$A \otimes_R A$-algèbres. La condition de cocycle sur $\varphi$ signifie que
le diagramme
$$
\xymatrix{
B \otimes_R A \otimes_R A & &
A \otimes_R A \otimes_R B \ar[ll] \ar[ld]\\
& A \otimes_R B \otimes_R A \ar[lu] &
}
$$
est commutatif. En inversant ces flèches, on obtient une donnée de descente
pour les modules relativement à $R \to A$, au sens de la
définition \ref{definition-descent-datum-modules}. On peut donc appliquer la
proposition \ref{proposition-descent-module} pour obtenir un $R$-module
$C = \Ker(B \to A \otimes_R B)$ et un isomorphisme
$A \otimes_R C \cong B$ compatible avec les données de descente. Pour toute
paire $c, c' \in C$, le produit $cc'$ dans $B$ appartient à $C$, puisque le
morphisme $\varphi$ est un homomorphisme d'algèbres. Ainsi, $C$ est une
$R$-algèbre dont le changement de base à $A$ est isomorphe à $B$ de manière
compatible avec les données de descente. En appliquant $\Spec$, on obtient un
schéma $U$ au-dessus de $S$ tel que
$(V, \varphi) \cong (X \times_S U, can)$, comme voulu.
\end{proof}

\begin{lemma}
\label{lemma-closed-immersion}
Soit $S$ un schéma. Soit $\{X_i \to S\}_{i\in I}$ un recouvrement fpqc ; voir
la définition \ref{topologies-definition-fpqc-covering} de Topologies. Soit
$(V_i/X_i, \varphi_{ij})$ une donnée de descente relativement à
$\{X_i \to S\}$. Si chaque morphisme $V_i \to X_i$ est une immersion fermée,
alors la donnée de descente est effective.
\end{lemma}

\begin{proof}
Cela résulte du fait qu'une immersion fermée est un morphisme affine
(lemme \ref{morphisms-lemma-closed-immersion-affine} de Morphismes), si bien que
le lemme \ref{lemma-affine} s'applique.
\end{proof}


\section{Descente des morphismes quasi-affines}
\label{section-quasi-affine}

\noindent
Dans cette section, nous montrons que
``les morphismes quasi-affines satisfont à la descente pour les recouvrements
pour la topologie fpqc''. En voici l'énoncé formel.

\begin{lemma}
\label{lemma-quasi-affine}
Soit $S$ un schéma. Soit $\{X_i \to S\}_{i\in I}$ un recouvrement fpqc ; voir
la définition \ref{topologies-definition-fpqc-covering} de Topologies. Soit
$(V_i/X_i, \varphi_{ij})$ une donnée de descente relativement à
$\{X_i \to S\}$. Si chaque morphisme $V_i \to X_i$ est quasi-affine, alors la
donnée de descente est effective.
\end{lemma}

\begin{proof}
La propriété d'être quasi-affine pour un morphisme de schémas est préservée
par tout changement de base ; voir les lemmes
\ref{morphisms-lemma-characterize-quasi-affine} et
\ref{morphisms-lemma-base-change-quasi-affine} de Morphismes. Le
lemme \ref{lemma-descending-types-morphisms} s'applique donc, et il suffit de
démontrer l'énoncé lorsque le recouvrement fpqc est donné par un unique
morphisme plat et surjectif de schémas affines $\{X \to S\}$. Écrivons
$X = \Spec(A)$ et $S = \Spec(R)$, de sorte que $R \to A$ soit un
homomorphisme d'anneaux fidèlement plat. Soit $(V, \varphi)$ une donnée de
descente relativement à $X$ au-dessus de $S$, et supposons que
$\pi : V \to X$ soit quasi-affine.

\medskip\noindent
D'après le
lemme \ref{morphisms-lemma-characterize-quasi-affine} de Morphismes, cela signifie que
$$
V \longrightarrow \underline{\Spec}_X(\pi_*\mathcal{O}_V) = W
$$
est une immersion ouverte quasi-compacte de schémas au-dessus de $X$. Les
projections $\text{pr}_i : X \times_S X \to X$ sont plates, et l'on a donc
$$
\text{pr}_0^*\pi_*\mathcal{O}_V =
(\pi \times \text{id}_X)_*\mathcal{O}_{V \times_S X}, \quad
\text{pr}_1^*\pi_*\mathcal{O}_V =
(\text{id}_X \times \pi)_*\mathcal{O}_{X \times_S V}
$$
par changement de base plat (lemme
\ref{coherent-lemma-flat-base-change-cohomology} de Cohomologie des schémas). Ainsi,
l'isomorphisme $\varphi : V \times_S X \to X \times_S V$ (qui est un
isomorphisme au-dessus de $X \times_S X$) induit un isomorphisme de faisceaux
quasi-cohérents d'algèbres
$$
\varphi^\sharp :
\text{pr}_0^*\pi_*\mathcal{O}_V
\longrightarrow
\text{pr}_1^*\pi_*\mathcal{O}_V
$$
sur $X \times_S X$. La condition de cocycle pour $\varphi$ implique la
condition de cocycle pour $\varphi^\sharp$. Autrement dit, on obtient une
donnée de descente $\varphi'$ sur le schéma affine $W$ relativement à $X$
au-dessus de $S$, avec la propriété supplémentaire que le morphisme $V \to W$
est un morphisme de données de descente. Le lemme \ref{lemma-affine} (ou
l'effectivité de la descente pour les algèbres quasi-cohérentes) fournit donc
un schéma $U' \to S$ et un isomorphisme
$(W, \varphi') \cong (X \times_S U', can)$ de données de descente. Notons au
passage que $U'$ est affine d'après le
lemme \ref{lemma-descending-property-affine}.

\medskip\noindent
On peut maintenant considérer $V$ comme l'ouvert quasi-compact
$V \subset X \times_S U'$ qui est stable par la donnée de descente
$$
can : X \times_S U' \times_S X \to X \times_S X \times_S U',
(x_0, u', x_1) \mapsto (x_0, x_1, u').
$$
Autrement dit, $(x_0, u') \in V \Rightarrow (x_1, u') \in V$ pour tous
$x_0, x_1, u'$ ayant la même image dans $S$. Comme $X \to S$ est surjectif,
on voit immédiatement que $V$ est l'image inverse d'un sous-ensemble
$U \subset U'$ par le morphisme $X \times_S U' \to U'$. Comme $X \to S$ est
quasi-compact, plat et surjectif, le morphisme $X \times_S U' \to U'$ est lui
aussi quasi-compact, plat et surjectif. D'après le
lemme \ref{morphisms-lemma-fpqc-quotient-topology} de Morphismes, le sous-ensemble
$U \subset U'$ est donc ouvert, ce qui achève la preuve.
\end{proof}












\section{Données de descente en termes de faisceaux}
\label{section-descent-data-sheaves}


\noindent
Voici une autre manière d'envisager les données de descente dans le cas d'un
recouvrement sur un site.

\begin{lemma}
\label{lemma-descent-data-sheaves}
Soit $\tau \in \{Zariski, fppf, \etale, lisse, syntomique\}$\footnote{L'absence
de fpqc n'est pas une coquille. Voir la discussion à la
section \ref{topologies-section-fpqc} de Topologies.}. Soit $\Sch_\tau$ un grand
site pour la topologie $\tau$. Soit $S \in \Ob(\Sch_\tau)$. Soit
$\{S_i \to S\}_{i \in I}$ un recouvrement du site $(\Sch/S)_\tau$. Il existe
une équivalence de catégories
$$
\left\{
\begin{matrix}
\text{données de descente }(X_i, \varphi_{ii'})\text{ telles que}\\
\text{chaque }X_i \in \Ob((\Sch/S)_\tau)
\end{matrix}
\right\}
\leftrightarrow
\left\{
\begin{matrix}
\text{faisceaux }F\text{ sur }(\Sch/S)_\tau\text{ tels que}\\
\text{chaque }h_{S_i} \times F\text{ est représentable}
\end{matrix}
\right\}.
$$
De plus,
\begin{enumerate}
\item les objets représentant $h_{S_i} \times F$ dans le membre de droite
correspondent aux schémas $X_i$ dans le membre de gauche, et
\item le faisceau $F$ est représentable si et seulement si la donnée de
descente correspondante $(X_i, \varphi_{ii'})$ est effective.
\end{enumerate}
\end{lemma}

\begin{proof}
Nous avons vu à la section \ref{section-fpqc-universal-effective-epimorphisms}
que les préfaisceaux représentables sont des faisceaux sur le site
$(\Sch/S)_\tau$. De plus, le lemme de Yoneda (lemme
\ref{categories-lemma-yoneda} de Catégories) garantit que les morphismes entre
faisceaux représentables correspondent bijectivement aux morphismes entre les
objets qui les représentent. Nous utiliserons ces remarques sans plus les
mentionner au cours de la démonstration.

\medskip\noindent
Construisons le foncteur de droite à gauche. Soit $F$ un faisceau sur
$(\Sch/S)_\tau$ tel que chaque $h_{S_i} \times F$ soit représentable. Dans ce
cas, prenons pour $X_i$ un objet représentant dans $(\Sch/S)_\tau$. Il est muni
d'un morphisme $X_i \to S_i$. Alors $X_i \times_S S_{i'}$ et
$S_i \times_S X_{i'}$ représentent tous deux le faisceau
$h_{S_i} \times F \times h_{S_{i'}}$ ; on obtient donc un isomorphisme
$$
\varphi_{ii'} : X_i \times_S S_{i'} \to S_i \times_S X_{i'}.
$$
On vérifie immédiatement que les morphismes $\varphi_{ii'}$ sont définis
au-dessus de $S_i \times_S S_{i'}$ et satisfont à la condition de cocycle. Le
foncteur de droite à gauche est donné par la construction
$F \mapsto (X_i, \varphi_{ii'})$.

\medskip\noindent
Construisons un foncteur de gauche à droite. Pour chaque $i$, notons $F_i$ le
faisceau $h_{X_i}$. Les isomorphismes $\varphi_{ii'}$ donnent des isomorphismes
$$
\varphi_{ii'} :
F_i \times h_{S_{i'}}
\longrightarrow
h_{S_i} \times F_{i'}
$$
au-dessus de $h_{S_i} \times h_{S_{i'}}$. Posons $F$ égal au coégalisateur du
diagramme suivant
$$
\xymatrix{
\coprod_{i, i'} F_i \times h_{S_{i'}}
\ar@<1ex>[rr]^-{\text{pr}_0}
\ar@<-1ex>[rr]_-{\text{pr}_1 \circ \varphi_{ii'}}
& &
\coprod_i F_i \ar[r]
&
F
}
$$
La condition de cocycle garantit que $h_{S_i} \times F$ est isomorphe à
$F_i$, donc représentable. Le foncteur de gauche à droite est donné par la
construction $(X_i, \varphi_{ii'}) \mapsto F$.

\medskip\noindent
Nous omettons de vérifier que ces constructions donnent des foncteurs
quasi-inverses l'un de l'autre. Les assertions finales (1) et (2) découlent
des constructions.
\end{proof}

\begin{remark}
\label{remark-what-product-means}
Dans l'énoncé du lemme \ref{lemma-descent-data-sheaves}, la condition que
$h_{S_i} \times F$ soit représentable équivaut à ce que la restriction de $F$
à $(\Sch/S_i)_\tau$ soit représentable.
\end{remark}
















\begin{multicols}{2}[\section{Autres chapitres}]
\noindent
Préliminaires
\begin{enumerate}
\item \hyperref[introduction-section-phantom]{Introduction}
\item \hyperref[conventions-section-phantom]{Conventions}
\item \hyperref[sets-section-phantom]{Théorie des ensembles}
\item \hyperref[categories-section-phantom]{Catégories}
\item \hyperref[topology-section-phantom]{Topologie}
\item \hyperref[sheaves-section-phantom]{Faisceaux sur les espaces}
\item \hyperref[sites-section-phantom]{Sites et faisceaux}
\item \hyperref[stacks-section-phantom]{Champs}
\item \hyperref[fields-section-phantom]{Corps}
\item \hyperref[algebra-section-phantom]{Algèbre commutative}
\item \hyperref[brauer-section-phantom]{Groupes de Brauer}
\item \hyperref[homology-section-phantom]{Algèbre homologique}
\item \hyperref[derived-section-phantom]{Catégories dérivées}
\item \hyperref[simplicial-section-phantom]{Méthodes simpliciales}
\item \hyperref[more-algebra-section-phantom]{Compléments d'algèbre}
\item \hyperref[smoothing-section-phantom]{Lissage des morphismes d'anneaux}
\item \hyperref[modules-section-phantom]{Faisceaux de modules}
\item \hyperref[sites-modules-section-phantom]{Modules sur les sites}
\item \hyperref[injectives-section-phantom]{Objets injectifs}
\item \hyperref[cohomology-section-phantom]{Cohomologie des faisceaux}
\item \hyperref[sites-cohomology-section-phantom]{Cohomologie sur les sites}
\item \hyperref[dga-section-phantom]{Algèbre différentielle graduée}
\item \hyperref[dpa-section-phantom]{Algèbre à puissances divisées}
\item \hyperref[sdga-section-phantom]{Faisceaux différentiels gradués}
\item \hyperref[hypercovering-section-phantom]{Hyperrecouvrements}
\end{enumerate}
Schémas
\begin{enumerate}
\setcounter{enumi}{25}
\item \hyperref[schemes-section-phantom]{Schémas}
\item \hyperref[constructions-section-phantom]{Constructions de schémas}
\item \hyperref[properties-section-phantom]{Propriétés des schémas}
\item \hyperref[morphisms-section-phantom]{Morphismes de schémas}
\item \hyperref[coherent-section-phantom]{Cohomologie des schémas}
\item \hyperref[divisors-section-phantom]{Diviseurs}
\item \hyperref[limits-section-phantom]{Limites de schémas}
\item \hyperref[varieties-section-phantom]{Variétés}
\item \hyperref[topologies-section-phantom]{Topologies sur les schémas}
\item \hyperref[descent-section-phantom]{Descente}
\item \hyperref[perfect-section-phantom]{Catégories dérivées de schémas}
\item \hyperref[more-morphisms-section-phantom]{Compléments sur les morphismes}
\item \hyperref[flat-section-phantom]{Compléments sur la platitude}
\item \hyperref[groupoids-section-phantom]{Schémas en groupoïdes}
\item \hyperref[more-groupoids-section-phantom]{Compléments sur les schémas en groupoïdes}
\item \hyperref[etale-section-phantom]{Morphismes étales de schémas}
\end{enumerate}
Thèmes de théorie des schémas
\begin{enumerate}
\setcounter{enumi}{41}
\item \hyperref[chow-section-phantom]{Homologie de Chow}
\item \hyperref[intersection-section-phantom]{Théorie de l'intersection}
\item \hyperref[pic-section-phantom]{Schémas de Picard des courbes}
\item \hyperref[weil-section-phantom]{Théories cohomologiques de Weil}
\item \hyperref[adequate-section-phantom]{Modules adéquats}
\item \hyperref[dualizing-section-phantom]{Complexes dualisants}
\item \hyperref[duality-section-phantom]{Dualité pour les schémas}
\item \hyperref[discriminant-section-phantom]{Discriminants et différentes}
\item \hyperref[derham-section-phantom]{Cohomologie de de Rham}
\item \hyperref[local-cohomology-section-phantom]{Cohomologie locale}
\item \hyperref[algebraization-section-phantom]{Géométrie algébrique et formelle}
\item \hyperref[curves-section-phantom]{Courbes algébriques}
\item \hyperref[resolve-section-phantom]{Résolution des surfaces}
\item \hyperref[models-section-phantom]{Réduction semi-stable}
\item \hyperref[functors-section-phantom]{Foncteurs et morphismes}
\item \hyperref[equiv-section-phantom]{Catégories dérivées de variétés}
\item \hyperref[pione-section-phantom]{Groupes fondamentaux de schémas}
\item \hyperref[etale-cohomology-section-phantom]{Cohomologie étale}
\item \hyperref[crystalline-section-phantom]{Cohomologie cristalline}
\item \hyperref[proetale-section-phantom]{Cohomologie pro-étale}
\item \hyperref[relative-cycles-section-phantom]{Cycles relatifs}
\item \hyperref[more-etale-section-phantom]{Compléments de cohomologie étale}
\item \hyperref[trace-section-phantom]{La formule des traces}
\end{enumerate}
Espaces algébriques
\begin{enumerate}
\setcounter{enumi}{64}
\item \hyperref[spaces-section-phantom]{Espaces algébriques}
\item \hyperref[spaces-properties-section-phantom]{Propriétés des espaces algébriques}
\item \hyperref[spaces-morphisms-section-phantom]{Morphismes d'espaces algébriques}
\item \hyperref[decent-spaces-section-phantom]{Espaces algébriques décents}
\item \hyperref[spaces-cohomology-section-phantom]{Cohomologie des espaces algébriques}
\item \hyperref[spaces-limits-section-phantom]{Limites d'espaces algébriques}
\item \hyperref[spaces-divisors-section-phantom]{Diviseurs sur les espaces algébriques}
\item \hyperref[spaces-over-fields-section-phantom]{Espaces algébriques sur des corps}
\item \hyperref[spaces-topologies-section-phantom]{Topologies sur les espaces algébriques}
\item \hyperref[spaces-descent-section-phantom]{Descente et espaces algébriques}
\item \hyperref[spaces-perfect-section-phantom]{Catégories dérivées d'espaces}
\item \hyperref[spaces-more-morphisms-section-phantom]{Compléments sur les morphismes d'espaces}
\item \hyperref[spaces-flat-section-phantom]{Platitude sur les espaces algébriques}
\item \hyperref[spaces-groupoids-section-phantom]{Groupoïdes dans les espaces algébriques}
\item \hyperref[spaces-more-groupoids-section-phantom]{Compléments sur les groupoïdes dans les espaces}
\item \hyperref[bootstrap-section-phantom]{Amorçage}
\item \hyperref[spaces-pushouts-section-phantom]{Sommes amalgamées d'espaces algébriques}
\end{enumerate}
Thèmes de géométrie
\begin{enumerate}
\setcounter{enumi}{81}
\item \hyperref[spaces-chow-section-phantom]{Groupes de Chow des espaces}
\item \hyperref[groupoids-quotients-section-phantom]{Quotients de groupoïdes}
\item \hyperref[spaces-more-cohomology-section-phantom]{Compléments sur la cohomologie des espaces}
\item \hyperref[spaces-simplicial-section-phantom]{Espaces simpliciaux}
\item \hyperref[spaces-duality-section-phantom]{Dualité pour les espaces}
\item \hyperref[formal-spaces-section-phantom]{Espaces algébriques formels}
\item \hyperref[restricted-section-phantom]{Algébrisation des espaces formels}
\item \hyperref[spaces-resolve-section-phantom]{Résolution des surfaces revisitée}
\end{enumerate}
Théorie des déformations
\begin{enumerate}
\setcounter{enumi}{89}
\item \hyperref[formal-defos-section-phantom]{Théorie formelle des déformations}
\item \hyperref[defos-section-phantom]{Théorie des déformations}
\item \hyperref[cotangent-section-phantom]{Le complexe cotangent}
\item \hyperref[examples-defos-section-phantom]{Problèmes de déformation}
\end{enumerate}
Champs algébriques
\begin{enumerate}
\setcounter{enumi}{93}
\item \hyperref[algebraic-section-phantom]{Champs algébriques}
\item \hyperref[examples-stacks-section-phantom]{Exemples de champs}
\item \hyperref[stacks-sheaves-section-phantom]{Faisceaux sur les champs algébriques}
\item \hyperref[criteria-section-phantom]{Critères de représentabilité}
\item \hyperref[artin-section-phantom]{Axiomes d'Artin}
\item \hyperref[quot-section-phantom]{Espaces de Quot et de Hilbert}
\item \hyperref[stacks-properties-section-phantom]{Propriétés des champs algébriques}
\item \hyperref[stacks-morphisms-section-phantom]{Morphismes de champs algébriques}
\item \hyperref[stacks-limits-section-phantom]{Limites de champs algébriques}
\item \hyperref[stacks-cohomology-section-phantom]{Cohomologie des champs algébriques}
\item \hyperref[stacks-perfect-section-phantom]{Catégories dérivées de champs}
\item \hyperref[stacks-introduction-section-phantom]{Introduction aux champs algébriques}
\item \hyperref[stacks-more-morphisms-section-phantom]{Compléments sur les morphismes de champs}
\item \hyperref[stacks-geometry-section-phantom]{La géométrie des champs}
\end{enumerate}
Thèmes de théorie des espaces de modules
\begin{enumerate}
\setcounter{enumi}{107}
\item \hyperref[moduli-section-phantom]{Champs de modules}
\item \hyperref[moduli-curves-section-phantom]{Espaces de modules de courbes}
\end{enumerate}
Divers
\begin{enumerate}
\setcounter{enumi}{109}
\item \hyperref[examples-section-phantom]{Exemples}
\item \hyperref[exercises-section-phantom]{Exercices}
\item \hyperref[guide-section-phantom]{Guide bibliographique}
\item \hyperref[desirables-section-phantom]{Desiderata}
\item \hyperref[coding-section-phantom]{Style de programmation}
\item \hyperref[obsolete-section-phantom]{Obsolète}
\item \hyperref[fdl-section-phantom]{Licence de documentation libre GNU}
\item \hyperref[index-section-phantom]{Index généré automatiquement}
\end{enumerate}
\end{multicols}


\bibliography {my} \bibliographystyle {amsalpha}


\end{document}
